% Content of the blueprint, shared by the web and the print version.
% Each chapter follows one section of the paper; the chapter files are in chapters/.
% The paper is cited by its own statement numbers and LaTeX labels.

\section*{The paper}

This blueprint follows the paper \emph{A fast dynamo with zero topological entropy} by Lukas
Niebel (arXiv:2608.02586v3) section by section: Section~1 (introduction and main theorem),
Section~2 (the local seed), Section~3 (the velocity field on $\T^3$), Section~4 (particle dynamics
and the absence of exponential stretching), Section~5 (the horizontal operator on $\T^2$),
Section~6 (the background resolvent), Section~7 (the approximate inverse and the unstable
eigenvalue), Section~8 (recovery of the magnetic field), and Appendix~A (the semigroup and its spectral growth), the only appendix.  The paper
is cited by its own statement numbers and LaTeX labels --- ``Theorem~1.2 of the paper'',
``the paper's \texttt{lem:torus-sector-basic}'', ``\texttt{eq:global-u}'' --- and unlabelled
passages by their position relative to a label or a section.  A \texttt{\textbackslash ref}
inside the blueprint prints the \emph{blueprint's} number, which is in general not the paper's.

\section*{Reading guide}

\textbf{What the green marks mean.}  \texttt{\textbackslash leanok} on a \emph{statement}
means that the Lean declarations named in that node's \texttt{\textbackslash lean} list exist
in this repository, elaborate, and together say what the statement says; for a
\texttt{definition} node it means that the Lean definition exists and formalises the object
described.  \texttt{\textbackslash leanok} inside a \emph{proof} means that the Lean proofs of
those declarations are complete --- no \texttt{sorry}, and no hypothesis that the paper
does not also have or discharge somewhere in the chain.  Every statement of the blueprint
carries a \texttt{\textbackslash lean} list and \texttt{\textbackslash leanok}, and so does
every proof: there are no non-green statements, and no statement carries
\texttt{\textbackslash notready}.  This includes every numbered statement of the paper except
Conjecture~1.1 (\texttt{conj:smooth-fast-dynamo}) and the remarks that discuss the literature
or alternatives: Remark~1.3 (\texttt{rem:subsequent-resolution}), Remark~7.5
(\texttt{rem:flexible-buffer-scales}) and Remark~8.4.  Those four have no node and no Lean
counterpart.  The other remarks of the paper are nodes: Remark~1.4
(Lemma~\ref{lem:finish:real-part-not-eigenmode}), Remark~1.5
(Lemmas~\ref{lem:finish:zero-mean} and~\ref{lem:main:velocity-regularity}), Remark~1.6
(Lemma~\ref{lem:main:normalisation}) and Remark~1.7
(Theorem~\ref{thm:vanishing-diffusivity}).  Theorem~\ref{thm:main} is
\texttt{FastDynamo.MainTheoremComplete}, and Lemma~\ref{lem:semigroup-spectral-growth} is
\texttt{FastDynamo.SemigroupSpectralGrowth}.

\textbf{What the \texttt{\textbackslash uses} arrows mean.}  An arrow records a mathematical
dependency of the paper's argument, and it is \emph{not} always an edge of the Lean import
graph.  Five proof arrows of this blueprint point at a node none of whose Lean modules is
imported, even indirectly, by a module of the citing node --- Lemma~\ref{lem:local:isolation} to
Lemma~\ref{lem:jump-horizontal-mode}, Lemma~\ref{lem:exact-renormalization} to
Definition~\ref{def:velocity:scale-family},
Corollary~\ref{cor:scaled-local-resolvents} and
Lemma~\ref{lem:transfer:contour-identity} to Proposition~\ref{prop:local-seed}, and
Theorem~\ref{thm:main} to Lemma~\ref{lem:main:growth-rate-limit}.  In the first four the Lean
statement is quantified over abstract data (a certified interface record, a set of scaling
parameters, a resolvent bound) whose concrete inhabitant is built in the cited node's module; in
the last, the Lean statement of Theorem~\ref{thm:main} is about the growth bound
\texttt{gamma}, and the theorem that this growth bound is the limit of the paper's definition is
proved in a module that imports it.  So the arrow is the right mathematical one while the Lean
dependency runs the other way or not at all.  Conversely, where the Lean dependency does run the
other way a proof arrow is \emph{omitted}: the proofs of
Lemma~\ref{lem:dynamics:conditional-growth-bound} and
Lemma~\ref{lem:dynamics:semigroup-unique} take forward uniqueness of strong solutions as a
hypothesis of an abstract statement, and that hypothesis is discharged, for the reduced
operator, for the all-frequency auxiliary operator and for the curl-form induction operator, by
Lemmas~\ref{lem:dynamics:strong-solution-uniqueness}, \ref{lem:dynamics:product-dissipative}
and~\ref{lem:dynamics:physical-uniqueness}, whose modules import the abstract ones.

\textbf{The velocity class.}  Throughout the blueprint the velocity $u$ ranges over the
axially independent (transverse) class \texttt{TransverseRealW1Inf}: real
$\W^{1,\infty}$ maps $\T^2\to\R^3$, viewed on $\T^3=\T^2\times\T$ as fields independent of
$x_3$.  The velocity constructed in the paper lies in that class, so the existential
statements about the velocity (Theorem~\ref{thm:main}, Theorem~\ref{thm:main-spectral},
Theorem~\ref{thm:main-strengthened} and the nodes that feed them) are \emph{stronger} than
the paper's Theorem~1.2, which asks only for $u\in\W^{1,\infty}(\T^3;\R^3)$; the universally
quantified statements over velocities (Lemma~\ref{lem:positive-eigenfunction-solenoidal},
Proposition~\ref{prop:finish:real-strong-mode}, Lemma~\ref{lem:semigroup-spectral-growth},
Lemma~\ref{lem:dynamics:physical-uniqueness}) are correspondingly \emph{narrower} than the
paper's ``for every divergence-free $u\in\W^{1,\infty}(\T^3;\R^3)$'', and are marked
\texttt{\textbackslash leanok} in that reading.  Clauses that the paper states on
$\T^3$ are formalised on the product $\T^2\times\T$; the literal three-torus
$(\R/\Z)^3$ is reached through the measure-preserving identification and the $\Lp^2$
isometries of \texttt{FastDynamo/Analysis/LiteralThreeTorus.lean}, which is how membership in
$\Lp^2(\T^3;\R^3)$ and the norm law of the real solution curve are read on $\T^3$
(Lemma~\ref{lem:main:literal-torus}); that module is where a reader should look for the
bridge.

\textbf{Where the differences are recorded.}  Every node whose Lean form differs from the
paper's --- a different route, a different shape of an estimate, an explicit numerical
threshold in place of ``sufficiently large'', a clause proved chart-side rather than on the
torus, a numeral chosen by the formalisation --- is followed by a short paragraph, outside
the theorem environment, that says exactly how, usually naming the Lean module.  Those
paragraphs, not the statements, are where the differences between the paper and the
formalisation are recorded; a green statement with no such paragraph is formalised as
written.

\chapter{Introduction: the theorem and its Lean statement}
This chapter fixes the objects of Section~1 of the paper -- the induction
operator on the flat three-torus and the growth rate $\gamma(u,\eps)$ --
together with the axial Fourier reduction used by the formalisation, and states
the main theorem, Theorem~1.2 of the paper (\texttt{thm:main}), together with
the Lean statement layer that carries it.

Theorem~1.2 of the paper is, up to the deviations recorded after Theorem~\ref{thm:main}, the
Lean predicate \texttt{FastDynamo.MainTheoremComplete}, proved by \texttt{FastDynamo.mainTheoremComplete\_holds}
(Theorem~\ref{thm:main-complete}); it has no hypothesis.  Its \emph{spectral half} is: one
real, divergence-free, axially independent
$\W^{1,\infty}$ velocity field, two constants $\eps_0,\gamma_0>0$ fixed before
the diffusivity, and, for every $0<\eps\leq\eps_0$, a non-zero axial Fourier
index $K_\eps$, an eigenvalue $\lambda_\eps$ with $\re\lambda_\eps\geq\gamma_0$
and a complex magnetic eigenmode
$V^\eps=\e^{2\pi iK_\eps x_3}b^\eps$ of the induction operator, together with
the non-zero real datum $\re V^\eps$ and an explicit real, divergence-free
strong solution obeying the exact exponential norm law.  Its existential form is the Lean
theorem \texttt{FastDynamo.mainPhysicalConclusion\_holds} (alias
\texttt{FastDynamo.main\_theorem\_spectral}).

The route by which that theorem is reached is the paper's, and it is short:
a two-component local seed on the plane (Proposition~\ref{prop:local-seed}), a multiscale
velocity on the torus (Lemma~\ref{lem:global-velocity}), the horizontal operator
$A_{\eps,K}$ on $\Hs^2(\T^2;\CC^2)$ with \emph{no} divergence constraint
(Lemma~\ref{lem:torus-sector-basic}) and its rescaling to the seed
(Lemma~\ref{lem:exact-renormalization}), one weighted-test coercivity estimate for the
background (Proposition~\ref{prop:background-resolvent}), a two-summand parametrix
(Proposition~\ref{prop:global-spectral-transfer}), and the algebraic recovery of the third
magnetic component (Lemma~\ref{lem:axial-recovery}).

The \emph{dynamical} half of Theorem~\ref{thm:main} is the last sentence of the paper's
statement: the topological entropy of every time map of the particle flow vanishes, and the
ideal induction group has exponential growth rate zero in $\Lp^2_\sigma$ operator norm.  It is
the Lean predicate \texttt{FastDynamo.MainTheoremDynamicalClauses}, discharged at the
constructed velocity by \texttt{FastDynamo.mainTheoremDynamical\_holds}; the flow is written
down in closed form and the ideal group is built from the Cauchy formula, in the
particle-dynamics chapter.  The paper proves those two clauses
from Proposition~\ref{prop:particle-flow} and Corollary~\ref{cor:ideal-growth}, which say
more: the linear Lipschitz bound on the flow maps and the two-sided per-datum bounds on the
ideal group.  Lean keeps that surplus in the larger bundle
Proposition~\ref{prop:main-lagrangian}, from which the theorem's two clauses are a projection.
The endpoint that joins the two halves is Theorem~\ref{thm:main-complete}.

The growth-rate clause of the paper's statement, $\gamma(u,\eps)\geq\gamma_0$, is
formalised by the paper's argument (the growing mode is an exponential minorant of the
semigroup norm).  The semigroup on
$\Lp^2_\sigma(\T^3)$ is constructed in Lean from a Hille--Yosida ray, one axial sector at a
time and then by restriction to the solenoidal subspace
(Theorem~\ref{thm:dynamics:solenoidal-semigroup}), and its generator is identified with the
paper's $L_{\eps,u}$ with domain $\Hs^2\cap\Lp^2_\sigma$
(Theorem~\ref{thm:dynamics:solenoidal-generator}), so it is the paper's $\e^{tL_{\eps,u}}$; its
growth bound is \texttt{FastDynamo.gamma}, which is a limit for every $(u,\eps)$
(Lemma~\ref{lem:main:growth-rate-limit}), and the clause is
Theorem~\ref{thm:main-growth-rate}.  The whole of the paper's Lemma~A.1
(\texttt{lem:semigroup-spectral-growth}) --- analyticity, compact resolvent, immediate
compactness, spectral mapping, and $\gamma(u,\eps)$ as the attained spectral bound --- is
Lemma~\ref{lem:semigroup-spectral-growth}, although Theorem~\ref{thm:main} uses none of it.
``Its solution'' is backed by a uniqueness theorem for the curl-form equation
(Lemma~\ref{lem:dynamics:physical-uniqueness}), and $\mathcal T_0$ is \emph{the} solution group
of the ideal equation in weak form (Corollary~\ref{cor:ideal-growth}).

The remarks and definitions of Section~1 of the paper have nodes: the definitions of a
dynamo, a fast dynamo and a slow dynamo (Definition~\ref{def:main:dynamo}, with
Corollary~\ref{cor:main:fast-dynamo}); Remark~1.4, that the real part of the eigenmode is not
an eigenmode (Lemma~\ref{lem:finish:real-part-not-eigenmode}); Remark~1.5
(\texttt{rem:scope-intro}), the zero mean of the magnetic fields
(Lemma~\ref{lem:finish:zero-mean}) and the regularity of the velocity
(Lemma~\ref{lem:main:velocity-regularity}); Remark~1.6, the normalisation
(Lemma~\ref{lem:main:normalisation}); Remark~1.7, the vanishing-diffusivity limit
(Theorem~\ref{thm:vanishing-diffusivity}); and the energy identity
\texttt{eq:energy-competition} (Lemma~\ref{lem:dynamics:energy-identity}).  Conjecture~1.1
(\texttt{conj:smooth-fast-dynamo}) and Remark~1.3 (\texttt{rem:subsequent-resolution}), which
discuss the literature, have no node and no Lean counterpart; the contrast that Remark~1.3
draws for the velocity constructed here is the ideal--resistive gap recorded after
Theorem~\ref{thm:main}.

\section{The induction operator and the statement layer}

\begin{definition}[The induction operator and its axial reduction]
  \label{def:main:induction-operator}
  \lean{FastDynamo.TransverseRealW1Inf, FastDynamo.IsTransverseDivergenceFree,
    FastDynamo.TransverseC3H2, FastDynamo.productPhysicalInductionOperator,
    FastDynamo.productInductionOperator, FastDynamo.reducedInductionOperator,
    FastDynamo.axialLiftH2, FastDynamo.productDivergenceFromH2,
    FastDynamo.DomainOperator.IsEigenpair,
    FastDynamo.DomainOperator.IsStrongSolution}
  \leanok
  For a divergence-free $u\in\W^{1,\infty}(\T^3;\R^3)$ and $\eps>0$, the
  induction operator is
  \[
    L_{\eps,u}B:=\curl(u\times B)+\eps\Delta B
  \]
  on $\Lp^2_\sigma(\T^3)=\{B\in\Lp^2(\T^3;\CC^3):\Div B=0\}$, with domain
  $\Hs^2(\T^3;\CC^3)\cap\Lp^2_\sigma(\T^3)$.  The formalisation also uses the
  auxiliary operator on all vector fields,
  \[
    \mathcal{A}_\eps B
    =\eps\Delta B-(u\cdot\nabla)B+(B\cdot\nabla)u,
    \qquad
    D(\mathcal{A}_\eps)=\Hs^2(\T^3;\CC^3),
  \]
  which agrees with $L_{\eps,u}$ on divergence-free fields.  For a velocity
  independent of $x_3$ every fixed-$K$ axial Fourier subspace is invariant
  under $\mathcal{A}_\eps$: writing $B(y,x_3)=\e^{2\pi iKx_3}b(y)$ and
  $u=(u_\perp,u_3)$,
  \begin{equation}
    A_{\eps,K}b
    =\eps\bigl(\Delta_\perp-(2\pi K)^2\bigr)b
    -u_\perp\cdot\nabla_\perp b-2\pi iKu_3b
    +(b_\perp\cdot\nabla_\perp)u,
    \qquad
    D(A_{\eps,K})=\Hs^2(\T^2;\CC^3).
    \label{eq:main:torus-sector}
  \end{equation}
  The horizontal block of \eqref{eq:main:torus-sector}, which acts on $b_\perp$
  and does not involve $b_3$, is the paper's horizontal operator
  $A_{\eps,K}$ on $\Lp^2(\T^2;\CC^2)$ with domain $\Hs^2(\T^2;\CC^2)$
  (\texttt{eq:torus-sector}).
\end{definition}

In Lean the three operators are
\texttt{FastDynamo.productPhysicalInductionOperator u eps} ($L_{\eps,u}$),
\texttt{FastDynamo.productInductionOperator u eps} ($\mathcal{A}_\eps$) and
\texttt{FastDynamo.reducedInductionOperator u eps K} ($A_{\eps,K}$), and the
lift $b\mapsto\e^{2\pi iKx_3}b$ is
\texttt{FastDynamo.axialLiftH2 K}.  Five points of the encoding matter for
reading every later node.

\begin{itemize}
  \item \emph{The auxiliary operator and the three-component reduction are the
    formalisation's.}  The paper has no auxiliary operator: it works with
    $L_{\eps,u}$ on $\Lp^2_\sigma$ (\texttt{eq:physical-generator}) and, for the
    spectral construction, with the two-component horizontal operator $A_{\eps,K}$ of
    its \texttt{eq:torus-sector}, recovering the axial magnetic component afterwards
    (Lemma~8.1 of the paper, \texttt{lem:axial-recovery}).  Lean's
    \texttt{reducedInductionOperator} is the three-component operator
    \eqref{eq:main:torus-sector}, whose horizontal block is the paper's $A_{\eps,K}$,
    and the passage from the reduced to the physical eigenpair goes through
    $\mathcal{A}_\eps$.
  \item Each operator is a \texttt{FastDynamo.DomainOperator}: a pair of
    bounded maps out of a Banach \emph{graph} domain, an injective inclusion
    into the ambient space and the operator itself.  Thus
    $\W^{1,\infty}(\T^2;\R^3)$, $\Hs^2(\T^2;\CC^3)$ and $\Hs^2(\T^3;\CC^3)$ are
    realised as the closed Fourier-graph submodules
    \texttt{FastDynamo.TransverseRealW1Inf},
    \texttt{FastDynamo.TransverseC3H2} and \texttt{FastDynamo.ProductC3H2},
    complete by construction.  An eigenpair is
    \texttt{DomainOperator.IsEigenpair}, that is, the ambient image of the
    eigenvector is non-zero and the operator acts on it by the scalar; a
    strong solution is \texttt{DomainOperator.IsStrongSolution}, that is, the
    ambient curve is differentiable in time with derivative the operator
    applied to the state.
  \item The velocity type is \emph{transverse}: it is a $\W^{1,\infty}$ field
    on $\T^2$ with values in $\R^3$, so axial independence is part of the type
    rather than an extra hypothesis.  Since the paper's construction also
    produces an $x_3$-independent field, the Lean existential ranges over a
    smaller class than the paper's and is therefore slightly
    \emph{stronger}.  The identification of $\W^{1,\infty}$ with the Lipschitz
    class is not itself formalised; $\W^{1,\infty}$ here means the Fourier
    graph of an essentially bounded value together with two essentially bounded
    weak derivatives, and divergence-freeness is the equation
    \texttt{IsTransverseDivergenceFree u} on the two diagonal weak
    derivatives (there is no axial term because $\partial_3u=0$).
  \item The statement layer has no space $\Lp^2_\sigma$ (the closed solenoidal
    subspace of $\Lp^2(\T^3;\CC^3)$ enters only with the semigroup,
    Definition~\ref{def:main:growth-rate}).  Solenoidality of a field is the
    side condition \texttt{productDivergenceFromH2 V = 0}, and the operators are
    defined on all vector fields, as $\mathcal{A}_\eps$ is.
  \item $L_{\eps,u}$ is realised weakly: \texttt{productPhysicalInductionOperator}
    maps $\Hs^2$ into a Fourier $\Hs^{-2}$ space, because $\curl(u\times B)$ with
    $u\in\W^{1,\infty}$ is a priori only a distribution.  The statement layer
    therefore asserts \emph{both} the $\Lp^2$-valued eigenvalue equation for
    $\mathcal{A}_\eps$ and the $\Hs^{-2}$-valued one for $L_{\eps,u}$; together
    with solenoidality these are equivalent to the paper's $\Lp^2$
    display, which the physical clause alone would not be.
\end{itemize}

\begin{definition}[The maximal sustained growth rate]
  \label{def:main:growth-rate}
  \lean{FastDynamo.gamma, FastDynamo.literalSolenoidalSubspace,
    FastDynamo.literalInductionSemigroup,
    FastDynamo.StronglyContinuousSemigroup.growthBound}
  \leanok
  \uses{def:dynamics:growth-bound, thm:dynamics:solenoidal-semigroup}
  For a time-independent divergence-free $u\in\W^{1,\infty}(\T^3;\R^3)$ and
  $\eps>0$, set
  \begin{equation}
    \gamma(u,\eps)
    :=\lim_{t\to\infty}\frac1t
    \log\norm{\e^{tL_{\eps,u}}}_{\Lp^2_\sigma\to\Lp^2_\sigma}.
    \label{eq:main:semigroup-growth-rate}
  \end{equation}
  The limit exists because $t\mapsto\log\norm{\e^{tL_{\eps,u}}}$ is subadditive
  and every $\C_0$-semigroup is locally bounded in time
  [Engel--Nagel, Prop.~IV.2.2].
\end{definition}

The Lean counterpart is \texttt{FastDynamo.gamma}
(\texttt{FastDynamo/Dynamics/Assembly/MainGrowthRate.lean}): the growth bound
(Definition~\ref{def:dynamics:growth-bound}) of
\texttt{literalInductionSemigroup} on the closed solenoidal subspace
\texttt{literalSolenoidalSubspace} of $\Lp^2(\T^3;\CC^3)$.  The semigroup is a
$\C_0$-semigroup, constructed and not assumed
(Theorem~\ref{thm:dynamics:solenoidal-semigroup}).  Three points of the encoding.

\begin{itemize}
  \item \emph{Which semigroup.}  \texttt{literalInductionSemigroup} is the isometric transport
    to the literal three-torus of \texttt{solenoidalInductionSemigroup}, whose generator is the
    paper's $L_{\eps,u}$ with domain $\Hs^2\cap\Lp^2_\sigma$, in the equivalent form recorded at
    Theorem~\ref{thm:dynamics:solenoidal-generator}: the values are the auxiliary action
    $\mathcal A_\eps B$ of Definition~\ref{def:main:induction-operator}, which on
    divergence-free $\Hs^2$ fields is the curl form computed in $\Hs^{-2}$.  A
    $\C_0$-semigroup is determined by its generator
    (\texttt{eq\_solenoidalInductionSemigroup\_of\_generator\_eq},
    \texttt{gamma\_eq\_growthBound\_of\_generator\_eq}), so \texttt{gamma} is the growth bound
    of the paper's $\e^{tL_{\eps,u}}$.  The route to the semigroup differs from the paper's:
    Lean generates on all of $\Lp^2(\T^3;\CC^3)$ with the auxiliary operator, restricts to
    $\Lp^2_\sigma$ and identifies the generator afterwards, where the paper generates on
    $\Lp^2_\sigma$ by perturbation of $\eps\Delta$ (Lemma~A.1 of the paper).
  \item \emph{$\limsup$ and $\lim$.}  Lean's \texttt{growthBound} is defined as a
    $\limsup$ in $[-\infty,\infty]$, so that it is defined for every semigroup with no
    preliminary work; Lemma~\ref{lem:dynamics:growth-bound-limit} proves that it is the limit
    for every $\C_0$-semigroup on a Banach space, by the paper's argument, and
    Lemma~\ref{lem:main:growth-rate-limit} records this for $\gamma(u,\eps)$, which is
    moreover a real number in $[0,M_u]$ for every $(u,\eps)$.
  \item \emph{Model.}  The torus is the product model $\T^2\times\T$, transported
    isometrically to the literal three-torus, and the velocity is axially independent;
    \texttt{literalSolenoidalSubspace} is the transported Fourier-defined solenoidal subspace
    of the product model.  The analytic and spectral properties of the semigroup are
    Lemma~\ref{lem:semigroup-spectral-growth}.
\end{itemize}

\begin{definition}[Dynamo, fast dynamo, slow dynamo]
  \label{def:main:dynamo}
  \lean{FastDynamo.gammaOf, FastDynamo.gammaOf_eq_gamma, FastDynamo.IsDynamoAt,
    FastDynamo.IsFastDynamo, FastDynamo.IsSlowDynamo}
  \leanok
  \uses{def:main:growth-rate}
  Let $u\in\W^{1,\infty}(\T^3;\R^3)$ be time-independent and divergence-free.  We call $u$ a
  \emph{dynamo at diffusivity $\eps$} when $\gamma(u,\eps)>0$.  It is a \emph{fast dynamo} if
  \[
    \liminf_{\eps\downarrow0}\gamma(u,\eps)>0.
  \]
  If $\gamma(u,\eps)>0$ for every sufficiently small $\eps$, but $\gamma(u,\eps)\to0$, we call
  it a \emph{slow dynamo}.
\end{definition}

These are the definitions of §1.1 of the paper, verbatim in Lean
(\texttt{FastDynamo/Dynamics/Assembly/FastDynamo.lean}) for a divergence-free velocity of
the axially independent class: \texttt{IsDynamoAt} is $0<\gamma(u,\eps)$, \texttt{IsFastDynamo}
is $0<\liminf$ along $\eps\to0^+$, and \texttt{IsSlowDynamo} is ``eventually positive along
$\eps\to0^+$ and tending to $0$''.  The limits in $\eps$ are taken of \texttt{gammaOf}, the
extension of $\gamma(u,\cdot)$ to all real $\eps$ (by $0$ for $\eps\leq0$), because
\texttt{gamma} takes the proof of $0<\eps$ as an argument; along $\eps\to0^+$ only positive
$\eps$ matter.  A fast dynamo is not a slow dynamo (\texttt{IsFastDynamo.not\_isSlowDynamo}),
which the paper leaves implicit; it is recorded with
Corollary~\ref{cor:main:fast-dynamo}.

\begin{lemma}[The growth rate is a limit]
  \label{lem:main:growth-rate-limit}
  \lean{FastDynamo.gamma_tendsto, FastDynamo.gamma_ne_bot_iff, FastDynamo.gamma_tendsto_real,
    FastDynamo.gamma_eq_lim, FastDynamo.zero_le_gamma,
    FastDynamo.mainGrowthFloor_le_gamma_lim}
  \leanok
  \uses{def:main:growth-rate, lem:dynamics:growth-bound-limit,
    lem:dynamics:mean-conservation, thm:main-growth-rate, def:main:witnesses}
  For every divergence-free $u$ and every $\eps>0$ the limit
  \eqref{eq:main:semigroup-growth-rate} exists and is a real number: there is $L\in\R$ with
  \[
    \frac1t\log\norm{\e^{tL_{\eps,u}}}_{\Lp^2_\sigma\to\Lp^2_\sigma}
    \xrightarrow[t\to\infty]{}L,
    \qquad
    L=\gamma(u,\eps),
    \qquad
    0\leq L\leq M_u .
  \]
  At the explicit witnesses of Definition~\ref{def:main:witnesses} and for every
  $0<\eps\leq\eps_0$, moreover $\gamma_0\leq L$.
\end{lemma}

This is the sentence after \texttt{eq:semigroup-growth-rate} in §1.1 of the paper (the limit
exists), together with the facts that it is finite and non-negative, which the paper does not
state.  In Lean: \texttt{gamma\_tendsto} (convergence in $[-\infty,\infty)$, from
Lemma~\ref{lem:dynamics:growth-bound-limit}), \texttt{gamma\_ne\_bot\_iff} and
\texttt{gamma\_tendsto\_real} (finiteness exactly under an exponential minorant, and then
convergence in $\R$), \texttt{gamma\_eq\_lim} and \texttt{zero\_le\_gamma} (the minorant
$\norm{\e^{tL_{\eps,u}}}\geq1$ of Lemma~\ref{lem:dynamics:mean-conservation}), and
\texttt{mainGrowthFloor\_le\_gamma\_lim} for the witnesses
(\texttt{FastDynamo/Dynamics/Assembly/MainGrowthRateLimit.lean} and
\texttt{FastDynamo/Dynamics/Assembly/InductionMeanConservation.lean}).  The convergence takes
place in \texttt{EReal}, with $\log0$ read as $-\infty$, which is the paper's $[-\infty,\infty)$
reading; the real-valued form is \texttt{gamma\_eq\_lim}.

\begin{proof}
  \leanok
  \uses{lem:dynamics:growth-bound-limit, lem:dynamics:mean-conservation,
    thm:main-growth-rate}
  Lemma~\ref{lem:dynamics:growth-bound-limit} applies to the semigroup on the closed subspace
  $\Lp^2_\sigma$ of a Hilbert space.  The minorant $\norm{\e^{tL_{\eps,u}}}\geq1$ makes the
  limit finite and non-negative, and the Hille--Yosida bound
  $\norm{\e^{tL_{\eps,u}}}\leq\e^{M_ut}$ bounds it by $M_u$.  At the witnesses,
  Theorem~\ref{thm:main-growth-rate} gives $\gamma_0\leq\gamma(u,\eps)$.
\end{proof}


\begin{definition}[The statement layer]
  \label{def:main:physical-clauses}
  \uses{def:main:induction-operator}
  \lean{FastDynamo.MainReducedClauses, FastDynamo.MainReducedTarget,
    FastDynamo.MainPhysicalClausesAt,
    FastDynamo.MainPhysicalClauses, FastDynamo.MainPhysicalConclusion,
    FastDynamo.mainPhysicalClauses_of_mainReducedClauses,
    FastDynamo.mainPhysicalConclusion_of_mainReducedTarget}
  \leanok
  \emph{The reduced clauses} $\mathrm{MRC}(u,\eps_0,\gamma_0)$ say: for every
  $0<\eps\leq\eps_0$ there are $K\in\Z\setminus\{0\}$, $\lambda\in\CC$ with
  $\re\lambda\geq\gamma_0$ and $b\in\Hs^2(\T^2;\CC^3)$ with
  $A_{\eps,K}b=\lambda b$ and $b\neq0$ in $\Lp^2$.  \emph{The reduced target}
  is
  \[
    \exists\,u\in\W^{1,\infty},\ \Div u=0,\ u\neq0,\ \
    \exists\,\eps_0,\gamma_0>0:\ \mathrm{MRC}(u,\eps_0,\gamma_0).
  \]
  \emph{The physical clauses} $\mathrm{MPC}(u,\eps_0,\gamma_0)$ are the
  sixteen conjuncts of the spectral half of Theorem~\ref{thm:main} at a fixed
  velocity and fixed constants, and \emph{the physical conclusion} is the same
  outer existential over $u,\eps_0,\gamma_0$ with $\mathrm{MPC}$ in place of
  $\mathrm{MRC}$.  The reduction
  \[
    \mathrm{MRC}(u,\eps_0,\gamma_0)\ \Longrightarrow\ \mathrm{MPC}(u,\eps_0,\gamma_0)
    \qquad(\Div u=0,\ \gamma_0>0)
  \]
  is the content of the finish layer.
\end{definition}

\begin{proof}
  \leanok
  \uses{def:finish:real-solution, prop:finish:real-strong-mode,
    lem:positive-eigenfunction-solenoidal, lem:real-axial-mode}
  Given a reduced eigenpair, the real strong-mode certificate of the finish
  layer supplies every physical clause at once: the lift
  $V=\e^{2\pi iKx_3}b$ is an eigenmode of $\mathcal{A}_\eps$ by the exact
  sector identity, it is divergence-free by
  Lemma~\ref{lem:positive-eigenfunction-solenoidal} (which needs
  $\re\lambda>0$ and $\Div u=0$) and hence an eigenmode of $L_{\eps,u}$, and
  Lemma~\ref{lem:real-axial-mode} turns the non-zero axial index into the
  non-vanishing of $\re V$ and the exact norm law for the real curve.
\end{proof}

The sixteen clauses of $\mathrm{MPC}(u,\eps_0,\gamma_0)$, in the order in
which \texttt{FastDynamo.MainPhysicalClauses} lists them, are the following.
Clauses (1)--(3) are the paper's
$K_\eps\in\Z\setminus\{0\}$, $\re\lambda_\eps\geq\gamma_0$ and
$0\neq b^\eps\in\Hs^2(\T^2;\CC^3)$.  Clause (4) is the reduced eigenpair
$A_{\eps,K}b=\lambda b$ of \eqref{eq:main:torus-sector}, kept as the reference
datum.  Clauses (5)--(7) carry the display
\eqref{eq:main-spectral-conclusion} of Theorem~\ref{thm:main}: the lift
$V^\eps=\e^{2\pi iK_\eps x_3}b^\eps$ lies in $\Hs^2(\T^3;\CC^3)$ by typing and
satisfies the auxiliary-operator equation $\mathcal{A}_\eps V^\eps=\lambda_\eps V^\eps$
(5), which the paper does not state, it is
divergence-free, hence lies in $\Lp^2_\sigma(\T^3)$ (6), and it satisfies the
curl-form equation $L_{\eps,u}V^\eps=\lambda_\eps V^\eps$ (7).  Clause (8) is
$B_{\rm in}^\eps=\re V^\eps\neq0$.  Clauses (9)--(11) and (15) are the
real-solution law of Theorem~\ref{thm:main}: the explicit real $\Lp^2$ curve
$B^\eps(t)=\re(\e^{t\lambda_\eps}V^\eps)$ never vanishes (9), obeys
$\norm{B^\eps(t)}_2=\e^{t\re\lambda_\eps}\norm{B_{\rm in}^\eps}_2$ (10) and the
uniform bound $\e^{t\gamma_0}\norm{B_{\rm in}^\eps}_2\leq\norm{B^\eps(t)}_2$
for $t\geq0$ (11), and is the real field represented by the $\Hs^2$ state
curve (15).  Clauses (12)--(14) and (16) say that this $\Hs^2$ state curve
starts at $\re V^\eps$, is fixed by conjugation (that is, real) at every time,
is divergence-free at every time, and is a strong solution of the induction
equation.

Two differences from the paper are worth recording.  The norm law (10) is
proved for \emph{every} real time, not only for $t\geq0$ as in
\eqref{eq:main-thm-growth}, and the uniform bound (11) with the
$\eps$-independent rate $\gamma_0$ is an extra clause, an immediate consequence of
\eqref{eq:main-thm-growth} and $\re\lambda_\eps\geq\gamma_0$ that the paper does not
state.  Clause (16) asserts that $B^\eps$ is a strong solution; that it is \emph{the}
solution, as the paper's ``its solution'' says, is
Lemma~\ref{lem:dynamics:physical-uniqueness}: forward strong solutions of the curl-form
equation in $\Hs^{-2}$ are unique.  The semigroup on $\Lp^2_\sigma(\T^3)$ carries the datum
$B_{\rm in}^\eps$ along exactly the curve $B^\eps$, and every forward strong solution with
datum $B_{\rm in}^\eps$ is that orbit; this is \texttt{FastDynamo.MainSemigroupOrbitClauses},
the field \texttt{semigroupOrbit} of Theorem~\ref{thm:main-strengthened}, which also asserts
the sixteen clauses of $\mathrm{MPC}$ at its own triple $(K,\lambda,b)$, so that the eigenmode,
the real datum and its solution belong to one triple.  That clause cannot
live in $\mathrm{MPC}$ itself: \texttt{FastDynamo/Main.lean} is in the import closure of the
semigroup construction, so naming the semigroup there would close a cycle.  The Lean modules
are \texttt{FastDynamo/Main.lean} (the four predicates and the reduction; the body of
$\mathrm{MPC}$ at one triple $(K,\lambda,b)$ is \texttt{MainPhysicalClausesAt}),
\texttt{FastDynamo/Finish/RealSolution.lean} (the certificate that discharges it) and
\texttt{FastDynamo/Dynamics/Assembly/SemigroupOrbit.lean} (the orbit clause).

\section{The main theorem}

\begin{theorem}[A fast dynamo with zero topological entropy]
  \label{thm:main}
  \lean{FastDynamo.MainTheoremComplete, FastDynamo.mainTheoremComplete_holds}
  \leanok
  \uses{def:main:induction-operator, def:main:growth-rate, thm:main-growth-rate,
    prop:particle-flow, cor:ideal-growth}
  There exist a real-valued autonomous velocity field
  \[
    u\in\W^{1,\infty}(\T^3;\R^3),\qquad\Div u=0,
  \]
  and constants $\eps_0>0$, $\gamma_0>0$, such that
  $\gamma(u,\eps)\geq\gamma_0$ for every $0<\eps\leq\eps_0$.  In particular,
  \begin{equation}
    \liminf_{\eps\downarrow0}\gamma(u,\eps)\geq\gamma_0.
    \label{eq:main-growth-conclusion}
  \end{equation}
  More precisely, for every $0<\eps\leq\eps_0$ there are
  $K_\eps\in\Z\setminus\{0\}$, $\lambda_\eps\in\CC$ and
  $0\neq b^\eps\in\Hs^2(\T^2;\CC^3)$ such that the complex magnetic eigenmode
  \begin{equation}
    V^\eps(y,x_3):=\e^{2\pi iK_\eps x_3}b^\eps(y)
    \in\Hs^2(\T^3;\CC^3)\cap\Lp^2_\sigma(\T^3),
    \quad
    L_{\eps,u}V^\eps=\lambda_\eps V^\eps,
    \quad
    \re\lambda_\eps\geq\gamma_0.
    \label{eq:main-spectral-conclusion}
  \end{equation}
  Moreover the real-valued datum
  $B_{\rm in}^\eps:=\re V^\eps\in\Hs^2(\T^3;\R^3)\cap\Lp^2_\sigma(\T^3)$ is
  non-zero, and its solution satisfies
  \begin{equation}
    \norm{B^\eps(t)}_{\Lp^2(\T^3)}
    =\e^{t\re\lambda_\eps}\norm{B_{\rm in}^\eps}_{\Lp^2(\T^3)}
    \qquad\text{for every }t\geq0.
    \label{eq:main-thm-growth}
  \end{equation}

  Moreover, $h_{\mathrm{top}}(\Phi_t)=0$ for all $t\in\R$ and
  \[
    \lim_{t\to\infty}\frac1t
    \log\norm{\mathcal T_0(t)}_{\Lp^2_\sigma\to\Lp^2_\sigma}=0,
  \]
  where $\Phi_t$ is the particle flow of $u$
  (Definition~\ref{def:dynamics:particle-flow}) and $\mathcal T_0(t)$ the solution
  group on $\Lp^2_\sigma(\T^3)$ of the ideal induction equation
  (Definition~\ref{def:dynamics:ideal-induction-group}).
\end{theorem}

\begin{proof}
  \leanok
  \uses{thm:main-spectral, thm:main-growth-rate, lem:main:growth-rate-limit,
    lem:real-axial-mode, prop:particle-flow, cor:ideal-growth, thm:main-strengthened,
    lem:dynamics:physical-uniqueness}
  Everything except the growth clause and \eqref{eq:main-growth-conclusion} is
  Theorem~\ref{thm:main-spectral}; the dynamical half is
  Proposition~\ref{prop:particle-flow} and Corollary~\ref{cor:ideal-growth}, both of which
  are formalised (Proposition~\ref{prop:main-lagrangian}).  The growth clause is
  Theorem~\ref{thm:main-growth-rate}: the eigenvalue $\lambda_\eps$ produced
  for each $0<\eps\leq\eps_0$ lies in the state space $\Lp^2_\sigma(\T^3)$ and
  the semigroup orbit of its eigenmode grows exactly like
  $\e^{t\re\lambda_\eps}$, which gives
  $\gamma(u,\eps)\geq\re\lambda_\eps\geq\gamma_0$; taking the lower limit as
  $\eps\downarrow0$ finishes.

  The paper's own route to the growth clause is the same in shape.  It reads the
  exact growth law \eqref{eq:main-thm-growth} of the \emph{real} datum as an exponential
  minorant of the operator norm,
  $\norm{\e^{tL_{\eps,u}}}\geq\norm{B^\eps(t)}_2/\norm{B_{\rm in}^\eps}_2
  =\e^{t\re\lambda_\eps}$, and concludes
  $\gamma(u,\eps)\geq\re\lambda_\eps\geq\gamma_0$; the semigroup appendix,
  Lemma~\ref{lem:semigroup-spectral-growth}, is not invoked here at all.  Lean divides by the
  norm of the \emph{complex} eigenmode $V^\eps$ rather than of the real datum
  $B_{\rm in}^\eps$ --- the two give the same inequality, $V^\eps\neq0$ being part of the
  definition of an eigenpair --- and \eqref{eq:main-thm-growth} itself, which the paper
  obtains from Lemma~\ref{lem:real-axial-mode}, is clause (10) of
  Definition~\ref{def:main:physical-clauses}.
\end{proof}

The Lean statement is \texttt{FastDynamo.MainTheoremComplete}, proved by
\texttt{mainTheoremComplete\_holds} (\texttt{FastDynamo/MainTheoremStrengthened.lean},
\texttt{FastDynamo/Dynamics/MainLagrangianClauses.lean}); it is
Theorem~\ref{thm:main-complete}, whose fields are the strengthened spectral statement
(Theorem~\ref{thm:main-strengthened}), the two dynamical clauses
(\texttt{MainTheoremDynamical}) and the larger Lagrangian bundle
(Proposition~\ref{prop:main-lagrangian}).  Clause by clause: the growth clause and
\eqref{eq:main-growth-conclusion} are Theorem~\ref{thm:main-growth-rate}
(\texttt{gamma\_ge\_gamma0}, \texttt{mainGrowthFloor\_le\_liminf\_gammaOf}), with
$\gamma(u,\eps)$ the limit of Lemma~\ref{lem:main:growth-rate-limit};
\eqref{eq:main-spectral-conclusion}, the non-vanishing of $B_{\rm in}^\eps$ and the exact norm
law \eqref{eq:main-thm-growth} are Theorem~\ref{thm:main-spectral}; ``its solution'' is the
semigroup-orbit clause of Theorem~\ref{thm:main-strengthened}, which says that $B_{\rm in}^\eps$
lies in $\Lp^2_\sigma$, that its orbit under $\e^{tL_{\eps,u}}$ is the explicit real curve
$B^\eps(t)=\re(\e^{t\lambda_\eps}V^\eps)$ for $t\geq0$, and that \emph{every} forward strong
solution of the curl-form equation with datum $B_{\rm in}^\eps$ is that orbit and obeys
\eqref{eq:main-thm-growth} (Lemma~\ref{lem:dynamics:physical-uniqueness}); and the two
dynamical clauses are \texttt{FastDynamo.mainTheoremDynamical\_holds}, with $\mathcal T_0$ the
unique solution group of the ideal equation (Corollary~\ref{cor:ideal-growth}).

The Lean statement differs from the paper's in these respects, each recorded at the node
named.
\begin{itemize}
  \item \emph{Shape of the statement.}  The clauses are fields and theorems at the explicit
    witnesses of Definition~\ref{def:main:witnesses}, not one proposition of the form
    $\exists u,\eps_0,\gamma_0$, and the per-diffusivity predicates
    \texttt{MainPhysicalClauses}, \texttt{MainSemigroupOrbitClauses} and
    \texttt{MainLiteralTorusClauses} each carry their own existential over $(K,\lambda,b)$.
    \texttt{MainSemigroupOrbitClauses} contains the sixteen clauses of
    Definition~\ref{def:main:physical-clauses} (\texttt{MainPhysicalClausesAt}) at its own
    triple, so the eigenmode, the real datum, the strong solution, the semigroup orbit and the
    uniqueness clause hold at one $(K,\lambda,b)$.  The triple of
    \texttt{MainLiteralTorusClauses} is the one of \texttt{mainSemigroupOrbitClauses\_holds}
    by its proof, not by its statement, and the field \texttt{physicalClauses} of
    Theorem~\ref{thm:main-strengthened} repeats the spectral clauses with a separately chosen
    triple (the paragraph after Theorem~\ref{thm:main-strengthened}).
  \item \emph{Velocity class and torus.}  The velocity ranges over the axially independent
    class \texttt{TransverseRealW1Inf}, so the existential is over a smaller set than the
    paper's and is stronger; the magnetic clauses live on the product model $\T^2\times\T$,
    transported isometrically to the literal three-torus (reading guide,
    Lemma~\ref{lem:main:literal-torus}).
  \item \emph{The operator} $L_{\eps,u}$ is realised as the curl form with values in
    $\Hs^{-2}$ in the statement layer and, as generator, in the equivalent auxiliary form on
    $\Hs^2\cap\Lp^2_\sigma$ (Definition~\ref{def:main:induction-operator},
    Theorem~\ref{thm:dynamics:solenoidal-generator}).
  \item \emph{The solution class} of ``its solution'' is not named by the paper; Lean takes the
    forward strong solutions of the curl-form equation in $\Hs^{-2}$
    (Lemma~\ref{lem:dynamics:physical-uniqueness}), and the clause adds the lower bound
    $\e^{t\gamma_0}\norm{B_{\rm in}^\eps}_2\leq\norm{B^\eps(t)}_2$, which the theorem does not
    state.
  \item \emph{The ideal solution group} is unique among the solution groups of the weak form
    of Definition~\ref{def:dynamics:ideal-weak-form}, whose test fields and time integration
    differ from the paper's distributional notion as recorded there.
\end{itemize}

\textbf{The ideal--resistive growth gap.}  The paper folds this contrast into
Theorem~\ref{thm:main}, whose closing line is its left member and
whose growth clause is its right member: for the velocity constructed here the ideal
induction group has exponential growth rate zero, while every sufficiently small positive
diffusivity grows at least at the rate $\gamma_0>0$,
\[
  \lim_{t\to\infty}\frac1t
  \log\norm{\mathcal T_0(t)}_{\Lp^2_\sigma\to\Lp^2_\sigma}
  =0<\gamma_0\leq\liminf_{\eps\downarrow0}\gamma(u,\eps).
\]
The chain is not a statement of the paper, and therefore not a node of this
blueprint; the paper's own qualitative reading of it is the second paragraph of its
Remark~1.3 (\texttt{rem:subsequent-resolution}).  In Lean,
\texttt{FastDynamo.ideal\_resistive\_gap\_holds} is the display itself, and
\texttt{FastDynamo.ideal\_resistive\_gap\_at\_diffusivity\_holds} is the sharper fixed-$\eps$
packaging, with $\gamma_0\leq\gamma(u,\eps)$ at every $0<\eps\leq\eps_0$ in place of the lower
limit.  Its left end is \texttt{HasVanishingIdealOperatorGrowthRate}
(Corollary~\ref{cor:ideal-growth}), its middle is $\gamma_0>0$, and its right end is
\texttt{mainGrowthFloor\_le\_liminf\_gammaOf}, the lower limit being taken along the filter
$\mathcal N_{>0}$ of the total extension \texttt{gammaOf} of $\gamma(u,\cdot)$.

\begin{theorem}[The spectral half of Theorem~\ref{thm:main}]
  \label{thm:main-spectral}
  \uses{def:main:physical-clauses}
  \lean{FastDynamo.MainPhysicalConclusion,
    FastDynamo.mainPhysicalConclusion_holds,
    FastDynamo.main_theorem_spectral, FastDynamo.mainReducedTarget_holds}
  \leanok
  There exist a real-valued autonomous $u\in\W^{1,\infty}(\T^3;\R^3)$ with
  $\Div u=0$ and $u\neq0$, and constants $\eps_0,\gamma_0>0$, such that for
  every $0<\eps\leq\eps_0$ the sixteen clauses
  $\mathrm{MPC}(u,\eps_0,\gamma_0)$ of
  Definition~\ref{def:main:physical-clauses} hold: an axial index
  $K_\eps\neq0$, an eigenvalue with $\re\lambda_\eps\geq\gamma_0$ and a
  non-zero $b^\eps\in\Hs^2(\T^2;\CC^3)$ producing the solenoidal $\Hs^2$
  eigenmode $V^\eps=\e^{2\pi iK_\eps x_3}b^\eps$ of
  \eqref{eq:main-spectral-conclusion}, the non-zero real datum
  $B_{\rm in}^\eps=\re V^\eps$, and the explicit real, divergence-free strong
  solution $B^\eps$ with
  $\norm{B^\eps(t)}_2=\e^{t\re\lambda_\eps}\norm{B_{\rm in}^\eps}_2$ for every
  $t\in\R$ and $\norm{B^\eps(t)}_2\geq\e^{t\gamma_0}\norm{B_{\rm in}^\eps}_2$
  for $t\geq0$.

  Equivalently: Theorem~\ref{thm:main} holds with its growth clause and its two dynamical
  clauses deleted, and ``its solution'' read as ``an explicit real, divergence-free strong
  solution with that datum''.
\end{theorem}

\begin{proof}
  \leanok
  \uses{def:transfer:parameter-choice, prop:global-spectral-transfer,
    lem:global-riesz-nonzero, lem:axial-recovery,
    lem:real-axial-mode, def:finish:real-solution, prop:finish:real-strong-mode}
  Fix the global parameter chain of
  Definition~\ref{def:transfer:parameter-choice}; it pins the velocity $u$,
  the threshold $\eps_0$ and, for each $0<\eps\leq\eps_0$, a matching cell and
  its axial index $K_\eps\neq0$.  Proposition~\ref{prop:global-spectral-transfer}
  gives, at that chain and for every such $\eps$, a resolvent of
  $A_{\eps,K_\eps}$ on the selected contour whose Riesz projection is non-zero;
  this is the step whose non-vanishing is
  Lemma~\ref{lem:global-riesz-nonzero}.  The projection has finite rank, so a
  non-zero Riesz projection yields an eigenpair of $A_{\eps,K_\eps}$
  inside the contour, whose real part is at least the contour's growth floor
  $\gamma_0=g_*$ -- a closed term of the fixed local seed, independent of
  $\eps$, of the velocity, of the sector and of the placement.  The eigenpair is
  \emph{horizontal}; Lemma~\ref{lem:axial-recovery} recovers the third component and
  delivers the reduced target.  The reduction of
  Definition~\ref{def:main:physical-clauses}, which uses
  Lemma~\ref{lem:real-axial-mode} through the certificate of
  Definition~\ref{def:finish:real-solution}, then produces all sixteen physical
  clauses.
\end{proof}

The Lean proof follows this route.  \texttt{mainReducedTarget\_holds} and
\texttt{mainPhysicalConclusion\_holds} in \texttt{FastDynamo/MainTheorem.lean} are obtained
from \texttt{mainReducedClauses\_holds}
(\texttt{FastDynamo/Global/MainTheoremWitnesses.lean}), which is three lines through the
modules
\texttt{Global/GlobalParameterChoiceConcrete.lean} $\to$
\texttt{Global/GlobalSpectralTransferConcrete.lean} $\to$
\texttt{Global/TorusGlobalSpectralTransfer.lean}; the finish layer is entered exactly once,
through \texttt{realStrongModeCertificate\_of\_reducedEigenpair}.
Corollary~\ref{cor:magnetic-scale} (the two-sided bounds
$\abs{K_\eps}\asymp\eps^{-1/2}$ and $L_B^\eps\asymp\sqrt\eps$) is \emph{not} used anywhere
in this chain: it is a separate quantitative statement about the modes once they have been
produced, and no step of the Lean proof invokes it.

\begin{theorem}[The growth-rate clause of Theorem~\ref{thm:main}]
  \label{thm:main-growth-rate}
  \uses{def:main:growth-rate, thm:main-spectral, thm:dynamics:solenoidal-semigroup,
    lem:dynamics:solenoidal-subspace, def:main:witnesses}
  \lean{FastDynamo.gamma_ge_gamma0, FastDynamo.mainGrowthFloor_le_gamma_le_stretching,
    FastDynamo.gamma_le_stretching,
    FastDynamo.gamma_ne_bot, FastDynamo.gamma_ne_top,
    FastDynamo.gammaOf, FastDynamo.gammaOf_eq_gamma,
    FastDynamo.mainGrowthFloor_le_liminf_gammaOf}
  \leanok
  At the explicit witnesses $u$, $\eps_0$, $\gamma_0$ of
  Definition~\ref{def:main:witnesses}, for every $0<\eps\leq\eps_0$,
  \[
    \gamma_0\ \leq\ \gamma(u,\eps)\ \leq\ M_u,
    \qquad
    M_u=\norm{\nabla_\perp u}_\infty,
  \]
  $\gamma(u,\eps)$ being the growth bound (Definition~\ref{def:main:growth-rate})
  of the semigroup on $\Lp^2_\sigma(\T^3;\CC^3)$ of
  Theorem~\ref{thm:dynamics:solenoidal-semigroup}.  In particular
  $\gamma(u,\eps)$ is neither $-\infty$ nor $+\infty$.  Since $\gamma_0$ does not
  depend on $\eps$, this is the growth clause of Theorem~\ref{thm:main}, and with it
  \eqref{eq:main-growth-conclusion}.
\end{theorem}

The lower bound is \texttt{FastDynamo.gamma\_ge\_gamma0}, the upper bound
\texttt{gamma\_le\_stretching}, and the pair
\texttt{mainGrowthFloor\_le\_gamma\_le\_stretching}
(\texttt{FastDynamo/Dynamics/Assembly/MainGrowthRate.lean}); both are fields of
the strengthened theorem (Theorem~\ref{thm:main-strengthened}).  With
Lemma~\ref{lem:main:growth-rate-limit}, which shows that $\gamma(u,\eps)$ is the paper's limit,
this is \eqref{eq:main-growth-conclusion} in the form displayed there; the packaged
statement is \texttt{mainGrowthFloor\_le\_gamma\_lim}.  Two smaller points.  The
paper's trailing restatement
$\liminf_{\eps\downarrow0}\gamma(u,\eps)\geq\gamma_0$ is
\texttt{mainGrowthFloor\_le\_liminf\_gammaOf}, the $\liminf$ along
$\mathcal N_{>0}$ of the total extension \texttt{gammaOf} of $\gamma(u,\cdot)$ ---
\texttt{gamma} carries the proof $0<\eps$ as an argument, so a filter limit in $\eps$ needs
that extension, which agrees with $\gamma$ on $(0,\infty)$ by
\texttt{gammaOf\_eq\_gamma}.  It is weaker than the per-diffusivity bound and is recorded
because it is the form of \eqref{eq:main-growth-conclusion} and of the paper's definition of
a fast dynamo.  The interval is the paper's: \texttt{gamma\_ge\_gamma0} and
\texttt{MainTheoremStrengthened.growthRate} carry the hypotheses $0<\eps$ and
$\eps\leq\eps_0$, the closed interval $0<\eps\leq\eps_0$ on which the paper states the growth
clause.  The semigroup is the paper's $\e^{tL_{\eps,u}}$, obtained by restriction and
identified by its generator (Theorem~\ref{thm:dynamics:solenoidal-generator}).  The upper bound is not in the
paper; it is
recorded because it is what makes the sandwich informative --- without it the
lower bound would be compatible with $\gamma=+\infty$ --- and because it is
uniform in the diffusivity, which is the paper's energy competition
(\texttt{eq:energy-competition}) read at the level of the semigroup.  The floor $\gamma_0$ is
a closed term of the fixed local seed with $\gamma_0\geq\tfrac34(448-4\pi^2)\geq288$; these
numerals are choices of the formalisation, not of the paper
(Definition~\ref{def:main:witnesses}).

\begin{proof}
  \leanok
  \uses{thm:main-spectral, thm:dynamics:solenoidal-semigroup,
    lem:dynamics:solenoidal-subspace, lem:positive-eigenfunction-solenoidal}
  This is the paper's route, with one extra Lean-side obligation.  Fix
  $0<\eps\leq\eps_0$.  The reduced clauses of Theorem~\ref{thm:main-spectral}
  supply an axial index $K_\eps\neq0$, an eigenvalue $\lambda_\eps$ with
  $\re\lambda_\eps\geq\gamma_0$ and a non-zero $b^\eps\in\Hs^2(\T^2;\CC^3)$ with
  $A_{\eps,K_\eps}b^\eps=\lambda_\eps b^\eps$; the exact sector identity promotes
  this to an eigenpair of $\mathcal{A}_\eps$ at the lift
  $V^\eps=\e^{2\pi iK_\eps x_3}b^\eps$.  The extra obligation is that $V^\eps$
  lie in the state space: since $\gamma_0>0$ forces $\re\lambda_\eps>0$, the
  Lemma~\ref{lem:positive-eigenfunction-solenoidal}, in the membership
  form of Lemma~\ref{lem:dynamics:solenoidal-subspace}, puts $V^\eps$ in
  $\Lp^2_\sigma$.  Theorem~\ref{thm:dynamics:solenoidal-semigroup} then converts
  the eigenvalue into $\re\lambda_\eps\leq\gamma(u,\eps)$ and supplies the
  ceiling $\gamma(u,\eps)\leq M_u$.  Since the lower bound holds for every
  $0<\eps\leq\eps_0$ with an $\eps$-independent $\gamma_0$, the lower limit as
  $\eps\downarrow0$ is at least $\gamma_0$.
\end{proof}

\begin{corollary}[The constructed velocity is a fast dynamo]
  \label{cor:main:fast-dynamo}
  \lean{FastDynamo.isDynamoAt_mainVelocity, FastDynamo.isFastDynamo_mainVelocity,
    FastDynamo.IsFastDynamo.not_isSlowDynamo}
  \leanok
  \uses{def:main:dynamo, thm:main-growth-rate, def:main:witnesses}
  The velocity $u$ of Definition~\ref{def:main:witnesses} is a dynamo at every diffusivity
  $0<\eps\leq\eps_0$ and a fast dynamo in the sense of Definition~\ref{def:main:dynamo}; in
  particular it is not a slow dynamo.
\end{corollary}

The paper draws this conclusion in its title and abstract and through
\eqref{eq:main-growth-conclusion}, which is its definition of a fast dynamo; it does not state
it as a separate result.  Module \texttt{FastDynamo/Dynamics/Assembly/FastDynamo.lean}.

\begin{proof}
  \leanok
  \uses{thm:main-growth-rate}
  $\gamma(u,\eps)\geq\gamma_0>0$ for $0<\eps\leq\eps_0$ by Theorem~\ref{thm:main-growth-rate},
  and $\liminf_{\eps\downarrow0}\gamma(u,\eps)\geq\gamma_0>0$.  If $\gamma(u,\eps)\to0$ the
  lower limit would be $0$.
\end{proof}

\section{The strengthened statement}

The proof produces closed terms for the velocity and both constants, and it
proves several statements that are strictly stronger than
Theorem~\ref{thm:main-spectral}.  They are collected in one structure.

\begin{definition}[The explicit witnesses]
  \label{def:main:witnesses}
  \uses{def:transfer:parameter-choice, def:velocity:global-velocity,
    def:local:seed}
  \lean{FastDynamo.mainVelocity,
    FastDynamo.mainEps0,
    FastDynamo.mainGrowthFloor,
    FastDynamo.mainEps0_eq,
    FastDynamo.mainEps0_pos,
    FastDynamo.mainGrowthFloor_eq,
    FastDynamo.mainGrowthFloor_pos,
    FastDynamo.mainGrowthFloor_ge_three_quarters,
    FastDynamo.mainGrowthFloor_ge_numeric,
    FastDynamo.mainVelocity_isTransverseDivergenceFree,
    FastDynamo.mainVelocity_ne_zero,
    FastDynamo.seed_kStar, FastDynamo.seed_wavenumber}
  \leanok
  Let $\Theta$ be the global parameter chain of
  Definition~\ref{def:transfer:parameter-choice} produced by the closure, let
  $u:=u_\Theta$ be its global velocity
  (Definition~\ref{def:velocity:global-velocity}), and put
  \[
    \eps_0:=\delta_+\frac{k_*^2}{K_0^2},
    \qquad
    \gamma_0:=g_*=\re(\text{centre }\Gamma)-\text{radius }\Gamma .
  \]
  Here $\delta_+$ is the seed's upper rescaled-diffusivity endpoint, $k_*=1$ the seed's
  integer wavenumber (so $N=2\pi$), $K_0$ the base axial index of the chain and $\Gamma$ the
  seed's isolating circle, so that $g_*=\min_{\zeta\in\cG}\re\zeta$.
\end{definition}

The threshold $\eps_0=\delta_+/K_0^2$ (at $k_*=1$) is the paper's value, the display
after \texttt{eq:diffusivity-intervals} in §3.2 of the paper, with no further cap;
\texttt{mainEps0\_eq} is that formula, and $\gamma_0=g_*$ is the choice made in the paper's
proof of Theorem~1.2.  The growth floor is a closed term of the fixed local seed alone: it
depends on neither $\eps$, nor the velocity, nor the sector, nor the axial shift, nor the
placement.

The numerals that follow are choices of the formalisation, not of the paper.  The paper
takes the radial scale $\varrho$ of the segment $\kappa=\varrho(1+it)$ ``sufficiently
large'', obtains $\re\lambda_*\geq\tfrac7{16}\varrho^2-4\pi^2$
(\texttt{eq:local-eigenvalue}), and fixes the isolating circle $\Gamma$ only qualitatively,
so it gives no numerical value of $\gamma_0=g_*$.  Lean takes $\varrho=32$ and $N=2\pi$,
which gives $\re\lambda_*\geq448-4\pi^2$, and a circle centred at $\lambda_*$ of radius at
most $\re\lambda_*/4$, which gives $g_*\geq\tfrac34\re\lambda_*$
(Lemmas~\ref{lem:local:segment-root} and~\ref{lem:local:isolation}); together they give
\[
  \gamma_0\ \geq\ \tfrac34\bigl(448-4\pi^2\bigr)\ \approx\ 306.4,
  \qquad\text{and in particular}\qquad
  \gamma_0\ \geq\ 288 .
\]
In Lean these are \texttt{FastDynamo.mainEps0\_eq},
\texttt{FastDynamo.mainGrowthFloor\_eq},
\texttt{FastDynamo.mainGrowthFloor\_ge\_three\_quarters} and
\texttt{FastDynamo.mainGrowthFloor\_ge\_numeric}
(\texttt{FastDynamo/Global/MainTheoremWitnesses.lean}); the numeral $288$ comes from
$\pi\leq4$, and the sharp value is $\approx306.4$.  Note that $\eps_0$ is explicit only as a
formula: $K_0$ is obtained from an existence statement about the closure thresholds and is
not computed, so no numeral for $\eps_0$ is claimed.

\begin{lemma}[Normalisation]
  \label{lem:main:normalisation}
  \uses{def:main:induction-operator, def:main:physical-clauses,
    def:main:witnesses, def:main:growth-rate, def:dynamics:particle-flow,
    def:dynamics:ideal-induction-group}
  \lean{FastDynamo.reducedInductionAction_smul,
    FastDynamo.reducedInductionOperator_isEigenpair_smul_iff,
    FastDynamo.mainReducedClauses_smul,
    FastDynamo.MainPhysicalConclusionNormalized,
    FastDynamo.mainPhysicalConclusionNormalized_of_mainReducedTarget,
    FastDynamo.mainPhysicalClauses_normalizedVelocity,
    FastDynamo.productInductionAction_smul,
    FastDynamo.productInductionAction_normalizedVelocity,
    FastDynamo.StronglyContinuousSemigroup.timeRescale,
    FastDynamo.StronglyContinuousSemigroup.growthBound_timeRescale,
    FastDynamo.literalInductionSemigroup_normalizedVelocity,
    FastDynamo.gamma_normalizedVelocity,
    FastDynamo.normalizedGrowthFloor_le_gamma,
    FastDynamo.normalizedGrowthFloor_le_liminf_gammaOf,
    FastDynamo.mainSemigroupOrbitClauses_normalizedVelocity,
    FastDynamo.IsParticleFlow.timeRescale, FastDynamo.IsIdealInductionGroup.timeRescale,
    FastDynamo.normalizedLiteralVelocity,
    FastDynamo.isIdealInductionGroup_normalizedLiteralVelocity,
    FastDynamo.normalizedLiteralFlow_unique,
    FastDynamo.MainTheoremNormalized, FastDynamo.mainTheoremNormalized_holds}
  \leanok
  Put $M=\max\{1,\norm{u}_{\W^{1,\infty}}\}$ and $\widetilde u=u/M$.  Then
  \[
    L_{\eps,\widetilde u}=M^{-1}L_{M\eps,u},
    \qquad
    \gamma(\widetilde u,\eps)=M^{-1}\gamma(u,M\eps).
  \]
  Hence Theorem~\ref{thm:main} remains valid for $\widetilde u$, with threshold
  $\widetilde\eps=\eps_0/M$ and lower bound $\widetilde\gamma=\gamma_0/M$ on the growth rate.
  Moreover, the particle flow and ideal induction group for $\widetilde u$ are the
  corresponding objects for $u$ evaluated at time $t/M$.  This preserves the ideal growth rate
  and the vanishing of topological entropy.  In particular, the velocity may be required to
  satisfy $\norm{\widetilde u}_{\W^{1,\infty}}\leq1$ without changing any of the conclusions of
  Theorem~\ref{thm:main}.
\end{lemma}

This is Remark~1.6 of the paper (unlabelled), whose wording the statement follows.  In Lean it
is the structure \texttt{MainTheoremNormalized} at $\widetilde u=$
\texttt{normalizedVelocity mainVelocity}, proved by \texttt{mainTheoremNormalized\_holds}
(\texttt{FastDynamo/MainTheoremNormalized.lean}): $\norm{\widetilde u}\leq1$; the two
identities (\texttt{action\_eq}, and \texttt{gamma\_eq} from the identity of semigroups
$\e^{tL_{\eps,\widetilde u}}=\e^{(t/M)L_{M\eps,u}}$, \texttt{semigroup\_eq}); every clause of
Theorem~\ref{thm:main} for $\widetilde u$ with $M^{-1}\eps_0$ and $M^{-1}\gamma_0$ (the spectral
clauses, the growth-rate clause and its $\liminf$ form, the semigroup-orbit and solution
clause, and the two dynamical clauses); the flow $t\mapsto\Phi_{t/M}$, which is the only
particle flow of $\widetilde u$, and the ideal solution group $t\mapsto\mathcal T_0(t/M)$; and
Proposition~\ref{prop:main-lagrangian} for $\widetilde u$.  The deviations:
\begin{itemize}
  \item The first identity is stated for the auxiliary all-frequency action $\mathcal A_\eps$
    on all of $\Hs^2$ (\texttt{productInductionAction\_normalizedVelocity}), the action of the
    generator, which on divergence-free fields is the curl form
    (Theorem~\ref{thm:dynamics:solenoidal-generator}); for the reduced operator the same
    homogeneity $A_{c\eps,cu,K}=c\,A_{\eps,u,K}$ is \texttt{reducedInductionAction\_smul}.
  \item The norm in $M$ is the norm of \texttt{TransverseRealW1Inf},
    $\norm u=\max(\norm u_\infty,\norm{\partial_1u}_\infty,\norm{\partial_2u}_\infty)$ with
    $\norm\cdot_\infty$ the essential supremum of the Euclidean norm; the paper does not
    specify the norm $\norm\cdot_{\W^{1,\infty}}$, and every identity holds for any $M>0$.
  \item Proposition~\ref{prop:main-lagrangian} for $\widetilde u$ is proved with the constant
    $(1+M^{-1})C_{\rm L}$; the paper does not restate the constants.
\end{itemize}
The normalised witness is \emph{named}, not merely some velocity in the unit ball.  Modules
\texttt{FastDynamo/Dynamics/SemigroupTimeRescaling.lean},
\texttt{FastDynamo/Dynamics/Assembly/NormalizedGrowthRate.lean},
\texttt{FastDynamo/Dynamics/NormalizedLagrangianClauses.lean} and
\texttt{FastDynamo/Operator/ReducedInductionScaling.lean}.

\begin{proof}
  \leanok
  \uses{thm:main-spectral, thm:main-growth-rate, prop:particle-flow, cor:ideal-growth,
    lem:dynamics:semigroup-unique, lem:dynamics:physical-uniqueness}
  The three velocity-dependent terms of the induction operator are $\R$-linear in $u$ and the
  diffusion term is linear in $\eps$, which gives $L_{c\eps,cu}=c\,L_{\eps,u}$.  The semigroup
  of $cu$ at $c\eps$ and the semigroup of $u$ at $\eps$ run at time $ct$ are realised by
  operators that differ by the factor $c$, so they coincide by uniqueness of the realising
  semigroup; the growth bound of $t\mapsto S(ct)$ is $c\,\omega_0(S)$.  For the clauses, given
  $0<\eps\leq M^{-1}\eps_0$ apply the unnormalised clauses at $M\eps\leq\eps_0$ and rescale
  the eigenvalue.  If $\Phi$ is the flow of $u$, then $t\mapsto\Phi_{t/M}$ solves the flow
  equation of $u/M$, and the Cauchy formula, the weak form and uniqueness rescale with it; the
  entropy of $\Phi_{t/M}$ is zero and the logarithmic rates are unchanged by a linear change of
  time.
\end{proof}

\begin{lemma}[Qualitative regularity of the velocity]
  \label{lem:main:velocity-regularity}
  \uses{def:main:witnesses, def:velocity:global-velocity,
    def:velocity:scale-family}
  \lean{FastDynamo.ScaleFamily.TorusCellPlacement.FullVelocityConcrete.GlobalVelocityQualitative,
    FastDynamo.mainVelocityChartRepresentative_qualitative,
    FastDynamo.differentiable_mainVelocityPeriodizedRepresentative,
    FastDynamo.contDiffAt_mainVelocityPeriodizedRepresentative,
    FastDynamo.fderiv_mainVelocityPeriodizedRepresentative_accumulationPoint,
    FastDynamo.not_continuousAt_fderiv_mainVelocityPeriodizedRepresentative,
    FastDynamo.lipschitzWith_mainVelocityPeriodizedRepresentative,
    FastDynamo.lipschitzWith_mainVelocityTorusRepresentative,
    FastDynamo.mainVelocityTorusRepresentative_transverseTorusPoint_eq_periodized}
  \leanok
  Let $u$ be the velocity of Definition~\ref{def:main:witnesses} and let
  $U\colon\R^2\to\R^3$ be its $\Z^2$-periodic lift, so that $u$ is $U$ read on
  $\T^2$ and, being independent of $x_3$, on $\T^3$.  Write
  $p_\infty=(0,\tfrac12)$ for the accumulation point.  Then:
  \begin{itemize}
    \item $U$ is Lipschitz, and so are $u$ on $\T^2$ for the quotient metric
      and its axially independent extension to $\T^3$;
    \item $U$ is differentiable at every point of $\R^2$;
    \item $U$ is smooth of every order at every point outside
      $p_\infty+\Z^2$;
    \item $DU(p_\infty)=0$, yet $DU$ is discontinuous at $p_\infty$.
  \end{itemize}
  Thus $u$ is differentiable everywhere and smooth away from the accumulation
  circle $\{p_\infty\}\times\T$, its derivative vanishes on that circle, and
  $u\notin\C^1$.
\end{lemma}

\begin{proof}
  \leanok
  The local profile $v$ is smooth, and regularity is lost only because the
  cell sum has non-vanishing gradients at infinitely many shrinking horizontal
  scales.  The rescaled cells have amplitude $O(\ell_n)$ and gradients $O(1)$,
  so the disjoint sum is Lipschitz and its derivative tends to $0$ at
  $p_\infty$, which gives differentiability with $Du(p_\infty,\cdot)=0$.
  Both cut-offs are constant near each cell centre, so there the amplitude
  $\ell_n$ of the local copy cancels its dilation exactly and
  \[
    Du(p_n,x_3)=
    \begin{pmatrix}
      0      & -\Omega & 0 \\
      \Omega & 0       & 0 \\
      0      & 0       & 0
    \end{pmatrix}
    \neq0
    \qquad(x_3\in\T),
  \]
  one and the same non-zero matrix at every $n$.  Since $p_n\to p_\infty$ while
  $Du(p_\infty,\cdot)=0$, the derivative $Du$ is discontinuous at every point of the
  accumulation circle; a continuous representative of the weak gradient would have to agree
  with $Du$ everywhere, so none exists and $u\notin\C^1$.
\end{proof}

The lift $U$ is the Lean declaration
\texttt{mainVelocityPeriodizedRepresentative}, the $\Z^2$-periodisation of the
chart-side representative; the theorem
\texttt{mainVelocityTorusRepresentative\_transverseTorusPoint\_eq\_periodized}
identifies it with the pullback of the torus-side representative along the
quotient map \emph{at every point of the plane}, including the accumulation
point, which the placement chart itself misses.  The four bullets are Lean
theorems verbatim, the accumulation circle appearing as the plane point
\texttt{accumulationPoint} $=(0,\tfrac12)$ together with its integer
translates.  The proof is the paper's (§3.4 of the paper), along the cell centres:
\texttt{fderiv\_\allowbreak parameterizedFullGlobalEuclideanRealSource\_\allowbreak cellCenter\_\allowbreak eq} gives one and the
same derivative at every centre,
\texttt{lineDeriv\_parameterizedSmoothedVelocity\_zero\_apply\_one} is the entry
$(\partial_{e_1}v)(0)(e_2)=\Omega$ that makes it non-zero, and
\texttt{not\_\allowbreak continuousAt\_\allowbreak fderiv\_\allowbreak parameterizedFullGlobalEuclideanRealSource} concludes along
\texttt{tendsto\_cellCenter\_accumulationPoint}; for the lift $U$ this is
\texttt{not\_\allowbreak continuousAt\_\allowbreak fderiv\_\allowbreak mainVelocityPeriodizedRepresentative}.  The matrix itself is not assembled in
Lean; only the one entry is, which is all the non-vanishing needs.  Two caveats.  First, the final sentence of the lemma -- the
passage from the lift to a statement about a function on $\T^3$ -- is
\emph{not} itself formalised for differentiability, smoothness or the failure
of $\C^1$: the Lean development contains no \texttt{ContDiff} statement about a
map defined on $\T^3$, only the standard chart form above.  Lipschitz
continuity is the exception: it is proved torus-side as well, by
\texttt{lipschitzWith\_mainVelocityTorusRepresentative} for the quotient metric
of $\T^2$ and by \texttt{lipschitzWith\_mainLiteralVelocity} on the literal
three-torus, and the torus-side representative is almost everywhere the
$\Lp^\infty$ value of $u$.  Second, the Lipschitz constant is explicit but not
claimed optimal: it is twice the chart-side bound
\texttt{mainVelocityLipschitzBound}.  The modules are
\texttt{FastDynamo/Global/TorusGlobalVelocityRegularity.lean} and
\texttt{FastDynamo/Global/TorusGlobalVelocityLipschitz.lean}.

\begin{lemma}[The literal three-torus]
  \label{lem:main:literal-torus}
  \uses{def:main:witnesses, def:main:physical-clauses, def:finish:real-solution}
  \lean{FastDynamo.toLiteralR3L2, FastDynamo.literalRealTrajectory,
    FastDynamo.literalVelocity, FastDynamo.mainLiteralVelocity,
    FastDynamo.lipschitzWith_mainLiteralVelocity}
  \leanok
  The magnetic clauses of Theorem~\ref{thm:main-spectral} and the Lipschitz
  bound on the velocity hold on the literal three-torus
  $\T^3=(\R/\Z)^3$: for every $0<\eps\leq\eps_0$ the transported real
  trajectory $B^\eps$ lies in $\Lp^2(\T^3;\R^3)$, never vanishes, obeys
  $\norm{B^\eps(t)}_2=\e^{t\re\lambda_\eps}\norm{B^\eps(0)}_2$ for every
  $t\in\R$ and $\norm{B^\eps(t)}_2\geq\e^{t\gamma_0}\norm{B^\eps(0)}_2$ for
  $t\geq0$, and $u$ is an axially independent Lipschitz function on $\T^3$.
\end{lemma}

\begin{proof}
  \leanok
  The development works in the product model $\T^2\times\T$ with the product
  of the transverse volume and the Haar measure of the circle.  That model is
  identified with the literal $\T^3=(\R/\Z)^3$, realised in Lean as
  \texttt{Torus 3}, carrying its own volume, by a
  measure-preserving equivalence, and the induced map on $\Lp^2$ is a linear
  isometry.  Transporting the real trajectory along it preserves every norm,
  so the exact law, the non-vanishing and the growth bound transfer verbatim;
  the axially independent extension of a Lipschitz transverse field is
  Lipschitz for the sup metric of $\T^3$.
\end{proof}

\begin{lemma}[The real solution on the literal three-torus]
  \label{lem:main:literal-torus-solution}
  \uses{lem:main:literal-torus, def:main:growth-rate, lem:dynamics:physical-uniqueness}
  \lean{FastDynamo.complexifyLiteralR3L2, FastDynamo.toLiteralC3L2_complexifyProductR3L2,
    FastDynamo.MainLiteralTorusClauses,
    FastDynamo.mainLiteralTorusClauses_holds}
  \leanok
  For every $0<\eps\leq\eps_0$ there are $K\neq0$, $\lambda$ with $\re\lambda\geq\gamma_0$ and
  $b\neq0$ such that the transported real trajectory $B^\eps$ of
  Lemma~\ref{lem:main:literal-torus} obeys the norm law, the non-vanishing and the growth
  bound stated there, and there is a datum $B_{\rm in}^\eps\in\Lp^2_\sigma(\T^3;\CC^3)$, the
  complexification of $B^\eps(0)$, whose orbit under the induction semigroup
  $\e^{tL_{\eps,u}}$ on $\Lp^2_\sigma(\T^3;\CC^3)$ (the semigroup of
  Definition~\ref{def:main:growth-rate}) is the complexification of $B^\eps(t)$ for every
  $t\geq0$.
\end{lemma}

\begin{proof}
  \leanok
  The triple and the product-model datum are those of the orbit clause
  \texttt{MainSemigroupOrbitClauses} at the witnesses
  (\texttt{mainSemigroupOrbitClauses\_holds}, from
  \texttt{FastDynamo/Dynamics/Assembly/SemigroupOrbit.lean}), where the orbit of $B_{\rm in}^\eps$ under the semigroup on $\Lp^2_\sigma(\T^2\times\T)$ is the
  complexified real curve.  The semigroup on $\Lp^2_\sigma(\T^3;\CC^3)$ is that semigroup
  conjugated by the isometric transport, and the transport commutes with complexification
  (\texttt{toLiteralC3L2\_complexifyProductR3L2}: both are composition with a pointwise map,
  up to null sets), so the orbit identity transfers.
\end{proof}

The literal-torus statements have two limitations.  No differential operator is transported, so on the literal
$\T^3$ divergence-freeness is membership of the datum in the transported subspace
$\Lp^2_\sigma$, and the uniqueness of the solution with datum $B_{\rm in}^\eps$ is stated in
the product model only (Lemma~\ref{lem:dynamics:physical-uniqueness}); and the norm statements
are stated against the trajectory's own value at time zero rather than against a separately
transported initial field, so that every norm in them is a norm of
$\Lp^2(\T^3;\R^3)$
(\texttt{FastDynamo/Analysis/LiteralThreeTorus.lean}).

\begin{theorem}[The strengthened main theorem]
  \label{thm:main-strengthened}
  \uses{thm:main-spectral, thm:main-growth-rate, def:main:witnesses,
    lem:main:normalisation, lem:main:velocity-regularity, lem:main:literal-torus,
    lem:main:literal-torus-solution,
    prop:global-spectral-transfer, lem:transfer:all-cells}
  \lean{FastDynamo.MainTheoremStrengthened,
    FastDynamo.main_theorem_strengthened,
    FastDynamo.MainTheoremStrengthened.growthRate,
    FastDynamo.MainTheoremStrengthened.growthRate_le,
    FastDynamo.MainEveryAdmissibleCellClauses,
    FastDynamo.mainEveryAdmissibleCellClauses_holds,
    FastDynamo.MainSemigroupOrbitClauses,
    FastDynamo.mainSemigroupOrbitClauses_holds}
  \leanok
  At the explicit witnesses $u$, $\eps_0$, $\gamma_0$ of
  Definition~\ref{def:main:witnesses} the following hold.
  \begin{enumerate}
    \item \emph{Explicit constants.}  $u$ is divergence-free and non-zero,
      $\eps_0=\delta_+k_*^2/K_0^2$ with $k_*=1$, and $\gamma_0>0$ satisfies both
      $\gamma_0\geq\tfrac34(448-4\pi^2)$, with $N=2\pi$, and the numeral bound
      $\gamma_0\geq288$.
    \item \emph{The sixteen physical clauses} $\mathrm{MPC}(u,\eps_0,\gamma_0)$
      of Definition~\ref{def:main:physical-clauses}.
    \item \emph{The growth-rate clause} \eqref{eq:main-growth-conclusion}, in
      the form of Theorem~\ref{thm:main-growth-rate}: for every
      $0<\eps\leq\eps_0$ the growth rate of the induction semigroup on
      $\Lp^2_\sigma(\T^3;\CC^3)$ satisfies $\gamma(u,\eps)\geq\gamma_0$, and at
      every positive diffusivity $\gamma(u,\eps)\leq\norm{\nabla_\perp u}_\infty$.
    \item \emph{Every admissible cell.}
      Proposition~\ref{prop:global-spectral-transfer} holds not only at the
      canonically chosen matching cell but at \emph{every} one: for every
      $0<\eps\leq\eps_0$ and every cell index $n$ with
      $\eps/\ell_n^2\in I$, the axial index $K_n$ is non-zero and carries an
      eigenpair, of the reduced and of the auxiliary operator, with
      $\re\lambda\geq\gamma_0$; and the Riesz-projection form of the
      proposition holds at every such cell as well.
    \item \emph{Normalisation} of the spectral clauses: the sixteen physical clauses hold
      for the named witness $\widetilde u=u/\max\{1,\norm u\}$, with
      $\norm{\widetilde u}\leq1$ and thresholds $M^{-1}\eps_0$, $M^{-1}\gamma_0$ (the whole of
      Remark~1.6 of the paper is Lemma~\ref{lem:main:normalisation}).
    \item \emph{Velocity regularity}, in the form of
      Lemma~\ref{lem:main:velocity-regularity}, together with the
      identification of the torus-side representative with the
      $\Z^2$-periodisation of the chart representative at every point of the
      plane.
    \item \emph{The literal three-torus}, in the form of
      Lemmas~\ref{lem:main:literal-torus} and~\ref{lem:main:literal-torus-solution}.
    \item \emph{The semigroup orbit of the real datum.}  For every
      $0<\eps\leq\eps_0$ the datum $B_{\rm in}^\eps$ lies in $\Lp^2_\sigma$, its orbit
      under the semigroup of Theorem~\ref{thm:dynamics:solenoidal-semigroup} is the
      explicit curve $B^\eps$ at every $t\geq0$, and that orbit obeys
      \eqref{eq:main-thm-growth}; every forward strong solution of the curl-form equation
      with datum $B_{\rm in}^\eps$ is that orbit, obeys \eqref{eq:main-thm-growth} and is
      bounded below by $\e^{t\gamma_0}\norm{B_{\rm in}^\eps}_2$; and at the same triple
      $(K,\lambda,b)$ the sixteen physical clauses of
      Definition~\ref{def:main:physical-clauses} hold.  This is the definite
      article of ``\emph{its} solution'' in Theorem~\ref{thm:main}.
  \end{enumerate}
\end{theorem}

\begin{proof}
  \leanok
  \uses{def:transfer:parameter-choice}
  Each item is the corresponding statement above, read at the named witnesses
  instead of at anonymous existentials.  Item (4) is the same closure
  certificate as in Theorem~\ref{thm:main-spectral}, except that the four
  defining properties of the parameter chain
  (Definition~\ref{def:transfer:parameter-choice}) are recorded one by one, so
  the certificate is available at every admissible diffusivity and every
  matching cell rather than only at the canonically chosen one.
\end{proof}

The structure is \texttt{FastDynamo.MainTheoremStrengthened} and the proof is
\texttt{FastDynamo.main\_theorem\_strengthened}
(\texttt{FastDynamo/MainTheoremStrengthened.lean}); the witnesses are named in
\texttt{FastDynamo/Global/MainTheoremWitnesses.lean}; the growth-rate item is
the pair of fields \texttt{growthRate} and \texttt{growthRate\_le}, and the semigroup-orbit
item is the field \texttt{semigroupOrbit}.  The
growth-rate clause, which Theorem~\ref{thm:main-spectral} does not contain, is included here,
and so is the uniqueness of the solution with datum $B_{\rm in}^\eps$ (item 8).  The dynamical
statements are Proposition~\ref{prop:main-lagrangian}; they are in a separate structure,
\texttt{MainTheoremComplete} (Theorem~\ref{thm:main-complete}), because the module that names
the velocity cannot import the flow without a cycle.

The items are separate fields, and the items stated for every $0<\eps\leq\eps_0$ choose their
triples $(K,\lambda,b)$ separately: items (2), (5), (7) and (8) each carry their own
existential over $(K,\lambda,b)$, item (5) for the normalised velocity.  Item (8) contains the sixteen physical clauses at its own triple, so the
eigenmode, the real datum, the strong solution, the semigroup orbit and the uniqueness clause
hold at one $(K,\lambda,b)$.  The triple of item (7) is the one of item (8) by the proof of
\texttt{mainLiteralTorusClauses\_holds}, not by its statement; item (2) is the existential
\texttt{MainPhysicalClauses} with a separately chosen triple; and item (4) is a statement about
every admissible cell, not about a chosen triple.


\chapter{The local seed: a helical velocity with unstable horizontal spectrum}
% Chapter: the local seed.  Section 2 of the paper (sec:local-model).
% Owned shared labels: def:local:seed.
% Paper labels: lem:jump-horizontal-mode (Lemma 2.1), prop:local-seed (Proposition 2.2).

Section~2 of the paper builds the seed on \emph{two}-component horizontal fields.  At unit
diffusivity the horizontal equations in the axial Fourier mode $\e^{2\pi ix_3}$ close on
$b\colon\R^2\to\CC^2$, and the weak form of the resulting generator
\eqref{eq:local:horizontal-induction} carries no derivative of the velocity once the
stretching term is integrated by parts.  This is what allows the seed to be built for a
\emph{discontinuous} velocity, the Ponomarenko screw flow, whose unstable eigenvalue is
produced by an intermediate-value argument on an explicit segment of the Bessel dispersion
relation.  The eigenvalue is isolated by the analytic Fredholm theorem, the velocity is
smoothed and the diffusivity is varied in \emph{one} form estimate
\eqref{eq:local:form-small}, and a fixed compactly supported test vector $f_*$ detects the
spectral subspace throughout.  The output is the record of Definition~\ref{def:local:seed},
inhabited at $k_*=1$ by Proposition~\ref{prop:local-seed}; the axial component of the
magnetic field is not part of the seed at all and is recovered only in the finish chapter
(Lemma~\ref{lem:axial-recovery}).

Throughout this chapter $\Lp^2$ inner products are linear in the first argument, as in the
paper; the Lean convention is the opposite one, and every statement below is the
transcription $\inner f\phi=\texttt{inner}\ \CC\ (j\phi)\ f$ of the paper's pairing
(the module docstring of \texttt{FastDynamo/Local/HorizontalWeakForm.lean} records it).

\section{The horizontal operator on $\R^2$ and its weak form}

\begin{definition}[The horizontal generator and its sesquilinear form]
  \label{def:local:horizontal-operator}
  \lean{FastDynamo.HorizontalWeakForm.horizontalDiffusionForm,
    FastDynamo.HorizontalWeakForm.horizontalTransportForm,
    FastDynamo.HorizontalWeakForm.horizontalAxialForm,
    FastDynamo.HorizontalWeakForm.horizontalStretchDivergenceForm,
    FastDynamo.HorizontalWeakForm.horizontalStretchGradientForm,
    FastDynamo.HorizontalWeakForm.horizontalWeakForm,
    FastDynamo.HorizontalWeakForm.HorizontalFormDomain,
    FastDynamo.HorizontalWeakForm.horizontalFormOperator,
    FastDynamo.HorizontalWeakForm.horizontalFormOperator_exists_iff,
    FastDynamo.HorizontalWeakForm.horizontalFormOperator_isEigenpair_iff}
  \leanok
  Fix $N=2\pi$ (the paper writes $4\pi^2$ and $2\pi i$ for $N^2$ and $iN$).  For a velocity $w=(w_\perp,w_3)$ independent of $x_3$ with
  $\Div_\perp w_\perp=0$ the horizontal generator at unit diffusivity is
  \begin{equation}
    A_wb=(\Delta_\perp-N^2)b-w_\perp\cdot\nabla_\perp b-iNw_3b
      +(b\cdot\nabla_\perp)w_\perp,
    \qquad b\colon\R^2\to\CC^2 .
    \label{eq:local:horizontal-induction}
  \end{equation}
  Integrating the stretching term by parts removes every derivative of $w$, so for
  $w\in\Lp^\infty(\R^2;\R^3)$ the sesquilinear form
  \begin{align}
    \mathfrak l_w(b,\phi)
    ={}&\inner{\nabla_\perp b}{\nabla_\perp\phi}+N^2\inner b\phi
       +\inner{w_\perp\cdot\nabla_\perp b}\phi+iN\inner{w_3b}\phi\notag\\
      &+\int_{\R^2}w_\perp\cdot
        \bigl((\Div_\perp b)\overline\phi+(b\cdot\nabla_\perp)\overline\phi\bigr)\dd y
    \label{eq:local:sectorial-form}
  \end{align}
  is defined for $b,\phi\in\Hs^1(\R^2;\CC^2)$.  We write $-A_w$ for the operator on
  $\Lp^2(\R^2;\CC^2)$ associated with $\mathfrak l_w$ in Kato's sense: $b\in D(A_w)$ with
  $A_wb=f$ exactly when $b\in\Hs^1$ and $\mathfrak l_w(b,\phi)=-\inner f\phi$ for every
  $\phi\in\Hs^1$.
\end{definition}

In Lean the five summands of \eqref{eq:local:sectorial-form} are the five declarations
listed above, all on \texttt{WholePlaneC2H1} $=\Hs^1(\R^2;\CC^2)$ with a velocity of type
\texttt{WholePlaneC3LInf}; \texttt{horizontalFormOperator} is the generic
\texttt{associatedFormProdOperator} of \texttt{FastDynamo/Operator/AssociatedForm.lean}
applied to $\mathfrak l_w$ with the pivot $\Hs^1\hookrightarrow\Lp^2$, so the domain is a
graph subspace of $\Hs^1\times\Lp^2$ rather than a subset of $\Lp^2$.  The two
characterisations \texttt{horizontalFormOperator\_exists\_iff} and
\texttt{\_isEigenpair\_iff} are the paper's domain description and the eigenvalue
equation $\mathfrak l_w(b,\phi)=-\lambda\inner b\phi$ verbatim.  The module is
\texttt{FastDynamo/Local/HorizontalWeakForm.lean} and
\texttt{FastDynamo/Local/HorizontalFormOperator.lean}.

\begin{lemma}[Boundedness, Gårding inequality and sectoriality]
  \label{lem:local:garding}
  \lean{FastDynamo.HorizontalWeakForm.norm_horizontalWeakForm_apply_le,
    FastDynamo.HorizontalWeakForm.horizontalWeakForm_garding,
    FastDynamo.HorizontalWeakForm.horizontalWeakForm_abs_im_le,
    FastDynamo.HorizontalWeakForm.horizontalGardingConstant,
    FastDynamo.HorizontalWeakForm.horizontalShiftedWeakForm,
    FastDynamo.HorizontalWeakForm.horizontalShiftedWeakForm_hasComplexCoercivityConstant,
    FastDynamo.HorizontalWeakForm.horizontalShiftedWeakForm_sectorial}
  \leanok
  \uses{def:local:horizontal-operator}
  For $w\in\Lp^\infty(\R^2;\R^3)$ the form $\mathfrak l_w$ is bounded on
  $\Hs^1\times\Hs^1$, and
  \[
    \re\mathfrak l_w(b,b)\geq\tfrac12\norm{\nabla_\perp b}_2^2-C\norm b_2^2,
    \qquad
    \abs{\im\mathfrak l_w(b,b)}\leq C'\bigl(\norm{\nabla_\perp b}_2^2+\norm b_2^2\bigr),
  \]
  with $C,C'$ depending only on $N$ and $\norm w_\infty$.  Consequently, for
  $\alpha=C+\tfrac12$, the shifted form $\mathfrak l_w+\alpha\inner\cdot\cdot$ is coercive on
  $\Hs^1$ with constant $\tfrac12$, and its imaginary part on the diagonal is at most $2C'$
  times its real part; so $\mathfrak l_w$ is closed and sectorial.
\end{lemma}

The Lean constants are explicit: $C=50\norm w_\infty^2+\abs N\norm w_\infty$
(\texttt{horizontalGardingConstant}), $C'=(5+\abs N)\norm w_\infty$, and the sectoriality
slope is $2C'$; all are generous, and no later step uses their value.  The first inequality
is the paper's \texttt{eq:sectorial-estimate-b}.  \textbf{Deviation (shape of the
imaginary-part bound).}  The paper's second display is
$\abs{\im\mathfrak l_w(b,b)}\leq\re\mathfrak l_w(b,b)+C\norm b_2^2$ for every $b\in\Hs^1$,
after increasing $C$ if necessary: a sector of slope one about the shifted real part.  Lean
bounds the imaginary part by the full $\Hs^1$ norm instead and obtains, after the shift by
$\alpha$, a sector of slope $2C'$, which depends on $N$ and $\norm w_\infty$.  Both say that
$\mathfrak l_w$ is sectorial, which is all the paper uses (Kato's representation theorem);
the slope-one form is not formalised.  The coercivity after the shift is stated in the
$\Hs^1$ graph norm rather than in the paper's split form, which is the same statement because
$\norm b_{\Hs^1}^2=\norm b_2^2+\norm{\nabla_\perp b}_2^2$ by construction of the graph
space.

\begin{proof}
  \leanok
  \uses{def:local:horizontal-operator}
  The three velocity terms carrying a derivative are bounded by
  $5\norm w_\infty\norm b_2\bigl(\norm{\partial_1b}_2+\norm{\partial_2b}_2\bigr)$ and the
  axial term by $\abs N\norm w_\infty\norm b_2^2$; Young's inequality
  $5Mm(G_0+G_1)\leq\tfrac14(G_0^2+G_1^2)+50M^2m^2$ absorbs the first into the diffusion
  term and leaves $\tfrac34$ of it, which is more than the claimed $\tfrac12$.  The
  imaginary part is estimated by the same three bounds without absorption, and
  $2\norm{\partial_jb}_2\norm b_2\leq\norm{\partial_jb}_2^2+\norm b_2^2$.  Adding a large
  real multiple of the $\Lp^2$ inner product makes the real part dominate, which is
  coercivity, and the quotient of imaginary by real part is then bounded, which is
  sectoriality.
\end{proof}

\section{An unstable eigenvalue for the discontinuous screw flow}

\begin{definition}[The Ponomarenko screw flow]
  \label{def:local:jump-flow}
  \lean{FastDynamo.HorizontalMode.horizontalJumpInteriorVelocity,
    FastDynamo.HorizontalMode.horizontalJumpVelocity,
    FastDynamo.HorizontalMode.horizontalJumpVelocityLInf,
    FastDynamo.HorizontalMode.horizontalModeVelocity,
    FastDynamo.PonomarenkoFreeOmegaInterface.jumpAxialSpeed,
    FastDynamo.PonomarenkoFreeOmegaInterface.scale_mul_jumpAxialSpeed}
  \leanok
  \uses{def:local:horizontal-operator}
  For $\Omega>0$ put $U=-\Omega/N$ and, in polar coordinates,
  \begin{equation}
    V^0(y)=\one_{\{r<1\}}\bigl(\Omega re_\theta+Ue_3\bigr).
    \label{eq:local:jump-flow}
  \end{equation}
  The normal component vanishes on both sides of $r=1$, so
  $\Div_\perp V^0_\perp=0$ in the sense of distributions and $V^0\in\Lp^\infty$ has
  compact support.
\end{definition}

In Lean $\Omega$ is a \emph{free} real parameter: \texttt{horizontalJumpVelocity Omega U}
is the Cartesian field $(-\Omega y_2,\Omega y_1,U)$ cut off by the indicator of the closed
unit disc, and \texttt{horizontalModeVelocity N Omega} pins $U=-\Omega/N$ through
\texttt{jumpAxialSpeed}.  The paper's open disc $\{r<1\}$ and the Lean closed disc
differ on a null set, hence define the same $\Lp^\infty$ class; this is recorded in the
module docstring.  Making $\Omega$ free is the change that lets the dispersion relation be
read as a \emph{condition on $\kappa$ alone}, which is what the intermediate-value
argument needs.

\begin{definition}[The dispersion relation and its certificate]
  \label{def:local:dispersion}
  \lean{FastDynamo.ponomarenkoDispersion_analyticOnNhd,
    FastDynamo.ModifiedBessel.ponomarenkoDispersion_eq_neg_two_div_mul_deriv_I1_mul_K1,
    FastDynamo.PonomarenkoFreeOmegaInterface.TransverseInterfaceLaws,
    FastDynamo.PonomarenkoFreeOmegaInterface.CertifiedInterfaceData,
    FastDynamo.PonomarenkoFreeOmegaInterface.transverseInterfaceLaws_of_dispersion,
    FastDynamo.ponomarenkoLocalEigenvalue}
  \leanok
  \uses{def:local:jump-flow}
  Separating the angular index $j=1$ and writing $B_\pm=B_r\pm iB_\theta$, the bounded
  solutions of the radial system are the modified Bessel branches $I_2,I_0$ inside and
  $K_2,K_0$ outside, at the transverse wave number $\kappa$ with
  $\lambda_*=\kappa^2-N^2$.  The four transmission conditions
  $\jump{B_r}=\jump{B_\theta}=\jump{B_r'}=0$ and $\jump{B_\theta'}=\Omega B_r(1)$ have a
  non-trivial solution exactly when
  \begin{equation}
    \mathcal D(\kappa):=I_0(\kappa)K_0(\kappa)-I_2(\kappa)K_2(\kappa)
    =-\frac{2i}{\Omega}.
    \label{eq:local:dispersion}
  \end{equation}
  A \emph{certificate} is a triple $(N,\Omega,\kappa)$ with $N>0$, $\re\kappa>0$,
  $\Omega>0$, \eqref{eq:local:dispersion}, the four transmission laws,
  $\re(\kappa^2-N^2)>0$, and the non-vanishing of $I_\nu(\kappa)$ and $K_\nu(\kappa)$ for
  $\nu=0,1,2$ and of the interior radial component at $r=1$.
\end{definition}

\texttt{CertifiedInterfaceData N Omega kappa} is exactly that record; the transmission laws
are a separate record \texttt{TransverseInterfaceLaws}, which
\texttt{transverseInterfaceLaws\_of\_dispersion} derives from
\eqref{eq:local:dispersion} and $\re\kappa>0$, so the record is carried only for
convenience.  The Lean branches are the \emph{cross-normalised} pair
$B_+=K_2(\kappa)I_2(\kappa r)$, $B_-=-K_0(\kappa)I_0(\kappa r)$ inside (and the mirrored
pair outside), a non-zero constant multiple of the paper's solution
(\texttt{eq:bessel-solutions} with $S=C_++C_-=1$) written without denominators; with that
normalisation the two Wronskian jumps
$\jump{B_+'}=-1$, $\jump{B_-'}=+1$ hold identically and
\eqref{eq:local:dispersion} is equivalent to $\Omega B_r(1)=i$.

\begin{lemma}[An admissible root on the segment $\kappa=\varrho(1+it)$]
  \label{lem:local:segment-root}
  \lean{FastDynamo.norm_ponomarenkoSegmentDispersion_sub_inv_cube_le,
    FastDynamo.exists_ponomarenkoSegmentDispersionRe_eq_zero,
    FastDynamo.modifiedBessel_denominators_ne_zero_on_segment,
    FastDynamo.exists_certifiedInterfaceData_of_segment,
    FastDynamo.exists_certifiedInterfaceData_two_pi}
  \leanok
  \uses{def:local:dispersion}
  For $\varrho\geq32$ and $\abs t\leq3/4$,
  \[
    \abs{\varrho^3\mathcal D(\varrho(1+it))-(1+it)^{-3}}\leq\frac{225}{4\varrho^2}
    \leq\frac1{16},
  \]
  and $\re(1+it)^{-3}=(1-3t^2)(1+t^2)^{-3}$ changes sign between $t=\tfrac12$ and
  $t=\tfrac34$ while $\im(1+it)^{-3}\leq-\tfrac13$ there.  Hence there is
  $t_\varrho\in(\tfrac12,\tfrac34)$ with $\re\mathcal D(\varrho(1+it_\varrho))=0$; setting
  $\kappa=\varrho(1+it_\varrho)$ and $\Omega=-2/\im\mathcal D(\kappa)>0$ produces a
  certificate, with
  \[
    \re\lambda_*=\varrho^2(1-t_\varrho^2)-N^2\geq\tfrac7{16}\varrho^2-N^2 .
  \]
  In particular at $N=2\pi$ and $\varrho=32$ a certificate exists with
  $\re\lambda_*\geq448-4\pi^2>0$.
\end{lemma}

\textbf{Deviation (explicit numerals).}  Every threshold here is explicit and is proved in
Lean at the stated numerals; the numerals are choices of the formalisation, not of the
paper.  The paper restricts to $\kappa=\varrho(1+it)$ with $t\in[\tfrac12,\tfrac34]$ and
$\varrho>0$ ``sufficiently large'', uses \texttt{eq:bessel-product-expansion} with an
unquantified $O(\varrho^{-5})$ remainder, and concludes
$\re\lambda_*=\varrho^2(1-t_\varrho^2)-4\pi^2\geq\tfrac7{16}\varrho^2-4\pi^2>0$ after
``increasing $\varrho$ if necessary'' (\texttt{eq:local-eigenvalue}); it fixes no value of
$\varrho$ and no numerical lower bound for $\re\lambda_*$.  Lean proves the expansion with
the explicit error $225/(4\varrho^2)$ for $\varrho\geq32$ ($\varrho=32$ is the smallest power
of two above the binding constraint $\varrho\geq30$), carries the hypothesis
$N^2<\tfrac7{16}\varrho^2$ in place of ``increasing $\varrho$'', and discharges it at
$N=2\pi$, $\varrho=32$ by $(2\pi)^2<448$; hence the numeral $448-4\pi^2$.  The
existence theorem also exposes $\re\kappa=\varrho$, $\im\kappa\in(\varrho/2,3\varrho/4)$
and $\norm\kappa\leq\tfrac54\varrho$, which the later sector estimates consume.

\begin{proof}
  \leanok
  \uses{def:local:dispersion}
  The product $zI_1(z)K_1(z)$ obeys $\abs{zI_1K_1-\tfrac12}\leq3(\re z)^{-2}$ in the
  sector $\norm z\leq2\re z$, by the Gaussian representations of $I_1$ and $K_1$; since
  $(I_1K_1)'(z)=-z\mathcal D(z)/2$, so that $\mathcal D$ is $-2/z$ times the derivative of
  that product, one gets the expansion above on a disc of radius $1/5$
  around each segment point by Cauchy's estimate, and the stated constants follow by
  summing $2\delta+2\eta$.  The sign pattern of $(1+it)^{-3}$ at $t=\tfrac12,\tfrac34$
  survives the error $1/16$, so the intermediate value theorem gives $t_\varrho$.  That the
  four Bessel values $I_0,I_2,K_0,K_2$ do not vanish at $\kappa$ follows from the four
  collapse identities for $z^3\mathcal D(z)$ together with
  $\norm{\kappa^3\mathcal D(\kappa)}\leq2$ and $\norm{\kappa I_1K_1-\tfrac12}<\tfrac12$,
  which force $\norm\kappa\leq3$, contradicting $\re\kappa=\varrho\geq32$.  Finally
  $\im\mathcal D(\kappa)<0$ on the segment makes $\Omega$ positive, and
  $t_\varrho^2\leq9/16$ gives the growth floor.
\end{proof}

\begin{definition}[The Cartesian eigenmode]
  \label{def:local:cartesian-mode}
  \lean{FastDynamo.HorizontalMode.horizontalMode,
    FastDynamo.HorizontalMode.horizontalEigenvalue,
    FastDynamo.HorizontalMode.horizontalMode_ne_zero,
    FastDynamo.HorizontalMode.horizontalMode_tail_le,
    FastDynamo.HorizontalMode.HorizontalModeWeakEquation,
    FastDynamo.HorizontalMode.horizontalModeWeakEquation_iff_distributional}
  \leanok
  \uses{def:local:dispersion, def:local:horizontal-operator}
  Given a certificate, the field $b$ with
  \[
    b_1+ib_2=\e^{2i\theta}B_+(r),\qquad b_1-ib_2=B_-(r)
  \]
  is smooth on each side of $r=1$, continuous across it, decays exponentially, and lies in
  $\Hs^1(\R^2;\CC^2)$.  Its eigenvalue is $\lambda_*=\kappa^2-N^2$.
\end{definition}

The Lean mode is built branch by branch: the origin regularity comes from the
$r^2$-analyticity of $I_0$ and $I_2$, the exterior decay from the Gaussian representation
of $K_\nu$, and the $\Hs^1$ membership from a slice fundamental theorem of calculus.  The
weak eigenvalue equation is stated in two equivalent forms, against all
$\phi\in\Hs^1$ and as an identity of tempered distributions; the equivalence rests on the
injectivity of $\Hs^{-1}(\R^2;\CC^2)\to\cS'(\R^2;\CC^2)$ and is a theorem, not a
definitional move.

\begin{lemma}[An unstable eigenvalue of the discontinuous flow]
  \label{lem:jump-horizontal-mode}
  \lean{FastDynamo.HorizontalMode.horizontalModeWeakIdentity_holds,
    FastDynamo.HorizontalMode.horizontalMode_isEigenpair_of_certified,
    FastDynamo.HorizontalMode.exists_horizontalMode_eigenpair_two_pi}
  \leanok
  \uses{def:local:jump-flow, def:local:cartesian-mode, lem:local:segment-root,
    def:local:horizontal-operator}
  There exists $\Omega>0$ such that, for the velocity $V^0$ of
  \eqref{eq:local:jump-flow}, the operator $A_{V^0}$ has an eigenvalue $\lambda_*$ with
  $\re\lambda_*>0$.  An associated non-zero eigenfunction lies in the angular mode
  $b=\e^{i\theta}(B_re_r+B_\theta e_\theta)$, is smooth on each side of $r=1$ and decays
  exponentially at infinity.
\end{lemma}

\textbf{Deviation: cross-normalised branches instead of the pair $(C_+,C_-)$.}  The
paper writes the interior and exterior solutions as
$C_\pm I_{\nu_\pm}(\kappa r)/I_{\nu_\pm}(\kappa)$ and
$C_\pm K_{\nu_\pm}(\kappa r)/K_{\nu_\pm}(\kappa)$ (\texttt{eq:bessel-solutions}), and
therefore needs $I_\nu(\kappa)K_\nu(\kappa)\neq0$ for its denominators; it obtains this
from the expansion \texttt{eq:bessel-product-expansion},
$I_\nu(\kappa)K_\nu(\kappa)=\tfrac1{2\kappa}-\tfrac{4\nu^2-1}{16\kappa^3}+O(\varrho^{-5})$
for $\nu\in\{0,2\}$, for $\varrho$ sufficiently large, and then normalises by
$S=C_++C_-=1$.  Lean cross-normalises instead -- $K_\nu(\kappa)I_\nu(\kappa r)$ inside and
$I_\nu(\kappa)K_\nu(\kappa r)$ outside -- so the two branches match at $r=1$ identically,
the Wronskian gives the derivative jumps $\mp1$ with no denominator, and the eigenfunction
obtained is a non-zero constant multiple of the paper's eigenfunction with $S=1$.  The
Bessel values do not vanish: $I_\nu(\kappa),K_\nu(\kappa)\neq0$
for $\nu=0,1,2$ on the whole segment for $\varrho\geq32$
(\texttt{modifiedBessel\_denominators\_ne\_zero\_on\_segment}, by the argument given in the
proof of Lemma~\ref{lem:local:segment-root} rather than from the expansion); these facts are
fields of the certificate, and $K_0(\kappa)\neq0$ is how $b\neq0$ is read off, the first
Cartesian component of the Lean field at the origin being $-K_0(\kappa)/2$.  Lean also
routes the expansion through the single order-one product $zI_1(z)K_1(z)$, using the exact
algebraic identity $z^3\mathcal D(z)=2z-4zI_1(z)K_1(z)$, rather than expanding $\nu=0$ and
$\nu=2$ separately: the sectorised estimate $\abs{zI_1(z)K_1(z)-\tfrac12}\leq3/(\re z)^2$
(\texttt{FastDynamo/\allowbreak Local/\allowbreak PonomarenkoSectorOrderOneProduct.lean}) is the effective form of
\texttt{eq:bessel-product-expansion}, and the segment statement
$\varrho^3\mathcal D(\varrho(1+it))=(1+it)^{-3}+O(\varrho^{-2})$ carries the explicit
constant $225/(4\varrho^2)$.  The root itself is produced, in the paper and here, by the
intermediate value theorem in $t\in[\tfrac12,\tfrac34]$.

At $N=2\pi$ the Lean theorem has no hypothesis and also returns the numeral
$448-4\pi^2\leq\re\lambda_*$, which comes from the formalisation's choice $\varrho=32$
(Lemma~\ref{lem:local:segment-root}); the paper asserts only $\re\lambda_*>0$, via
$\re\lambda_*\geq\tfrac7{16}\varrho^2-4\pi^2$.  The auxiliary clauses of the paper's
statement -- the angular form of the eigenfunction, smoothness on each side of the circle,
exponential decay -- are separate declarations (\texttt{horizontalMode},
\texttt{horizontalInteriorField\_contDiff}, \texttt{horizontalMode\_tail\_le}) rather than
conjuncts of the headline theorem.  The paper's sentence ``the transmission conditions
cancel the interface distributions'' is, in Lean, a \emph{pointwise} identity between the
two one-sided jump densities on the unit circle, so the result does not depend on the
normalisation of the circle measure.

\begin{proof}
  \leanok
  \uses{def:local:cartesian-mode, lem:local:segment-root, def:local:jump-flow}
  Writing $D_j=\partial_jb+w_\perp b_j$ for the pivot field, the weak equation
  $\mathfrak l_{V^0}(b,\phi)=-\lambda_*\inner b\phi$ is equivalent to the distributional
  identity $\sum_j\partial_jD_j=M+\lambda_*b$ with
  $M=N^2b+w_\perp\cdot\nabla_\perp b+iNw_3b+w_\perp\Div_\perp b$.  Away from $r=1$ this is
  the bulk identity: the Cartesian Laplacian acts on the two circular combinations
  $b_1\pm ib_2$ with angular indices $(j\pm1)^2=4,0$, giving $\Delta_\perp b=\kappa^2b$,
  while the
  transport and stretching terms cancel because $(b\cdot\nabla_\perp)V^0_\perp
  =V^0_\perp\cdot\nabla_\perp b-i\Omega b$ and $iNU=-i\Omega$.  Across $r=1$ the
  distributional derivative of $D_j$ acquires a surface term with density
  $n_j\jump{D_j}$; the outward normal is the position vector and
  $\jump{\partial_rb}=\Omega B_re_\theta$ is the fourth transmission law, which is
  \eqref{eq:local:dispersion}.  The two one-sided densities are therefore pointwise
  opposite and the surface term vanishes.  Lemma~\ref{lem:local:segment-root} supplies a
  certificate, and $b\neq0$ because its first component at the origin is
  $-K_0(\kappa)/2\neq0$.
\end{proof}

\section{Spectral isolation and regularisation}

\begin{lemma}[Compactness of the velocity form and the analytic Fredholm family]
  \label{lem:local:compactness}
  \lean{FastDynamo.HorizontalWeakForm.freeHorizontalForm,
    FastDynamo.HorizontalWeakForm.horizontalVelocityForm,
    FastDynamo.HorizontalWeakForm.horizontalPencilForm,
    FastDynamo.HorizontalWeakForm.freeHorizontalForm_hasComplexCoercivityConstant,
    FastDynamo.HorizontalWeakForm.freeHorizontalForm_bijective,
    FastDynamo.HorizontalWeakForm.isCompactOperator_horizontalVelocityZeroOrder_of_ae_zero_off,
    FastDynamo.isCompactOperator_adjoint}
  \leanok
  \uses{def:local:horizontal-operator, lem:local:garding}
  Let $w\in\Lp^\infty(\R^2;\R^3)$ vanish almost everywhere off a set of finite measure.
  On $V=\Hs^1(\R^2;\CC^2)$ with antidual $V^*=\Hs^{-1}$, write
  $\mathbb A_w(z)=F_z+K_w$ with $F_z=z-(\Delta_\perp-N^2)$.  Then $F_z\colon V\to V^*$ is
  an isomorphism for $\re z>0$, $K_w\colon V\to V^*$ is compact, and
  $\mathbb A_w(z)=F_z(\id+F_z^{-1}K_w)$ is an analytic Fredholm family on $\{\re z>0\}$.
\end{lemma}

Three deviations.  First, the hypothesis on the velocity: the paper takes $w$ with compact
support, $\supp w\subset D(0,R_w)$, and restricts to that disc; Lean assumes only that $w$
vanishes almost everywhere off a measurable set $s$ of finite measure and restricts to $s$,
local Rellich compactness being proved for finite-measure sets
(\texttt{FastDynamo/Analysis/WholePlaneBandLimitedRellichUnconditional.lean}).  This is more
general, and is applied to $V^0$, which vanishes off the unit disc.  Second, the paper
inverts $F_z$ by the Fourier multiplier $(z+\abs\xi^2+4\pi^2)^{-1}$; Lean instead observes that
$\re F_z(b,b)=\norm{\nabla b}_2^2+(N^2+\re z)\norm b_2^2$ and applies complex
Lax--Milgram, with coercivity constant $\min(1,\re z)$.  A form is determined by its
values, so the operator is the same.  Third, Mathlib has no Schauder theorem, so the
compactness of the adjoint $R^*$ needed for the ``derivative on the input'' half of the
argument is \texttt{isCompactOperator\_adjoint}, proved here from scratch in the
Hilbert-space shortcut $\norm{S^\dagger y}^2=\re\inner y{SS^\dagger y}$.  The argument itself
is the paper's: §2.3 of the paper takes $R\colon\Hs^1(\R^2)\to\Lp^2(D(0,R_w))$ compact by
Rellich, $R^*\colon\Lp^2(D(0,R_w))\to\Hs^{-1}(\R^2)$ compact as its adjoint, and decomposes
each of the three kinds of term through $R$, multiplication and $R^*$; only the proof of the
adjoint step is supplied here.

\begin{proof}
  \leanok
  \uses{lem:local:garding}
  Let $s$ be the finite-measure set off which $w$ vanishes.  Terms with a derivative on
  the test function factor through the compact restriction $\Hs^1(\R^2)\to\Lp^2(s)$;
  terms with a derivative on the input factor through its adjoint, a compact map
  $\Lp^2(s)\to\Hs^{-1}(\R^2)$, after the coefficient has been moved across the
  inner product (multiplication by a real $\Lp^\infty$ coefficient is self-adjoint); the
  zeroth-order term factors through the same restriction.  Multiplication by a bounded
  coefficient is continuous in each case, so $K_w$ is compact by Rellich's theorem.  Since
  $F_z$ is affine in $z$, $z\mapsto F_1^{-1}F_z$ is entire and its inverse is analytic on
  the invertible locus by the Neumann series, which makes the family analytic Fredholm.
\end{proof}

\begin{lemma}[Discreteness of the right-half-plane spectrum and the isolating contour]
  \label{lem:local:isolation}
  \lean{FastDynamo.HorizontalWeakForm.horizontalSpectrum,
    FastDynamo.HorizontalWeakForm.isDiscrete_horizontalSpectrum,
    FastDynamo.HorizontalWeakForm.exists_ne_zero_horizontalPencilForm_eq_zero,
    FastDynamo.HorizontalWeakForm.exists_isEigenpair_of_mem_horizontalSpectrum,
    FastDynamo.HorizontalWeakForm.mem_horizontalSpectrum_iff,
    FastDynamo.HorizontalWeakForm.exists_punctured_ball_isInvertible_horizontalPencilForm,
    FastDynamo.HorizontalWeakForm.HorizontalIsolation,
    FastDynamo.HorizontalWeakForm.HorizontalIsolation.projection,
    FastDynamo.HorizontalWeakForm.HorizontalIsolation.isIdempotentElem_projection,
    FastDynamo.HorizontalWeakForm.HorizontalIsolation.finiteDimensional_projection_range,
    FastDynamo.HorizontalWeakForm.HorizontalIsolation.projection_apply_eigenvector,
    FastDynamo.HorizontalWeakForm.HorizontalIsolation.projection_ne_zero,
    FastDynamo.HorizontalMode.modeIsolation,
    FastDynamo.HorizontalMode.modeIsolation_spec}
  \leanok
  \uses{lem:local:compactness, lem:jump-horizontal-mode}
  Let $w$ be as in Lemma~\ref{lem:local:compactness}.  The set of $z$ with $\re z>0$ at
  which $\mathbb A_w(z)$ is not invertible is discrete, every point of $\{\re z>0\}$ has a
  punctured neighbourhood on which $\mathbb A_w(z)$ is invertible, and that set is exactly the
  set of eigenvalues of $A_w$ with positive real part; so the spectrum of $A_w$ in the open
  right half-plane consists of isolated eigenvalues.  Applied to $w=V^0$ this fixes a
  positively oriented circle $\Gamma$, centred at $\lambda_*$ and of radius at most
  $\tfrac14\re\lambda_*$, whose closed interior $\cG$ lies in $\{\re z>0\}$ and meets that
  set only in $\lambda_*$; the Riesz projection
  $P^0=(2\pi i)^{-1}\int_\Gamma(z-A_{V^0})^{-1}\dd z$ is idempotent, of finite rank, and
  satisfies $P^0b=b\neq0$ for the eigenmode $b$.
\end{lemma}

The paper's analytic Fredholm paragraph in §2.3 concludes that $\mathbb A_w(z)$ is invertible
except at a discrete set of points of the open right half-plane, that at each exceptional point
Fredholm index zero gives a non-zero kernel, which the definition of $A_w$ identifies with
$\ker(z-A_w)$, and hence that the spectrum of $A_w$ in the open right half-plane consists of
isolated eigenvalues.  Lean formalises all of it: the discreteness
(\texttt{isDiscrete\_horizontalSpectrum}, where \texttt{horizontalSpectrum} is \emph{defined}
as the non-invertibility set of $\mathbb A_w(z)$ in $\{\re z>0\}$), the punctured
neighbourhoods, and the identification of that set with the eigenvalues
(\texttt{mem\_horizontalSpectrum\_iff}, \texttt{FastDynamo/Local/HorizontalSpectrumEigenvalue.lean}).
The non-zero kernel at an exceptional point comes from the factorisation
$\mathbb A_w(z)=F_z(\id-K(z))$ and the Fredholm alternative for the compact $K(z)$, which is
the paper's index-zero statement in the form ``injective implies invertible''.  The eigenvector
lies in the form-operator domain of $A_w$, which is the paper's $D(A_w)$ for $w\in\Lp^\infty$ with
compact support; the hypotheses are those of \texttt{isDiscrete\_horizontalSpectrum}
($w$ vanishing off a set of finite measure).  In the other direction, Lean proves that the
Riesz projection of every isolating circle has finite rank
(\texttt{HorizontalIsolation.finiteDimensional\_projection\_range}), which the paper does not
state.  The radius bound $r\leq\tfrac14\re\lambda_*$, which gives
$g_*\geq\tfrac34\re\lambda_*$, is a choice of the formalisation; the paper chooses $\Gamma$
only qualitatively.

The Lean contour is a \texttt{SpectralCircle}: a centre, a radius, and the proof that the
closed disc lies in $\{\re z>0\}$; its \texttt{growthFloor} is
$\re(\text{centre})-\text{radius}$, which \emph{is} $\min_{\zeta\in\cG}\re\zeta$.  The
resolvent delivered by the isolation bundle is domain valued, that is $\Lp^2\to D(A_w)$
with the $\Hs^1$ graph domain; the upgrade to $\Lp^2\to\Hs^2$ is
Lemma~\ref{lem:local:h2-resolvent} and is where \eqref{eq:local:continuation-bound} is
proved.

\begin{proof}
  \leanok
  \uses{lem:local:compactness, lem:jump-horizontal-mode}
  The analytic Fredholm theorem applied to $\id+F_z^{-1}K_w$ on the connected open set
  $\{\re z>0\}$, at the invertible real base point supplied by the Gårding estimate of
  Lemma~\ref{lem:local:garding} (the paper uses \texttt{eq:sectorial-estimate-b} in the
  same way), shows that the non-invertibility set is discrete.  An eigenvector of $A_w$
  makes $\mathbb A_w$ singular at its eigenvalue; conversely, at a point of non-invertibility
  the Fredholm alternative for the compact part gives a non-zero null vector, which solves the
  weak eigenvalue identity and is therefore an eigenvector.  Around $\lambda_*$ pick a punctured ball
  of invertibility; a circle of radius at most $\tfrac14\re\lambda_*$ inside it is $\Gamma$.
  Idempotence of the contour integral is the annular resolvent identity, finite rank
  follows from the compactness of the state resolvent, and $P^0b=b$ is Cauchy's formula
  for $(z-A_{V^0})^{-1}b=(z-\lambda_*)^{-1}b$.
\end{proof}

\begin{definition}[The smoothed velocity and the diffusivity family]
  \label{def:local:smoothed-velocity}
  \lean{FastDynamo.HorizontalMode.horizontalSmoothedVelocityLInf,
    FastDynamo.HorizontalMode.horizontalSmoothedModeVelocity,
    FastDynamo.HorizontalMode.horizontalSmoothingErrorL4,
    FastDynamo.HorizontalWeakForm.horizontalWeakFormDelta,
    FastDynamo.HorizontalWeakForm.horizontalPencilFormDelta,
    FastDynamo.HorizontalWeakForm.horizontalFormDeltaOperator,
    FastDynamo.horizontalLocalOperator}
  \leanok
  \uses{def:local:jump-flow, def:local:horizontal-operator}
  For $0<h<1/2$ let $\chi_h$ be a smooth radial cut-off with $0\leq\chi_h\leq1$, equal to
  one on $[0,1-h]$ and zero on $[1+h,\infty)$, and set
  \begin{equation}
    v_h(y)=\chi_h(\abs y)(-\Omega y_2,\Omega y_1,U)\in\C_c^\infty(\R^2;\R^3),
    \qquad\Div_\perp(v_h)_\perp=0 .
    \label{eq:local:smooth-velocity}
  \end{equation}
  This display is the paper's \texttt{eq:smooth-local-velocity}; the
  solenoidality is immediate from $\nabla\chi_h(\abs y)\cdot y^\perp=0$.
  For $\delta>0$ put
  $\mathfrak l_{h,\delta}=\mathfrak l_{v_h}
   +(\delta-1)\bigl(\inner{\nabla_\perp\cdot}{\nabla_\perp\cdot}+N^2\inner\cdot\cdot\bigr)$,
  with shifted form operator $\mathbb A_{h,\delta}(z)$ and generator
  \begin{equation}
    \cL_\delta b=\delta(\Delta_\perp-N^2)b-v_\perp\cdot\nabla_\perp b-iNv_3b
      +(b\cdot\nabla_\perp)v_\perp,
    \qquad D(\cL_\delta)=\Hs^2(\R^2;\CC^2).
    \label{eq:local:generator}
  \end{equation}
\end{definition}

A note on $N$.  The paper has no axial wave number parameter: it works in the axial mode
$\e^{2\pi ix_3}$ and writes $4\pi^2$ and $2\pi i$ in \texttt{eq:horizontal-induction},
\texttt{eq:local-generator} and everywhere else.  This chapter keeps the symbol $N$, because
the Lean
development does: \texttt{LocalSeedParameters.wavenumber} is $2\pi k_*$ with $k_*$ a
positive natural-number field, and the concrete inhabitant is taken at $k_*=1$
(\texttt{exists\_localSeedC2\_two\_pi}), so every statement below is the paper's at
$k_*=1$ and is more general otherwise.  Keeping $k_*$ free costs nothing and is what makes
the rescaling (Lemma~5.2 of the paper) in a later chapter a statement about an arbitrary scale family.

The whole $\delta$-dependence is the single affine term $(\delta-1)\cdot$(diffusion form),
which is what makes the constant of \eqref{eq:local:form-small} independent of $z$: a
difference of two members of the family is $z$-free, the free form $F_z$ cancelling.
\texttt{horizontalLocalOperator N delta w} is the \emph{strong} operator
\eqref{eq:local:generator} on $\Hs^2$; that it agrees with the form operator on the form
domain is \texttt{FastDynamo/Local/HorizontalFormStrongIdentification.lean} and is the
bridge the resolvent bound below needs.

\begin{lemma}[The one form perturbation estimate]
  \label{lem:local:form-small}
  \lean{FastDynamo.HorizontalWeakForm.horizontalPencilFormDelta_sub,
    FastDynamo.HorizontalWeakForm.complexFormOpNorm_horizontalPencilFormDelta_sub_le,
    FastDynamo.HorizontalMode.exists_norm_horizontalSmoothingErrorL4_le,
    FastDynamo.HorizontalWeakForm.exists_complexFormOpNorm_horizontalModePencil_sub_le}
  \leanok
  \uses{def:local:smoothed-velocity, def:local:jump-flow}
  There is $C$, independent of $z$, of $h\in(0,\tfrac12)$ and of $\delta$, with
  \begin{equation}
    \norm{\mathbb A_{h,\delta}(z)-\mathbb A_{0,1}(z)}_{V\to V^*}
    \leq C\bigl(h^{1/4}+\abs{\delta-1}\bigr).
    \label{eq:local:form-small}
  \end{equation}
\end{lemma}

\begin{proof}
  \leanok
  \uses{def:local:smoothed-velocity}
  The difference splits as $(K_{v_h}-K_{V^0})+(\delta-1)\cdot$(diffusion form) with no $z$
  in it.  The velocity difference is uniformly bounded and supported in an annulus of area
  $O(h)$, so $\norm{v_h-V^0}_4\leq Ch^{1/4}$; each first-order term is then estimated with
  H\"older exponents $(4,4,2)$ or $(4,2,4)$ according to which field carries the
  derivative, using $\Hs^1(\R^2)\hookrightarrow\Lp^4(\R^2)$, and the coefficient is moved
  onto the test function by self-adjointness of multiplication by a real $\Lp^\infty$
  field.  The zeroth-order term is bounded the same way and the diffusion difference
  directly on $V\times V$.
\end{proof}

\begin{lemma}[The Riesz projections move by less than a quarter]
  \label{lem:local:riesz-perturbation}
  \lean{FastDynamo.HorizontalWeakForm.HorizontalNeumannData,
    FastDynamo.HorizontalWeakForm.HorizontalNeumannData.projection,
    FastDynamo.HorizontalWeakForm.HorizontalNeumannData.norm_stateResolvent_sub_le_two,
    FastDynamo.HorizontalWeakForm.HorizontalNeumannData.norm_projection_sub_le_two,
    FastDynamo.HorizontalMode.exists_seedNeumannData,
    FastDynamo.HorizontalMode.exists_seedThresholds}
  \leanok
  \uses{lem:local:form-small, lem:local:isolation}
  The inverses $\mathbb A_{0,1}(z)^{-1}$ are uniformly bounded on $\Gamma$, so for
  $h^{1/4}+\abs{\delta-1}$ small a Neumann series gives
  $\sup_{z\in\Gamma}\norm{\mathbb A_{h,\delta}(z)^{-1}-\mathbb A_{0,1}(z)^{-1}}
  \leq C(h^{1/4}+\abs{\delta-1})$; restricting to $\Lp^2$ gives the resolvents of
  $\cL_{h,\delta}$ on $\Gamma$, and integrating gives
  $\norm{P_{h,\delta}-P^0}_{2\to2}\leq C(h^{1/4}+\abs{\delta-1})$.  There are thresholds
  $h_1,\eta_1>0$ with $\norm{P_{h,\delta}-P^0}<\tfrac14$ whenever $h\leq h_1$ and
  $\abs{\delta-1}\leq\eta_1$.
\end{lemma}

The Lean perturbation is run on the open annulus $B(\lambda_*,3r/2)\setminus\bar
B(\lambda_*,r/4)$ rather than on $\Gamma$ alone; this costs nothing, because
\eqref{eq:local:form-small} is $z$-free, and it buys idempotence of the perturbed
projection from the annular Riesz calculus.  The explicit thresholds are
$h_1=\min(h_0,(8(K+1))^{-4})$ and $\eta_1=\min(\eta_0,(8(K+1))^{-1})$, which give
$K(h^{1/4}+\abs{\delta-1})\leq K/(4(K+1))<\tfrac14$.

\begin{proof}
  \leanok
  \uses{lem:local:form-small, lem:local:isolation}
  Uniform boundedness of $\mathbb A_{0,1}(z)^{-1}$ on the compact annulus is continuity of
  the inverse plus compactness.  With $\eta=C(h^{1/4}+\abs{\delta-1})$ bounding the
  perturbation, the algebraic identity $S^{-1}=\id-(TD)S^{-1}$ gives
  $\norm{S^{-1}}\leq(1-B\eta)^{-1}$ and hence the resolvent difference; the contour
  integral of a bound is the bound times the length of $\Gamma$ divided by $2\pi$, that is
  times the radius.  The thresholds are the stated explicit minima.
\end{proof}

\begin{lemma}[A compactly supported spectral witness]
  \label{lem:local:witness}
  \lean{FastDynamo.rayleighRe,
    FastDynamo.wholePlaneC2DiskCutoff,
    FastDynamo.exists_diskSupported_witness,
    FastDynamo.HorizontalMode.exists_horizontalSeedWitness}
  \leanok
  \uses{lem:local:isolation}
  There are $L_*>0$ and $f_*\in\Lp^2(\R^2;\CC^2)$ with $\norm{f_*}_2=1$, $f_*=0$ almost
  everywhere off $D(0,L_*)$, and $\re\inner{P^0f_*}{f_*}>\tfrac34$.
\end{lemma}

\textbf{Deviation.}  The paper asks for $f_*\in\C_c^\infty(\R^2;\CC^2)$ in Proposition~2.2
and obtains it by density in the proof; Lean produces
the indicator truncation $\one_{D(0,L)}x$ of a normalised element of $\Ran P^0$,
renormalised, which is an $\Lp^2$ field with compact support and not a smooth one.  Nothing
downstream uses smoothness: the global layer consumes only $\norm{f_*}_2=1$, the almost
everywhere support statement and $\re\inner{P_\delta f_*}{f_*}>\tfrac12$.  The support
statement is also realised on the \emph{closed} disc ($L_*<\norm y$ implies $f_*(y)=0$
a.e.), which is the weaker and therefore the safe direction, since
\eqref{eq:transfer:scaled-witness-support} is strict.

\begin{proof}
  \leanok
  \uses{lem:local:isolation}
  Take a unit $x\in\Ran P^0$, which exists because $P^0\neq0$.  Then
  $\re\inner{P^0x}x=1$, and $\one_{D(0,L)}x\to x$ in $\Lp^2$ as $L\to\infty$ by dominated
  convergence, so the normalised Rayleigh quotient
  $y\mapsto\re\inner{P^0y}y/\norm y^2$, which is continuous away from the origin, exceeds
  $\tfrac34$ for $L$ large.
\end{proof}

\begin{lemma}[The uniform $\Lp^2\to\Hs^2$ resolvent bound]
  \label{lem:local:h2-resolvent}
  \lean{FastDynamo.exists_wholePlaneH2_of_laplacian_eq_l2,
    FastDynamo.HorizontalWeakForm.horizontalStrongPairMap,
    FastDynamo.HorizontalWeakForm.norm_horizontalStrongPairMap_le_bound,
    FastDynamo.HorizontalWeakForm.horizontalLocalResolventBoundOfBounds,
    FastDynamo.HorizontalWeakForm.horizontalStrongResolvent_stateResolvent,
    FastDynamo.HorizontalWeakForm.horizontalStrongResolvent_rieszProjection}
  \leanok
  \uses{def:local:smoothed-velocity, lem:local:riesz-perturbation}
  For $\delta$ in a compact interval $I=[\delta_-,\delta_+]$ with $\delta_->0$ the form
  domain of $\mathfrak l_{h,\delta}$ is $\Hs^2(\R^2;\CC^2)$, the strong operator
  $\cL_\delta$ of \eqref{eq:local:generator} has the same resolvent as the form operator,
  hence the same Riesz projection, and there is one constant $L$ with
  \begin{equation}
    \sup_{\delta\in I,\ z\in\Gamma}
      \norm{(z-\cL_\delta)^{-1}}_{\Lp^2\to\Hs^2}\leq L .
    \label{eq:local:continuation-bound}
  \end{equation}
\end{lemma}

\textbf{Lean proves more than the paper asks.}  The paper's
\texttt{eq:local-continuation-bound} in Proposition~2.2 (\texttt{prop:local-seed}) is the
$\Lp^2\to\Hs^1$ bound; what is stated and proved here is the $\Lp^2\to\Hs^2$ bound, and the
paper's bound is its corollary through the embedding $\Hs^2\hookrightarrow\Hs^1$.  This is not extra work.  The Lean
$\cL_\delta$ is a \texttt{DomainOperator} with bundled domain $\Hs^2(\R^2;\CC^2)$ -- the
domain identification that the paper's proof of Proposition~2.2 performs too -- and
\texttt{HasResolventOn} stores a domain-valued inverse $\Lp^2\to\Hs^2$, so the $\Hs^2$
operator norm is the only norm available at that type.  Only the $\Lp^2$ and gradient halves
are ever consumed, by \texttt{norm\_embedding\_inv\_apply\_le} and
\texttt{norm\_deriv\_inv\_apply\_le}, which is exactly the $\Hs^1$ content that
Corollary~\ref{cor:scaled-local-resolvents} needs.

The $\Hs^2$ estimate $\norm b_{\Hs^2}\leq C(\norm b_2+\norm{\Delta_\perp b}_2)$ is obtained
from Banach's isomorphism theorem applied to $h\mapsto(h,\Delta h)$ onto the closed pair
space, not from the Fourier symbol; the constant is therefore unspecified, which is all
\eqref{eq:local:continuation-bound} asks for.  The identification of the two Riesz
projections is the step that transports the witness clause, proved for the form operator,
to the strong generator that the global layer uses.

\begin{proof}
  \leanok
  \uses{def:local:smoothed-velocity, lem:local:riesz-perturbation}
  An element of the form-operator domain lies in $\Hs^1$; since $v_h$ is smooth with
  bounded derivatives the lower-order terms lie in $\Lp^2$, so
  $\delta\Delta_\perp b\in\Lp^2$ and $b\in\Hs^2$, with
  $\delta\norm{\Delta_\perp b}_2\leq\norm f_2+C_v\norm b_{\Hs^1}$.  The form inverse gives
  $\norm b_{\Hs^1}\leq C\norm f_2$ uniformly on $I\times\Gamma$ by
  Lemma~\ref{lem:local:riesz-perturbation}, and $\delta\geq\delta_->0$ turns the two into
  one $\delta$-free bound.  The two resolvents have the same ambient value because the
  $\Hs^2$ lift of a graph element has the same $\Lp^2$ value, and Riesz projections of
  operators with equal state resolvents on the same circle agree.
\end{proof}

\section{The local seed package}

\begin{definition}[The local seed]
  \label{def:local:seed}
  \lean{FastDynamo.LocalSeedParameters, FastDynamo.LocalSeedC2,
    FastDynamo.LocalSeedC2.velocity,
    FastDynamo.LocalSeedC2.coefficients,
    FastDynamo.LocalSeedC2.resolver,
    FastDynamo.LocalSeedC2.norm_resolver_le,
    FastDynamo.LocalSeedC2.projection,
    FastDynamo.LocalSeedC2.re_inner_projection_witness_gt,
    FastDynamo.LocalSeedC2.growthFloor,
    FastDynamo.LocalSeedC2.growthFloor_pos}
  \leanok
  \uses{def:local:smoothed-velocity, lem:local:h2-resolvent, lem:local:witness}
  A \emph{local seed} is the data
  \[
    \bigl(N=2\pi k_*,\ v=v_h,\ I=[\delta_-,\delta_+],\ \Gamma,\ \cG,\ g_*,\ f_*,\ L_*,\ L\bigr)
  \]
  with $0<\delta_-<1<\delta_+$, $v$ helical with $\Div_\perp v_\perp=0$, $\Gamma$ a
  positively oriented circle with $\cG\subset\{\re z>0\}$,
  $\Gamma\subset\rho(\cL_\delta)$ for every $\delta\in I$ with the uniform bound
  \eqref{eq:local:continuation-bound}, the witness clause
  \eqref{eq:local:spectral-witness} and the gap
  $g_*=\min_{\zeta\in\cG}\re\zeta>0$.
\end{definition}

The Lean record is \texttt{FastDynamo.LocalSeedC2}
(\texttt{FastDynamo/Local/LocalSeedC2.lean}), extending
\texttt{LocalSeedParameters} by the velocity data, the family of resolvents and the
witness.  The velocity is a \emph{definition} of the record, not a field, so the record
cannot be inhabited by a non-helical field, and $\Div_\perp v_\perp=0$ is a theorem about
it.  The growth floor is \texttt{contour.growthFloor}$=\re(\text{centre})-\text{radius}$,
which is exactly $\min_{\zeta\in\cG}\re\zeta$ over the closed disc.

\begin{proposition}[Uniform instability of the reference model on $\R^2$]
  \label{prop:local-seed}
  \lean{FastDynamo.HorizontalMode.exists_localSeedC2_two_pi,
    FastDynamo.HorizontalMode.localSeedC2,
    FastDynamo.HorizontalMode.localSeedC2_kStar,
    FastDynamo.HorizontalMode.localSeedC2_wavenumber,
    FastDynamo.HorizontalMode.seedGrowthFloorBound,
    FastDynamo.HorizontalMode.seedGrowthFloorBound_le_localSeedC2_growthFloor,
    FastDynamo.HorizontalMode.localSeedC2_angularSpeed_pos}
  \leanok
  \uses{def:local:seed, lem:local:riesz-perturbation, lem:local:witness,
    lem:local:h2-resolvent, lem:jump-horizontal-mode}
  With $N=2\pi$ there are $h\in(0,\tfrac12]$, an interval $I=[\delta_-,\delta_+]$ with
  $0<\delta_-<1<\delta_+$, a compactly supported unit $f_*$ and $L_*>0$ with the following
  properties.  For the velocity $v=v_h$ of \eqref{eq:local:smooth-velocity} and
  $\delta\in I$, the generator $\cL_\delta$ of \eqref{eq:local:generator} satisfies
  $\Gamma\subset\rho(\cL_\delta)$ together with \eqref{eq:local:continuation-bound}, and the
  Riesz projections $P_\delta=(2\pi i)^{-1}\int_\Gamma(z-\cL_\delta)^{-1}\dd z$ obey
  \begin{equation}
    \norm{f_*}_2=1,\qquad\supp f_*\subset D(0,L_*),\qquad
    \re\inner{P_\delta f_*}{f_*}>\tfrac12\quad(\delta\in I).
    \label{eq:local:spectral-witness}
  \end{equation}
  In particular each $\cL_\delta$ has spectrum inside $\Gamma$ and
  $g_*=\min_{\zeta\in\cG}\re\zeta>0$.
\end{proposition}

The Lean inhabitant is at $k_*=1$, hence $N=2\pi$, with $h=\min(h_0,h_1,\tfrac12)$ and
$I=[1-\eta,1+\eta]$, $\eta=\min(\eta_0,\eta_1,\tfrac12)$, so
$0<\tfrac12\leq\delta_-<1<\delta_+\leq\tfrac32$; the growth floor obeys
\[
  g_*\ \geq\ \tfrac34\bigl(448-4\pi^2\bigr)\ \approx\ 306.4,
  \qquad\text{and in particular}\qquad g_*\geq288 .
\]
These numerals come from the formalisation's choices $\varrho=32$
(Lemma~\ref{lem:local:segment-root}) and $r\leq\tfrac14\re\lambda_*$
(Lemma~\ref{lem:local:isolation}); the paper asserts only $g_*>0$
(\texttt{eq:spectral-gap}).  Three points of shape.  First, the clauses
``$v_h\in\C_c^\infty$ and $\Div_\perp(v_h)_\perp=0$'' are not part of the proposition: the
paper states them with \texttt{eq:smooth-local-velocity} itself, just before the
proposition, and the proposition only says ``for the reference velocity profile on $\R^2$,
$v=v_h$, from equation \eqref{eq:local:smooth-velocity}''.  The Lean record has the same shape: the velocity
is stored as the triple $(h,\Omega,U)$, and the vanishing of $v_h$ outside $D(0,1+h)$ is a
theorem about it (\texttt{parameterizedSmoothedVelocity\_eq\_zero\_of\_radius\_le\_norm}),
not a field; the horizontal divergence-freeness is proved for the assembled global velocity
(Lemma~\ref{lem:global-velocity}), not for $v_h$ on its own.  Second,
the resolvent bound quoted above is the $\Lp^2\to\Hs^2$ one of
Lemma~\ref{lem:local:h2-resolvent}; the paper asks only for the $\Lp^2\to\Hs^1$ bound,
which follows from it, and the deviation is recorded at that node.  Third, two small
differences of the Lean inhabitant: the smoothing width is built as a minimum with
$\tfrac12$, so $h=\tfrac12$ is not excluded where the paper writes
$h\in(0,\tfrac12)$ (hence $h\in(0,\tfrac12]$ above), and $f_*$ is an $\Lp^2$ field with compact support rather than a
$\C_c^\infty$ one (Lemma~\ref{lem:local:witness}); nothing downstream uses either.

\begin{proof}
  \leanok
  \uses{lem:jump-horizontal-mode, lem:local:isolation, lem:local:riesz-perturbation,
    lem:local:witness, lem:local:h2-resolvent, def:local:seed}
  Lemma~\ref{lem:jump-horizontal-mode} supplies $\Omega$, $\kappa$ and $\lambda_*$ at
  $N=2\pi$, and Lemma~\ref{lem:local:isolation} the contour $\Gamma$ with
  $g_*\geq\tfrac34\re\lambda_*$ and the projection $P^0$.
  Lemma~\ref{lem:local:witness} gives $f_*$, $L_*$ with $\re\inner{P^0f_*}{f_*}>\tfrac34$,
  and Lemma~\ref{lem:local:riesz-perturbation} thresholds with
  $\norm{P_\delta-P^0}<\tfrac14$; subtracting gives
  $\re\inner{P_\delta f_*}{f_*}>\tfrac12$, so $P_\delta\neq0$ and the spectrum inside
  $\Gamma$ is non-empty.  Lemma~\ref{lem:local:h2-resolvent} supplies the single constant
  $L$ of \eqref{eq:local:continuation-bound} and identifies the projection of the strong
  generator with the one just estimated.  Assembling the pieces into the record of
  Definition~\ref{def:local:seed} finishes.
\end{proof}


\chapter{The global velocity}
% Chapter: The global velocity.  Section 3 of the paper (sec:velocity).
% Owned shared labels: def:velocity:scale-family, def:velocity:placement,
% def:velocity:global-velocity.  Paper label: lem:global-velocity (Lemma 3.1 of the paper).

The local seed of the previous chapter is a single smooth, compactly supported,
divergence-free velocity $v $ on $\R^2 $ together with one axial wave number
$N=2\pi k_* $, one interval $I $ of admissible rescaled diffusivities, one contour
$\Gamma $ and the growth floor $g_*>0 $ (node \texttt{prop:local-seed}).  It produces
instability at one scale only.  This chapter installs rescaled copies of that seed at a
geometric sequence of horizontal scales $\ell_n\downarrow0 $, packs them into pairwise
disjoint discs of a fixed Euclidean chart of $\T^2 $, adds in each disc a small axial
velocity whose only role is to shift that cell's spectrum vertically by the fixed amount
$-i\tau $, and sums everything into one autonomous velocity field $u $ on $\T^3 $.

The point of the construction is a quantifier order: nothing here depends on the physical
diffusivity $\eps $.  In this blueprint's shorthand the order of the choices is
\[
  (v,I,\Gamma,\cG,g_*,f_*,L_*,q,\vartheta)
  \quad\longrightarrow\quad\tau
  \quad\longrightarrow\quad K_0 :
\]
the local seed, the single geometric ratio $q $ and the axial cut-off
$\vartheta $ come first, then the axial shift $\tau $ of the background chapter, then the
base Fourier index $K_0 $, which fixes the field $u $ and the threshold $\eps_0 $; only
afterwards is a cell matched to a given $\eps $.  The paper states this order in prose, in
Section~3.1 (``Choice of the parameters'') and in the paragraph collecting the parameter choices at
the end of Section~7.2, just before Remark~7.5 of the paper (\texttt{rem:flexible-buffer-scales});
it does not display the chain.  Every estimate that chooses $K_0 $ is uniform once the
earlier parameters are fixed.  The chapter ends
with the field's regularity (Lemma~\ref{lem:global-velocity}): $u $ is Lipschitz but not
$\C^1 $, and the failure of $\C^1 $ is exactly the accumulation of nonvanishing gradients at
infinitely many shrinking scales.  The dynamical consequences of the Lipschitz bound
(nodes \texttt{prop:particle-flow} and \texttt{cor:ideal-growth}) are in the particle
chapter.

Terminology: the paper calls the $y $-directions \emph{horizontal} and this chapter
follows it; the Lean identifiers use the adjective \texttt{transverse}
(\texttt{TransverseRealW1Inf}, \texttt{transverseTorusPoint},
\texttt{IsTransverseDivergenceFree}, \dots).  The correspondence is recorded once in
\texttt{FastDynamo/Conventions.lean}.

\section{Diffusivity intervals and horizontal scales}

\begin{definition}[The scale family]
  \label{def:velocity:scale-family}
  \lean{FastDynamo.ScaleRatios, FastDynamo.ScaleFamily,
    FastDynamo.ScaleFamily.fourierIndex, FastDynamo.ScaleFamily.ell,
    FastDynamo.ScaleFamily.bufferRadius, FastDynamo.ScaleFamily.cellCenter,
    FastDynamo.ScaleFamily.axialAmplitude,
    FastDynamo.ScaleFamily.diffusivityInterval, FastDynamo.ScaleFamily.epsilon0,
    FastDynamo.ScaleFamily.MatchingCell,
    FastDynamo.ScaleRatios.q_lt_scaleRatioCeiling, FastDynamo.ScaleFamily.rounding_bound,
    FastDynamo.ScaleFamily.TorusCellPlacement.FullVelocityConcrete.accumulationPoint}
  \leanok
  \uses{prop:local-seed, def:local:seed}
  Fix a ratio $1<q<\sqrt{\delta_+/\delta_-} $ and an integer $K_0 $ so large that
  $q+K_0^{-1}<\sqrt{\delta_+/\delta_-} $.  Set
  \[
    K_{n+1}=\lceil qK_n\rceil,
    \qquad
    \ell_n=\frac{k_*}{K_n},
    \qquad
    R_n=\sqrt{\ell_n},
    \qquad
    a_n=\frac{\tau}{2\pi K_n}=\frac{\tau}{2\pi k_*}\ell_n
  \]
  for a parameter $\tau\geq0 $ fixed later.  Here $\ell_n $ is the horizontal scale of the
  $n $-th rescaled copy of the local velocity (the \emph{core} scale), $R_n $ is the radius
  reserved around it (the \emph{buffer} scale) and $a_n $ is the amplitude of the added
  axial velocity.  The centres and the accumulation point are
  \[
    p_n=\Bigl(4\sum_{m\geq n}R_m,\tfrac12\Bigr),
    \qquad
    p_\infty=\bigl(0,\tfrac12\bigr),
  \]
  and the corresponding accumulation set in $\T^3 $ is the circle $\{p_\infty\}\times\T $.
  The diffusivity intervals and the threshold are
  \[
    \cI_n=\Bigl[\delta_-\frac{k_*^2}{K_n^2},\;\delta_+\frac{k_*^2}{K_n^2}\Bigr],
    \qquad
    \eps_0=\delta_+\frac{k_*^2}{K_0^2}.
  \]
  An index $n $ is \emph{admissible for} $\eps $ when $\delta_n:=\eps/\ell_n^2\in I $; an
  admissible index together with that membership is a \emph{matching cell}.
\end{definition}

In Lean the two rounds of choice are two structures: \texttt{ScaleRatios p} carries the
single ratio $q $ together with $q<\sqrt{\delta_+/\delta_-} $
(\texttt{ScaleRatios.q\_lt\_scaleRatioCeiling}), and \texttt{ScaleFamily r} carries only
$K_0 $ with the rounding bound $q+K_0^{-1}<\sqrt{\delta_+/\delta_-} $
(\texttt{ScaleFamily.rounding\_bound}), so that later stages may still enlarge $K_0 $:
the bound is monotone in $K_0 $, which is the paper's remark after
\texttt{eq:ratio-choice} that every subsequent enlargement of $K_0 $ preserves it
(\texttt{ScaleFamily.rounding\_bound\_of\_le}).
All of $K_n,\ell_n,R_n,a_n,p_n,\cI_n,\eps_0 $ are then functions of a
\texttt{ScaleFamily} (modules \texttt{FastDynamo/Global/ScaleSelection.lean} and
\texttt{FastDynamo/Global/CellGeometry.lean}); the point $p_\infty $ is the constant
\texttt{accumulationPoint} of
\texttt{FastDynamo/Global/TorusGlobalVelocityQualitative.lean}.  A
\texttt{MatchingCell f eps} is an index bundled with the proof that
$\eps/\ell_n^2\in I $; uniqueness of the index is neither claimed nor needed, exactly as
in the paper.

\textbf{Deviation: the free wave number $k_* $.}  The paper sets $\ell_n=1/K_n $, so that
$K_n\ell_n=1 $, $a_n=\tau\ell_n/(2\pi) $, $\cI_n=[\delta_-/K_n^2,\delta_+/K_n^2] $ and
$\eps_0=\delta_+/K_0^2 $.  Lean keeps the local axial wave number $N=2\pi k_* $ of the seed
as a parameter ($k_* $ a positive integer, the field \texttt{LocalSeedParameters.kStar}) and
sets $\ell_n=k_*/K_n $; the displays above are written in that generality.  The concrete
seed has $k_*=1 $ (\texttt{globalSeedInterface\_kStar}), where every formula of the definition is the
paper's.

\begin{lemma}[Exact identities and geometric decay of the scale family]
  \label{lem:velocity:scale-arithmetic}
  \lean{FastDynamo.ScaleFamily.adjacent_ratio_bounds,
    FastDynamo.ScaleFamily.adjacent_ratio_sq_lt,
    FastDynamo.ScaleFamily.axial_wavenumber_match,
    FastDynamo.ScaleFamily.bufferRadius_sq,
    FastDynamo.ScaleFamily.ell_div_bufferRadius,
    FastDynamo.ScaleFamily.axialAmplitude_match,
    FastDynamo.ScaleFamily.summable_bufferRadius,
    FastDynamo.ScaleFamily.tendsto_bufferRadius_div_ell_atTop,
    FastDynamo.ScaleFamily.epsilon0_formula,
    FastDynamo.ScaleFamily.iUnion_diffusivityInterval,
    FastDynamo.ScaleFamily.exists_matchingCell,
    FastDynamo.ScaleFamily.TorusCellPlacement.FullVelocityConcrete.tendsto_cellCenter_accumulationPoint}
  \leanok
  \uses{def:velocity:scale-family}
  For every scale family:
  \begin{enumerate}
    \item $q\leq K_{n+1}/K_n\leq q+1/K_0<\sqrt{\delta_+/\delta_-} $, the paper's
      \texttt{eq:ratio-choice};
    \item $2\pi K_n\ell_n=N $ and $2\pi K_na_n=\tau $;
    \item $R_n^2=\ell_n $ and $\ell_n/R_n=R_n $, hence $R_n/\ell_n\to\infty $;
    \item $\sum_{n\geq0}R_n<\infty $, and $p_n\to p_\infty $;
    \item $\bigcup_{n\geq0}\cI_n=(0,\eps_0] $ with $\eps_0=\delta_+k_*^2/K_0^2 $; hence
      every $0<\eps\leq\eps_0 $ admits at least one index $n $ with $\eps/\ell_n^2\in I $.
  \end{enumerate}
\end{lemma}

Items (i)--(v) are separate Lean theorems, each holding for every
\texttt{ScaleRatios} and \texttt{ScaleFamily}.  Item (i) is the full three-part chain
\texttt{eq:ratio-choice}, one Lean theorem,
\texttt{ScaleFamily.adjacent\_ratio\_bounds}.  In the paper its upper bound gives the
overlap of consecutive diffusivity intervals (Section~3.2), and its lower bound gives $r_m\geq q $
for $m<n $ and $r_m\leq q^{-1} $ for $m>n $ in the proof of Proposition~6.1 of the paper
(\texttt{prop:background-resolvent}), where the axial shifts of distinct cells are separated
through $c_q=(q-1)/(q+1) $.  Item (ii) is the exact frequency matching
that makes each cell a literal dilate of the local problem, together with the identity that
makes the added axial velocity act as the fixed spectral translation $-i\tau $ on the
fixed-$K_n $ sector.  The paper sets $\ell_n=1/K_n $ and uses $K_n\ell_n=1 $ directly,
without a label.  Lean keeps $k_* $ free, so \texttt{ScaleFamily.axial\_wavenumber\_match}
reads $2\pi K_n\ell_n=N=2\pi k_* $ and is the paper's identity at $k_*=1 $.

\begin{proof}
  \leanok
  \uses{def:velocity:scale-family}
  (i) is the rounding estimate $qK_n\leq\lceil qK_n\rceil<qK_n+1 $ combined with
  $K_0\leq K_n $ and the rounding bound $q+K_0^{-1}<\sqrt{\delta_+/\delta_-} $.  (ii) is
  immediate from $\ell_n=k_*/K_n $ and
  $a_n=\tau\ell_n/(2\pi k_*) $; (iii) from $R_n=\sqrt{\ell_n} $.  (iv) follows from the
  geometric bound $\ell_n\leq\ell_0q^{-n} $, hence $R_n\leq R_0q^{-n/2} $.  For (v), the
  ratio condition $(q+1/K_0)^2<\delta_+/\delta_- $
  (\texttt{ScaleFamily.adjacent\_ratio\_sq\_lt}) makes consecutive intervals overlap, and
  both endpoints tend to zero, so their union is the half-open interval down from $\eps_0 $.
\end{proof}

\begin{lemma}[The base index may be chosen above any later bound]
  \label{lem:velocity:base-index}
  \lean{FastDynamo.ScaleFamily.exists_ge, FastDynamo.exists_scale_selection,
    FastDynamo.ScaleFamily.exists_largeK0_chartPacking_and_supportFit,
    FastDynamo.ScaleFamily.rounding_bound_of_le,
    FastDynamo.ScaleFamily.TorusCellPlacement.nonempty}
  \leanok
  \uses{def:velocity:scale-family}
  Admissible ratios exist, and for every $M\in\Z_{\geq0} $ there is a scale family with
  $K_0\geq M $.  Moreover, for every $R_v\geq0 $ a large enough $K_0 $ satisfies both
  \[
    4\sum_{n\geq0}R_n+R_0<\tfrac14
    \qquad\text{and}\qquad
    R_v\ell_n<\frac{R_n}{10}\quad\text{for every }n .
  \]
\end{lemma}

The \texttt{exists\_ge} shape ($\forall M,\ \exists f,\ M\leq f.K_0 $) is what implements
the paper's dependency order: every later stage of the proof imposes one further
monotone lower bound on $K_0 $, and all of them are absorbed into a single maximum.  The
two displayed conditions are \texttt{eq:chart-packing} and
\texttt{eq:local-support-in-buffer}.  The single object that funnels every later bound
into one choice is \texttt{GlobalParameterChoice}, the node
\texttt{def:transfer:parameter-choice} of the transfer chapter.

\textbf{Deviation: explicit thresholds on $K_0 $.}  The paper asks for
\texttt{eq:chart-packing} ``for $K_0 $ sufficiently large'' and derives
\texttt{eq:local-support-in-buffer} from the explicit sufficient condition
$K_0>(10R_v)^2 $, via $R_v\ell_n/R_n=R_v/\sqrt{K_n}\leq R_v/\sqrt{K_0} $ (Section~3.3).  Lean
proves the same criterion in the form $10R_vR_0<1 $
(\texttt{ScaleFamily.supportScale\_lt\_tenth\_buffer}), which at $k_*=1 $ is exactly
$K_0>(10R_v)^2 $; the funnel
\texttt{ScaleFamily.exists\_ge\_chartPacking\_supportFit\_bufferRadius\_zero\_lt} behind
\texttt{exists\_largeK0\_chartPacking\_and\_supportFit} then chooses one explicit integer
$K_0>k_*/t^2 $ with
$t=\min\bigl(\tfrac14/(4/(1-q^{-1/2})+1),\ 1/(10(R_v+1)),\ \text{target}\bigr) $.  So the
packing threshold is the explicit $R_0\bigl(4/(1-q^{-1/2})+1\bigr)<\tfrac14 $ (via
\texttt{chartPacking\_of\_geometric\_bound}), and the support-fit threshold used
is $K_0>k_*(10(R_v+1))^2 $, slightly larger than the paper's $(10R_v)^2 $ (the shift by one
keeps the bound meaningful at $R_v=0 $).  Both are sufficient conditions of the paper's kind;
nothing downstream depends on the particular values.

\begin{proof}
  \leanok
  \uses{def:velocity:scale-family, lem:velocity:scale-arithmetic}
  Both conditions are controlled by $R_0=\sqrt{k_*/K_0} $ alone.  By the geometric bound
  $R_n\leq R_0(\sqrt{q^{-1}})^n $ the packing quantity is at most
  $R_0\bigl(4/(1-\sqrt{q^{-1}})+1\bigr) $, and $R_v\ell_n<R_n/10 $ reduces to
  $10R_vR_0<1 $ because $\ell_n/R_n=R_n\leq R_0 $.  Choosing $K_0 $ beyond the corresponding
  threshold, and beyond the prescribed $M $, gives both.
\end{proof}

\section{Disjoint placement and added axial velocity profiles on $\T^2$}

\begin{definition}[Placement of the cells in the chart and in the torus]
  \label{def:velocity:placement}
  \lean{FastDynamo.ScaleFamily.placementChart, FastDynamo.ScaleFamily.bufferDisk,
    FastDynamo.ScaleFamily.localCoreDisk,
    FastDynamo.ScaleFamily.TorusCellPlacement,
    FastDynamo.transverseTorusPoint,
    FastDynamo.ScaleFamily.TorusCellPlacement.torusBufferDisk,
    FastDynamo.ScaleFamily.TorusCellPlacement.torusLocalCoreDisk,
    FastDynamo.GlobalSeedInterface.placementRadius,
    FastDynamo.GlobalSeedInterface.profileSupportRadius_le_placementRadius,
    FastDynamo.GlobalSeedInterface.witnessSupportRadius_le_placementRadius,
    FastDynamo.ScaleFamily.TorusCellPlacement.restrictSupportRadius}
  \leanok
  \uses{def:velocity:scale-family, lem:velocity:base-index}
  All cells are placed in one fixed Euclidean coordinate square of $\T^2 $.  The $n $-th
  \emph{buffer disc} is $D(p_n,R_n) $ --- what the paper calls the $n $-th \emph{cell}
  --- and the $n $-th \emph{local-flow core} --- the paper's \emph{core} --- is the disc
  $D(p_n,R_v\ell_n) $ carrying the support of the rescaled reference profile, where $R_v $ is
  chosen with
  \[
    R_v\geq L_*
    \qquad\text{and}\qquad
    \supp v\cup\supp\nabla_\perp v\subset D(0,R_v)
  \]
  (Section~3.3 of the paper).  The corresponding subsets
  of $\T^3 $ are $D(p_n,R_n)\times\T $ and $D(p_n,R_v\ell_n)\times\T $.  A \emph{placement}
  is a scale family together with the two certificates of
  Lemma~\ref{lem:velocity:base-index}: chart packing and support fit.
\end{definition}

In Lean the chart is \texttt{ScaleFamily.placementChart}, the discs are
\texttt{bufferDisk} and \texttt{localCoreDisk} (Euclidean, in $\R\times\R $), and the
placement certificate is the structure \texttt{ScaleFamily.TorusCellPlacement
supportRadius}, whose two fields are exactly \texttt{chartPacking} and
\texttt{supportFit} (module \texttt{FastDynamo/Global/TorusCellPlacement.lean}).  The
quotient map $\R^2\to\T^2 $ is \texttt{transverseTorusPoint}
(\texttt{FastDynamo/Analysis/TransverseChart.lean}), and \texttt{torusBufferDisk},
\texttt{torusLocalCoreDisk} are the images of the Euclidean discs
(\texttt{FastDynamo/Global/TorusChartEmbedding.lean}).

Two differences from the paper.  First, the Lean chart is
$(0,\tfrac14)\times(\tfrac14,\tfrac34) $, the right half of the paper's
$(-\tfrac14,\tfrac14)\times(\tfrac14,\tfrac34) $; since every centre has positive first
coordinate and $p_\infty $ is the left edge, every disc still lies in it and nothing about
the accumulation changes.  Second, the concrete local profile is the smoothed Ponomarenko
velocity of the seed chapter, whose own support radius is $1+h $ for the smoothing width
$h $; nothing in the development bounds the witness radius $L_* $ above, so the paper's
$R_v $ is realised as
\[
  R_v=\texttt{GlobalSeedInterface.placementRadius}=\max(1+h,L_*),
\]
which is the smallest choice satisfying both requirements of the definition.  The two
inequalities are the Lean theorems
\texttt{GlobalSeedInterface.profileSupportRadius\_le\_placementRadius} and
\texttt{GlobalSeedInterface.witnessSupportRadius\_le\_placementRadius}, the second being
exactly $R_v\geq L_* $.  The placement is carried at this joint radius (the record field
\texttt{GlobalParameterChoice.jointPlacement}) and read back at $1+h $ for the velocity
constructors through \texttt{TorusCellPlacement.restrictSupportRadius}, which leaves the
scale family untouched; that reading is the derived accessor
\texttt{GlobalParameterChoice.placement}.  The cost of the joint radius is nothing: it is
one call of the funnel of Lemma~\ref{lem:velocity:base-index} instead of two, and it is what
makes \texttt{eq:scaled-witness-support} a consequence rather than a further threshold
(node \texttt{lem:transfer:witness}).

\begin{lemma}[The placed cells are disjoint, in the chart and in the torus]
  \label{lem:velocity:disjoint-placement}
  \lean{FastDynamo.ScaleFamily.bufferDisk_subset_placementChart,
    FastDynamo.ScaleFamily.disjoint_bufferDisk,
    FastDynamo.ScaleFamily.localCoreDisk_subset_bufferDisk,
    FastDynamo.transverseTorusPoint_injectiveOn_placementChart,
    FastDynamo.ScaleFamily.TorusCellPlacement.disjoint_torusBufferDisk,
    FastDynamo.ScaleFamily.TorusCellPlacement.torusLocalCoreDisk_subset_torusBufferDisk,
    FastDynamo.ScaleFamily.TorusCellPlacement.exists_matchingCell_with_torusPlacedCore}
  \leanok
  \uses{def:velocity:placement}
  For a placement: every $D(p_n,R_n) $ lies in the fixed chart; the discs $D(p_n,R_n) $ are
  pairwise disjoint; and $D(p_n,R_v\ell_n)\subset D(p_n,R_n) $ for every $n $.  The quotient
  map is injective on the chart, so the images in $\T^2 $ are again pairwise disjoint, with
  each core image inside its own buffer image.  Consequently at most one cell is active at
  each point of $\T^2 $.
\end{lemma}

\begin{proof}
  \leanok
  \uses{def:velocity:placement, lem:velocity:base-index}
  Chart containment: the first coordinate of a point of $D(p_n,R_n) $ lies between
  $4\sum_{m\geq n}R_m-R_n>0 $ and $4\sum_{m\geq0}R_m+R_0<\tfrac14 $, while its second
  coordinate differs from $\tfrac12 $ by less than $R_n\leq R_0<\tfrac14 $.  Disjointness:
  the centres are collinear with $\abs{p_n-p_{n+1}}=4R_n>R_n+R_{n+1} $, and summing adjacent
  gaps gives $\abs{p_n-p_m}>R_n+R_m $ for $n<m $.  The core inclusion is the support fit
  $R_v\ell_n<R_n/10 $.  Injectivity on the chart is uniqueness of representatives in
  $[0,1) $ for $\R/\Z $, applied coordinatewise; it transports both statements to the torus.
\end{proof}

\begin{lemma}[The added axial velocity and its cut-off gradient]
  \label{lem:velocity:axial-plateau}
  \lean{FastDynamo.ScaleFamily.TorusCellPlacement.axialPlateauBump,
    FastDynamo.ScaleFamily.TorusCellPlacement.axialPlateau_eq_one_of_norm_pairToWholePlane_le,
    FastDynamo.ScaleFamily.TorusCellPlacement.axialOffsetEuclidean,
    FastDynamo.ScaleFamily.TorusCellPlacement.axialOffsetEuclidean_eq_axialAmplitude_of_core,
    FastDynamo.ScaleFamily.TorusCellPlacement.axialOffsetEuclidean_eq_zero_of_not_mem_bufferDisk,
    FastDynamo.ScaleFamily.TorusCellPlacement.norm_axialOffsetEuclideanDerivative_le,
    FastDynamo.ScaleFamily.axialGradientScale,
    FastDynamo.ScaleFamily.axialGradientScale_formula,
    FastDynamo.ScaleFamily.axialGradientScale_le_zero}
  \leanok
  \uses{def:velocity:placement, lem:velocity:scale-arithmetic}
  Let $\vartheta\in\C^\infty_c([0,\infty)) $ equal one on $[0,3/4] $ with
  $\supp\vartheta\subset[0,1) $, and put
  \[
    u_n^\mathrm{ax}(y)=a_n\vartheta\Bigl(\frac{\abs{y-p_n}}{R_n}\Bigr)e_3 .
  \]
  Then $u_n^\mathrm{ax}=a_ne_3 $ on the \emph{constant axial region} $D(p_n,3R_n/4) $,
  $\supp u_n^\mathrm{ax}\subset D(p_n,R_n) $, the field is independent of $x_3 $ and
  divergence-free, and
  \[
    \norm{\nabla u_n^\mathrm{ax}}_\infty\leq\norm{\vartheta'}_\infty\frac{a_n}{R_n},
    \qquad
    \frac{a_n}{R_n}=\frac{\tau}{N}R_n=\frac{\tau}{2\pi\sqrt{k_*K_n}}
    \leq\frac{a_0}{R_0} .
  \]
  In particular the supremum
  $\eta_\mathrm{ax}=\sup_n\norm{\vartheta'}_\infty a_n/R_n $ is attained at $n=0 $ and tends to
  $0 $ as $K_0\to\infty $.
\end{lemma}

The Lean plateau is a \texttt{ContDiffBump} at the origin of the pair chart, equal to one
on the closed ball of radius $3/4 $ and vanishing outside the open unit ball;
$\norm{\vartheta'}_\infty $ is realised by its Lipschitz constant
\texttt{axialPlateauLipschitzConstant}, and the quantity $a_n/R_n $ is
\texttt{ScaleFamily.axialGradientScale}.  The closed expression
$\eta_\mathrm{ax}=\norm{\vartheta'}_\infty\tau/(2\pi\sqrt{k_*K_0}) $ is not a single Lean
lemma: what is proved is the monotonicity $a_n/R_n\leq a_0/R_0 $ together with
$a_0/R_0=(\tau/N)R_0 $ and $R_0=\sqrt{k_*/K_0} $, which is the same statement after
substitution.  The paper uses $\eta_\mathrm{ax} $ only as a gradient bound (in the proof of
Lemma~3.1 of the paper and in \texttt{eq:axial-shear-bound} of the proof of
Proposition~4.1 of the paper); its smallness for large $K_0 $ is noted in Section~3.3 but is not one
of the thresholds on $K_0 $, neither in the paper nor in
\texttt{GlobalParameterChoice}.  The axial taper is \emph{radial} --
$\vartheta(\abs{y-p_n}/R_n)$ in the Euclidean metric of the chart, not in the sup metric --
which is what makes the plateau $D(p_n,3R_n/4)$ a Euclidean disc and lets the support fit of
the placement be read off without losing a factor two.

\begin{proof}
  \leanok
  \uses{def:velocity:placement, lem:velocity:scale-arithmetic}
  The plateau statements are the defining properties of $\vartheta $, transported by
  $y\mapsto(y-p_n)/R_n $.  The chain rule produces the factor $R_n^{-1} $ in the gradient,
  and $a_n/R_n=(\tau/N)\ell_n/R_n=(\tau/N)R_n $ by
  Lemma~\ref{lem:velocity:scale-arithmetic}(iii); monotonicity in $n $ follows from
  $R_n\leq R_0 $.  Divergence-freeness holds because the only nonzero component is the third
  one and it does not depend on $x_3 $.
\end{proof}

\begin{definition}[The global velocity]
  \label{def:velocity:global-velocity}
  \lean{FastDynamo.ScaleFamily.TorusCellPlacement.parameterizedHelicalCellSource,
    FastDynamo.ScaleFamily.TorusCellPlacement.FullVelocityConcrete.parameterizedFullEuclideanCellRealSource,
    FastDynamo.ScaleFamily.TorusCellPlacement.FullVelocityConcrete.parameterizedFullGlobalEuclideanRealSource,
    FastDynamo.ScaleFamily.TorusCellPlacement.FullVelocityConcrete.parameterizedFullGlobalPeriodizedRealSource,
    FastDynamo.ScaleFamily.TorusCellPlacement.FullVelocityConcrete.parameterizedFullGlobalRealTorusSource,
    FastDynamo.ScaleFamily.TorusCellPlacement.FullVelocityConcrete.parameterizedFullGlobalW1Inf,
    FastDynamo.TransverseRealW1Inf,
    FastDynamo.literalVelocity}
  \leanok
  \uses{def:velocity:placement, lem:velocity:disjoint-placement, lem:velocity:axial-plateau,
    def:local:seed}
  With $u_n^\mathrm{loc}(y)=\ell_nv\bigl((y-p_n)/\ell_n\bigr) $ and $u_n^\mathrm{ax} $ as above,
  the \emph{global velocity} is
  \[
    u(y,x_3)=\sum_{n\geq0}\bigl(u_n^\mathrm{loc}(y)+u_n^\mathrm{ax}(y)\bigr).
  \]
  It is independent of $x_3 $, and by Lemma~\ref{lem:velocity:disjoint-placement} at most
  one summand is nonzero at each point, so the sum is defined pointwise without any
  convergence hypothesis.
\end{definition}

Lean builds this in three stages (modules
\texttt{FastDynamo/Global/TorusVelocityConcrete.lean},
\texttt{FastDynamo/Global/TorusAxialVelocityConcrete.lean},
\texttt{FastDynamo/Global/TorusFullVelocityConcrete.lean},
\texttt{FastDynamo/Global/TorusGlobalVelocityLipschitz.lean}):
\begin{itemize}
  \item the \emph{chart-side representative}
    \texttt{parameterizedFullGlobalEuclideanRealSource}, the literal pointwise sum
    $\sum_n(\re u_n^\mathrm{loc}+u_n^\mathrm{ax}) $ on $\R^2 $, supported in the placement
    chart;
  \item its \emph{$\Z^2 $-periodisation}
    \texttt{parameterizedFullGlobalPeriodizedRealSource}, which again has at most one
    nonzero translate at each point of the plane and descends to the torus-side
    representative \texttt{parameterizedFullGlobalRealTorusSource};
  \item the \emph{torus-side $\W^{1,\infty} $ class} \texttt{parameterizedFullGlobalW1Inf},
    an element of \texttt{TransverseRealW1Inf}: a real $\Lp^\infty $ field on $\T^2 $
    bundled with two $\Lp^\infty $ derivative fields subject to the Fourier
    weak-derivative identities
    $\widehat{\partial_ju}(m)=2\pi i\,m_j\,\widehat u(m) $.  Membership in
    $\W^{1,\infty} $ is thus carried by the type; the class is the sum of a pointwise
    disjoint-support limit (the helical part, whose cell derivatives are \emph{not}
    $\Lp^\infty $-summable) and an absolutely convergent $\Lp^\infty $ series
    (the axial part).
\end{itemize}
The axial shift is the parameter \texttt{tau}; the value used downstream is the chain's
\texttt{axialShift}.  The velocity of the main theorem is
\texttt{GlobalParameterChoice.velocity}, that is
\texttt{parameterizedFullGlobalW1Inf} with the helical
speeds pinned to the seed's own $(\Omega,U) $, so that the on-core identity of
Lemma~\ref{lem:velocity:matching-cell} produces the seed velocity $v $.  Finally
\texttt{literalVelocity} is the axially independent extension of the torus-side
representative to $\T^3 $, which is the field the paper writes as $u $.

\begin{lemma}[On a cell core the global field is a rescaled seed copy plus $a_ne_3 $]
  \label{lem:velocity:matching-cell}
  \lean{FastDynamo.ScaleFamily.TorusCellPlacement.FullVelocityConcrete.parameterizedFullGlobalTorusSource_eq_rescaledLocalCopy_add_axialShift_of_mem_core,
    FastDynamo.ScaleFamily.TorusCellPlacement.FullVelocityConcrete.parameterizedFullGlobalTorusFirstDerivative_eq_rescaledLocalCopy_of_mem_core,
    FastDynamo.ScaleFamily.TorusCellPlacement.FullVelocityConcrete.parameterizedFullGlobalRealTorusSource_eq_axialShift_of_mem_plateau_of_not_mem_core}
  \leanok
  \uses{def:velocity:global-velocity, lem:velocity:disjoint-placement,
    lem:velocity:axial-plateau, lem:velocity:scale-arithmetic}
  For every $n $ and every $y $ in the $n $-th local-flow core,
  \[
    u(y)=\ell_nv\Bigl(\frac{y-p_n}{\ell_n}\Bigr)+a_ne_3,
    \qquad
    \nabla_\perp u(y)=(\nabla_\perp v)\Bigl(\frac{y-p_n}{\ell_n}\Bigr),
  \]
  and for $y\in D(p_n,3R_n/4) $ outside the core, $u(y)=a_ne_3 $.  Since $2\pi K_na_n=\tau $,
  the added axial velocity acts on the fixed-$K_n $ Fourier subspace
  $B(y,x_3)=\e^{2\pi iK_nx_3}b(y) $ as the constant translation
  $-(a_ne_3\cdot\nabla)B=-i\tau B $, which changes no spectral real part.
\end{lemma}

The last two Lean names state the same identities almost everywhere for the stored
$\Lp^\infty $ classes at the selected matching cell, in the form the parametrix chapter
consumes (module \texttt{FastDynamo/Global/TorusSelectedCellChartAgreement.lean}); there
the local profile is exactly the seed's, because the helical speeds of the global
velocity are pinned to the seed's profile data (the fields \texttt{smoothingWidth},
\texttt{angularSpeed} and \texttt{axialSpeed} of \texttt{GlobalSeedInterface}).  The consequent operator identity
$U_n^{-1}A_nU_n=L_{\delta_n}-i\tau $ is the exact-renormalisation node
\texttt{lem:exact-renormalization} of the operator chapter (Lemma~5.2 of the paper,
``Scaling'', identity \texttt{eq:scaling}), which is where the paper carries out the
$-i\tau $ computation; it is not restated here.

\begin{proof}
  \leanok
  \uses{lem:velocity:disjoint-placement, lem:velocity:axial-plateau,
    lem:velocity:scale-arithmetic}
  By disjointness only the $n $-th summand is active on $D(p_n,R_n) $.  The support fit puts
  the core well inside $D(p_n,3R_n/4) $, where the axial cut-off is identically one, so the
  axial summand is the constant $a_ne_3 $ there and contributes nothing to the derivative.
  The helical summand is by definition $\ell_nv((y-p_n)/\ell_n) $, whose derivative is
  $(\nabla_\perp v)((y-p_n)/\ell_n) $ because the amplitude $\ell_n $ exactly cancels the
  dilation.  Outside the core the helical summand vanishes.  The spectral statement is the
  identity $2\pi K_na_n=\tau $ of Lemma~\ref{lem:velocity:scale-arithmetic}(ii).
\end{proof}

\section{Regularity and incompressibility of the velocity field on $\T^3$}

\begin{lemma}[Reality, $x_3 $-independence and incompressibility]
  \label{lem:velocity:divergence-free}
  \lean{FastDynamo.IsTransverseDivergenceFree,
    FastDynamo.ScaleFamily.TorusCellPlacement.FullVelocityConcrete.parameterizedFullGlobalTorusSource_pointwiseReal,
    FastDynamo.ScaleFamily.TorusCellPlacement.FullVelocityConcrete.parameterizedFullGlobalTorusFirstDerivative_transverseDivergence,
    FastDynamo.ScaleFamily.TorusCellPlacement.FullVelocityConcrete.parameterizedFullGlobalW1Inf_isTransverseDivergenceFree}
  \leanok
  \uses{def:velocity:global-velocity}
  The global velocity is real valued, independent of $x_3 $, and $\Div u=0 $: pointwise
  $\partial_1u_1+\partial_2u_2=0 $ for the displayed representatives, hence $\Div u=0 $ in
  the $\W^{1,\infty} $ graph sense.
\end{lemma}

In Lean, $x_3 $-independence is by type (the class lives on the transverse torus $\T^2 $),
reality is the vanishing of the imaginary part of every component of the complexified
representative, and $\Div u=0 $ is the predicate \texttt{IsTransverseDivergenceFree}: the
vanishing in $\Lp^\infty $ of the sum of the two stored diagonal derivative components.
Because the stored derivatives satisfy the Fourier weak-derivative identities, this is the
paper's distributional statement and not merely a pointwise one.

\begin{proof}
  \leanok
  \uses{def:velocity:global-velocity, lem:velocity:disjoint-placement,
    lem:velocity:axial-plateau}
  Each rescaled local copy is divergence-free because $\Div_\perp v_\perp=0 $ and dilation
  preserves that identity; each added axial velocity is divergence-free because it is
  parallel to $e_3 $ and independent of $x_3 $.  The supports are pairwise disjoint, so the
  identity for the sum is the identity for the unique active summand at each point.  Passing
  from the pointwise identity to the graph identity uses that the stored derivative fields
  are represented almost everywhere by the displayed pointwise sums.
\end{proof}

\begin{lemma}[Explicit uniform derivative bound and Lipschitz constant]
  \label{lem:velocity:w1inf-bound}
  \lean{FastDynamo.ScaleFamily.TorusCellPlacement.FullVelocityConcrete.norm_parameterizedFullEuclideanCellRealSource_le,
    FastDynamo.ScaleFamily.TorusCellPlacement.FullVelocityConcrete.globalVelocityCoordinateDerivativeBound,
    FastDynamo.ScaleFamily.TorusCellPlacement.FullVelocityConcrete.globalVelocityLipschitzBound,
    FastDynamo.ScaleFamily.TorusCellPlacement.FullVelocityConcrete.norm_fderiv_parameterizedFullGlobalEuclideanRealSource_le,
    FastDynamo.ScaleFamily.TorusCellPlacement.FullVelocityConcrete.lipschitzWith_parameterizedFullGlobalEuclideanRealSource,
    FastDynamo.ScaleFamily.TorusCellPlacement.FullVelocityConcrete.lipschitzWith_parameterizedFullGlobalRealTorusSource,
    FastDynamo.ScaleFamily.TorusCellPlacement.FullVelocityConcrete.lipschitzWith_literalVelocity_parameterizedFullGlobalRealTorusSource}
  \leanok
  \uses{def:velocity:global-velocity, lem:velocity:axial-plateau,
    lem:velocity:disjoint-placement}
  Each combined cell obeys
  $\norm{u_n^\mathrm{loc}+u_n^\mathrm{ax}}_\infty\leq(\norm v_\infty+\tau/N)\ell_n $, and there is
  an explicit constant
  \[
    L=2\Bigl(\norm{\nabla v}_\infty+\norm{\vartheta'}_\infty\frac{a_0}{R_0}\Bigr)
  \]
  with $\norm{Du(x)}\leq L $ at every point of the chart.  Consequently $u $ is Lipschitz
  with constant $2L $ on $\T^2 $ and, after the axial extension, on $\T^3 $.
\end{lemma}

This is the quantitative half of the paper's $\W^{1,\infty} $ claim, in slightly
different bookkeeping.  The proof of Lemma~3.1 of the paper bounds
$\norm{Du}\leq\norm{\nabla v}_\infty+\eta_\mathrm{ax} $ at every point; Lean proves
an everywhere bound on the chart carrying two extra factors of $2 $: one because the
operator norm on the pair chart is estimated by the sum of the two coordinate derivatives,
one because the periodisation estimate compares two competing integer translates.  The
role of $\norm{\nabla v}_\infty $ is played by the sum of the two first-order Schwartz
seminorms of the smoothed profile.  The constant is explicit but not claimed optimal, and
no later step uses its value.

\begin{proof}
  \leanok
  \uses{def:velocity:global-velocity, lem:velocity:axial-plateau,
    lem:velocity:disjoint-placement}
  Directly from the definitions, $\norm{u_n^\mathrm{loc}}_\infty\leq\ell_n\norm v_\infty $ and
  $\norm{u_n^\mathrm{ax}}_\infty\leq a_n=(\tau/N)\ell_n $.  For the derivative, the amplitude
  $\ell_n $ of the local copy exactly cancels its dilation, so
  $\norm{\nabla u_n^\mathrm{loc}}_\infty\leq\norm{\nabla v}_\infty $ uniformly in $n $, while
  $\norm{\nabla u_n^\mathrm{ax}}_\infty\leq\norm{\vartheta'}_\infty a_n/R_n
  \leq\norm{\vartheta'}_\infty a_0/R_0 $ by Lemma~\ref{lem:velocity:axial-plateau}.  Since
  the supports are pairwise disjoint, only one cell contributes at each point, which gives
  the bound for the whole sum; the mean value inequality turns it into a Lipschitz
  statement, and the descent lemma for $\Z^2 $-periodic Lipschitz maps transports it to
  $\T^2 $ and then to $\T^3 $.
\end{proof}

\begin{lemma}[Quadratic decay and differentiability at the accumulation circle]
  \label{lem:velocity:accumulation-decay}
  \lean{FastDynamo.ScaleFamily.TorusCellPlacement.FullVelocityConcrete.three_mul_bufferRadius_lt_dist_accumulationPoint_of_mem_bufferDisk,
    FastDynamo.ScaleFamily.TorusCellPlacement.FullVelocityConcrete.parameterizedFullGlobalEuclideanRealSource_eq_cell_of_mem_bufferDisk,
    FastDynamo.ScaleFamily.TorusCellPlacement.FullVelocityConcrete.parameterizedFullGlobalEuclideanRealSource_eq_zero_of_outside_buffers,
    FastDynamo.ScaleFamily.TorusCellPlacement.FullVelocityConcrete.norm_parameterizedFullGlobalEuclideanRealSource_le_dist_sq,
    FastDynamo.ScaleFamily.TorusCellPlacement.FullVelocityConcrete.parameterizedFullGlobalEuclideanRealSource_hasFDerivAt_accumulationPoint,
    FastDynamo.ScaleFamily.TorusCellPlacement.FullVelocityConcrete.parameterizedFullGlobalEuclideanRealSource_differentiableAt_accumulationPoint}
  \leanok
  \uses{def:velocity:global-velocity, lem:velocity:disjoint-placement,
    lem:velocity:scale-arithmetic}
  With $C=\norm v_\infty+\tau/N $,
  \[
    \abs{u(y,x_3)}\leq C\abs{y-p_\infty}^2 ,
  \]
  so $u $ vanishes at $p_\infty $ and is differentiable there with derivative zero.  The
  bound does not involve $x_3 $, so the same holds at every point of
  $\{p_\infty\}\times\T $.
\end{lemma}

The Lean statements are for the chart-side representative.  The only geometric inputs are
$\ell_n=R_n^2 $ and the separation $\abs{y-p_\infty}\geq3R_n $ for $y\in D(p_n,R_n) $.

\begin{proof}
  \leanok
  \uses{def:velocity:global-velocity, lem:velocity:disjoint-placement,
    lem:velocity:scale-arithmetic}
  If $y\in D(p_n,R_n) $ then
  $\abs{y-p_\infty}\geq\abs{p_n-p_\infty}-R_n=4\sum_{m\geq n}R_m-R_n\geq3R_n $, and there the
  field is the single $n $-th cell, of size at most $C\ell_n=CR_n^2 $; hence
  $\abs{u(y)}\leq C\abs{y-p_\infty}^2/9\leq C\abs{y-p_\infty}^2 $.  Outside all buffer discs
  the field vanishes, so the bound holds everywhere.  A field that is
  $O(\abs{y-p_\infty}^2) $ at $p_\infty $ is Fr\'echet differentiable there with derivative
  zero.
\end{proof}

\begin{lemma}[The global velocity is non-zero for a positive axial shift]
  \label{lem:velocity:nonvanishing}
  \lean{FastDynamo.ScaleFamily.TorusCellPlacement.FullVelocityConcrete.parameterizedFullGlobalTorusSource_ne_zero,
    FastDynamo.ScaleFamily.TorusCellPlacement.FullVelocityConcrete.parameterizedFullGlobalW1Inf_ne_zero}
  \leanok
  \uses{def:velocity:global-velocity, lem:velocity:matching-cell, lem:velocity:axial-plateau}
  If $\tau>0 $ then $u\neq0 $: on the nonempty set
  $D(p_n,3R_n/4)\setminus\supp u_n^\mathrm{loc} $ the field equals $a_ne_3\neq0 $.  The axial
  shift used in the construction is strictly positive, so the velocity of the main
  theorem is non-zero.
\end{lemma}

Because the torus-side object is an $\Lp^\infty $ equivalence class, the Lean proof needs
a set of positive measure rather than one point: it uses the explicit annular witness at
distance $R_n/2 $ from $p_n $ (outside the core by the support fit, inside the constant
axial plateau) together with a ball of radius $R_n/20 $ around it.  At $\tau=0 $ nothing
is proved; a helical-core witness would be needed, and the library contains none.
This is harmless: the shift used downstream is strictly positive.

\begin{proof}
  \leanok
  \uses{lem:velocity:matching-cell, lem:velocity:axial-plateau}
  Take the point at distance $R_n/2 $ from $p_n $.  The support fit $R_v\ell_n<R_n/10 $ puts
  it outside the local-flow core, and $R_n/2<3R_n/4 $ puts it inside the constant axial
  region, so by Lemma~\ref{lem:velocity:matching-cell} the field equals $a_ne_3 $ there,
  with $a_n=\tau\ell_n/(2\pi k_*)>0 $.  A ball of radius $R_n/20 $ around that point stays
  in the same region and has positive measure, which upgrades the pointwise statement to the
  $\Lp^\infty $ class.
\end{proof}

\begin{lemma}[Regularity and incompressibility of the velocity field on $\T^3$]
  \label{lem:global-velocity}
  \lean{FastDynamo.ScaleFamily.TorusCellPlacement.FullVelocityConcrete.parameterizedFullGlobalW1Inf,
    FastDynamo.ScaleFamily.TorusCellPlacement.FullVelocityConcrete.parameterizedFullGlobalW1Inf_isTransverseDivergenceFree,
    FastDynamo.ScaleFamily.TorusCellPlacement.FullVelocityConcrete.parameterizedFullGlobalW1Inf_ne_zero,
    FastDynamo.ScaleFamily.TorusCellPlacement.FullVelocityConcrete.lipschitzWith_literalVelocity_parameterizedFullGlobalRealTorusSource}
  \leanok
  \uses{def:velocity:global-velocity, lem:velocity:divergence-free, lem:velocity:w1inf-bound,
    lem:velocity:accumulation-decay, lem:velocity:nonvanishing}
  The field $u $ is a non-zero real-valued autonomous vector field belonging to
  $\W^{1,\infty}(\T^3;\R^3) $.  It is independent of $x_3 $, and $\Div u=0 $.
\end{lemma}

The paper's lemma ends with one further clause: \emph{$u $ is smooth away from
$\{p_\infty\}\times\T $ and is differentiable at every point of this circle, with derivative
zero.}  That clause is not part of the statement above, and it is not proved on $\T^3 $:
it is formalised for the chart-side representative and its periodisation in
Lemma~\ref{lem:main:velocity-regularity}.

Four remarks on the Lean form.  (a) Membership in $\W^{1,\infty} $ is carried by the type
\texttt{TransverseRealW1Inf}, so it is true by construction rather than proved; the
content is the construction itself --- that the infinite sum does define such an element
--- together with the explicit bound of Lemma~\ref{lem:velocity:w1inf-bound}.  The
$\W^{1,\infty} $ class lives on $\T^2 $ and reaches $\T^3 $ through the axially
independent extension \texttt{literalVelocity}, which is proved Lipschitz; there is no
separate $\W^{1,\infty}(\T^3) $ space in the development.  (b) Reality,
$x_3 $-independence and $\Div u=0 $ are Lemma~\ref{lem:velocity:divergence-free}.
(c) Non-vanishing holds for $\tau>0 $ (Lemma~\ref{lem:velocity:nonvanishing}).
(d) The smoothness and differentiability clauses are proved for the chart-side
representative and its $\Z^2 $-periodisation on the plane, not as smoothness in a manifold
structure on $\T^2 $ (the development installs none); they are collected, with the failure
of $\C^1 $, in the node \texttt{lem:main:velocity-regularity} and its Lean predicate
\texttt{GlobalVelocityQualitative}.

\begin{proof}
  \leanok
  \uses{lem:velocity:divergence-free, lem:velocity:w1inf-bound,
    lem:velocity:accumulation-decay, lem:velocity:nonvanishing,
    lem:velocity:disjoint-placement}
  Incompressibility and reality are Lemma~\ref{lem:velocity:divergence-free}, the uniform
  gradient bound is Lemma~\ref{lem:velocity:w1inf-bound}, non-vanishing is
  Lemma~\ref{lem:velocity:nonvanishing}, and differentiability with vanishing derivative on
  the accumulation circle is Lemma~\ref{lem:velocity:accumulation-decay}.  Away from
  $p_\infty $ the family of supports is locally finite, so the sum is locally a finite sum
  of smooth functions; with the previous sentence this makes $u $ differentiable at every
  point of $\T^3 $, with derivative bounded by $\norm{\nabla v}_\infty+\eta_\mathrm{ax} $
  (in Lean by the constant $L $ of Lemma~\ref{lem:velocity:w1inf-bound}).  Applying the mean value theorem to the periodic
  lift of $u $, we deduce that $u $ is globally Lipschitz with this bound.  Its classical
  divergence vanishes everywhere.  Since the weak derivative of a Lipschitz function agrees
  almost everywhere with its classical derivative, $u\in\W^{1,\infty} $ and $\Div u=0 $ in
  distributions.
\end{proof}

\textbf{Deviation: a Fourier dominated limit, not Rademacher.}  The route just narrated is
the paper's, and the two steps it rests on --- global
Lipschitz-ness from the mean value theorem on the periodic lift, and the vanishing of the
classical divergence --- are proved in Lean exactly that way
(\texttt{lipschitzWith\_parameterizedFullGlobalEuclideanRealSource}, obtained from
\texttt{norm\_\allowbreak fderiv\_\allowbreak parameterizedFullGlobalEuclideanRealSource\_\allowbreak le} through Mathlib's mean value inequality, and
\texttt{parameterizedFullGlobalTorusFirstDerivative\_transverseDivergence}, which is a
pointwise identity on the displayed representatives).  The last step is not: membership in
$\W^{1,\infty} $ is, in this development, membership in the closed Fourier graph
\texttt{TransverseRealW1Inf}, and no bridge \emph{Lipschitz $\Rightarrow$ Fourier graph}
--- a Rademacher theorem on $\T^2 $ together with an integration by parts for Lipschitz maps
--- exists in this library or in Mathlib in the required form.  Lean therefore builds the
graph element directly, by a route that needs neither a $\W^{1,1} $ summation of the cells
nor the paper's a.e.\ identification of weak and classical derivatives: the helical part is a dominated
limit of the Fourier coefficients of the partial sums (its cell derivatives are uniformly
bounded but \emph{not} $\Lp^\infty $-summable, so the graph norms cannot simply be added),
the axial part is an absolutely convergent $\Lp^\infty $ series, and the two graph elements
are added
(\texttt{integral\_parameterizedHelicalGlobal\_weakDerivative},
\texttt{parameterizedHelicalGlobalW1Inf}, \texttt{axialGlobalW1Inf} in
\texttt{FastDynamo/Global/TorusFullVelocityConcrete.lean}).  The conclusion --- that the
stored derivative fields are the Fourier weak derivatives of the stored value, and that
their diagonal sum vanishes --- is the paper's $u\in\W^{1,\infty} $, $\Div u=0 $.

\section{Why the construction stops at Lipschitz regularity}

The local field $v $ is smooth; regularity is lost only because the global field has
nonvanishing gradients at infinitely many shrinking horizontal scales.  The estimate of
Lemma~\ref{lem:velocity:accumulation-decay} shows that $u $ is nevertheless differentiable on
the accumulation circle, with $Du(p_\infty,x_3)=0 $ for every $x_3\in\T $.  At the other end,
both cut-offs are constant near each cell centre: the support fit puts $p_n $ well inside the
axial plateau $D(p_n,3R_n/4) $, and the rescaled local copy has amplitude $\ell_n $ and
dilation $\ell_n $, which cancel.  Hence, by Lemma~\ref{lem:velocity:matching-cell},
\[
  Du(p_n,x_3)=
  \begin{pmatrix}
    0      & -\Omega & 0 \\
    \Omega & 0       & 0 \\
    0      & 0       & 0
  \end{pmatrix}
  \neq0
  \qquad(x_3\in\T),
\]
that is, $Du(p_n,x_3)e_1=\Omega e_2 $, $Du(p_n,x_3)e_2=-\Omega e_1 $ and $Du(p_n,x_3)e_3=0 $;
it is non-zero because $\Omega\neq0 $.  It is one and the same map for every $n $, namely the
derivative at the origin of the unrescaled profile $v $.  Since $p_n\to p_\infty $, the
classical derivative $Du $ is discontinuous at every point of $\{p_\infty\}\times\T $.  A
continuous representative of the weak gradient would have to agree everywhere with $Du $, so
no such representative exists and
$u\notin\C^1 $, even though $u $ is differentiable everywhere (Section~3.4 of the paper).

These three statements --- differentiable everywhere, smooth off the accumulation point,
derivative discontinuous there --- are the content of the node
\texttt{lem:main:velocity-regularity} and of the Lean predicate
\texttt{GlobalVelocityQualitative} in
\texttt{FastDynamo/Global/TorusGlobalVelocityRegularity.lean}, and the Lean proof is the
argument just given, along the cell centres:
\texttt{fderiv\_\allowbreak parameterizedFullGlobalEuclideanRealSource\_\allowbreak cellCenter\_\allowbreak eq} gives one
and the same map at every centre,
\texttt{lineDeriv\_parameterizedSmoothedVelocity\_zero\_apply\_one} identifies its $(2,1) $
entry as the angular speed $\Omega $, whence
\texttt{fderiv\_\allowbreak parameterizedFullGlobalEuclideanRealSource\_\allowbreak cellCenter\_\allowbreak ne\_\allowbreak zero} under
$\Omega\neq0 $, and
\texttt{not\_\allowbreak continuousAt\_\allowbreak fderiv\_\allowbreak parameterizedFullGlobalEuclideanRealSource} concludes from
\texttt{tendsto\_cellCenter\_accumulationPoint} and
\texttt{fderiv\_\allowbreak parameterizedFullGlobalEuclideanRealSource\_\allowbreak accumulationPoint}, which says that
the derivative is $0 $ there.  Lean extracts only the entry $\Omega $ of the
matrix rather than the matrix itself, which is all the non-vanishing needs.  The global
failure of $\C^1 $ is the field \texttt{periodized\_not\_continuousAt\_fderiv} of the
strengthened main theorem (Theorem~\ref{thm:main-strengthened}): the derivative of the
$\Z^2$-periodised plane representative \texttt{mainVelocityPeriodizedRepresentative} is
discontinuous at the accumulation point
(\texttt{not\_\allowbreak continuousAt\_\allowbreak fderiv\_\allowbreak mainVelocityPeriodizedRepresentative}), a statement on
the plane, not on the manifold $\T^3 $ (Lemma~\ref{lem:main:velocity-regularity}).


\chapter{Particle dynamics, the ideal induction group and vanishing diffusivity}
This chapter follows Section~4 of the paper, \emph{Particle dynamics and absence of
exponential stretching} (\texttt{sec:particle}).  It has two halves: the
particle flow of the global velocity, written down in closed form, with its linear Lipschitz
bound and its zero topological entropy (Proposition~\ref{prop:particle-flow}); and the ideal
induction group built from the Cauchy formula, with its two-sided polynomial bounds and its
vanishing exponential growth rate (Corollary~\ref{cor:ideal-growth}).  Together they give
the two dynamical clauses of Theorem~\ref{thm:main}, and rather more
(Proposition~\ref{prop:main-lagrangian}).

The predicates are stated in
\texttt{FastDynamo/Dynamics/ParticleFlow.lean}; the flow is constructed in closed form, cell
by cell, in \texttt{FastDynamo/Dynamics/Flow/*}; the ideal group is constructed from the
Cauchy formula in \texttt{FastDynamo/Dynamics/Ideal/*}; and
\texttt{FastDynamo/Dynamics/MainLagrangianClauses.lean} assembles them at the constructed
velocity (\texttt{mainTheoremLagrangian\_holds}), which discharges the hypothesis of
\texttt{mainTheoremComplete\_of\_lagrangian} and of \texttt{ideal\_resistive\_gap}.  Almost
no ODE theory is used: Picard--Lindel\"of, integral curves and Rademacher are avoided
altogether by writing the flow down.  The steps of the paper's \emph{proof} of Corollary~4.2
(\texttt{cor:ideal-growth}) are formalised: the variational equation
\texttt{eq:flow-variational-equation} (Lemma~\ref{lem:dynamics:variational-equation}), the
weak form of the ideal equation and the fact that the Cauchy formula solves it
(Definition~\ref{def:dynamics:ideal-weak-form}, Lemma~\ref{lem:dynamics:ideal-weak-existence}),
and the uniqueness of weak solutions in $\C([-T,T];\Lp^2)$, which makes $\mathcal T_0$
\emph{the} solution group (Lemma~\ref{lem:dynamics:ideal-weak-uniqueness}).  The last section
is Remark~1.7 of the paper, the convergence of the resistive evolution to the ideal one as
$\eps\downarrow0$ (Theorem~\ref{thm:vanishing-diffusivity}).

\section{The particle flow and its entropy}

\begin{definition}[Particle flow and Lipschitz constants]
  \label{def:dynamics:particle-flow}
  \lean{FastDynamo.literalTorusPoint, FastDynamo.IsLiteralTorusLift,
    FastDynamo.IsParticleFlow,
    FastDynamo.ScaleFamily.TorusCellPlacement.literalFlow,
    FastDynamo.ScaleFamily.TorusCellPlacement.literalFlowLift,
    FastDynamo.mainLiteralFlow,
    FastDynamo.IsParticleFlow.unique,
    FastDynamo.ScaleFamily.TorusCellPlacement.literalFlow_unique,
    FastDynamo.mainLiteralFlow_unique,
    FastDynamo.existsUnique_isParticleFlow_mainLiteralVelocity}
  \leanok
  \uses{def:velocity:global-velocity, lem:global-velocity}
  The \emph{particle flow} of the global velocity $u $ is the family
  $\Phi_t:\T^3\to\T^3 $ solving
  \[
    \frac{\dd}{\dd t}\Phi_t(x)=u(\Phi_t(x)),
    \qquad
    \Phi_0(x)=x.
  \]
  Existence and uniqueness for every $t\in\R $ follow from the global Lipschitz regularity
  of $u $ (Lemma~\ref{lem:global-velocity}), and $\Phi_t^{-1}=\Phi_{-t} $.  For a
  Lipschitz map $F $, $\operatorname{Lip}(F) $ denotes its optimal Lipschitz constant.
\end{definition}

There is both a Lean statement and a Lean flow.  \texttt{IsParticleFlow v} $\Phi$ says
that for every $x$ there is a single curve $\gamma\colon\R\to\R^3$ with
$\Phi_t(x)=\pi(\gamma(t))$ and $\gamma'(t)=v(\Phi_t(x))$, together with $\Phi_0=\id$; the
inverse law $\Phi_{-t}\circ\Phi_t=\id$ (used in the proof of Proposition~4.1 of the paper
as $\Phi_t^{-1}=\Phi_{-t}$) is a separate conjunct of
\texttt{MainLagrangianClauses}, and the full group law $\Phi_s\circ\Phi_t=\Phi_{s+t}$ is proved
for the periodised data from which the flow is built
(\texttt{PeriodizedPlaneFlowData.literalFlow\_comp}) but is not part of any clause.  The paper's own
device of working with periodic lifts is
\texttt{IsLiteralTorusLift}.  The flow itself is
\texttt{ScaleFamily.TorusCellPlacement.literalFlow}
(\texttt{FastDynamo/Dynamics/Flow/TorusFlow.lean}), with periodic lift
\texttt{literalFlowLift}, and \texttt{mainLiteralFlow} is its instance at the constructed
velocity.  Its \emph{existence} is built \emph{in closed form}, not by an existence theorem:
the library never invokes Picard--Lindel\"of or integral curves, where the paper obtains
existence and uniqueness of the flow from the global Lipschitz bound of
Lemma~3.1 of the paper (opening paragraph of Section~4) and then derives the closed formula
\texttt{eq:explicit-particle-flow} for that flow.  The closed formula itself is the
paper's -- the transverse rotation of
\texttt{Global/PairChartRotation.lean}, the radial angular and axial profiles of
\texttt{Dynamics/Flow/CellProfiles.lean}, the single-cell map of
\texttt{Dynamics/Flow/CellFlowMap.lean}, the sum over the accumulating cells and its
$\Z^2$-periodisation in \texttt{Dynamics/Flow/PlaneFlow\{,Regularity\}.lean}, and the
descent through the chart in \texttt{Dynamics/Flow/TorusFlow.lean}.  That the closed
formula solves the ODE is Proposition~\ref{prop:particle-flow}.

\emph{Uniqueness} is proved, by Gronwall's inequality.  The paper's
``the unique particle flow'' (opening paragraph of Section~4) is
\texttt{IsParticleFlow.unique}
(\texttt{FastDynamo/Dynamics/Flow/ParticleFlowUniqueness.lean}): for a Lipschitz $v$ there is at
most one $\Phi$ with $\texttt{IsParticleFlow}\ v\ \Phi$.  The argument runs on the universal
cover, in the spirit of the paper's periodic lifts.  Two flows
give, through a fixed point $x$, two lifted trajectories $\gamma,\gamma'$ solving the same
\emph{autonomous} ODE $\dot r=\tilde v(r)$ with $\tilde v=v\circ\pi$, which is Lipschitz with
the same constant because $\pi$ is $1$-Lipschitz
(\texttt{lipschitzWith\_literalTorusPoint}).  Their starting points agree modulo $\Z^3$, so the
translate $\gamma'+(\gamma(0)-\gamma'(0))$ solves the same equation with the same initial value,
and Mathlib's \texttt{dist\_le\_of\_trajectories\_ODE} (Gronwall with $\delta=0$), applied
forwards and to the time-reversed system, identifies the two curves; projecting back gives
$\Phi=\Phi'$.  Hence \texttt{literalFlow} is \emph{the} flow at every admissible parameter
tuple (\texttt{literalFlow\_unique}), and \texttt{mainLiteralFlow} is the paper's
$\Phi_t$ at the constructed velocity (\texttt{mainLiteralFlow\_unique},
\texttt{existsUnique\_isParticleFlow\_mainLiteralVelocity}).  Nothing downstream needs it:
$\Phi_{-t}\circ\Phi_t=\id$ is a clause of \texttt{MainLagrangianClauses} in its own right, not a
consequence of uniqueness.

One deviation, forced by the geometry: the chart-side plane flow is \emph{not}
$\Z^2$-equivariant (its displacement is supported in the placement chart, which is
disjoint from its nontrivial integer translates), exactly as the velocity is not, so the
object that descends to the torus is the periodisation
\texttt{planePeriodizedTransverseFlow} built with the same
\texttt{UnitBoxPeriodization.periodization} the velocity already uses.  On the placement
chart the two agree, so no constant is lost.

\begin{definition}[Topological entropy]
  \label{def:dynamics:topological-entropy}
  \lean{FastDynamo.HasZeroTopologicalEntropy,
    FastDynamo.coverEntropy_eq_zero_of_lipschitz_iterate}
  \leanok
  Let $F $ be a continuous self-map of a compact metric space $(X,d) $, and write
  $h_{\mathrm{top}}(F) $ for its topological entropy.  The paper fixes no display for
  it: it refers to [Hasselblatt--Katok, \emph{Principal structures}, Section~2.5e]
  (Section~1.2 of the paper, before Theorem~1.2).  That section introduces $h_{\mathrm{top}} $ through the dynamical metrics
  $d_m(x,y)=\max_{0\leq j<m}d(F^j(x),F^j(y)) $ and the exponential growth rate in $m $ of
  the minimal cardinality of an $(m,\rho) $-cover, equivalently of the maximal cardinality
  of an $(m,\rho) $-separated set; on a compact metric space the cover, spanning and
  separated forms are one and the same invariant.
\end{definition}

\textbf{Which invariant.}  The Lean predicate \texttt{HasZeroTopologicalEntropy F} is
$\texttt{Dynamics.coverEntropy}\ F\ \texttt{univ}=0$, Mathlib's \emph{cover} entropy.  On
the compact metric space $\T^3$ the proof reduces it through
\texttt{Metric.uniformity\_basis\_dist} (inside
\texttt{coverEntropy\_eq\_zero\_of\_lipschitz\_iterate}) to the growth rate of the minimal
$(m,\rho)$-cover, so the quantity proved to vanish is the Bowen--Dinaburg invariant of the
cited section, not a weaker substitute.  Mathlib's own equivalence with the separated (net)
form, \texttt{Dynamics.coverEntropy\_eq\_iSup\_netEntropyEntourage}, exists in the
pinned Mathlib but is not imported here, because no clause needs it
(\texttt{Mathlib/Dynamics/TopologicalEntropy/NetEntropy.lean} is not built in this tree).

\textbf{The proof route differs.}  The paper (last paragraph of the proof of
Proposition~4.1) combines two facts quoted
from Hasselblatt--Katok: the power rule $h_{\mathrm{top}}(F^m)=m\,h_{\mathrm{top}}(F)$
(Section~2.5f(3)) and the Lipschitz bound
$h_{\mathrm{top}}(F)\leq3\log\operatorname{Lip}(F)$ (Proposition~2.5.5), applied to
$F=\Phi_s$, $F^m=\Phi_{ms}$ and $\operatorname{Lip}(\Phi_{ms})\leq C_{\rm L}(1+m\abs s)$.
Mathlib has neither fact -- there is no entropy power rule in
\texttt{Mathlib/Dynamics/TopologicalEntropy/} and no entropy bound for a Lipschitz self-map
-- so the development proves the \emph{combination} directly instead of its two factors.
\texttt{coverEntropy\_eq\_zero\_of\_lipschitz\_iterate}
(\texttt{FastDynamo/Analysis/EntropyCriterion.lean}) says that if every iterate $F^{[m]}$ is
$C(1+m)$-Lipschitz on $\T^3$ then $\texttt{coverEntropy}\,F\,\texttt{univ}=0$, and
\texttt{lipschitzWith\_literalFlow\_iterate} supplies that hypothesis from the group law
$\Phi_s^{[m]}=\Phi_{ms}$ and the Lipschitz bound of
Proposition~\ref{prop:particle-flow}.  Its proof is the classical Bowen--Dinaburg covering
argument that underlies the Lipschitz bound of Hasselblatt--Katok:
the $N^3$-point grid of \texttt{Analysis/TorusNet.lean} is a $1/(2N)$-net of $\T^3$, so it
is a dynamical $(m,\rho)$-cover for $N=\lceil C(1+m)/\rho\rceil+1$, the covering number
grows polynomially in $m$, and the exponential growth rate of a polynomial is zero.  The
input is the paper's (the linear Lipschitz bound and the group law); only the two quoted
theorems are bypassed, because formalising either of them would be strictly more work than
the combination.

\begin{proposition}[Explicitly solvable particle dynamics and zero entropy]
  \label{prop:particle-flow}
  \lean{FastDynamo.IsParticleFlow, FastDynamo.HasZeroTopologicalEntropy,
    FastDynamo.MainLagrangianClauses,
    FastDynamo.ScaleFamily.TorusCellPlacement.isParticleFlow_literalFlow,
    FastDynamo.ScaleFamily.TorusCellPlacement.forall_literalFlow_neg_literalFlow,
    FastDynamo.ScaleFamily.TorusCellPlacement.lipschitzWith_literalFlow,
    FastDynamo.ScaleFamily.TorusCellPlacement.exists_lipschitzWith_literalFlow,
    FastDynamo.ScaleFamily.TorusCellPlacement.coverEntropy_literalFlow_eq_zero,
    FastDynamo.ScaleFamily.TorusCellPlacement.forall_hasZeroTopologicalEntropy_literalFlow,
    FastDynamo.lipConst,
    FastDynamo.ScaleFamily.TorusCellPlacement.lipConst_literalFlow_le,
    FastDynamo.ScaleFamily.TorusCellPlacement.inv_flowLipschitzBound_le_lipConst_literalFlow,
    FastDynamo.ScaleFamily.TorusCellPlacement.tendsto_abs_inv_mul_log_lipConst_literalFlow,
    FastDynamo.tendsto_abs_inv_mul_log_lipConst_mainLiteralFlow,
    FastDynamo.ScaleFamily.TorusCellPlacement.literalFlow_eq_cell_of_mem_bufferDisk,
    FastDynamo.mainLiteralFlow_eq_cell_of_mem_bufferDisk,
    FastDynamo.ScaleFamily.TorusCellPlacement.literalFlow_eq_self_of_forall_not_mem_bufferDisk,
    FastDynamo.ScaleFamily.TorusCellPlacement.literalFlow_eq_self_of_accumulationPoint,
    FastDynamo.mainLiteralFlow_eq_self_of_accumulationPoint,
    FastDynamo.ScaleFamily.TorusCellPlacement.cellAngularSpeed_cellCenter,
    FastDynamo.ScaleFamily.TorusCellPlacement.cellAxialSpeed_cellCenter}
  \leanok
  \uses{def:dynamics:particle-flow, def:dynamics:topological-entropy, def:velocity:global-velocity, def:velocity:scale-family, def:velocity:placement}
  Let $u $ be the global velocity field.  There is a constant $C_{\rm L}\geq1 $ with
  \[
    \operatorname{Lip}(\Phi_t)+\operatorname{Lip}(\Phi_t^{-1})
    \leq C_{\rm L}(1+\abs t)
    \qquad(t\in\R).
  \]
  In particular $\lim_{\abs t\to\infty}\abs t^{-1}\log\operatorname{Lip}(\Phi_t)=0 $, and
  every time map has zero topological entropy:
  \[
    h_{\mathrm{top}}(\Phi_t)=0
    \qquad(t\in\R).
  \]

  The flow is moreover explicit.  With $\mathsf R_\alpha $ the rotation by the angle
  $\alpha $ in the horizontal plane and, for the $n $-th cell,
  \[
    \omega_n(r)=\Omega\chi_h\!\left(\frac r{\ell_n}\right),
    \qquad
    g_n(r)=\ell_nU\chi_h\!\left(\frac r{\ell_n}\right)
    +a_n\vartheta\!\left(\frac r{R_n}\right),
  \]
  one has, for $y\in D(p_n,R_n) $, $r=\abs{y-p_n} $ and $x_3\in\T $,
  \begin{equation}
    \Phi_t(y,x_3)
    =\left(p_n+\mathsf R_{t\omega_n(r)}(y-p_n),\ x_3+tg_n(r)\right),
    \label{eq:dynamics:explicit-flow}
  \end{equation}
  the second component modulo $1 $.  At $y=p_n $ the horizontal component is $p_n $ while
  the axial component translates at the speed $g_n(0)=\ell_nU+a_n $.  Every point outside
  the cells, on their boundaries, and of the accumulation circle
  $\{p_\infty\}\times\T $ is fixed.
\end{proposition}

Proved, in every assertion.  The ODE clause is
\texttt{isParticleFlow\_literalFlow}, the inverse clause
$\Phi_{-t}\circ\Phi_t=\id$ is \texttt{forall\_literalFlow\_neg\_literalFlow}, the Lipschitz
bound is \texttt{exists\_lipschitzWith\_literalFlow} with the explicit constant
$C_{\rm L}=\texttt{flowCL}=4(1+s)$, $s=\texttt{flowLipschitzSlope}$, the log-Lipschitz limit
is \texttt{tendsto\_abs\_inv\_mul\_log\_lipConst\_literalFlow}, and the entropy clause is
\texttt{coverEntropy\_literalFlow\_eq\_zero}.  All of them except the log-Lipschitz limit are
conjuncts of the Lean predicate \texttt{MainLagrangianClauses}, and
Proposition~\ref{prop:main-lagrangian} is that predicate at the constructed velocity; the
log-Lipschitz limit is a theorem beside it
(\texttt{tendsto\_abs\_inv\_mul\_log\_lipConst\_mainLiteralFlow}), because no clause of
Theorem~\ref{thm:main} consumes it.

Four statements are rendered differently from the paper.  First,
$\operatorname{Lip}$ is \texttt{FastDynamo.lipConst}, the infimum of the set of admissible
Lipschitz constants.  That this infimum is itself admissible is \emph{not} proved, and is not
needed: the upper bound is an \texttt{sInf} below an admissible constant and the lower bound
holds for every admissible one, so the squeeze goes through without it.  The bundle's own
Lipschitz clause quantifies existentially over admissible constants instead, which is the
same statement, the optimal constant being the least admissible one.  Second, $\abs t\to\infty$
is read as \texttt{Filter.cocompact}~$\R$, the standard Mathlib reading of that limit.
Third, \eqref{eq:dynamics:explicit-flow} is realised twice: as the \emph{definition} of
\texttt{literalFlow} -- the horizontal component is \texttt{planeRotate} at the angle
$t\,\omega_n(r)$ about the cell centre and the axial component the shift by $t\,g_n(r)$,
with the radius a constant of the motion
(\texttt{cellTransverseFlow\_mem\_bufferDisk}), so that the closed form is what the ODE
verification consumes -- and, as a theorem about the torus flow,
\texttt{literalFlow\_eq\_cell\_of\_mem\_bufferDisk}, whose hypotheses pin $y$ as a
fundamental-box representative of the horizontal block of $x$; that is what makes the
chart-side cell formula visible through the $\Z^2$-periodisation.  Fourth, the fixed-point
clause is stated translate-wise --
\texttt{literalFlow\_eq\_self\_of\_forall\_not\_mem\_bufferDisk} assumes that no integer
translate of the representative lies in a buffer disc -- because the torus flow is the
$\Z^2$-periodisation of the plane flow and the naive form would not suffice; points on the
boundaries are covered because \texttt{bufferDisk} is open, and the accumulation circle is
\texttt{literalFlow\_eq\_self\_of\_accumulationPoint}, no translate of
$p_\infty=(0,\tfrac12)$ lying in a buffer disc.  The two centre speeds are
\texttt{cellAngularSpeed\_cellCenter} ($\omega_n(0)=\Omega$) and
\texttt{cellAxialSpeed\_cellCenter} ($g_n(0)=\ell_nU+a_n$).

The uniform bound across the accumulating cells is carried out in
\texttt{Dynamics/Flow/PlaneFlowLipschitz.lean}, where a sum over cells is replaced by a
supremum: the buffer discs are disjoint and the flow is the identity off their union.  This
is not a deviation: the paper argues the same way in the proof of Proposition~4.1: the
natural lift of $\Phi_t$ is globally Lipschitz by
Lemma~\ref{lem:global-velocity}, its derivative is the cell formula on each disc and the
identity almost everywhere outside their union.

\begin{proof}
  \leanok
  \uses{def:dynamics:particle-flow, def:dynamics:topological-entropy, def:velocity:global-velocity, lem:global-velocity}
  On the $n $-th disc the velocity is
  $u(y,x_3)=(\omega_n(r)(y-p_n)^\perp,g_n(r)) $ with $r=\abs{y-p_n} $, so
  $\frac{\dd}{\dd t}\abs{y-p_n}^2=2\omega_n(r)(y-p_n)\cdot(y-p_n)^\perp=0 $: the radius is
  a constant of the motion, and the angular and axial coordinates move at the constant
  speeds $\omega_n(r) $ and $g_n(r) $.  This gives the closed formula, shows that no
  trajectory crosses a cell boundary, and matches the identity smoothly near
  $\partial D(p_n,R_n) $.  Differentiating the formula and using the two uniform shear
  bounds $\sup_n r\abs{\omega_n'(r)}\leq\abs\Omega\sup_s s\abs{\chi_h'(s)} $ and
  $\sup_n\abs{g_n'(r)}\leq\abs U\norm{\chi_h'}_\infty+\eta_{\rm ax} $ gives
  $\norm{D\Phi_t}\leq C(1+\abs t) $ on each cell with $C $ independent of $n $.  That the
  accumulation of the cells costs nothing is Lemma~\ref{lem:global-velocity}: the periodic
  lift of $\Phi_t $ is globally Lipschitz, its derivative is given by the cell formula on
  each disc and equals the identity almost everywhere outside their union, so the uniform
  cell bound is a bound on $\operatorname{Lip}(\Phi_t) $.  Applying the same estimate at
  $-t $ and using $\Phi_t^{-1}=\Phi_{-t} $ gives the two-sided bound after increasing the
  constant, and dividing its logarithm by $\abs t $ gives the vanishing log-Lipschitz rate.
  For the entropy, fix $s\in\R $ and an integer $m\geq1 $; the entropy formula for iterates
  and the Lipschitz entropy bound
  [Hasselblatt--Katok, Section~2.5f(3) and Proposition~2.5.5] give
  \[
    m\,h_{\mathrm{top}}(\Phi_s)
    =h_{\mathrm{top}}(\Phi_{ms})
    \leq3\log\bigl(C_{\rm L}(1+m\abs s)\bigr),
  \]
  and dividing by $m $ and letting $m\to\infty $ gives $h_{\mathrm{top}}(\Phi_s)=0 $.
\end{proof}

\section{The ideal induction group and the growth gap}

\begin{definition}[Ideal induction group]
  \label{def:dynamics:ideal-induction-group}
  \lean{FastDynamo.IsIdealInductionGroup, FastDynamo.HasVanishingIdealGrowthRate,
    FastDynamo.HasVanishingIdealOperatorGrowthRate,
    FastDynamo.ScaleFamily.TorusCellPlacement.idealTransport,
    FastDynamo.ScaleFamily.TorusCellPlacement.idealGroup,
    FastDynamo.mainIdealGroup}
  \leanok
  \uses{def:dynamics:particle-flow, def:velocity:global-velocity}
  The \emph{ideal induction equation} for the global velocity $u $ is
  \[
    \partial_tB+(u\cdot\nabla)B-(B\cdot\nabla)u=0,
    \qquad\Div B=0,
  \]
  that is, the induction equation of the introduction at $\eps=0 $.  Its solution group
  $\mathcal T_0(t) $ on $\Lp^2_\sigma(\T^3) $ is defined by the Cauchy formula
  \begin{equation}
    \bigl(\mathcal T_0(t)B_{\rm in}\bigr)(\Phi_t(a))
    =D\Phi_t(a)B_{\rm in}(a)
    \qquad\text{for a.e. }a,
    \label{eq:dynamics:ideal-cauchy}
  \end{equation}
  the magnetic field being transported and deformed by the particle flow.  Its
  exponential growth rate in operator norm is
  \begin{equation}
    \gamma_{\rm id}(u)
    :=\lim_{t\to\infty}\frac1t
    \log\norm{\mathcal T_0(t)}_{\Lp^2_\sigma\to\Lp^2_\sigma}.
    \label{eq:dynamics:ideal-rate}
  \end{equation}
\end{definition}

In the paper, $\mathcal T_0(t) $ is introduced in Section~1.2 as ``the solution group on
$\Lp^2_\sigma(\T^3) $ of the ideal induction equation''; the Cauchy formula is
\texttt{eq:ideal-cauchy-formula} of the \emph{proof} of Corollary~4.2, where it is derived from
the Lipschitz flow and shown to give that group, and $\gamma_{\rm id} $ is defined in
\texttt{eq:zero-ideal-growth} of the statement of Corollary~4.2.  The definition above takes
the Cauchy formula as the definition, which is the Lean model.

Stated \emph{and} constructed.  \texttt{IsIdealInductionGroup v} $\Phi$ $D\Phi$ $T$ bundles the
$\C_0$-group $T$ with the Cauchy representation \eqref{eq:dynamics:ideal-cauchy}, the
almost-everywhere Fr\'echet differentiability of a periodic lift of $\Phi_t$ (Rademacher),
the measure preservation of $\Phi_t$, and the two fields \texttt{solutionGroup} and
\texttt{unique}: $T$ is a solution group of the ideal equation of $v$ in the weak sense of
Definition~\ref{def:dynamics:ideal-weak-form}, and every solution group of that equation is
$T$.  Measure preservation replaces Liouville's formula, whose statement would need
$\det D\Phi_t$; with \texttt{unique} the group is determined by $v$ alone, which is the
paper's ``\emph{the} solution group'' of §1.2.

The group itself is built in two steps.  \texttt{idealTransport}
(\texttt{FastDynamo/Dynamics/Ideal/CauchyOperator.lean}) is
\eqref{eq:dynamics:ideal-cauchy} read as an operator on all of $\Lp^2(\T^3;\C^3)$, solved for
$\mathcal T_0(t)B$: the substitution $B\mapsto B\circ\Phi_{-t}$, which is a linear isometry of
$\Lp^2$ because $\Phi_{-t}$ preserves the volume, followed by pointwise multiplication by the
$\Lp^\infty$ matrix symbol $D\Phi_t(\Phi_{-t}\,\cdot)$ (the complexification
\texttt{flowFDerivC}, of norm at most $24(1+\abs t\,s)$ almost everywhere).  Then
\texttt{idealGroup} (\texttt{FastDynamo/Dynamics/Ideal/SolenoidalInvariance.lean})
corestricts it to $\Lp^2_\sigma$, which is the substantial step: see
Corollary~\ref{cor:ideal-growth}.  \texttt{mainIdealGroup} is the instance at the constructed
velocity.

The state space $\Lp^2_\sigma(\T^3) $ is \texttt{literalSolenoidalSubspace}
(\texttt{FastDynamo/Dynamics/Assembly/MainGrowthRate.lean}), which is \emph{defined} as the
image, under the isometric transport \texttt{toLiteralC3L2}, of the Fourier-defined solenoidal
subspace \texttt{productSolenoidalSubspace} of the spectral half.  Its identification with an
intrinsically defined subspace $\{B\in\Lp^2(\T^3;\C^3):\Div B=0\} $ of the literal torus
model is not formalised, no differential operator being transported to the literal torus.

The weak form of the ideal equation and the variational equation
\eqref{eq:dynamics:variational} (the paper's \texttt{eq:flow-variational-equation}), which the
paper uses in the proof of Corollary~4.2, are the nodes that follow.  Both need $\nabla u$ on
the literal torus, which is the derivative of the periodic lift of $u$ at every point
(Lemma~\ref{lem:dynamics:velocity-gradient}).  Remark~1.7 of the paper, the
vanishing-diffusivity limit, is Theorem~\ref{thm:vanishing-diffusivity} at the end of this
chapter.

\begin{definition}[The weak form of the ideal induction equation]
  \label{def:dynamics:ideal-weak-form}
  \lean{FastDynamo.literalFDeriv, FastDynamo.literalAdvection, FastDynamo.literalStretch,
    FastDynamo.literalFDeriv_literalTorusPoint, FastDynamo.literalAdvection_literalTorusPoint,
    FastDynamo.IsIdealTestField, FastDynamo.IsIdealWeakSolution,
    FastDynamo.IsIdealSolutionGroup, FastDynamo.IdealWeakIdentity}
  \leanok
  \uses{def:dynamics:ideal-induction-group}
  For a velocity $v$ on $\T^3$ whose periodic lift is differentiable, write $\nabla v$ for the
  derivative of the lift, $(v\cdot\nabla)\varphi$ and $(B\cdot\nabla)v=\nabla v\,B$.  A curve
  $W\in\C([-T,T];\Lp^2(\T^3;\CC^3))$ is a \emph{weak solution} of the ideal induction
  equation $\partial_tB+(v\cdot\nabla)B-(B\cdot\nabla)v=0$ on $[-T,T]$ if for every smooth
  periodic test field $\varphi$ and all $s,t\in[-T,T]$
  \[
    \inner{W(t)}{\varphi}-\inner{W(s)}{\varphi}
    =\int_s^t\Bigl(\inner{(W(\tau)\cdot\nabla)v}{\varphi}
      +\inner{W(\tau)}{(v\cdot\nabla)\varphi}\Bigr)\dd\tau .
  \]
  A \emph{solution group} is a one-parameter group $G(t)$ of bounded operators on
  $\Lp^2_\sigma(\T^3)$ all of whose orbits $t\mapsto G(t)B$ are weak solutions on every
  $[-T,T]$.
\end{definition}

The paper (proof of Corollary~4.2) understands solutions in the distributional sense in
$\C([-T,T];\Lp^2)$ and writes the weak form as
$\frac{\dd}{\dd t}\inner{B}{\varphi}=\inner{(B\cdot\nabla)u}{\varphi}+\inner{B}{(u\cdot\nabla)\varphi}$
for almost every $t$.  The Lean notion (\texttt{FastDynamo/Dynamics/IdealWeakForm.lean})
differs in these respects: the identity is tested against smooth spatial fields and integrated
in time (the equivalence with space--time test functions is not formalised); $\Div B=0$ is
not part of the notion of solution, only of the state space of a solution group; $\nabla v$
is the pointwise derivative of the lift; the integrals are Bochner integrals; and pairings are
$\int\langle\varphi,B\rangle$, conjugate-linear in the test field.  A solution group is
required to satisfy the group law, which the paper's ``solution group'' includes.

\begin{lemma}[The gradient of the constructed velocity]
  \label{lem:dynamics:velocity-gradient}
  \lean{FastDynamo.ScaleFamily.TorusCellPlacement.FullVelocityConcrete.fderiv_parameterizedFullGlobalPeriodizedRealSource_apply_pairBasis,
    FastDynamo.ScaleFamily.TorusCellPlacement.FullVelocityConcrete.transverseW1InfDeriv_parameterizedFullGlobalW1Inf_ae_eq_fderiv,
    FastDynamo.ScaleFamily.TorusCellPlacement.FullVelocityConcrete.fderiv_parameterizedFullGlobalPeriodizedRealSource_transverseDivergence,
    FastDynamo.ScaleFamily.TorusCellPlacement.hasFDerivAt_literalVelocity_lift,
    FastDynamo.ScaleFamily.TorusCellPlacement.norm_literalFDeriv_literalVelocity_le,
    FastDynamo.ScaleFamily.TorusCellPlacement.sum_literalFDeriv_literalVelocity_diag,
    FastDynamo.ScaleFamily.TorusCellPlacement.continuous_literalFDeriv_literalVelocity_literalFlow,
    FastDynamo.literalStretchL2,
    FastDynamo.ScaleFamily.TorusCellPlacement.literalStretchL2_coe_ae,
    FastDynamo.ScaleFamily.TorusCellPlacement.norm_literalStretchL2_le}
  \leanok
  \uses{def:velocity:global-velocity, lem:global-velocity, lem:main:velocity-regularity,
    def:dynamics:ideal-weak-form}
  The periodic lift of the global velocity $u$ is differentiable at every point of $\R^3$,
  with derivative bounded uniformly and divergence $\sum_i\partial_iu_i=0$ at every point; its
  transverse derivatives are, almost everywhere, the weak derivatives of $u$ as an element of
  $\W^{1,\infty}$; $t\mapsto\nabla u(\Phi_t(x))$ is continuous for every $x$; and
  $B\mapsto(B\cdot\nabla)u$ is a bounded operator on $\Lp^2(\T^3;\CC^3)$.
\end{lemma}

The paper's $\nabla u$ is the weak gradient of a $\W^{1,\infty}$ field and $\Div u=0$ holds
almost everywhere (Lemma~3.1 of the paper); here both hold at every point of the plane for the
closed-form lift, and the identification with the weak derivatives of the Fourier-graph field
\texttt{TransverseRealW1Inf} is the second item.  Modules
\texttt{FastDynamo/Global/TorusGlobalVelocityDerivativeBridge.lean} and
\texttt{FastDynamo/Dynamics/Ideal/VelocityGradient.lean}.

\begin{proof}
  \leanok
  \uses{lem:global-velocity, lem:main:velocity-regularity}
  Reduce a point to its representative in the unit box.  In a buffer disc the periodised field
  coincides near the point with the field of that cell, whose derivative is the derivative of
  the rescaled seed copy plus the axial offset; outside every buffer the derivative vanishes.
  The divergence vanishes cell by cell, for the helical part and for the axial offset.  The
  bound is the explicit derivative bound of the global velocity, and the almost-everywhere
  identification follows from the representative of the weak derivatives through the
  measure-preserving quotient map.
\end{proof}

\begin{lemma}[The variational equation]
  \label{lem:dynamics:variational-equation}
  \lean{FastDynamo.flowDeriv_apply_eq_add_integral,
    FastDynamo.ScaleFamily.TorusCellPlacement.flowFDeriv_eq_one_add_integral,
    FastDynamo.ScaleFamily.TorusCellPlacement.hasDerivAt_flowFDeriv,
    FastDynamo.ScaleFamily.TorusCellPlacement.flowFDeriv_apply_literalVelocity}
  \leanok
  \uses{def:dynamics:particle-flow, prop:particle-flow, lem:dynamics:velocity-gradient}
  For the particle flow $\Phi_t$ of the global velocity and its derivative $D\Phi_t$,
  \[
    D\Phi_t(a)=\id+\int_0^t(\nabla u)(\Phi_s(a))D\Phi_s(a)\dd s
  \]
  for every $a$ and every $t\in\R$; in particular $t\mapsto D\Phi_t(a)$ is differentiable with
  derivative $(\nabla u)(\Phi_t(a))D\Phi_t(a)$, and $D\Phi_t(a)u(a)=u(\Phi_t(a))$.
\end{lemma}

This is the paper's \texttt{eq:flow-variational-equation}, stronger: it holds at every point
and every time, where the paper has almost every $a$.  The route differs: the paper obtains it
from Cauchy--Lipschitz, Rademacher and the chain rule [Ambrosio--Crippa, Section~2]; here the
integral form of the flow equation is differentiated in the initial point under the integral
sign (dominated convergence and mean-value bounds), using the everywhere Fr\'echet derivative
of the closed-form lift.  The last identity, which the paper does not state, comes from
differentiating the group law of the lift.  Module
\texttt{FastDynamo/Dynamics/Flow/VariationalEquation.lean}.

\begin{proof}
  \leanok
  \uses{prop:particle-flow, lem:dynamics:velocity-gradient}
  Write $\Phi_t(a)=a+\int_0^tu(\Phi_s(a))\dd s$ on the lift.  The integrand is differentiable
  in $a$ with derivative $(\nabla u)(\Phi_s(a))D\Phi_s(a)$, bounded uniformly for $s$ in a
  compact interval, so differentiation under the integral sign gives the identity at every
  point.  The integrand is continuous in time
  (Lemma~\ref{lem:dynamics:velocity-gradient}), so the fundamental theorem of calculus gives
  the derivative.
\end{proof}

\begin{lemma}[The Cauchy formula solves the ideal equation]
  \label{lem:dynamics:ideal-weak-existence}
  \lean{FastDynamo.ScaleFamily.TorusCellPlacement.hasDerivAt_inner_literalFlow_flowFDerivC,
    FastDynamo.ScaleFamily.TorusCellPlacement.hasDerivAt_integral_inner_idealTransport,
    FastDynamo.ScaleFamily.TorusCellPlacement.isIdealWeakSolution_idealTransport,
    FastDynamo.ScaleFamily.TorusCellPlacement.isIdealSolutionGroup_idealGroup,
    FastDynamo.isIdealSolutionGroup_mainIdealGroup}
  \leanok
  \uses{def:dynamics:ideal-weak-form, def:dynamics:ideal-induction-group,
    lem:dynamics:variational-equation, lem:dynamics:velocity-gradient}
  For every $B\in\Lp^2(\T^3;\CC^3)$ and every $T$, the curve $t\mapsto\mathcal T_0(t)B$ given
  by the Cauchy formula \eqref{eq:dynamics:ideal-cauchy} is a weak solution of the ideal
  induction equation on $[-T,T]$, and
  \[
    \frac{\dd}{\dd t}\inner{\mathcal T_0(t)B}{\varphi}
    =\inner{(\mathcal T_0(t)B\cdot\nabla)u}{\varphi}
    +\inner{\mathcal T_0(t)B}{(u\cdot\nabla)\varphi}
  \]
  at every $t$.  Hence $\mathcal T_0$ is a solution group.
\end{lemma}

This is the paragraph ``Formula \texttt{eq:ideal-cauchy-formula} solves the ideal induction
equation'' of the paper's proof of Corollary~4.2, stronger: for every $B\in\Lp^2$ (not only
$\Lp^2_\sigma$), every $T$, and with the derivative at every time where the paper has almost
every $t$.  The route is the paper's: change of variables $x=\Phi_t(a)$, differentiation under
the integral, and the variational equation.  Module
\texttt{FastDynamo/Dynamics/Ideal/WeakSolution.lean}.

\begin{proof}
  \leanok
  \uses{lem:dynamics:variational-equation, lem:dynamics:velocity-gradient}
  By measure preservation and the Cauchy formula,
  $\inner{\mathcal T_0(t)B}{\varphi}=\int\langle\varphi(\Phi_t(a)),D\Phi_t(a)B(a)\rangle\dd a$.
  The integrand is differentiable in $t$ at every point, by the chain rule on the smooth lift
  of $\varphi$ and Lemma~\ref{lem:dynamics:variational-equation}, with a derivative bounded by
  $C\abs{B(a)}$ near each $t_0$; differentiation under the integral and the change of
  variables back give the displayed derivative, which is continuous in $t$.  The fundamental
  theorem of calculus gives the integrated identity.
\end{proof}

\begin{lemma}[Lipschitz test fields]
  \label{lem:dynamics:lipschitz-test-fields}
  \lean{FastDynamo.dense_idealTestFieldL2,
    FastDynamo.ScaleFamily.TorusCellPlacement.idealWeakIdentity_of_isIdealTestField,
    FastDynamo.IdealWeakIdentity.of_tendsto,
    FastDynamo.ScaleFamily.TorusCellPlacement.IsIdealWeakSolution.weakIdentity_of_lipschitz}
  \leanok
  \uses{def:dynamics:ideal-weak-form, lem:dynamics:velocity-gradient}
  The smooth periodic test fields are dense in $\Lp^2(\T^3;\CC^3)$.  If $W$ is a weak
  solution of the ideal equation of the global velocity on $[-T,T]$, the weak identity holds
  also for every test field $\varphi$ whose periodic lift is Lipschitz and differentiable at
  every point, with $\varphi,(u\cdot\nabla)\varphi\in\Lp^2$.
\end{lemma}

The paper extends identities from smooth to $\Hs^1$ test functions ``by density''.  Here the
extension is to Lipschitz, everywhere differentiable lifts, through Fourier partial sums:
integration by parts against the characters, Parseval convergence of the partial sums of
$\varphi$ and of $\partial_j\varphi$, and multiplication by the bounded components of $u$; no
mollifier is used.  The fields needed are $\psi\circ\Phi_{-\sigma}$ with $\psi$ smooth, which
are of this kind.  Module \texttt{FastDynamo/Dynamics/Ideal/LipschitzTestFields.lean}.

\begin{proof}
  \leanok
  \uses{lem:dynamics:velocity-gradient}
  Trigonometric polynomials are smooth and dense.  For a Lipschitz field with everywhere
  differentiable lift, the Fourier coefficients of $\partial_j\varphi$ are $2\pi im_j$ times
  those of $\varphi$ (integration by parts against characters), so the Fourier partial sums
  $\varphi_N$ converge to $\varphi$ and $(u\cdot\nabla)\varphi_N$ to $(u\cdot\nabla)\varphi$
  in $\Lp^2$.  The weak identity is linear and continuous in the test pair
  $(\varphi,(u\cdot\nabla)\varphi)$, so it passes to the limit.
\end{proof}

\begin{lemma}[Uniqueness of weak solutions; the solution group]
  \label{lem:dynamics:ideal-weak-uniqueness}
  \lean{FastDynamo.eq_zero_of_le_mul_abs_intervalIntegral,
    FastDynamo.ScaleFamily.TorusCellPlacement.inner_flowCompL2_eq_intervalIntegral,
    FastDynamo.ScaleFamily.TorusCellPlacement.norm_le_mul_abs_intervalIntegral,
    FastDynamo.ScaleFamily.TorusCellPlacement.IsIdealWeakSolution.eq_zero,
    FastDynamo.ScaleFamily.TorusCellPlacement.IsIdealSolutionGroup.eq_idealGroup,
    FastDynamo.ScaleFamily.TorusCellPlacement.existsUnique_isIdealSolutionGroup,
    FastDynamo.existsUnique_isIdealSolutionGroup_mainLiteralVelocity}
  \leanok
  \uses{def:dynamics:ideal-weak-form, lem:dynamics:ideal-weak-existence,
    lem:dynamics:lipschitz-test-fields, prop:particle-flow}
  If $W\in\C([-T,T];\Lp^2)$ is a weak solution of the ideal induction equation of the global
  velocity with $W(0)=0$, then $W=0$ on $[-T,T]$.  Consequently every solution group of that
  equation on $\Lp^2_\sigma(\T^3)$ is $\mathcal T_0$, and $\mathcal T_0$ is the unique solution
  group.
\end{lemma}

The first sentence is the paragraph ``Uniqueness follows from the same characteristic
argument'' of the paper's proof of Corollary~4.2.  The route differs: the paper obtains
$W(t,\Phi_t(a))=\int_0^t(\nabla u)(\Phi_s(a))W(s,\Phi_s(a))\dd s$ almost everywhere from
[Ambrosio--Crippa, Proposition~2.3] and applies Gronwall's inequality; here that identity is
proved in tested form,
$\inner{W(t)\circ\Phi_t}{\psi}=\int_0^t\inner{(W(\tau)\cdot\nabla)u}{\psi\circ\Phi_{-\tau}}\dd\tau$
for smooth $\psi$, directly from the weak form with the Lipschitz test fields
$\psi\circ\Phi_{-\sigma}$ (Lemma~\ref{lem:dynamics:lipschitz-test-fields}); density of the test
fields and the isometry $B\mapsto B\circ\Phi_t$ give
$\norm{W(t)}_2\leq\norm{S_u}\,\abs{\int_0^t\norm{W(\tau)}_2\dd\tau}$, with $S_u$ the bounded
operator $B\mapsto(B\cdot\nabla)u$ on $\Lp^2$, and Gronwall's inequality,
forward and backward in time, gives $W=0$.  Modules
\texttt{FastDynamo/Dynamics/Ideal/WeakSolutionUniqueness.lean} and
\texttt{FastDynamo/Dynamics/MainLagrangianClauses.lean}.

\begin{proof}
  \leanok
  \uses{lem:dynamics:lipschitz-test-fields, lem:dynamics:ideal-weak-existence}
  Test the weak identity of $W$ at $\psi\circ\Phi_{-\sigma}$ and differentiate the diagonal
  $\sigma\mapsto\inner{W(\sigma)}{\psi\circ\Phi_{-\sigma}}$; the transport terms cancel because
  $\frac{\dd}{\dd\sigma}\psi\circ\Phi_{-\sigma}=-(u\cdot\nabla)\psi\circ\Phi_{-\sigma}$, and what
  remains is the stretching term.  Integrating gives the tested characteristic identity; taking
  the supremum over unit $\psi$ gives the norm bound, and Gronwall's inequality gives $W=0$.
  For two solution groups $G$ and $\mathcal T_0$, the difference of their orbits from the same
  $B$ is a weak solution vanishing at $0$, so $G=\mathcal T_0$; existence is
  Lemma~\ref{lem:dynamics:ideal-weak-existence}.
\end{proof}

\begin{corollary}[No exponential growth for ideal induction]
  \label{cor:ideal-growth}
  \lean{FastDynamo.IsIdealInductionGroup, FastDynamo.MainLagrangianClauses,
    FastDynamo.HasVanishingIdealGrowthRate,
    FastDynamo.HasVanishingIdealOperatorGrowthRate,
    FastDynamo.MainLagrangianClauses.hasVanishingIdealGrowthRate,
    FastDynamo.MainLagrangianClauses.hasVanishingIdealOperatorGrowthRate,
    FastDynamo.inner_gradient_eq_zero_of_mem_literalSolenoidalSubspace,
    FastDynamo.ScaleFamily.TorusCellPlacement.isIdealInductionGroup_idealGroup,
    FastDynamo.ScaleFamily.TorusCellPlacement.norm_idealGroup_bounds,
    FastDynamo.ScaleFamily.TorusCellPlacement.tendsto_inv_mul_log_norm_idealGroup,
    FastDynamo.ScaleFamily.TorusCellPlacement.existsIdealGroupClauses,
    FastDynamo.exists_ne_zero_mainIdealGroup_rate,
    FastDynamo.existsUnique_isIdealSolutionGroup_mainLiteralVelocity}
  \leanok
  \uses{def:dynamics:ideal-induction-group, def:dynamics:particle-flow, def:velocity:global-velocity,
    def:dynamics:ideal-weak-form}
  Let $\mathcal T_0(t) $ be the solution group on $\Lp^2_\sigma(\T^3) $ of the ideal
  induction equation for the global velocity $u $ (a strongly continuous group of bounded
  operators, given by the Cauchy formula of
  Definition~\ref{def:dynamics:ideal-induction-group}).  Then, with $C_{\rm L} $ the
  constant of Proposition~\ref{prop:particle-flow},
  \[
    \frac1{C_{\rm L}(1+\abs t)}
    \leq\norm{\mathcal T_0(t)}_{\Lp^2_\sigma\to\Lp^2_\sigma}
    \leq C_{\rm L}(1+\abs t)
    \qquad(t\in\R).
  \]
  Consequently $\gamma_{\rm id}(u)=0 $.  Moreover, for every
  $B_{\rm in}\in\Lp^2_\sigma(\T^3) $ and $t\in\R $,
  \[
    \frac{\norm{B_{\rm in}}_2}{C_{\rm L}(1+\abs t)}
    \leq\norm{\mathcal T_0(t)B_{\rm in}}_2
    \leq C_{\rm L}(1+\abs t)\norm{B_{\rm in}}_2 ;
  \]
  in particular every non-zero ideal solution has exponential growth rate zero.
\end{corollary}

Every bound and rate of the statement is proved, for $\mathcal T_0$ the unique solution group
of the ideal equation in the weak sense of Definition~\ref{def:dynamics:ideal-weak-form}: it is
the group built from the Cauchy formula, it solves the equation
(Lemma~\ref{lem:dynamics:ideal-weak-existence}), and every solution group equals it
(Lemma~\ref{lem:dynamics:ideal-weak-uniqueness}); the fields \texttt{solutionGroup} and
\texttt{unique} of \texttt{IsIdealInductionGroup} record both, and
\texttt{existsUnique\_isIdealSolutionGroup\_mainLiteralVelocity} states them as one existence
and uniqueness statement.  The two-sided bounds are proved in the
\emph{per-datum} form, the display after \texttt{eq:zero-ideal-growth} in the statement,
$\norm B/(C_{\rm L}(1+\abs t))\leq\norm{T_tB}\leq C_{\rm L}(1+\abs t)\norm B$
-- \texttt{norm\_idealGroup\_bounds}, a conjunct of \texttt{MainLagrangianClauses} -- from
which the operator-norm form \texttt{eq:ideal-group-growth} follows by taking the supremum
over unit data, which is how the paper's proof obtains it too.  Both vanishing rates follow
-- the \emph{per-datum} one
$t^{-1}\log(\norm{T_tB}_2/\norm B_2)\to0$ for every $B\neq0$
(\texttt{tendsto\_inv\_mul\_log\_norm\_idealGroup}, packaged as
\texttt{HasVanishingIdealGrowthRate}, the last sentence of the statement)
and the \emph{operator-norm} one $t^{-1}\log\norm{\mathcal T_0(t)}\to0$
(\texttt{HasVanishingIdealOperatorGrowthRate}, which is both \texttt{eq:zero-ideal-growth}
here and the second dynamical clause of Theorem~1.2 of the paper (\texttt{thm:main}),
obtained from the bounds by
\texttt{MainLagrangianClauses.hasVanishingIdealOperatorGrowthRate}, which is the paper's
own argument: taking the supremum over unit initial data gives two-sided polynomial bounds
on the ideal operator norm, and the rate is squeezed between the logarithms of two linear
functions -- the squeeze itself is the reusable
\texttt{tendsto\_inv\_mul\_log\_opNorm\_of\_two\_sided\_bounds}).
\texttt{existsIdealGroupClauses} is the whole last conjunct of
\texttt{MainLagrangianClauses} in one statement, at the single constant
$C_{\rm L}=\texttt{dynamicalCL}=24(1+s)$.

None of these clauses is vacuous.  Every one of them -- the Cauchy formula, the group laws,
the strong continuity, the two-sided bounds and the per-datum rate -- is quantified over
$B\in\Lp^2_\sigma(\T^3) $, and would hold word for word if that space were $0 $; it is not,
by \texttt{exists\_ne\_zero\_literalSolenoidalSubspace}
(Theorem~\ref{thm:dynamics:solenoidal-semigroup}), whose witness is the spectral half's own
unstable eigenfunction.  \texttt{exists\_ne\_zero\_mainIdealGroup\_rate}
(\texttt{FastDynamo/Dynamics/MainLagrangianClauses.lean}) records the conclusion next to the
closing statements: there is a $B\neq0 $ in the state space at which
$t^{-1}\log(\norm{\mathcal T_0(t)B}/\norm B)\to0 $.

\textbf{Deviation: the solution notion.}  The paper's proof of Corollary~4.2 shows that the
Cauchy formula solves the ideal equation in weak form and that a distributional solution
$W\in\C([-T,T];\Lp^2)$ with $W(0)=0$ vanishes, so that $\mathcal T_0$ is the solution group.
Lean proves the same two facts for the weak notion of
Definition~\ref{def:dynamics:ideal-weak-form}, whose differences from the paper's
distributional notion (spatial smooth test fields and time-integrated identity, Bochner
integrals, pointwise $\nabla u$) are recorded there; the uniqueness argument differs from the
paper's (Lemma~\ref{lem:dynamics:ideal-weak-uniqueness}).

The substantial step of the construction is the one the paper's proof settles in one display
(testing against $\nabla\psi$, moving $\Phi_t$ onto the test function, and extending to
$\psi\circ\Phi_t\in\Hs^1$ by density): that $\Lp^2_\sigma$ is \emph{invariant}.  Read on the
literal torus it is the orthogonality
$\int_{\T^3}\overline{\nabla w}\cdot B=0$ for every solenoidal $B$ and every $C^1$ periodic
test function $w$ (\texttt{inner\_gradient\_eq\_zero\_of\_mem\_literalSolenoidalSubspace},
\texttt{FastDynamo/Dynamics/Ideal/GradientOrthogonality.lean}), proved by integration by parts
against the characters $e_m$ and Parseval.  Applying it at $w=e_m\circ\Phi_t$ needs a $C^1$
lift of $\Phi_t$, which the accumulation circle forbids, so the invariance is proved first for
the \emph{truncated} flows $\Phi^{(N)}$ -- finitely many cells, hence a $C^\infty$ lift --
and then transferred by $\Lp^2$-convergence
$\mathcal T_0^{(N)}(t)B\to\mathcal T_0(t)B$ and closedness of $\Lp^2_\sigma$
(\texttt{FastDynamo/Dynamics/Ideal/Truncated*.lean}).

\begin{proof}
  \leanok
  \uses{prop:particle-flow, def:dynamics:ideal-induction-group, def:dynamics:particle-flow,
    lem:dynamics:variational-equation, lem:dynamics:ideal-weak-existence,
    lem:dynamics:ideal-weak-uniqueness}
  Work with the periodic lift of $u$ and the corresponding lift of $\Phi_t$.  Cauchy--Lipschitz, Rademacher's theorem and the chain rule
  give the variational equation
  \begin{equation}
    D\Phi_t(a)=\id+\int_0^t(\nabla u)(\Phi_s(a))D\Phi_s(a)\dd s
    \qquad\text{for almost every }a
    \label{eq:dynamics:variational}
  \end{equation}
  [Ambrosio--Crippa, Section 2]; Liouville's formula and $\Div u=0 $ give
  $\det D\Phi_t=1 $, so the bi-Lipschitz map $\Phi_t $ preserves Lebesgue measure by the
  area formula.  Measure preservation turns the Cauchy formula into a bounded operator
  with $\norm{\mathcal T_0(t)B_{\rm in}}_2\leq\operatorname{Lip}(\Phi_t)\norm{B_{\rm in}}_2 $;
  testing against $\nabla\psi $ and moving $\Phi_t $ onto the test function shows the
  divergence constraint is preserved; differentiating
  $\inner{\mathcal T_0(t)B_{\rm in}}{\varphi} $ and using the variational equation gives
  the weak form of the ideal equation.  Uniqueness follows from the same characteristic
  representation plus Gronwall's inequality forward and backward in time, and the chain
  rule $D\Phi_{t+s}(a)=D\Phi_t(\Phi_s(a))D\Phi_s(a) $ gives the group law
  $\mathcal T_0(t+s)=\mathcal T_0(t)\mathcal T_0(s) $, $\mathcal T_0(-t)=\mathcal T_0(t)^{-1} $.
  Strong continuity follows from $\norm{\Phi_t-\id}_\infty\leq\norm u_\infty\abs t $ and
  $\norm{D\Phi_t-\id}_\infty\leq\e^{\norm{\nabla u}_\infty\abs t}-1 $ on a dense set of
  smooth fields, together with the local operator bound.  The Cauchy formula and
  Proposition~\ref{prop:particle-flow} give the per-datum upper bound
  $\norm{\mathcal T_0(t)B_{\rm in}}_2\leq C_{\rm L}(1+\abs t)\norm{B_{\rm in}}_2 $, and
  $\mathcal T_0(-t)\mathcal T_0(t)=\id $ gives the matching per-datum lower bound
  $\norm{B_{\rm in}}_2\leq C_{\rm L}(1+\abs t)\norm{\mathcal T_0(t)B_{\rm in}}_2 $.  Taking
  the supremum over unit initial data turns them into the displayed operator bounds, and
  taking logarithms and dividing by $t $ gives $\gamma_{\rm id}(u)=0 $ and the per-datum
  rate.
\end{proof}

\textbf{How the Lean proof differs from this one.}  Several steps are replaced by cheaper ones
that the closed form makes available.  Cauchy--Lipschitz is not used at all (the flow is
written down); Rademacher is not used (the lift is differentiable at every point, so the
almost-everywhere clause is discharged by \texttt{Filter.Eventually.of\_forall}), and the
variational equation is proved at every point by differentiating the integral form of the flow
equation (Lemma~\ref{lem:dynamics:variational-equation}); and Liouville's formula is replaced
by a pointwise Jacobian computation on the plane, $\det D\Phi_t=1$ at every point from the
closed formula, followed by the change of variables and a skew-product argument for the axial
shear.  The group law comes from the group law of the flow, and strong continuity is proved at
$t=0$ and moved along the group, by density of \emph{continuous} data and dominated convergence
for the symbol $D\Phi_t(a) $ at fixed $a $
(\texttt{FastDynamo/Dynamics/Ideal/StrongContinuity.lean}), instead of the paper's smooth data
and the bounds $\norm{\Phi_t-\id}_\infty\leq\norm u_\infty\abs t $,
$\norm{D\Phi_t-\id}_{\Lp^\infty}\leq\e^{L\abs t}-1 $.  The weak form is derived as in the paper
(Lemma~\ref{lem:dynamics:ideal-weak-existence}); uniqueness is proved from the weak form in
tested form rather than through [Ambrosio--Crippa, Proposition~2.3]
(Lemma~\ref{lem:dynamics:ideal-weak-uniqueness}).  The upper bound uses the bound on the
symbol $D\Phi_t $ rather than $\operatorname{Lip}(\Phi_t) $, and the lower bound is
$1\leq\norm{\mathcal T_0(-t)}\norm{\mathcal T_0(t)}$ exactly as here.  The last step is the
paper's: the per-datum bounds come first, and the operator norm follows by the supremum over
unit data.

\begin{proposition}[Particle flow and ideal induction]
  \label{prop:main-lagrangian}
  \lean{FastDynamo.MainLagrangianClauses, FastDynamo.mainLiteralVelocity,
    FastDynamo.MainTheoremLagrangian,
    FastDynamo.ScaleFamily.TorusCellPlacement.mainLagrangianClauses_literalVelocity,
    FastDynamo.mainTheoremLagrangian_holds}
  \leanok
  \uses{thm:main, def:velocity:global-velocity, def:dynamics:particle-flow, def:dynamics:topological-entropy, def:dynamics:ideal-induction-group}
  The velocity field $u $ of Theorem~\ref{thm:main} can be chosen so that there is
  $C_{\rm L}\geq1 $ with
  \[
    \operatorname{Lip}(\Phi_t)+\operatorname{Lip}(\Phi_t^{-1})
    \leq C_{\rm L}(1+\abs t),
    \qquad
    h_{\mathrm{top}}(\Phi_t)=0
    \qquad(t\in\R),
  \]
  and so that the ideal induction equation defines a strongly continuous solution group
  $\mathcal T_0(t) $ on $\Lp^2_\sigma(\T^3) $ with
  \[
    \lim_{t\to\infty}\frac1t
    \log\norm{\mathcal T_0(t)}_{\Lp^2_\sigma\to\Lp^2_\sigma}=0.
  \]
\end{proposition}

Proved.  The paper has no proposition of this shape: what it has is
Proposition~\ref{prop:particle-flow}, Corollary~\ref{cor:ideal-growth}, and the last
sentence of Theorem~1.2 of the paper, which quotes them.  The node exists because it is the
packaging of those statements \emph{at the velocity that
Theorem~\ref{thm:main} produces}, not at an arbitrary one, and because that packaging is
exactly the Lean predicate
\texttt{MainTheoremLagrangian}~$:=$~\texttt{MainLagrangianClauses mainLiteralVelocity}.

The paper makes the difference between this bundle and the theorem's own dynamical half
visible, and the Lean statement layer mirrors it.
\texttt{MainTheoremDynamicalClauses}~$v$ is exactly the two dynamical clauses of
Theorem~1.2 of the paper: an
ideal induction group for $v$ together with $h_{\mathrm{top}}(\Phi_t)=0$ at every $t\in\R$
and $t^{-1}\log\norm{\mathcal T_0(t)}\to0$, the entropy clause carrying the continuity of
every $\Phi_t$ with it because the paper defines $h_{\mathrm{top}}$ for a continuous self-map
of a compact metric space (Section~1.2 of the paper); \texttt{MainTheoremDynamical} is that
predicate
at \texttt{mainLiteralVelocity}, discharged by
\texttt{mainTheoremDynamical\_holds}.  It is a projection of this node --
\texttt{MainLagrangianClauses.toMainTheoremDynamicalClauses} -- and the bundle keeps what
the theorem drops: the inverse law $\Phi_{-t}\circ\Phi_t=\id$, the linear Lipschitz
bound of Proposition~\ref{prop:particle-flow}, and the two-sided per-datum bounds of
Corollary~\ref{cor:ideal-growth}.  \texttt{MainTheoremComplete} therefore carries both, as
the fields \texttt{dynamical} and \texttt{lagrangian}.

\texttt{mainLagrangianClauses\_literalVelocity}
(\texttt{FastDynamo/Dynamics/MainLagrangianClauses.lean}) proves the predicate at the general
parameters $(h,P,\Omega,U,\tau)$ with $\tau\geq0$, and \texttt{mainTheoremLagrangian\_holds}
instantiates it at the main parameter choice, where $\tau\geq0$ is
\texttt{mainParameterChoice.tau\_nonneg}.  The single constant shared by the Lipschitz
clause and the ideal bounds is
$C_{\rm L}=\texttt{dynamicalCL}=6\,\texttt{flowCL}=24(1+s)$ with
$s=\texttt{flowLipschitzSlope}$, the sum of the horizontal and axial shear bounds of the cell
profiles; the factor $6$ is what the passage from the sup metric of $\T^3$ to the Euclidean
metric of $\R^3$ and the complexification of $D\Phi_t$ cost.  The normalisation remark of the
paper (the unlabelled Remark~1.6) notes that the conclusion is invariant under $u\mapsto u/M $,
since the flow and the ideal group of the rescaled field are those of $u $ evaluated at time
$t/M $; that is Lemma~\ref{lem:main:normalisation}, where the bundle is proved for $u/M$ with
the constant $(1+M^{-1})C_{\rm L}$.

\begin{proof}
  \leanok
  \uses{prop:particle-flow, cor:ideal-growth, def:velocity:global-velocity,
    lem:global-velocity}
  The velocity field constructed for Theorem~\ref{thm:main} is the global velocity of
  Lemma~\ref{lem:global-velocity}, so Proposition~\ref{prop:particle-flow} gives the
  Lipschitz bound and the vanishing of the topological entropy of every time map, and
  Corollary~\ref{cor:ideal-growth} gives the solution group and its zero exponential growth
  rate.
\end{proof}

\begin{theorem}[Theorem~\ref{thm:main} in full]
  \label{thm:main-complete}
  \lean{FastDynamo.MainTheoremLagrangian, FastDynamo.MainTheoremDynamical,
    FastDynamo.MainTheoremComplete,
    FastDynamo.mainTheoremComplete_of_lagrangian,
    FastDynamo.mainTheoremComplete_holds}
  \leanok
  \uses{thm:main-strengthened, prop:main-lagrangian, thm:main}
  The conjunction of the strengthened spectral statement
  (Theorem~\ref{thm:main-strengthened}) with the dynamical clauses
  (Proposition~\ref{prop:main-lagrangian}) is Theorem~\ref{thm:main} in full, and more.
  In Lean it is the structure \texttt{MainTheoremComplete}, inhabited by
  \texttt{mainTheoremComplete\_holds}.
\end{theorem}

The spectral half is
\texttt{main\_theorem\_strengthened} and the dynamical half
\texttt{mainTheoremLagrangian\_holds}, and
\texttt{mainTheoremComplete\_of\_lagrangian} is the implication that combines them.  The
structure has three fields, not two: \texttt{strengthened}, then
\texttt{dynamical}~$:$~\texttt{MainTheoremDynamical}, the theorem's own two clauses,
and \texttt{lagrangian}~$:$~\texttt{MainTheoremLagrangian}, the larger bundle from which the
first is projected.  The implication form exists because the
module that defines \texttt{mainLiteralVelocity} cannot import the flow without a cycle;
\texttt{FastDynamo/Dynamics/MainLagrangianClauses.lean} supplies the closed term.
\texttt{Dynamics/ParticleFlow.lean} is not a leaf of the development: the whole of
\texttt{Dynamics/Flow/*} and \texttt{Dynamics/Ideal/*} depends on it.

\begin{proof}
  \leanok
  \uses{thm:main-strengthened, prop:main-lagrangian}
  Immediate: the two halves are proved separately and the structure is their conjunction.
\end{proof}


\section{Vanishing diffusivity}

Remark~1.7 of the paper (unlabelled, the last remark of §1.2) relates the two halves of
Theorem~\ref{thm:main}: the resistive evolution converges to the ideal one as
$\eps\downarrow0$.  The paper derives it from the Trotter--Kato approximation theorem
[Engel--Nagel, Theorem~III.4.8] and the Friedrichs commutator lemma [DiPerna--Lions,
Lemma~II.1].  The formalisation proves it by a direct energy estimate against a mollified ideal
orbit; the first three nodes are its ingredients, none of which the paper states.

\begin{lemma}[The mollifier on the torus]
  \label{lem:vanishing:mollifier}
  \lean{FastDynamo.literalTranslate, FastDynamo.continuous_literalTranslate_apply,
    FastDynamo.mollifierKernel, FastDynamo.integral_mollifierKernel,
    FastDynamo.literalMollify, FastDynamo.literalMollifyDeriv, FastDynamo.mollifierSymbol,
    FastDynamo.inner_flatTorus3FourierBasis_literalMollify,
    FastDynamo.inner_flatTorus3FourierBasis_literalMollifyDeriv,
    FastDynamo.norm_literalMollify_le, FastDynamo.tendsto_literalMollify,
    FastDynamo.exists_bound_sq_mul_mollifierSymbol}
  \leanok
  Let $\rho_\delta(z)=\delta^{-3}\rho(z/\delta)$ with $\rho\geq0$ a smooth bump on $\R^3$
  supported in the unit ball, of integral one, and let $J_\delta f=\int\rho_\delta(z)f(\cdot-z)\dd z$
  on $\Lp^2(\T^3;\CC^3)$, and $\partial_jJ_\delta$ the same with $\partial_j\rho_\delta$.  Then
  $\norm{J_\delta}\leq1$, $J_\delta f\to f$ as $\delta\downarrow0$, $J_\delta$ and
  $\partial_jJ_\delta$ are Fourier multipliers, and for each $\delta$ the symbol of $J_\delta$
  decays like $\abs m^{-2}$.
\end{lemma}

The operators are Bochner integrals of translations of $\Lp^2$
(\texttt{FastDynamo/Analysis/LiteralTorusMollifier.lean}).  The decay of the symbol is what makes
$J_\delta$ map $\Lp^2$ into $\Hs^2$.

\begin{proof}
  \leanok
  Translations are isometries and strongly continuous, so the integral has norm at most
  $\int\rho_\delta=1$ and converges to $f$ as the kernel concentrates.  The Fourier coefficient
  of a translate is the coefficient times a phase, which gives the multipliers; integrating by
  parts twice in $z$ gives the decay of the symbol.
\end{proof}

\begin{lemma}[The Friedrichs commutator lemma on the torus]
  \label{lem:vanishing:friedrichs}
  \lean{FastDynamo.literalComponentMul, FastDynamo.friedrichsCommutator,
    FastDynamo.friedrichsCommutator_apply, FastDynamo.exists_norm_friedrichsCommutator_le,
    FastDynamo.sum_integral_fderiv_mollifierKernel_mul_eq_zero,
    FastDynamo.tendsto_friedrichsCommutator}
  \leanok
  \uses{lem:vanishing:mollifier}
  Let $v$ be an $L$-Lipschitz velocity on $\T^3$ and
  $K_\delta f=\sum_j\bigl(\partial_jJ_\delta(v_jf)-v_j\partial_jJ_\delta f\bigr)$.  Then
  $\norm{K_\delta}_{\Lp^2\to\Lp^2}\leq C$ uniformly in $\delta\in(0,1]$.  If moreover the
  periodic lift of $v$ is differentiable at every point with $\Div v=0$ at every point, then
  $K_\delta f\to0$ in $\Lp^2$ as $\delta\downarrow0$, uniformly for $f$ in a compact subset of
  $\Lp^2$.
\end{lemma}

DiPerna and Lions state the lemma on $\R^d$, for Sobolev velocities and $\Lp^p$ fields, with
convergence in $\Lp^1_{\rm loc}$; here it is proved on the torus for Lipschitz $v$ with
everywhere differentiable lift, in $\Lp^2$, uniformly on compact sets, which is the case the
energy estimate needs; the constructed velocity satisfies the hypotheses
(Lemma~\ref{lem:dynamics:velocity-gradient}).  Module
\texttt{FastDynamo/Analysis/LiteralFriedrichsCommutator.lean}.

\begin{proof}
  \leanok
  \uses{lem:vanishing:mollifier}
  $K_\delta f=\sum_j\int\partial_j\rho_\delta(z)(\tau_zv_j-v_j\tau_z)f\dd z$ and
  $\norm{\tau_zv_j-v_j\tau_z}\leq L\abs z$, so
  $\norm{K_\delta}\leq3L\int\abs z\abs{\nabla\rho_\delta(z)}\dd z$, which does not depend on
  $\delta$.  For a continuous field $g$, the divergence term
  $\sum_j\int\partial_j\rho_\delta(z)(v_j(x-z)-v_j(x))\dd z=\int\rho_\delta(z)(\Div v)(x-z)\dd z$
  vanishes, which leaves
  $K_\delta g(x)=\sum_j\int\partial_j\rho_\delta(z)(v_j(x-z)-v_j(x))(g(x-z)-g(x))\dd z$, of size
  $O(L\sup_{\abs z\leq\delta}\norm{g(\cdot-z)-g}_\infty)$.  Continuous fields are dense, and
  the uniform bound extends the convergence to all of $\Lp^2$, uniformly on compact sets.
\end{proof}

\begin{lemma}[The resistive operator on mollified fields]
  \label{lem:vanishing:bridge}
  \lean{FastDynamo.productMollifyH2, FastDynamo.productH2Embedding_productMollifyH2,
    FastDynamo.toLiteralC3L2_productInductionAction_productMollifyH2,
    FastDynamo.ScaleFamily.TorusCellPlacement.hasDerivAt_literalMollify_idealGroup,
    FastDynamo.hasDerivAt_literalMollify_mainIdealGroup,
    FastDynamo.literalInductionSemigroup_apply_coe}
  \leanok
  \uses{lem:vanishing:mollifier, def:main:induction-operator,
    lem:dynamics:ideal-weak-existence, lem:dynamics:velocity-gradient,
    thm:dynamics:solenoidal-semigroup}
  For every $\delta$, $J_\delta$ maps $\Lp^2(\T^3;\CC^3)$ continuously into $\Hs^2$, and on
  $J_\delta f$ the auxiliary induction operator of the constructed velocity
  (Definition~\ref{def:main:induction-operator}) reads
  \[
    \mathcal A_\eps J_\delta f=\eps\Delta J_\delta f-\sum_ju_j\,\partial_jJ_\delta f
    +(J_\delta f\cdot\nabla)u .
  \]
  For every $B\in\Lp^2_\sigma$ the curve $t\mapsto J_\delta\mathcal T_0(t)B$ is differentiable
  at every $t$, with derivative
  $J_\delta\bigl((\mathcal T_0(t)B\cdot\nabla)u\bigr)-\sum_j\partial_jJ_\delta\bigl(u_j\,\mathcal T_0(t)B\bigr)$.
\end{lemma}

The first display is stated for the auxiliary action $\mathcal A_\eps$ on the product model and
transported to the literal torus; the derivative of the mollified ideal orbit comes from the
weak form tested against the characters.  Module
\texttt{FastDynamo/Dynamics/VanishingDiffusivity/ProductBridge.lean}.

\begin{proof}
  \leanok
  \uses{lem:vanishing:mollifier, lem:dynamics:ideal-weak-existence}
  The Fourier certificate of $J_\delta f$ in $\Hs^2$ is $(1+4\pi^2\abs m^2)\sigma_\delta(m)\hat
  f(m)$, square summable by the decay of the symbol; the derivative identities hold coefficient
  by coefficient, the transport term is computed on the continuous representative of $u$, and the
  stretching term on its pointwise derivative.  For the ideal orbit, the weak form at the
  characters $e_m$, placed in each vector coordinate, gives
  $J_\delta\mathcal T_0(s)B=J_\delta B+\int_0^sF$ with $F$ continuous, and the fundamental
  theorem of calculus gives the derivative.
\end{proof}

\begin{theorem}[Vanishing diffusivity; Remark~1.7 of the paper]
  \label{thm:vanishing-diffusivity}
  \lean{FastDynamo.resistiveEvolution,
    FastDynamo.norm_productInductionSemigroup_sub_literalMollify_le,
    FastDynamo.norm_mainIdealGroup_sub_resistiveEvolution_le,
    FastDynamo.tendstoUniformlyOn_resistiveEvolution_mainIdealGroup}
  \leanok
  \uses{thm:main, def:dynamics:ideal-induction-group, def:main:growth-rate}
  For each fixed $B_{\rm in}\in\Lp^2_\sigma$, $\e^{tL_{\eps,u}}B_{\rm in}$ converges strongly in
  $\Lp^2$ to $\mathcal T_0(t)B_{\rm in}$ as $\eps\downarrow0$, uniformly on bounded time
  intervals in $[0,\infty)$.
\end{theorem}

The statement is the paper's, for the velocity $u$ of Theorem~\ref{thm:main}: in Lean,
$u$ is \texttt{mainVelocity}, read on the literal torus as \texttt{mainLiteralVelocity},
$\mathcal T_0$ is \texttt{mainIdealGroup}, and the uniformity is on $[0,T]$ for every $T$.  The
resistive evolution \texttt{resistiveEvolution} $\eps\,t\,B$ is
\texttt{literalInductionSemigroup} at $\eps>0$, the semigroup generated by $L_{\eps,u}$ on
$\Lp^2_\sigma$ (Theorem~\ref{thm:dynamics:solenoidal-generator}); it is extended by $0$ to
$\eps\leq0$ and clamps negative times to $0$, conventions that the limit along
$\eps\downarrow0$ on $[0,T]$ does not see.  The paper's remark follows Theorem~1.2 and
concerns its velocity and its group $\mathcal T_0$; the Lean statement is for that velocity
only, not for a general divergence-free Lipschitz velocity.  The route differs: not
Trotter--Kato, but the energy estimate below.  Module
\texttt{FastDynamo/Dynamics/VanishingDiffusivity/Limit.lean}.

\begin{proof}
  \leanok
  \uses{lem:vanishing:friedrichs, lem:vanishing:bridge, lem:vanishing:mollifier,
    lem:dynamics:product-dissipative, thm:dynamics:product-semigroup}
  Fix a mollification scale $\delta$ and start the resistive semigroup at $J_\delta B$; let $w$
  be its $\Hs^2$ orbit and $e(t)=w(t)-J_\delta\mathcal T_0(t)B$, so $e(0)=0$.  By
  Lemma~\ref{lem:vanishing:bridge}, $e$ solves the resistive equation up to the residual
  $\eps\Delta J_\delta\mathcal T_0(t)B+K_\delta\mathcal T_0(t)B
  +(S_uJ_\delta-J_\delta S_u)\mathcal T_0(t)B$, with $S_uB=(B\cdot\nabla)u$,
  and the dissipativity $\re\inner{e}{\mathcal A_\eps e}\leq M_u\norm e^2$ with Gronwall's
  inequality bounds $\norm{e(t)}$ by the residual times $\e^{(2M_u+1)T}$.  On the compact orbit
  $\{\mathcal T_0(t)B:t\in[0,T]\}$ the commutator $K_\delta$ and the mollification errors are
  small uniformly for small $\delta$ (Lemmas~\ref{lem:vanishing:friedrichs}
  and~\ref{lem:vanishing:mollifier}); for fixed $\delta$ the diffusion term is $O(\eps)$; and
  the initial error $\norm{J_\delta B-B}$ is propagated by
  $\norm{\e^{tL_{\eps,u}}}\leq\e^{M_ut}$, with $M_u$ independent of $\eps$.  Choose $\delta$
  first and then $\eps$.
\end{proof}


\chapter{The horizontal operator on $\T^2$ and the scaling identity}
% Chapter: the horizontal operator on T^2 and the scaling identity.  Section 5 of the paper
% (sec:axial).
% Paper labels: lem:torus-sector-basic (Lemma 5.1), lem:exact-renormalization (Lemma 5.2).

The paper's proof works with the two horizontal components only.  Because the global
velocity of the previous chapter is independent of $x_3$, the horizontal equations in the
axial Fourier mode $\e^{2\pi iKx_3}$ close on $h\colon\T^2\to\CC^2$, and the resulting
operator $A_{\eps,K}$ has domain $\Hs^2(\T^2;\CC^2)$ with \emph{no} constraint on $h$ --
in particular multiplication by a cut-off preserves the domain, which is what makes the
parametrix of the transfer chapter possible.  The axial velocity enters only through the
zeroth-order multiplication operator $-2\pi iKu_3$, so no derivative of $u_3$ occurs
anywhere.  This chapter states the basic operator-theoretic facts about $A_{\eps,K}$ and
the exact renormalisation that identifies the selected cell with the local seed of the
previous chapter; the missing third component is recovered only at the very end
(Lemma~\ref{lem:axial-recovery}).

The paper fixes its terminology at the start of its Section~5: in the spectral analysis it
distinguishes the reference model on $\R^2$, the model on $\R^2$ for the selected cell, and the
horizontal problem on $\T^2$ associated with the velocity profile $u$, and uses the same terms
for the corresponding operators, resolvents, approximate inverses and Riesz projections.  This
blueprint says \emph{local} for the two models on $\R^2$ and \emph{global} for the horizontal
problem on $\T^2$, and qualifies the corresponding objects in the same way.

\section{The horizontal operator on $\T^2$}

\begin{definition}[The horizontal induction operator in one axial sector]
  \label{def:operator:horizontal-induction}
  \lean{FastDynamo.horizontalInductionDiffusion,
    FastDynamo.horizontalInductionTransport,
    FastDynamo.horizontalInductionAxialTransport,
    FastDynamo.horizontalInductionStretching,
    FastDynamo.horizontalInductionAction,
    FastDynamo.horizontalInductionOperator,
    FastDynamo.horizontalInductionAction_apply,
    FastDynamo.transverseW1InfHorizontalDerivLInf,
    FastDynamo.horizontalProjH2}
  \leanok
  \uses{def:main:induction-operator, def:velocity:global-velocity}
  For $\eps>0$ and $K\in\Z$ define, on $\Lp^2(\T^2;\CC^2)$,
  \begin{equation}
    A_{\eps,K}h=\eps\bigl(\Delta_\perp-(2\pi K)^2\bigr)h
      -u_\perp\cdot\nabla_\perp h-2\pi iKu_3h+(h\cdot\nabla_\perp)u_\perp,
    \qquad D(A_{\eps,K})=\Hs^2(\T^2;\CC^2).
    \label{eq:operator:torus-sector}
  \end{equation}
  No constraint is imposed on $h$.
\end{definition}

The Lean operator is \texttt{FastDynamo.horizontalInductionOperator u eps K}, a
\texttt{DomainOperator TransverseC2L2 TransverseC2H2}
(\texttt{FastDynamo/Operator/HorizontalSector.lean}), assembled from the four displayed
summands.  Two points of the encoding.  The stretching coefficient is
\texttt{transverseW1InfHorizontalDerivLInf j u}, the \emph{horizontal projection}
$(\partial_ju)_\perp$ of the velocity gradient: this is exactly where the two-component
operator differs from the three-component $A_{\eps,K}^{(3)}$ of
Definition~\ref{def:main:induction-operator}, whose stretching term also produces a third
row $\sum_jh_j\partial_ju_3$; because that row is non-zero, the horizontal subspace is not
invariant under $A^{(3)}_{\eps,K}$.  Lifting a horizontal eigenpair back to three
components is the content of Lemma~\ref{lem:axial-recovery}.

\begin{lemma}[The horizontal operator on $\T^2$ associated with $w$]
  \label{lem:torus-sector-basic}
  \lean{FastDynamo.HorizontalTorusSectorBasic,
    FastDynamo.horizontalInductionOperator_basic,
    FastDynamo.horizontalInductionOperator_isClosedGraph,
    FastDynamo.horizontalInductionOperator_hasNonemptyResolvent,
    FastDynamo.norm_transverseC2H2_le_horizontalInduction_graph_of_hasResolventOn,
    FastDynamo.isCompactOperator_horizontalInduction_stateResolvent,
    FastDynamo.pencilFredholmIndexZeroOn_horizontalInductionOperator,
    FastDynamo.isFredholm_horizontalInductionPencil,
    FastDynamo.index_horizontalInductionPencil_eq_zero}
  \leanok
  \uses{def:operator:horizontal-induction}
  Let $\eps>0$, $K\in\Z$ and $w=(w_\perp,w_3)\in\W^{1,\infty}(\T^2;\R^3)$ with
  $\Div_\perp w_\perp=0$, and write $A_{\eps,K}[w]$ for the operator
  \eqref{eq:operator:torus-sector} with $u$ replaced by $w$.  Then $A_{\eps,K}[w]$ is closed
  and has non-empty resolvent set; its graph norm is equivalent to the $\Hs^2$ norm; its
  resolvent is compact; and
  \[
    z-A_{\eps,K}[w]\colon\Hs^2(\T^2;\CC^2)\longrightarrow\Lp^2(\T^2;\CC^2)
  \]
  is Fredholm of index zero for every $z\in\CC$.
\end{lemma}

The paper states the lemma (its Lemma~5.1) for an arbitrary such $w$ rather than for the
constructed velocity $u$, which is what lets the background chapter reuse it verbatim at
$u_n^{\rm bg}$.  \textbf{The Lean statement is stronger still}:
\texttt{horizontalInductionOperator\_basic} quantifies over an arbitrary
$u\in\W^{1,\infty}(\T^2;\R^3)$ -- an arbitrary element of \texttt{TransverseRealW1Inf} --
with \emph{no} hypothesis $\Div_\perp u_\perp=0$, because the route below never uses it.
The paper's lemma is the instance $u:=w$ with the hypothesis discarded.

\textbf{Deviation: a Neumann series and a compact base point, in place of coercivity and a
compact perturbation.}  The paper's proof first obtains Fredholm index zero for every $z$
from the fact that $z-A_{\eps,K}[w]$ differs from $\id-\eps\Delta_\perp$ by an operator
bounded $\Hs^1\to\Lp^2$, hence compact $\Hs^2\to\Lp^2$, while $\id-\eps\Delta_\perp$ is an
isomorphism $\Hs^2\to\Lp^2$; it then obtains injectivity of $z_0-A_{\eps,K}[w]$ at every
real $z_0>\norm{\nabla_\perp w_\perp}_\infty$ from a coercivity estimate that uses
$\Div_\perp w_\perp=0$ and the realness of $w_3$; surjectivity at $z_0$ comes from index
zero, and the bounded inverse theorem gives the graph-norm equivalence, closedness and,
with the compact embedding $\Hs^2\hookrightarrow\Lp^2$, the compact resolvent.  Lean
reverses the order and uses neither ingredient.  It inverts the diffusion part exactly as a Fourier multiplier on the graph
space and shows that the three lower-order terms composed with that inverse have operator
norm tending to zero along the real ray $z\to+\infty$; a Neumann series gives a resolvent
point directly, and the factorisation
$z-A=\bigl[\id+(z-z_0)(z_0-A)^{-1}\bigr](z_0-A)$ -- identity-plus-compact on the
\emph{ambient} Hilbert space $\Lp^2(\T^2;\CC^2)$ -- gives Fredholm index zero at every $z$.
The ambient side is forced: the graph domain \texttt{TransverseC2H2} is a closed submodule
of a sup-normed product and carries no inner product, so Atkinson's theorem is not available
on it (\texttt{FastDynamo/Operator/CompactBaseFredholm.lean}), and the paper's route --
index zero for $T+K$ with $T$ an isomorphism between two \emph{different} spaces -- is not
available in Mathlib in that form.  The pencil $z-A\colon\Hs^2\to\Lp^2$ itself is then
proved Fredholm with index zero
(\texttt{isFredholm\_horizontalInductionPencil},
\texttt{index\_horizontalInductionPencil\_eq\_zero}); only the argument passes through the
ambient space.  The conclusion is the paper's, and the Lean route
is \emph{stronger}, since it never uses $\Div_\perp w_\perp=0$.  The half of the
graph-norm equivalence that needs no hypothesis
($\norm h_{\Hs^2}\geq c(\norm h_2+\norm{A h}_2)$) is the continuity of the embedding
$\Hs^2\hookrightarrow\Lp^2$ and of the operator action $\Hs^2\to\Lp^2$, which the Lean
development does not bundle under a name; the other half is the listed
\texttt{norm\_transverseC2H2\_le\_horizontalInduction\_graph\_of\_hasResolventOn}.  The
Fredholm clause is stated in Lean for an arbitrary subset of $\CC$ and is therefore the
paper's ``for every $z\in\CC$'' (take the whole plane).

\begin{proof}
  \leanok
  \uses{def:operator:horizontal-induction}
  Write $z-A_{\eps,K}=\eps(z/\eps-\Delta_\perp+(2\pi K)^2)-B$ with $B$ the three
  lower-order terms.  The diffusion pencil is diagonal in the Fourier basis and invertible
  with $\Hs^2$-bounded inverse for $z$ large and real, and the norm of
  $B\circ(\text{that inverse})$ tends to zero, because the transport term gains one
  derivative and the axial and stretching terms two.  So the pencil is invertible for
  large real $z$ -- this is the deviation recorded above, in place of the paper's
  coercivity at $z_0>\norm{\nabla_\perp w_\perp}_\infty$ -- which gives simultaneously the
  non-empty resolvent set, the closedness of the graph and the graph-norm equivalence.
  Compactness of the resolvent is the Rellich embedding
  $\Hs^2(\T^2)\hookrightarrow\Lp^2(\T^2)$ composed with the bounded resolvent.  Finally
  $z-A=\bigl[\id+(z-z_0)(z_0-A)^{-1}\bigr](z_0-A)$ exhibits the pencil as an isomorphism
  composed with the identity plus a compact operator on the ambient $\Lp^2$, which is
  Fredholm of index zero; the paper instead compares $z-A_{\eps,K}[w]$ with
  $\id-\eps\Delta_\perp$.
\end{proof}

\section{Rescaling to the local model}

\begin{definition}[The selected whole-plane operator and the rescaling]
  \label{def:operator:selected-whole-plane}
  \lean{FastDynamo.wholePlaneC2LocalAction,
    FastDynamo.wholePlaneC2LocalOperator,
    FastDynamo.horizontalLocalOperator,
    FastDynamo.selectedWholePlaneC2AxialShift,
    FastDynamo.selectedWholePlaneC2InductionAction,
    FastDynamo.selectedWholePlaneC2InductionOperator,
    FastDynamo.WholePlaneExactScalingParameters,
    FastDynamo.shiftedWholePlaneC2LocalAction,
    FastDynamo.shiftedWholePlaneC2LocalOperator}
  \leanok
  \uses{def:local:smoothed-velocity, def:velocity:scale-family,
    def:velocity:global-velocity}
  Fix $0<\eps\leq\eps_0$ and an index $n$ with $\delta_n=\eps/\ell_n^2\in I$, and work at
  the axial index $K_n$.  The \emph{selected whole-plane model} keeps $u_n^{\rm loc}$ and
  extends the constant axial velocity $a_ne_3$ to all of $\R^2$:
  \begin{equation}
    A^{(n)}_{\eps,n}h=\eps\bigl(\Delta_\perp-(2\pi K_n)^2\bigr)h
      -(u_n^{\rm loc})_\perp\cdot\nabla_\perp h
      -2\pi iK_n\bigl((u_n^{\rm loc})_3+a_n\bigr)h
      +(h\cdot\nabla_\perp)(u_n^{\rm loc})_\perp
    \label{eq:operator:selected-whole-plane}
  \end{equation}
  on $\Lp^2(\R^2;\CC^2)$ with domain $\Hs^2(\R^2;\CC^2)$.  The rescaling is the unitary
  \[
    U_n\colon\Lp^2(\R^2_Y;\CC^2)\to\Lp^2(\R^2_y;\CC^2),
    \qquad(U_nf)(y)=\ell_n^{-1}f\Bigl(\frac{y-p_n}{\ell_n}\Bigr).
  \]
\end{definition}

The three exact identities $\delta_n=\eps/\ell_n^2$, $N=2\pi K_n\ell_n$ and
$\tau=2\pi K_na_n$ -- which the paper, having fixed $N=2\pi$ and $\ell_n=1/K_n$, writes
as $\delta_n=\eps/\ell_n^2$, $K_n\ell_n=1$ and $2\pi K_na_n=\tau$ in the proof of its
Lemma~5.2 -- are the
three fields of the Lean record
\texttt{WholePlaneExactScalingParameters}; they are the \emph{only} input of the next
lemma, which is therefore stated for an arbitrary such record rather than for the
particular cell.  $U_n$ is realised by two differently typed maps, a linear isometry
equivalence on $\Lp^2$ and a continuous linear equivalence on $\Hs^2$, so the conjugation
identity below is stated for the actions and the embeddings separately rather than as an
equality of \texttt{DomainOperator}s.

\begin{lemma}[Scaling]
  \label{lem:exact-renormalization}
  \lean{FastDynamo.exactScalingC2_diffusion,
    FastDynamo.exactScalingC2_transverseTransport,
    FastDynamo.exactScalingC2_localAxialTransport,
    FastDynamo.exactScalingC2_stretching,
    FastDynamo.exactScalingC2_constantAxialShift,
    FastDynamo.exactScalingC2_selectedAction_clm,
    FastDynamo.exactScalingC2_embedding,
    FastDynamo.exactScalingC2_selectedOperator_pencil_conjugated,
    FastDynamo.exactScalingC2_eigenpair_iff,
    FastDynamo.exactScalingC2_selectedEigenpair_to_local}
  \leanok
  \uses{def:operator:selected-whole-plane, def:local:smoothed-velocity}
  With $n=n(\eps)$ as above,
  \begin{equation}
    U_n^{-1}A^{(n)}_{\eps,n}U_n=\cL_{\delta_n}-i\tau
    \label{eq:operator:exact-scaling}
  \end{equation}
  as operators with domain $\Hs^2(\R^2;\CC^2)$.  Consequently $\lambda$ is an eigenvalue of
  $A^{(n)}_{\eps,n}$ with eigenfunction $U_nf$ exactly when $\lambda+i\tau$ is an
  eigenvalue of $\cL_{\delta_n}$ with eigenfunction $f$.
\end{lemma}

The paper's Lemma~5.2 consists of the identity \texttt{eq:scaling} alone; the eigenvalue
correspondence of the last sentence is the form in which the later chapters use it, and in
Lean it is a separate pair of theorems (\texttt{exactScalingC2\_eigenpair\_iff} and its
forward half \texttt{exactScalingC2\_selectedEigenpair\_to\_local}).  The identity itself is
formalised as the conjugation of each of the four summands, of the action and of the
embedding, and of the pencil at every spectral parameter
(\texttt{exactScalingC2\_selectedOperator\_pencil\_conjugated}), for the reason given
before the lemma, and for an arbitrary \texttt{WholePlaneExactScalingParameters} record.

\begin{proof}
  \leanok
  \uses{def:operator:selected-whole-plane, def:velocity:scale-family}
  Change variables by $y=p_n+\ell_nY$ and use
  $u_n^{\rm loc}(p_n+\ell_nY)=\ell_nv(Y)$.  Each of the four summands scales exactly: the
  diffusion picks up $\ell_n^{-2}$ and $\eps\ell_n^{-2}=\delta_n$, while
  $\eps(2\pi K_n)^2=\delta_nN^2$ by $2\pi K_n\ell_n=N$ (the paper's $K_n\ell_n=1$ at
  $N=2\pi$, giving $\eps(2\pi K_n)^2=4\pi^2\delta_n$); the transport and the stretching each
  pick up one factor $\ell_n$ from the amplitude and one $\ell_n^{-1}$ from the derivative,
  so they are unchanged; the axial multiplier splits as $-iNv_3-i\tau$ by
  $2\pi K_n\ell_n=N$ and $2\pi K_na_n=\tau$ (the paper's $-2\pi iv_3-i\tau$).  Lean keeps
  $N=2\pi k_*$ symbolic throughout, so the record is stated with $N$ and the paper's
  display \texttt{eq:scaling} is the case $k_*=1$.
\end{proof}

\begin{lemma}[The translated contour]
  \label{lem:operator:shifted-contour}
  \lean{FastDynamo.backgroundShiftedContour,
    FastDynamo.backgroundShiftedContour_growthFloor,
    FastDynamo.verticallyShiftedSet,
    FastDynamo.verticallyShiftedSet_boundary,
    FastDynamo.verticallyShiftedSet_closedRegion}
  \leanok
  \uses{lem:exact-renormalization, def:local:seed}
  Set $\Gamma_\tau=\{z-i\tau:z\in\Gamma\}$ and $\cG_\tau=\{z-i\tau:z\in\cG\}$.  Both are
  independent of $n$ and of $\eps$, and translation by $-i\tau$ leaves real parts
  unchanged, so $\min_{\zeta\in\cG_\tau}\re\zeta=g_*$.
\end{lemma}

The paper does not state this as a lemma: it is the display after the proof of its
Lemma~5.2, followed by the remark that both sets are independent of $n$ and $\eps$ and that
their real parts are unchanged.  In Lean $\Gamma_\tau$ is again a \texttt{SpectralCircle}, with centre
$\Gamma.\texttt{centre}-i\tau$ and the \emph{same} radius, so the growth floor is exactly
unchanged; \texttt{verticallyShiftedSet tau S} is the preimage $\{z:z+i\tau\in S\}$, and
the two lemmas identify its boundary and closed region with those of the shifted circle.
That is the bridge by which the rescaled seed resolvent of
Corollary~\ref{cor:scaled-local-resolvents} and the background resolvent of
Proposition~\ref{prop:background-resolvent} live on one and the same contour.

\begin{proof}
  \leanok
  \uses{lem:exact-renormalization}
  Immediate from \eqref{eq:operator:exact-scaling}: the spectrum of
  $U_n^{-1}A^{(n)}_{\eps,n}U_n$ is that of $\cL_{\delta_n}$ translated by $-i\tau$, and
  $\re(z-i\tau)=\re z$.
\end{proof}


\chapter{A uniform estimate for the background resolvent on $\T^2$}
% Chapter: the uniform estimate for the background resolvent.  Section 6 of the paper
% (sec:localization).
% Paper labels: prop:background-resolvent (Proposition 6.1),
% cor:scaled-local-resolvents (Corollary 6.2).

This is the analytic heart of the proof.  Remove the \emph{selected} local cell from the
global velocity and keep
everything else -- every other helical cell and every added axial velocity.  The claim is
that the resulting operator has \emph{no} spectrum in the translated region $\cG_\tau$, with
a resolvent bound uniform in the cell index, in the diffusivity and in the spectral
parameter.  The whole estimate comes from testing the resolvent equation against the
weighted field $(1+i\beta_n)h$, where $\beta_n=F(s_n)$ is a rational function of the axial
velocity profile: the axial shifts $\tau s_n$, which differ from $\tau$ by a definite
amount on every non-selected core, turn into a positive zeroth-order term there, while
$\re\zeta\geq g_*$ carries the rest of the torus.

\section{The background velocity and its operator}

\begin{definition}[The background velocity and the background operator]
  \label{def:background:periodic-background}
  \lean{FastDynamo.ScaleFamily.TorusCellPlacement.selectedLocalCell,
    FastDynamo.ScaleFamily.TorusCellPlacement.selectedBackgroundVelocity,
    FastDynamo.ScaleFamily.TorusCellPlacement.selectedBackgroundVelocity_add_selectedLocalCell,
    FastDynamo.ScaleFamily.TorusCellPlacement.selectedBackgroundVelocity_isTransverseDivergenceFree,
    FastDynamo.ScaleFamily.TorusCellPlacement.selectedBackgroundOperator,
    FastDynamo.ScaleFamily.TorusCellPlacement.selectedBackgroundOperator_eq}
  \leanok
  \uses{def:velocity:global-velocity, def:operator:horizontal-induction}
  Let $\eps>0$ and $n\geq0$ with $\delta_n=\eps/\ell_n^2\in I$.
  The \emph{background velocity} is
  \begin{equation}
    u_n^{\rm bg}=u-u_n^{\rm loc}
      =\sum_{m\neq n}u_m^{\rm loc}+\sum_{m\geq0}u_m^{\rm ax},
    \label{eq:background:velocity}
  \end{equation}
  and the \emph{background operator} is
  \begin{equation}
    A^{\rm bg}_{\eps,n}h=\eps\bigl(\Delta_\perp-(2\pi K_n)^2\bigr)h
      -(u_n^{\rm bg})_\perp\cdot\nabla_\perp h-2\pi iK_n(u_n^{\rm bg})_3h
      +(h\cdot\nabla_\perp)(u_n^{\rm bg})_\perp,
    \qquad D(A^{\rm bg}_{\eps,n})=\Hs^2(\T^2;\CC^2).
    \label{eq:background:operator}
  \end{equation}
  It agrees with $A_{\eps,K_n}$ outside $\supp u_n^{\rm loc}$.
\end{definition}

In Lean $A^{\rm bg}_{\eps,n}$ is \emph{exactly}
\texttt{horizontalInductionOperator (selectedBackgroundVelocity n) eps (K n)}, so the whole
operator theory of Lemma~\ref{lem:torus-sector-basic} applies to it verbatim with no new
work; the agreement outside the selected support is the subtraction identity
\texttt{selectedBackgroundVelocity\_add\_selectedLocalCell}.  This is also how the paper
argues: it states Lemma~5.1 (\texttt{lem:torus-sector-basic}) for an arbitrary
$w\in\W^{1,\infty}(\T^2;\R^3)$ with $\Div_\perp w_\perp=0$, and the proof of
Proposition~6.1 of the paper applies it to $u_n^{\rm bg}$, which is real valued, belongs to
$\W^{1,\infty}$ and is horizontally divergence-free
(\texttt{selectedBackgroundVelocity\_isTransverseDivergenceFree}).

\begin{definition}[The scalar multiplier $F$, the profile $s_n$ and the weight $\beta_n$]
  \label{def:background:multiplier}
  \lean{FastDynamo.backgroundMultiplierF, FastDynamo.backgroundMultiplierDeriv,
    FastDynamo.abs_backgroundMultiplierF_le_one,
    FastDynamo.abs_mul_backgroundMultiplierF_le_one,
    FastDynamo.abs_backgroundMultiplierDeriv_le_three,
    FastDynamo.backgroundRatioGap, FastDynamo.backgroundRatioGap_le_abs,
    FastDynamo.ScaleFamily.TorusCellPlacement.backgroundShiftedGlobalTorus,
    FastDynamo.ScaleFamily.TorusCellPlacement.backgroundMultiplier,
    FastDynamo.ScaleFamily.TorusCellPlacement.norm_backgroundMultiplier_value_le_one,
    FastDynamo.ScaleFamily.TorusCellPlacement.gradSupNormSq_backgroundMultiplier_le,
    FastDynamo.TransverseScalarW1Inf.gradSupNormSq}
  \leanok
  \uses{def:velocity:scale-family, def:velocity:placement, def:background:periodic-background}
  With $r_m=K_n/K_m=\ell_m/\ell_n$, $c_q=(q-1)/(q+1)$,
  $M=\max_{\zeta\in\cG}\abs{\im\zeta}$ and $d(K_0)=\sup_m\ell_m/R_m=\sqrt{\ell_0}$, set
  \begin{equation}
    s_n(y)=\sum_{m\geq0}r_m\vartheta\Bigl(\frac{\abs{y-p_m}}{R_m}\Bigr),
    \qquad\beta_n=F(s_n),\qquad F(s)=\frac{s-1}{(1+s)^2}.
    \label{eq:background:multiplier}
  \end{equation}
  Then $r_m\geq q$ for $m<n$ and $r_m\leq q^{-1}$ for $m>n$, and
  \begin{equation}
    \abs{\beta_n}\leq1,
    \qquad\sqrt\eps\,\norm{\nabla_\perp\beta_n}_\infty\leq C\sqrt{\delta_+}\,d(K_0).
    \label{eq:background:multiplier-bounds}
  \end{equation}
\end{definition}

Two encoding points.  The summands of $s_n$ have pairwise disjoint supports, but the series
$\sum_mr_m\vartheta_m$ has \emph{unbounded} coefficients, so Lean sums the shifted series
$\beta_n+1=\sum_m\Psi_{r_m}(\vartheta_m)$ with $\Psi_\rho(t)=F(\rho t)+1$ and
$\Psi_\rho(0)=0$, which is summable in $\Lp^\infty$, and defines $\beta_n$ as that sum minus
one.  And $\norm{\nabla_\perp\beta_n}_\infty$ is realised as
$\texttt{gradSupNormSq}=\sum_j\norm{\partial_j\beta}_\infty^2$, which dominates
$\norm{\abs{\nabla\beta}}_\infty^2$ and so appears only on the larger side of every
estimate; the proved bound is
$\texttt{gradSupNormSq}\,\beta_n\leq2(3\norm{\vartheta'}_\infty R_0/\ell_n)^2$, giving
\eqref{eq:background:multiplier-bounds} with $C=3\sqrt2\norm{\vartheta'}_\infty$.  (The paper
introduces $r_m,c_q,M,d(K_0),s_n,\beta_n$ at the start of the proof of Proposition~6.1 and
derives \eqref{eq:background:multiplier-bounds}, its \texttt{eq:background-multiplier-bounds},
at the start of the estimate of the diffusion term $I_1$.)

\begin{lemma}[The phase inequality]
  \label{lem:background:phase}
  \lean{FastDynamo.backgroundMultiplierF_phase_lower_bound_sharp,
    FastDynamo.backgroundMultiplierF_phase_lower_bound_ratioGap,
    FastDynamo.backgroundMultiplierF_phase_lower_bound,
    FastDynamo.backgroundTauThreshold, FastDynamo.exists_backgroundTau,
    FastDynamo.ScaleFamily.TorusCellPlacement.BackgroundParameters.phase_bound}
  \leanok
  \uses{def:background:multiplier}
  Write $\mathcal C_n=\bigcup_{m\neq n}\supp u_m^{\rm loc}$.  For $\abs t\leq M$ and
  $s\geq0$,
  \begin{equation}
    F(s)\bigl(t+\tau(s-1)\bigr)
      \geq\frac\tau2\Bigl(\frac{s-1}{s+1}\Bigr)^2-\frac{M^2}{2\tau(1+s)^2},
    \label{eq:background:phase-positive}
  \end{equation}
  and on $\mathcal C_n$ one has $\abs{(s_n-1)/(s_n+1)}\geq c_q$, so that
  $\beta_n(\im\zeta+\tau(s_n-1))\geq\tfrac{\tau c_q^2}2\one_{\mathcal C_n}-\tfrac{M^2}{2\tau}$
  throughout $\T^2$.
\end{lemma}

\begin{proof}
  \leanok
  \uses{def:background:multiplier}
  The difference of the two sides of \eqref{eq:background:phase-positive} is
  $\bigl[(\tau(s-1)+t)^2-t^2+M^2\bigr]/(2\tau(1+s)^2)$, which is non-negative when
  $\abs t\leq M$.  On the core of cell $m$ the plateau gives $s_n=r_m$, and
  $r\mapsto(r-1)/(1+r)$ is increasing, so $r_m\geq q$ gives $\geq c_q$ and
  $r_m\leq q^{-1}$ gives $\leq-c_q$.  Relaxing $(1+s)^2\geq1$ gives the second display.
\end{proof}

\section{The weighted energy estimate}

\begin{lemma}[The weighted test field]
  \label{lem:background:weighted-test}
  \lean{FastDynamo.weightedTransverseState,
    FastDynamo.transverseGradientEnergy,
    FastDynamo.transverseWeightedGradientDefect,
    FastDynamo.re_inner_weightedTransverseState_neg_laplacian_ge,
    FastDynamo.transverseWeightedGradientDefect_young_scaled,
    FastDynamo.re_inner_weightedTransverseState_neg_laplacian_ge_young,
    FastDynamo.re_inner_weightedTransverseState_smul,
    FastDynamo.re_inner_weightedTransverseState_axialMultiplier,
    FastDynamo.re_inner_weightedTransverseState_transport_eq,
    FastDynamo.norm_weightedTransverseState_le,
    FastDynamo.norm_le_norm_weightedTransverseState}
  \leanok
  \uses{def:background:multiplier}
  Let $\beta\in\W^{1,\infty}(\T^2)$ and $a\in\Lp^\infty(\T^2)$ be real,
  $w\in\W^{1,\infty}(\T^2;\R^2)$, $k\in\R$, $\eps\geq0$, $z\in\CC$ and $h\in\Hs^2(\T^2;\CC^2)$.  Then
  \[
    \re\inner{\nabla_\perp h}{\nabla_\perp((1+i\beta)h)}
      \geq\norm{\nabla_\perp h}_2^2-\mathcal E(\beta,h),
  \]
  the zeroth-order term contributes
  $\re\bigl[(z+\eps k^2)\inner h{(1+i\beta)h}\bigr]
   =\int\bigl[\re z+\eps k^2+\beta\im z\bigr]\abs h^2$, the axial multiplier contributes
  $k\int\beta a\abs h^2$, the unweighted transport term has zero real part when
  $\Div_\perp w=0$, and $\norm h_2\leq\norm{(1+i\beta)h}_2\leq\sqrt2\norm h_2$.
\end{lemma}

\textbf{Deviation: the rendering of $\int\abs{\nabla_\perp\beta}\abs h\abs{\nabla_\perp h}$.}
In its estimate of the diffusion term $I_1$ of \texttt{eq:background-testing}, the paper
majorises the exact defect $\sum_j\im\inner{(\partial_j\beta)h}{\partial_jh}$ by the
integral $\int\abs{\nabla_\perp\beta_n}\abs h\abs{\nabla_\perp h}$ (pointwise
Cauchy--Schwarz in $j$, then integration).  Lean majorises it instead by
$\mathcal E(\beta,h)=\sum_j\norm{(\partial_j\beta)h}_2\norm{\partial_jh}_2$
(Cauchy--Schwarz in $\Lp^2$, one $j$ at a time).  \emph{The two majorants are
incomparable}; neither dominates the other.  Nothing depends on the choice, because the
only use of either is the same Young absorption
$\eps\mathcal E\leq\tfrac\eps4\norm{\nabla_\perp h}_2^2+C\delta_+d(K_0)^2\norm h_2^2$, which
is proved directly for the Lean defect.  A second orientation point: the Lean library
carries the axial-multiplier identity in \emph{both} signs, and the proof of
Proposition~\ref{prop:background-resolvent} consumes the one with $+2\pi iK_n$, matching
$\mathfrak a_z(h,\phi)=\inner{(z-A^{\rm bg})h}\phi$; with the normalisation
$2\pi K_na^{\rm ax}_m=\tau r_m$, forced by $2\pi K_n\ell_m=Nr_m$ (the Lean keeps the seed
wavenumber $N=2\pi k_*$ symbolic; the paper's value is $2\pi$), this is what produces the
$+\tau s_n$ inside the phase bracket of $I_2$.

\begin{proof}
  \leanok
  \uses{def:background:multiplier}
  Expand $\nabla_\perp((1+i\beta)h)=(1+i\beta)\nabla_\perp h+i(\nabla_\perp\beta)h$ and take
  real parts: the first summand gives $\norm{\nabla_\perp h}_2^2$ exactly, the second the
  defect, which Cauchy--Schwarz bounds one direction at a time.  The zeroth-order and axial
  identities are direct computations with a real $\Lp^\infty$ coefficient, for which
  multiplication is self-adjoint.  The transport identity splits the weighted pairing into
  the unweighted one, whose real part vanishes for a divergence-free $w$, and a remainder
  carrying $\beta w$.  Finally $\abs{1+i\beta}^2=1+\beta^2\in[1,2]$ gives the two norm
  bounds.
\end{proof}

\begin{lemma}[The coercivity estimate]
  \label{lem:background:coercivity}
  \lean{FastDynamo.re_inner_weightedTransverseState_horizontalInductionPencil_ge,
    FastDynamo.BackgroundCoercivity,
    FastDynamo.BackgroundCoercivity.re_inner_pencil_ge,
    FastDynamo.BackgroundCoercivity.norm_add_sqrt_eps_gradientEnergy_le,
    FastDynamo.BackgroundCoercivity.pencil_injective,
    FastDynamo.BackgroundCoercivity.backgroundResolventConstant,
    FastDynamo.ScaleFamily.TorusCellPlacement.backgroundCoercivityConstant,
    FastDynamo.ScaleFamily.TorusCellPlacement.backgroundDiffusionConstant,
    FastDynamo.ScaleFamily.TorusCellPlacement.BackgroundParameters,
    FastDynamo.ScaleFamily.TorusCellPlacement.BackgroundParameters.coercivity,
    FastDynamo.ScaleFamily.TorusCellPlacement.BackgroundParameters.transport_bound}
  \leanok
  \uses{lem:background:weighted-test, lem:background:phase,
    def:background:periodic-background}
  With $z=\zeta-i\tau$, $\zeta\in\cG$, and $h\in\Hs^2(\T^2;\CC^2)$,
  \begin{align}
    \re\mathfrak a_z\bigl(h,(1+i\beta_n)h\bigr)
    \geq{}&\tfrac\eps2\norm{\nabla_\perp h}_2^2
      +\Bigl(g_*-\frac{M^2}{2\tau}-C_1\delta_+d(K_0)^2\Bigr)\norm h_2^2\notag\\
      &+\Bigl(\frac{\tau c_q^2}2-C_0\Bigr)\int_{\mathcal C_n}\abs h^2,
    \label{eq:background:coercivity}
  \end{align}
  with $C_0,C_1$ independent of $\tau$, $K_0$, $n$ and $\delta_n\in I$.  Choosing
  \begin{equation}
    \frac{M^2}{2\tau}\leq\frac{g_*}4,\qquad\frac{\tau c_q^2}2\geq C_0+1,
    \qquad\text{then}\qquad C_1\delta_+\ell_0\leq\frac{g_*}4,
    \label{eq:background:axial-separation-choice}
  \end{equation}
  (the paper's \texttt{eq:axial-separation-choice} and \texttt{eq:background-diffusion-small})
leaves $\tfrac\eps2\norm{\nabla_\perp h}_2^2+\tfrac{g_*}2\norm h_2^2
  \leq\re\mathfrak a_z(h,(1+i\beta_n)h)\leq\sqrt2\norm f_2\norm h_2$ whenever
  $(z-A^{\rm bg}_{\eps,n})h=f$; in particular the pencil is injective on $\cG_\tau$.
\end{lemma}

The Lean constants are closed terms in the seed data alone:
$C_0=C_vN+2(C_v/\sqrt{\delta_-})^2+4C_s$ and $C_1=18\norm{\vartheta'}_\infty^2$, both free
of $\tau$ and of $K_0$, which is what makes the order
$(\text{seed})\to q\to\tau\to K_0$ of the parameter chain non-circular (the order the
paper states in prose in Section~3.1 and in the summary paragraph at the end of Section~7.2).  The
stretching term is absorbed with the crude constant $4C_s$ where
$\abs{1+i\beta}\leq\sqrt2$ would give $2\sqrt2C_s$; that only enlarges $C_0$, hence
$\tau_0$.  The Lean estimate also keeps the term $\eps(2\pi K_n)^2\norm h_2^2$ of $I_2$,
which the paper drops when it passes to \texttt{eq:background-coercivity}; it is
non-negative and is discarded by the caller.  The Lean $M$ is any upper bound for
$\abs{\im\zeta}$ on $\cG$ (field \texttt{imBound}), the paper's $M$ the maximum.

\begin{proof}
  \leanok
  \uses{lem:background:weighted-test, lem:background:phase}
  The sketch follows the paper's identity \texttt{eq:background-testing},
  $\re\inner{(z-A^{\rm bg}_{\eps,n})h}{(1+i\beta_n)h}=I_1+I_2+I_3+I_4$, whose terms are the
  contributions of Lemma~\ref{lem:background:weighted-test} (the transport term without a
  factor $\beta_n$ has zero real part since $\Div_\perp(u_n^{\rm bg})_\perp=0$).  The
  diffusion term $I_1$ is $\eps\norm{\nabla_\perp h}_2^2$ minus $\eps$ times the defect; the
  Young absorption with \eqref{eq:background:multiplier-bounds} removes $\tfrac\eps4$ of the
  gradient and leaves $C_1\delta_+d(K_0)^2\norm h_2^2$.  The zeroth-order term
  $I_2=\int[\re\zeta+\eps(2\pi K_n)^2+\beta_n(\im\zeta+\tau(s_n-1))]\abs h^2$ collects
  $z=\zeta-i\tau$ and the added axial velocities; $\re\zeta\geq g_*$ and
  Lemma~\ref{lem:background:phase} turn it into the second and third brackets.  The weighted
  transport term $I_3$ obeys $\abs{\beta_n(u_n^{\rm bg})_\perp}\leq C\sqrt\eps\one_{\mathcal C_n}$,
  so a second Young absorption removes another $\tfrac\eps4$ and leaves
  $C\int_{\mathcal C_n}\abs h^2$.  The term $I_4$, stretching plus the local axial multiplier
  $2\pi K_n\int\beta_n\sum_{m\neq n}(u_m^{\rm loc})_3\abs h^2$, has coefficients bounded by
  $C\one_{\mathcal C_n}$ and contributes the same way.  For the last display test with
  $\phi=(1+i\beta_n)h$ and use $\norm{(1+i\beta_n)h}_2\leq\sqrt2\norm h_2$.
\end{proof}

\section{The resolvent}

\begin{proposition}[Uniform background resolvent on $\T^2$]
  \label{prop:background-resolvent}
  \lean{FastDynamo.ScaleFamily.TorusCellPlacement.BackgroundParameters.exists_hasResolventOn_shiftedRegion,
    FastDynamo.ScaleFamily.TorusCellPlacement.BackgroundParameters.resolvent,
    FastDynamo.ScaleFamily.TorusCellPlacement.BackgroundParameters.boundaryResolvent,
    FastDynamo.ScaleFamily.TorusCellPlacement.BackgroundParameters.norm_resolvent_le,
    FastDynamo.ScaleFamily.TorusCellPlacement.BackgroundParameters.sqrt_eps_mul_sqrt_gradientEnergy_resolvent_le,
    FastDynamo.ScaleFamily.TorusCellPlacement.BackgroundParameters.sqrt_eps_mul_norm_deriv_resolvent_le,
    FastDynamo.TorusBackgroundResolventData,
    FastDynamo.ScaleFamily.TorusCellPlacement.BackgroundParameters.backgroundResolventData,
    FastDynamo.ScaleFamily.TorusCellPlacement.BackgroundParameters.boundaryResolventData,
    FastDynamo.DomainOperator.analyticOnNhd_inverse_pencil,
    FastDynamo.DomainOperator.resolventSetResolvent,
    FastDynamo.DomainOperator.analyticOnNhd_resolventSetResolvent_inv,
    FastDynamo.ScaleFamily.TorusCellPlacement.BackgroundParameters.exists_isOpen_hasResolventOn_analyticOnNhd}
  \leanok
  \uses{lem:background:coercivity, lem:torus-sector-basic, lem:operator:shifted-contour,
    def:transfer:parameter-choice}
  There is $\tau_0>0$, depending only on $v$, $I$, $\cG$, $g_*$ and $q$,
  with the following property.  Fix $\tau\geq\tau_0$.  For all sufficiently large $K_0$, every $n\geq0$ and
  every $\eps=\delta\ell_n^2$ with $\delta\in I$ satisfy
  $\cG_\tau\subset\rho(A^{\rm bg}_{\eps,n})$, and for $z\in\cG_\tau$ and $f\in\Lp^2(\T^2;\CC^2)$,
  \begin{equation}
    \norm{(z-A^{\rm bg}_{\eps,n})^{-1}f}_2
      +\sqrt\eps\,\norm{\nabla_\perp(z-A^{\rm bg}_{\eps,n})^{-1}f}_2\leq C\norm f_2,
    \label{eq:background:bg-resolvent}
  \end{equation}
  with $C$ independent of $n$, $\delta$, $z$ and the final enlargement of $K_0$.
\end{proposition}

\textbf{The route is the paper's.}  The paper uses the coercivity only for
\emph{injectivity} of the pencil and takes surjectivity from Lemma~5.1 of the paper
(\texttt{lem:torus-sector-basic}): $z-A^{\rm bg}_{\eps,n}:\Hs^2\to\Lp^2$ has index zero, so
injective implies bijective, and the inverse lands in $\Hs^2$ by construction.  That is
exactly the Lean chain -- \texttt{BackgroundCoercivity.pencil\_injective}, then
\texttt{DomainOperator.hasResolventAt\_of\_compactStateResolvent\_of\_pencil\_injective},
then \texttt{hasResolventOnOfPointwiseTwoSidedInverses}.  Neither side needs an
elliptic-regularity theorem or any $\Hs^1$ multiplier theory.

\textbf{Additional clauses: continuity and holomorphy.}  The paper's statement says nothing about
the dependence of the resolvent on $z$.  The Lean resolvent of the estimate is a
\texttt{HasResolventOn} family, continuous on $\cG_\tau$.  In addition, and not stated by the
paper, Lean proves holomorphy: \texttt{exists\_isOpen\_hasResolventOn\_analyticOnNhd}
(\texttt{FastDynamo/Global/TorusBackgroundResolventHolomorphic.lean}): there is an open
$U\supset\cG_\tau$ carrying a two-sided resolvent $(z-A^{\rm bg}_{\eps,n})^{-1}:\Lp^2\to\Hs^2$,
analytic on $U$ in operator norm, that agrees on $\cG_\tau$ with the resolvent of the
estimate.  The set $U$ is the whole resolvent set, which is open, and the resolvent is analytic
on it by the local Neumann series
(\texttt{DomainOperator.analyticOnNhd\_inverse\_pencil},
\texttt{FastDynamo/Operator/ResolventSetHolomorphy.lean}).  The contour integral of Section~7
consumes only the continuity of $z\mapsto R_{\rm bg}(z)$ on $\cG_\tau$ that
\texttt{HasResolventOn} records.

\textbf{Deviation: explicit thresholds.}  Where the paper says ``there exists $\tau_0$'' and
``for all sufficiently large $K_0$'', the Lean proposition is stated for a
\texttt{BackgroundParameters} record whose fields are the explicit choices:
$\tau\geq\tau_0=\max\bigl(2M^2/g_*,\,2(C_0+1)/c_q^2\bigr)$
(\texttt{backgroundTauThreshold}, which gives both inequalities of
\texttt{eq:axial-separation-choice}) and $C_1\delta_+\ell_0\leq g_*/4$ (the paper's
\texttt{eq:background-diffusion-small}, satisfiable by enlarging $K_0$ since $\ell_0=1/K_0$),
with $C_0$, $C_1$ the closed constants above.  The explicit resolvent constant is
$C=\sqrt2/g+2/\sqrt g$ at $g=g_*/2$.

\begin{proof}
  \leanok
  \uses{lem:background:coercivity, lem:torus-sector-basic}
  Choose $\tau$ and then $K_0$ by \eqref{eq:background:axial-separation-choice}, so that
  the two brackets of \eqref{eq:background:coercivity} are at least $g_*/2$ and $1$.
  Injectivity at $z\in\cG_\tau$ is then immediate from
  $\tfrac{g_*}2\norm h_2^2\leq\re\mathfrak a_z(h,(1+i\beta_n)h)$.  By
  Lemma~\ref{lem:torus-sector-basic} the pencil is Fredholm of index zero, so it is
  bijective onto $\Lp^2$ with bounded inverse into $\Hs^2$; the pointwise inverses are
  continuous in $z$ because the pencil is affine in $z$, which assembles them into one
  family on $\cG_\tau$.  The bound follows from
  $\tfrac\eps2\norm{\nabla_\perp h}_2^2+\tfrac{g_*}2\norm h_2^2\leq\sqrt2\norm f_2\norm h_2$:
  the second term gives $\norm h_2\leq(\sqrt2/g)\norm f_2$ and, feeding that back,
  $\eps\norm{\nabla_\perp h}_2^2\leq(4/g)\norm f_2^2$.  For the additional holomorphy clause:
  the resolvent set of the pencil is open and the inverse is analytic on it, by the Neumann
  series around each of its points, and it contains $\cG_\tau$.
\end{proof}

\begin{corollary}[Resolvent on $\R^2$ for the selected cell]
  \label{cor:scaled-local-resolvents}
  \lean{FastDynamo.WholePlaneLocalResolventBound,
    FastDynamo.scaledLocalResolventInv, FastDynamo.scaledLocalResolvent,
    FastDynamo.norm_embedding_scaledLocalResolventInv_le,
    FastDynamo.sqrt_epsilon_mul_norm_deriv_scaledLocalResolventInv_le,
    FastDynamo.sqrt_epsilon_mul_norm_deriv_le_of_delta_le,
    FastDynamo.scaledLocalResolventOnShiftedContour,
    FastDynamo.LocalSeedC2.scaledResolventBoundData,
    FastDynamo.LocalSeedC2.scaledResolvent}
  \leanok
  \uses{lem:exact-renormalization, prop:local-seed, lem:operator:shifted-contour}
  Let $\eps=\delta_n\ell_n^2$ with $\delta_n\in I$.  For $z\in\Gamma_\tau$ the resolvent
  $\cR_n(z)=(z-A^{(n)}_{\eps,n})^{-1}$ exists and
  \begin{equation}
    \norm{\cR_n(z)f}_2+\sqrt\eps\,\norm{\nabla_\perp\cR_n(z)f}_2\leq C\norm f_2,
    \qquad f\in\Lp^2(\R^2;\CC^2),
    \label{eq:background:local-gradient}
  \end{equation}
  with $C$ independent of $n$, $\delta_n$ and $z$.
\end{corollary}

The Lean statement is per transverse coordinate,
$\sqrt\eps\norm{\partial_j\cR_n(z)f}_2\leq\sqrt{\delta_+}L\norm f_2$ for $j=0,1$; the
gradient-energy form that Section~7 consumes costs one extra $\sqrt2$, and the transfer
chapter
builds its constant as $\sqrt2\sqrt{\delta_+}L$ accordingly.

\begin{proof}
  \leanok
  \uses{lem:exact-renormalization, prop:local-seed}
  By Lemma~\ref{lem:exact-renormalization},
  $\cR_n(z)=U_n(z+i\tau-\cL_{\delta_n})^{-1}U_n^{-1}$, and $z\in\Gamma_\tau$ is exactly
  $z+i\tau\in\Gamma$, so Proposition~\ref{prop:local-seed} applies.  $U_n$ is unitary on
  $\Lp^2$, so the first term is bounded by the seed's own constant $L$; and
  $\norm{\partial_jU_ng}_2=\ell_n^{-1}\norm{\partial_jg}_2$, so the gradient term picks up
  $\sqrt\eps/\ell_n=\sqrt{\delta_n}\leq\sqrt{\delta_+}$.
\end{proof}


\chapter{An approximate inverse for the horizontal problem on $\T^2$ and an unstable eigenvalue}
% Chapter: the approximate inverse for the horizontal problem on T^2 and the unstable
% eigenvalue.  Section 7 of the paper (sec:parametrix).
% Owned shared labels: def:transfer:parameter-choice, def:transfer:structural-data.
% Paper labels: prop:global-spectral-transfer (Proposition 7.1), lem:parametrix-error
% (Lemma 7.2), lem:global-contour-resolvent (Lemma 7.3), lem:global-riesz-nonzero (Lemma 7.4),
% rem:flexible-buffer-scales (Remark 7.5, no node; see the end of this chapter's first section).

This chapter transfers the instability from the selected cell to the torus.  The
parametrix has \emph{two} summands: the background resolvent of the
previous chapter carries everything away from the selected cell, and the rescaled seed
resolvent carries the cell itself.  The error is a sum of two diffusion commutators, both
supported in annuli of width $\asymp R_n$ where the horizontal velocity vanishes, so it is
$O(\sqrt\eps/R_n+\eps/R_n^2)=O(\sqrt{\ell_n}+\ell_n)$ and is made small by one last
enlargement of the base index.  The fixed test vector $f_*$ of the seed, rescaled and
placed, then separates the global Riesz projection from zero.

\section{The ordered choice of parameters}

\begin{definition}[The global parameter chain]
  \label{def:transfer:parameter-choice}
  \lean{FastDynamo.GlobalSeedInterface,
    FastDynamo.GlobalSeedInterface.axialShift,
    FastDynamo.GlobalSeedInterface.errorConstant,
    FastDynamo.GlobalSeedInterface.approximateInverseConstant,
    FastDynamo.GlobalSeedInterface.errorTarget,
    FastDynamo.GlobalSeedInterface.mul_errorTarget_div_lt_quarter,
    FastDynamo.GlobalParameterChoice,
    FastDynamo.GlobalParameterChoice.velocity,
    FastDynamo.GlobalParameterChoice.jointPlacement,
    FastDynamo.GlobalParameterChoice.eps0,
    FastDynamo.GlobalParameterChoice.eps0_formula,
    FastDynamo.GlobalParameterChoice.matchingCell,
    FastDynamo.GlobalParameterChoice.axialAmplitude_match,
    FastDynamo.GlobalParameterChoice.errorSmall_cell,
    FastDynamo.exists_globalParameterChoice,
    FastDynamo.nonempty_globalParameterChoice,
    FastDynamo.LocalSeedC2.toGlobalSeedInterface,
    FastDynamo.HorizontalMode.globalSeedInterface,
    FastDynamo.HorizontalMode.mainParameterChoice,
    FastDynamo.HorizontalMode.mainParameterChoice_eps0_formula,
    FastDynamo.HorizontalMode.kStar_le_mainParameterChoice_baseIndex}
  \leanok
  \uses{def:local:seed, prop:local-seed, def:velocity:scale-family, def:velocity:placement}
  The parameters are fixed in the order
  \[
    (v,I,\Gamma,\cG,g_*,f_*,L_*,q,\vartheta)\ \longrightarrow\ \tau
    \ \longrightarrow\ K_0 ,
  \]
  each stage depending only on the previous ones.  (The paper states this order in prose,
  in Section~3.1 and in the paragraph after the proof of Proposition~7.1; it does not display
  it.)  The seed
  data are Proposition~\ref{prop:local-seed}; the single geometric ratio satisfies
  $1<q<\sqrt{\delta_+/\delta_-}$ and $\vartheta$ is the fixed axial cut-off; the axial shift
  is fixed by \eqref{eq:background:axial-separation-choice}; and the base index $K_0$ is
  enlarged once, to satisfy simultaneously the five conditions: the ratio bound
  $q+K_0^{-1}<\sqrt{\delta_+/\delta_-}$, the chart packing, the core containment
  $R_v\ell_n<R_n/10$ (for which the paper gives the sufficient threshold
  $K_0>(10R_v)^2$), \eqref{eq:background:axial-separation-choice}'s diffusion clause and
  \eqref{eq:transfer:error-small}.  One final $K_0$ fixes the velocity $u$ and the threshold
  $\eps_0=\delta_+k_*^2/K_0^2$ ($=\delta_+/K_0^2$ in the paper, where $k_*=1$); only the matching index $n=n(\eps)$ depends on $\eps$
  afterwards.  The witness-support bound \eqref{eq:transfer:scaled-witness-support} is
  \emph{not} a sixth threshold: it follows from the core containment because $R_v\geq L_*$
  (Lemma~\ref{lem:transfer:witness}).
\end{definition}

In Lean the chain is a record \texttt{GlobalParameterChoice S errorConstant errorTarget}
over a \texttt{GlobalSeedInterface} $S$, and
\texttt{exists\_globalParameterChoice} inhabits it for every seed interface $S$, every
$\texttt{errorConstant}\geq0$ and every $\texttt{errorTarget}>0$: the
$K_0$-thresholds are collected into a single monotone funnel of the shape ``$\ell_0<t$ or
$\sqrt{\ell_0}<t$'' and discharged at once.  The record has exactly four fields --
\texttt{ratios} (the single ratio $q$), \texttt{jointPlacement} (the chart packing and the
core containment, both read at $R_v=\max(1+h,L_*)$), \texttt{diffusionSmall} and
\texttt{errorSmall} -- matching the paper's five conditions, since the two placement
conditions travel together and the integer-ratio bound is a field of the
\texttt{ScaleFamily} inside the placement.  The axial shift
$\tau=\texttt{axialShift}$ is a closed term in $\tau$-free and $K_0$-free data, which is
what makes the order non-circular.  Note that the
interface record itself stores $L_*$ and $L$ as bare reals and carries no seed facts, so
\texttt{nonempty\_globalParameterChoice} on its own is a \emph{shape} statement; the
chain is instantiated only at \texttt{LocalSeedC2.toGlobalSeedInterface}, and every seed
fact the transfer consumes ($f_*$, its norm, its support, the resolvents, the projection
bound) is taken from the \texttt{LocalSeedC2} record, not from the interface.  The chosen
chain is \texttt{HorizontalMode.mainParameterChoice}.

One point about the shape of the dependency graph here.  The record
\texttt{GlobalParameterChoice S errorConstant errorTarget} takes the Section~6 and
Section~7 constants
as \emph{arguments}: $\tau$ is a closed term in the seed data and the ratio, and the
thresholds \eqref{eq:background:axial-separation-choice} and \eqref{eq:transfer:error-small}
enter as fields whose numerical content the caller supplies.  That is exactly the paper's
quantifier order, and it is why this node depends
only on the seed and on the placement in the graph, while the \emph{instantiation}
\texttt{mainParameterChoice} additionally feeds it
\texttt{GlobalSeedInterface.errorConstant} and \texttt{errorTarget}, closed terms built from
the constants of Proposition~\ref{prop:background-resolvent},
Corollary~\ref{cor:scaled-local-resolvents} and Lemma~\ref{lem:parametrix-error}.  Drawing
those two extra edges would make the graph cyclic; the paper avoids the circle in the
same way, by fixing the constants before the base index.

Remark~7.5 of the paper (\texttt{rem:flexible-buffer-scales}), which observes that the
arguments of Sections~6 and~7 also work for other non-increasing buffer radii $R_n$ with
$\sum_nR_n\to0$ and $\sup_n\ell_n/R_n\to0$ as $K_0\to\infty$, has no node and is not
formalised: the Lean fixes $R_n=\sqrt{\ell_n}$ (\texttt{ScaleFamily.bufferRadius}), the
paper's own choice.

\section{The two-resolvent approximate inverse}

\begin{definition}[The three cut-offs and the chart transports]
  \label{def:transfer:structural-data}
  \lean{FastDynamo.ParametrixCutoffs.backgroundCutoff,
    FastDynamo.ParametrixExplicitCutoff.inputCutoff,
    FastDynamo.ParametrixExplicitCutoff.outputCutoff,
    FastDynamo.ParametrixCutoffs.norm_cutoff_coordinateDerivative_sum_le,
    FastDynamo.ParametrixCutoffs.norm_cutoff_iteratedFDeriv_two_sum_le,
    FastDynamo.ParametrixCutoffs.backgroundCutoff_mul_one_sub_inputCutoff,
    FastDynamo.ParametrixCutoffs.outputCutoff_mul_inputCutoff,
    FastDynamo.ParametrixCutoffs.cutoff_partition,
    FastDynamo.ParametrixCutoffs.TransverseVecSecondOrderCutoff,
    FastDynamo.ParametrixCutoffs.HorizontalParametrixCutoffs,
    FastDynamo.ParametrixCutoffs.HorizontalParametrixCutoffs.partition,
    FastDynamo.TorusPlacedCutoffC2.placedCutoffs,
    FastDynamo.ScaleFamily.TorusCellPlacement.placedParametrixCutoffs,
    FastDynamo.TorusChartC2Transfer.restrictInputC2,
    FastDynamo.TorusChartC2Transfer.extendStateC2,
    FastDynamo.TorusChartC2Transfer.extendDomainC2,
    FastDynamo.TorusChartC2Transfer.extendStateC2_comp_restrictInputC2}
  \leanok
  \uses{def:velocity:placement, def:transfer:parameter-choice}
  Choose smooth scalar cut-offs with values in $[0,1]$, the \emph{input cut-off} $\phi_n$ and
  the \emph{output cut-offs} $\psi_n$, $\chi_n$,
  \[
    \phi_n=1\text{ on }D(p_n,R_n/5),\quad\supp\phi_n\subset D(p_n,R_n/4),
    \qquad
    \psi_n=1\text{ on }D(p_n,R_n/3),\quad\supp\psi_n\subset D(p_n,R_n/2),
  \]
  \[
    \chi_n=0\text{ on }D(p_n,R_n/10),\qquad\chi_n=1\text{ outside }D(p_n,R_n/6),
  \]
  obtained by rescaling fixed radial cut-offs, so that
  \begin{equation}
    \norm{D^j\psi_n}_\infty+\norm{D^j\chi_n}_\infty\leq CR_n^{-j},
    \qquad j=1,2,
    \label{eq:transfer:cutoff-bounds}
  \end{equation}
  and
  \begin{equation}
    \chi_n(1-\phi_n)=1-\phi_n,\qquad\psi_n\phi_n=\phi_n .
    \label{eq:transfer:cutoff-identities}
  \end{equation}
\end{definition}

Only the input cut-off $\phi_n$ splits the forcing, so only the two output cut-offs need
derivative bounds; the paper's \texttt{eq:parametrix-cutoff-bounds} accordingly carries no
summand $\norm{D^j\phi_n}_\infty$.  Lean matches
this: the consumed interface \texttt{HorizontalParametrixCutoffs} stores, for $\phi_n$, only
the two sup-norm bounds $\norm{\phi_n}_\infty\leq1$ and $\norm{1-\phi_n}_\infty\leq1$, and
the four derivative bounds only for $\psi_n$ and $\chi_n$.  The auxiliary lemmas
\texttt{norm\_cutoff\_coordinateDerivative\_sum\_le} and
\texttt{norm\_cutoff\_iteratedFDeriv\_two\_sum\_le} do bound all three cut-offs in a single
three-term sum -- that is how the explicit construction produces them -- but the $\phi_n$
summand is never read: the placed construction extracts the $\psi_n$ and $\chi_n$ halves,
which is exactly the paper's display.

In Lean a cut-off is packaged as a \texttt{TransverseVecSecondOrderCutoff}: the
multiplication operators by the value, by each first derivative and by the Laplacian, with
their $\Lp^\infty$ bounds and the identity that the value operator maps $\Hs^2$ into
$\Hs^2$.  \texttt{HorizontalParametrixCutoffs R} bundles the three of them together with
the operator form $\chi_n(1-\phi_n)+\psi_n\phi_n=\id$ of
\eqref{eq:transfer:cutoff-identities}.  The concrete placed triple is built on the plane
and transported to the torus through \texttt{placedCutoffs}, which needs the no-wrap
statement that the closed ball of radius $R_n/2$ around $p_n$ lies in one fundamental box;
the chart transports $\texttt{restrictInputC2}$ (torus $\to$ plane, multiply by $\phi_n$,
extend by zero) and $\texttt{extendStateC2}$ (plane $\to$ torus, multiply by $\psi_n$) are
the two maps that make the selected summand well defined.

\begin{definition}[The approximate inverse]
  \label{def:transfer:approximate-inverse}
  \lean{FastDynamo.TorusTwoSummandParametrix.twoSummandParametrix,
    FastDynamo.TorusTwoSummandParametrix.twoSummandError,
    FastDynamo.TorusTwoSummandParametrix.BackgroundParametrixData,
    FastDynamo.TorusTwoSummandParametrix.BackgroundParametrixData.ofTorusBackgroundResolventData,
    FastDynamo.TorusTwoSummandParametrix.SelectedLocalResolventData,
    FastDynamo.TorusTwoSummandParametrix.TwoSummandParametrixData,
    FastDynamo.TorusTwoSummandParametrix.TwoSummandParametrixData.approximateInverse,
    FastDynamo.TorusTwoSummandParametrix.TwoSummandParametrixData.continuousOn_approximateInverse,
    FastDynamo.TorusChartC2Transfer.selectedLocalResolventData,
    FastDynamo.TorusChartC2Transfer.selectedCellParametrixData,
    FastDynamo.TorusParametrixError.norm_stateParametrix_apply_le}
  \leanok
  \uses{def:transfer:structural-data, prop:background-resolvent,
    cor:scaled-local-resolvents}
  With $\cR_{\rm bg}(z)=(z-A^{\rm bg}_{\eps,n})^{-1}$ and
  $\cR_n(z)=(z-A^{(n)}_{\eps,n})^{-1}$ -- the first an abbreviation the paper introduces in
  running text just before \texttt{eq:Rapp}, the second the resolvent of Corollary~6.2 of
  the paper -- put
  \begin{equation}
    \cR_{\rm app}(z)=\chi_n\cR_{\rm bg}(z)(1-\phi_n)+\psi_n\cR_n(z)\phi_n,
    \qquad z\in\Gamma_\tau .
    \label{eq:transfer:Rapp}
  \end{equation}
  For the second summand, $\phi_nf$ is extended by zero from the fixed chart to $\R^2$, the
  whole-plane resolvent is applied, and the product with $\psi_n$ is restricted to the
  torus.  Both summands map $\Lp^2(\T^2;\CC^2)$ boundedly into $\Hs^2(\T^2;\CC^2)$, and
  \eqref{eq:background:bg-resolvent}, \eqref{eq:background:local-gradient} give
  \begin{equation}
    \sup_{z\in\Gamma_\tau}\norm{\cR_{\rm app}(z)}_{2\to2}\leq C_{\rm app},
    \label{eq:transfer:parametrix-uniform-bound}
  \end{equation}
  uniformly in $n$, in $\delta_n\in I$ and in the final enlargement of $K_0$.  The
  $\Lp^2\to\Hs^2$ norm need not be uniform.
\end{definition}

The Lean \texttt{TwoSummandParametrixData} is an interface: the two resolvent data, the
cut-off triple, the chart transports and the hypotheses that make the commutator
computation valid.  Its instance at the selected cell is
\texttt{selectedCellParametrixData}.  The uniform bound is assembled as
$C_{\rm app}=C_{\rm bg}+C_{\rm loc}$
(\texttt{TorusParametrixError.norm\_stateParametrix\_apply\_le}), using the sharp
$\norm{1-\phi_n}_{\Lp^\infty}\leq1$, which holds because $0\leq\phi_n\leq1$.

\begin{lemma}[The commutator of the operator with a cut-off]
  \label{lem:transfer:commutator}
  \lean{FastDynamo.HorizontalCutoffCommutator.commutator,
    FastDynamo.HorizontalCutoffCommutator.diffusion_commutator_apply,
    FastDynamo.HorizontalCutoffCommutator.axialTransport_commutes,
    FastDynamo.HorizontalCutoffCommutator.stretching_commutes,
    FastDynamo.HorizontalCutoffCommutator.commutator_apply,
    FastDynamo.HorizontalCutoffCommutator.VelocityVanishesOnGradient,
    FastDynamo.HorizontalCutoffCommutator.commutator_apply_of_velocityVanishesOnGradient,
    FastDynamo.HorizontalCutoffCommutator.norm_commutator_apply_le_explicitScale}
  \leanok
  \uses{def:operator:horizontal-induction, def:transfer:structural-data}
  Write $[A_{\eps,K_n},\chi]=A_{\eps,K_n}\chi-\chi A_{\eps,K_n}$.  For a scalar cut-off $\chi$
  whose gradient is supported where $u_\perp$ vanishes,
  \[
    [A_{\eps,K_n},\chi]h
      =\eps\bigl(2\nabla_\perp\chi\cdot\nabla_\perp h+(\Delta_\perp\chi)h\bigr),
    \qquad
    \norm{[A_{\eps,K_n},\chi]h}_2\leq C\Bigl(\frac\eps{R_n}\norm{\nabla_\perp h}_2
      +\frac\eps{R_n^2}\norm h_2\Bigr).
  \]
\end{lemma}

\begin{proof}
  \leanok
  The axial multiplier and the stretching term are of order zero and commute with
  multiplication by the scalar $\chi$.  The transport term contributes
  $(u_\perp\cdot\nabla_\perp\chi)h$, which vanishes because $u_\perp=0$ on the support of
  $\nabla_\perp\chi$.  The commutator is therefore that of $\eps\Delta_\perp$, which the
  Leibniz rule turns into the displayed identity, and the bound follows from
  $\norm{\nabla_\perp\chi}_{\Lp^\infty}\leq C/R_n$ and
  $\norm{\Delta_\perp\chi}_{\Lp^\infty}\leq C/R_n^2$.
\end{proof}

The axial multiplier and the stretching term are zeroth order, so they commute with a
scalar cut-off \emph{exactly}; this is proved in Lean rather than asserted.  The transport
term is killed by the hypothesis \texttt{VelocityVanishesOnGradient}, which is the only
geometric input and is discharged for $\chi_n$ (gradient in $R_n/10\leq\abs{y-p_n}\leq
R_n/6$) and for $\psi_n$ (gradient in $R_n/3\leq\abs{y-p_n}\leq R_n/2$) from the support
fit $\supp u_n^{\rm loc}\subset D(p_n,R_n/10)$, in one statement covering both.  The paper
proves this commutator identity and bound inside the proof of Lemma~7.2, for
$\chi\in\{\chi_n,\psi_n\}$.

\begin{lemma}[Identity and error for the approximate inverse]
  \label{lem:parametrix-error}
  \lean{FastDynamo.TorusTwoSummandParametrix.TwoSummandParametrixData.right_error_identity,
    FastDynamo.TorusTwoSummandParametrix.TwoSummandParametrixData.error_eq,
    FastDynamo.TorusTwoSummandParametrix.TwoSummandParametrixData.transfer_rightError_eq,
    FastDynamo.TorusParametrixError.scaleErrorMajorant,
    FastDynamo.TorusParametrixError.commutatorErrorConstant,
    FastDynamo.TorusParametrixError.parametrixErrorConstant,
    FastDynamo.TorusParametrixError.norm_error_apply_le,
    FastDynamo.TorusParametrixError.norm_error_le,
    FastDynamo.TorusParametrixError.scaleErrorMajorant_scaled_le}
  \leanok
  \uses{def:transfer:approximate-inverse, lem:transfer:commutator,
    prop:background-resolvent, cor:scaled-local-resolvents}
  For $z\in\Gamma_\tau$,
  \begin{equation}
    (z-A_{\eps,K_n})\cR_{\rm app}(z)=\id+E_n(z),
    \qquad
    E_n(z)=-[A_{\eps,K_n},\chi_n]\cR_{\rm bg}(z)(1-\phi_n)
      -[A_{\eps,K_n},\psi_n]\cR_n(z)\phi_n,
    \label{eq:transfer:parametrix}
  \end{equation}
  and
  \begin{equation}
    \norm{E_n(z)}_{2\to2}\leq C\Bigl(\frac{\sqrt\eps}{R_n}+\frac\eps{R_n^2}\Bigr)
      \leq C(\sqrt{\ell_n}+\ell_n),
    \label{eq:transfer:explicit-error}
  \end{equation}
  with $C$ independent of $n$, of $\delta_n\in I$, of $z$ and of the final enlargement of
  $K_0$.
\end{lemma}

The paper gives the formula for $E_n(z)$ in the proof; the statement above includes it.
The Lean constant is explicit: $C=(2\sqrt2C_1+2C_2)(C_{\rm bg}+C_{\rm loc})$, with $C_1$,
$C_2$ the first- and second-derivative constants of the rescaled cut-offs, and the scaling
identity $\eps=\delta\ell^2$, $R=\sqrt\ell$ turns $C(\sqrt\eps/R+\eps/R^2)$ into
$C(1+\delta_+)(\sqrt\ell+\ell)$.

\begin{proof}
  \leanok
  \uses{lem:transfer:commutator, def:transfer:approximate-inverse}
  The full and background coefficients agree wherever $\chi_n\neq0$, because
  $\supp u_n^{\rm loc}\subset D(p_n,R_n/10)$; the full and selected whole-plane
  coefficients agree on $\supp\psi_n$, because the added axial velocity is constant on
  $D(p_n,3R_n/4)$ and all other cells are disjoint from the buffer disc.  So
  \eqref{eq:transfer:cutoff-identities} gives
  $(z-A_{\eps,K_n})\chi_n\cR_{\rm bg}(1-\phi_n)
  =1-\phi_n-[A_{\eps,K_n},\chi_n]\cR_{\rm bg}(1-\phi_n)$ and
  $(z-A_{\eps,K_n})\psi_n\cR_n\phi_n=\phi_n-[A_{\eps,K_n},\psi_n]\cR_n\phi_n$; adding gives
  \eqref{eq:transfer:parametrix}.  Lemma~\ref{lem:transfer:commutator} bounds each
  commutator by $C(\eps/R_n)\norm{\nabla_\perp h}_2+C(\eps/R_n^2)\norm h_2$, and the two
  resolvent estimates \eqref{eq:background:bg-resolvent},
  \eqref{eq:background:local-gradient} convert
  $\eps\norm{\nabla_\perp h}_2\leq\sqrt\eps\cdot C\norm f_2$ and
  $\norm h_2\leq C\norm f_2$.
\end{proof}

\section{The global resolvent on the contour}

\begin{lemma}[The resolvent of the horizontal operator on $\T^2$ along the shifted contour]
  \label{lem:global-contour-resolvent}
  \lean{FastDynamo.TorusParametrixError.error_norm_lt_one_of_majorant,
    FastDynamo.TorusParametrixError.exists_scale_threshold,
    FastDynamo.TorusParametrixError.globalContourResolvent,
    FastDynamo.TorusParametrixError.globalContourResolvent_inv,
    FastDynamo.TorusChartC2Transfer.selectedCellGlobalContourResolvent,
    FastDynamo.TorusGlobalRieszNonzero.cellContourResolvent,
    FastDynamo.HorizontalMode.mainSelectedCellResolvent,
    FastDynamo.GlobalParameterChoice.scaleErrorMajorant_lt_errorTarget}
  \leanok
  \uses{lem:parametrix-error, lem:torus-sector-basic, def:transfer:parameter-choice}
  For all sufficiently large $K_0$, uniformly in $n\geq0$ and in
  $\eps=\delta_n\ell_n^2$ with $\delta_n\in I$, one has
  $\Gamma_\tau\subset\rho(A_{\eps,K_n})$ and
  \begin{equation}
    (z-A_{\eps,K_n})^{-1}=\cR_{\rm app}(z)\bigl(\id+E_n(z)\bigr)^{-1},
    \qquad z\in\Gamma_\tau .
    \label{eq:transfer:global-resolvent}
  \end{equation}
  The threshold is: choose $0<\eta_*<1/2$ with
  $\tfrac{\operatorname{length}(\Gamma)}{2\pi}C_{\rm app}\tfrac{\eta_*}{1-\eta_*}<\tfrac14$
  and enlarge $K_0$ until
  \begin{equation}
    \sup_{n\geq0}\sup_{\delta\in I}\sup_{z\in\Gamma_\tau}\norm{E_{n,\delta}(z)}_{2\to2}
      <\eta_* .
    \label{eq:transfer:error-small}
  \end{equation}
\end{lemma}

The paper chooses $\eta_*$ and states \texttt{eq:error-small} in the proof of Lemma~7.3; the
statement above includes them.  The strict inequality of \eqref{eq:transfer:error-small}
is the strict inequality in Lean as well.  The Lean threshold is
$\eta_*=\min\bigl(\tfrac12,(100(1+K))^{-1}\bigr)$ with $K=(1+r_\Gamma)(1+C_{\rm app})$
(\texttt{GlobalSeedInterface.errorTarget}), which gives $A\eta_*/(1-\eta_*)<\tfrac14$ for
every prefactor $0\leq A\leq10K$ (\texttt{mul\_errorTarget\_div\_lt\_quarter}), in
particular for $A=r_\Gamma C_{\rm app}$; the prefactor
$\operatorname{length}(\Gamma)/2\pi$ is exactly the radius $r_\Gamma$.

\begin{proof}
  \leanok
  \uses{lem:parametrix-error, lem:torus-sector-basic}
  Since $\ell_n\leq\ell_0=k_*/K_0$ ($=1/K_0$ in the paper, where $k_*=1$),
  \eqref{eq:transfer:explicit-error} tends to zero
  uniformly as $K_0\to\infty$, so \eqref{eq:transfer:error-small} can be arranged; the
  Neumann series then inverts $\id+E_n(z)$ and \eqref{eq:transfer:parametrix} exhibits
  $\cR_{\rm app}(\id+E_n)^{-1}$ as a right inverse with values in $\Hs^2$.  By
  Lemma~\ref{lem:torus-sector-basic} the pencil $z-A_{\eps,K_n}\colon\Hs^2\to\Lp^2$ is
  Fredholm of index zero, so surjectivity forces injectivity and the right inverse is the
  resolvent; continuity in $z$ is inherited from the two resolvent families and the
  Neumann series.
\end{proof}

\section{The Riesz projection of the horizontal operator on $\T^2$ is non-zero}

\begin{lemma}[The placed witness]
  \label{lem:transfer:witness}
  \lean{FastDynamo.TorusSelectedWitness.VanishesOutsideDisk,
    FastDynamo.TorusSelectedWitness.scaledWitness,
    FastDynamo.TorusSelectedWitness.placedWitness,
    FastDynamo.TorusSelectedWitness.norm_placedWitness,
    FastDynamo.TorusSelectedWitness.inputMul_stateMul_placedWitness,
    FastDynamo.TorusSelectedWitness.outputMul_stateMul_placedWitness,
    FastDynamo.TorusSelectedWitness.restrictInputC2_placedWitness,
    FastDynamo.TorusSelectedWitness.vanishesOutsideDisk_scaledWitness,
    FastDynamo.GlobalParameterChoice.witnessSupport}
  \leanok
  \uses{prop:local-seed, def:transfer:structural-data, def:transfer:parameter-choice}
  Since $L_*\leq R_v$, the core containment $R_v\ell_n<R_n/10$ of
  Definition~\ref{def:velocity:placement} gives
  \begin{equation}
    L_*\ell_n<R_n/5\qquad\text{for every }n\geq0 .
    \label{eq:transfer:scaled-witness-support}
  \end{equation}
  Read $f_n=U_nf_*$ as a torus field by restriction to the fixed chart and extension by
  zero.  Then $\norm{f_n}_2=1$ and $\phi_nf_n=\psi_nf_n=f_n$.
\end{lemma}

\begin{proof}
  \leanok
  The placement radius satisfies $R_v\geq L_*$ by construction
  (\texttt{GlobalSeedInterface.witnessSupportRadius\_le\_placementRadius}), so
  $L_*\ell_n\leq R_v\ell_n<R_n/10<R_n/5$; the support of $f_*$ lies in $D(0,L_*)$, so the
  support of the rescaled and placed $f_n$ lies in $D(p_n,L_*\ell_n)\subset D(p_n,R_n/5)$,
  where both $\phi_n$ and $\psi_n$ equal one, and unitarity of $U_n$ gives $\norm{f_n}_2=1$.
\end{proof}

The radius $R_n/5$ in \eqref{eq:transfer:scaled-witness-support} is exactly the plateau
radius of the input cut-off, so the constant is not slack.  The paper requires
$R_v\geq L_*$ when $R_v$ is chosen (Section~3.3, just before \texttt{eq:local-support-in-buffer})
and derives the bound, its \texttt{eq:scaled-witness-support}, from the core containment
at the start of the proof of Lemma~7.4 rather than from a further enlargement of $K_0$,
which keeps the list of $K_0$-thresholds at five.
Lean does the same: the placement is chosen once, at support radius
$R_v=\max(1+h,L_*)$, and \texttt{GlobalParameterChoice.witnessSupport} is proved
from the placement's support fit together with $R_n/10<R_n/5$.

\begin{lemma}[The contour identity at the witness]
  \label{lem:transfer:contour-identity}
  \lean{FastDynamo.TorusSelectedContourIdentity.SelectedCellData,
    FastDynamo.TorusSelectedContourIdentity.SelectedCellData.parametrix,
    FastDynamo.TorusSelectedContourIdentity.circleIntegral_translate,
    FastDynamo.TorusSelectedContourIdentity.inner_stateParametrix_eq_of_witness,
    FastDynamo.TorusSelectedContourIdentity.inner_stateParametrix_placedWitness,
    FastDynamo.TorusSelectedContourIdentity.inner_contourOperator_placedWitness,
    FastDynamo.TorusGlobalRieszNonzero.re_inner_contourOperator_placedWitness_gt_half}
  \leanok
  \uses{lem:transfer:witness, def:transfer:approximate-inverse, lem:exact-renormalization,
    prop:local-seed}
  For $z\in\Gamma_\tau$,
  $\inner{\cR_{\rm app}(z)f_n}{f_n}=\inner{(z+i\tau-\cL_{\delta_n})^{-1}f_*}{f_*}$, and
  therefore
  \begin{equation}
    \frac1{2\pi i}\int_{\Gamma_\tau}\inner{\cR_{\rm app}(z)f_n}{f_n}\dd z
      =\inner{P_{\delta_n}f_*}{f_*} .
    \label{eq:transfer:selected-contour-identity}
  \end{equation}
  In particular $\re$ of the left-hand side exceeds $\tfrac12$.
\end{lemma}

\begin{proof}
  \leanok
  \uses{lem:transfer:witness, lem:exact-renormalization, prop:local-seed}
  The input of the background summand is $(1-\phi_n)f_n=0$, and the output cut-off of the
  selected summand equals one on $\supp f_n$, so only the selected summand survives; the
  scaling identity \eqref{eq:operator:exact-scaling} and the unitarity of $U_n$ turn it into
  the seed resolvent at $z+i\tau$.  Integrating and translating the circle by $+i\tau$ --
  the parametrisation and its derivative are unchanged, so the translation is exact and
  orientation preserving -- gives the seed's Riesz projection, and
  \eqref{eq:local:spectral-witness} gives the bound.
\end{proof}

\begin{lemma}[The Riesz projection of the horizontal operator on $\T^2$ is non-zero]
  \label{lem:global-riesz-nonzero}
  \lean{FastDynamo.TorusGlobalRieszNonzero.norm_stateResolvent_sub_stateParametrix_le_of_inv,
    FastDynamo.TorusGlobalRieszNonzero.norm_rieszProjection_sub_contourOperator_le_cell,
    FastDynamo.TorusGlobalRieszNonzero.re_inner_rieszProjection_placedWitness_gt_quarter,
    FastDynamo.TorusGlobalRieszNonzero.rieszProjection_ne_zero_cell,
    FastDynamo.LocalSeedC2.projectionError_matchingCellData_lt_quarter}
  \leanok
  \uses{lem:transfer:contour-identity, lem:global-contour-resolvent,
    lem:parametrix-error}
  The projection
  $Q_{\eps,n}=(2\pi i)^{-1}\int_{\Gamma_\tau}(z-A_{\eps,K_n})^{-1}\dd z$ satisfies
  \begin{equation}
    \Bigl|\inner{Q_{\eps,n}f_n}{f_n}-\inner{P_{\delta_n}f_*}{f_*}\Bigr|
      \leq\frac{\operatorname{length}(\Gamma)}{2\pi}C_{\rm app}
        \frac{\eta_*}{1-\eta_*}<\frac14,
    \label{eq:transfer:global-projection-error}
  \end{equation}
  hence $\re\inner{Q_{\eps,n}f_n}{f_n}>\tfrac14$ and $Q_{\eps,n}\neq0$, for every $n\geq0$
  and every $\eps=\delta_n\ell_n^2$ with $\delta_n\in I$.
\end{lemma}

The paper's statement of Lemma~7.4 asserts only $Q_{\eps,n}\neq0$; the bound
\eqref{eq:transfer:global-projection-error} (the paper's \texttt{eq:global-projection-error})
and $\re\inner{Q_{\eps,n}f_n}{f_n}>\tfrac14$ are established in its proof, together with
the content of Lemmas~\ref{lem:transfer:witness} and~\ref{lem:transfer:contour-identity}.
In Lean the error is bounded in operator norm,
$\norm{Q_{\eps,n}-(2\pi i)^{-1}\int_{\Gamma_\tau}\cR_{\rm app}(z)\dd z}
\leq r_\Gamma C_{\rm app}e/(1-e)$ with $e<1$ the explicit majorant of
$\sup_{z}\norm{E_n(z)}$ (\texttt{norm\_rieszProjection\_sub\_contourOperator\_le\_cell}),
then $e<\eta_*$ is used and the result is paired with the unit vector $f_n$.

\begin{proof}
  \leanok
  \uses{lem:global-contour-resolvent, lem:transfer:contour-identity}
  From \eqref{eq:transfer:global-resolvent},
  $(z-A_{\eps,K_n})^{-1}-\cR_{\rm app}(z)=-\cR_{\rm app}(z)(\id+E_n(z))^{-1}E_n(z)$, whose
  norm is at most $C_{\rm app}\eta_*/(1-\eta_*)$ by
  \eqref{eq:transfer:parametrix-uniform-bound} and
  \eqref{eq:transfer:error-small}; integrating over a circle of radius $r_\Gamma$ and
  pairing with the unit vector $f_n$ gives
  \eqref{eq:transfer:global-projection-error}.  Subtracting from
  Lemma~\ref{lem:transfer:contour-identity} leaves
  $\re\inner{Q_{\eps,n}f_n}{f_n}>\tfrac12-\tfrac14$.
\end{proof}

\section{The spectral transfer}

\begin{lemma}[Extraction of an eigenpair from a non-zero Riesz projection]
  \label{lem:transfer:eigenpair-extraction}
  \lean{FastDynamo.HorizontalSpectralExtraction.isCompactOperator_rieszProjection_horizontalInduction,
    FastDynamo.HorizontalSpectralExtraction.finiteDimensional_rieszProjection_range_horizontalInduction,
    FastDynamo.HorizontalSpectralExtraction.exists_horizontalInduction_eigenpair_mem_closedRegion_of_rieszProjection_ne_zero,
    FastDynamo.HorizontalSpectralExtraction.exists_growthFloor_horizontalInduction_eigenpair_of_rieszProjection_ne_zero,
    FastDynamo.ParametrixNeumannCorrection.ne_zero_of_inner_contour_test}
  \leanok
  \uses{lem:torus-sector-basic, lem:global-riesz-nonzero,
    lem:operator:shifted-contour}
  If $A_{\eps,K}$ has a resolvent on a spectral circle and the associated Riesz projection
  is non-zero, then $A_{\eps,K}$ has an eigenvalue $\lambda$ in the closed interior of that
  circle, with an eigenfunction in $\Hs^2(\T^2;\CC^2)$; in particular
  $\re\lambda\geq$ the circle's growth floor.
\end{lemma}

\begin{proof}
  \leanok
  \uses{lem:torus-sector-basic}
  The state resolvent is compact at every resolvent point
  (Lemma~\ref{lem:torus-sector-basic}), so $Q$ is a compact idempotent and has
  finite-dimensional range; a non-zero finite-rank idempotent built from a contour integral
  forces a spectral point inside, and index zero turns that spectral point into an
  eigenvalue with an $\Hs^2$ eigenfunction.  The real-part bound is the definition of the
  growth floor, which is invariant under the translation by $-i\tau$.
\end{proof}

\begin{proposition}[Existence of an unstable eigenvalue of the horizontal operator on $\T^2$]
  \label{prop:global-spectral-transfer}
  \lean{FastDynamo.TorusGlobalSpectralTransfer.exists_growing_horizontalInduction_eigenpair_cell,
    FastDynamo.TorusGlobalSpectralTransfer.exists_growing_reducedInduction_eigenpair_cell,
    FastDynamo.LocalSeedC2.matchingCellData,
    FastDynamo.LocalSeedC2.exists_growing_reducedInduction_eigenpair_matchingCell,
    FastDynamo.HorizontalMode.exists_growing_reducedInduction_eigenpair_matchingCell,
    FastDynamo.mainReducedClauses_holds}
  \leanok
  \uses{def:transfer:parameter-choice, lem:global-contour-resolvent,
    lem:global-riesz-nonzero, lem:transfer:eigenpair-extraction, lem:axial-recovery}
  Fix $v,I,\Gamma,\cG,g_*,f_*,L_*$, and choose $q$, the axial cut-off $\vartheta$ and
  $\tau$ in the order of Definition~\ref{def:transfer:parameter-choice}.
  Then $K_0$ can be chosen so that for
  every $0<\eps\leq\eps_0$ and every $n$ with $\delta_n=\eps/\ell_n^2\in I$,
  \[
    \Gamma_\tau\subset\rho(A_{\eps,K_n}),\qquad
    Q_{\eps,n}=\frac1{2\pi i}\int_{\Gamma_\tau}(z-A_{\eps,K_n})^{-1}\dd z\neq0 .
  \]
  In particular $A_{\eps,K_n}$ has an eigenvalue inside $\Gamma_\tau$ with real part at least
  $g_*$:
  \[
    \lambda^\eps=\alpha_\eps+i\omega_\eps,\qquad\alpha_\eps\geq g_* .
  \]
\end{proposition}

The paper's statement is existential (``has an eigenvalue inside $\Gamma_\tau$ with real
part at least $g_*$''), and so is the Lean,
\texttt{exists\_growthFloor\_horizontalInduction\_eigenpair\_of\_rieszProjection\_ne\_zero};
the paper does not name the eigenvalue, and the display above only names it for later use.
The paper's proof also notes the universal form, ``every eigenvalue in this region has real
part at least $g_*$, as translation by $-i\tau$ preserves real parts''.  In Lean that is
one application of \texttt{SpectralCircle.growthFloor\_le\_re\_of\_mem\_ball} -- the very
lemma the existential proof applies -- and nothing downstream needs it, so it is not stated
separately.

In Lean the proposition is used in its \emph{eigenpair} form and immediately composed with
the axial recovery, so the two published endpoints are
\texttt{TorusGlobalSpectralTransfer.exists\_growing\_reducedInduction\_eigenpair\_cell}
(one selected cell, three components) and, at the chosen chain,
\texttt{HorizontalMode.exists\_growing\_reducedInduction\_eigenpair\_matchingCell}, whose
packaging over all $0<\eps\leq\eps_0$ is \texttt{mainReducedClauses\_holds}.  The last of
these is the exact content of this node at the constructed velocity; it is the theorem the
whole development exists to prove.

\begin{proof}
  \leanok
  \uses{lem:global-contour-resolvent, lem:global-riesz-nonzero,
    lem:transfer:eigenpair-extraction, lem:torus-sector-basic}
  Lemma~\ref{lem:global-contour-resolvent} gives the contour resolvent and
  Lemma~\ref{lem:global-riesz-nonzero} the non-zero projection, at every admissible cell
  and every admissible diffusivity.  By Lemma~\ref{lem:torus-sector-basic} the resolvent is
  compact, so the spectrum inside $\Gamma_\tau$ consists of isolated eigenvalues of finite
  algebraic multiplicity, and the non-zero projection shows that it is non-empty;
  Lemma~\ref{lem:transfer:eigenpair-extraction} extracts the eigenpair.  Every eigenvalue in
  the region has $\re\lambda\geq g_*$ because translation by $-i\tau$ preserves real
  parts.
\end{proof}

\begin{lemma}[Every admissible cell]
  \label{lem:transfer:all-cells}
  \lean{FastDynamo.exists_growing_eigenpair_of_matchingCell,
    FastDynamo.exists_growing_eigenpair_of_admissible_index,
    FastDynamo.MainEveryAdmissibleCellClauses,
    FastDynamo.mainEveryAdmissibleCellClauses_holds}
  \leanok
  \uses{prop:global-spectral-transfer, def:transfer:parameter-choice}
  Proposition~\ref{prop:global-spectral-transfer} holds not only at a canonically chosen
  matching cell but at every cell index $n$ with $\eps/\ell_n^2\in I$: the axial index
  $K_n$ is non-zero and carries an eigenpair of $A_{\eps,K_n}$, hence of the
  three-component operator, with $\re\lambda\geq\gamma_0$.
\end{lemma}

This is stronger than Theorem~\ref{thm:main} requires and is recorded as item (4) of
Theorem~\ref{thm:main-strengthened}.  The Lean predicate quantifies over
$0<\eps\leq\eps_0$, although the paper's proof needs no upper bound on $\eps$ once an
admissible index is given; the threshold enters only through the \emph{existence} of such
an index.

\begin{proof}
  \leanok
  \uses{prop:global-spectral-transfer}
  Every step of the chain is stated at an arbitrary admissible index, so nothing has to be
  redone; only the choice of the index is dropped.
\end{proof}


\chapter{Recovery of the magnetic field and the main theorem}
The previous chapters produce, for every small diffusivity $\eps$, an eigenpair
$(\lambda^\eps,h^\eps)$ of the \emph{horizontal} operator $A_{\eps,K}$ on
$\Lp^2(\T^2;\CC^2)$, with $\re\lambda^\eps\geq g_*>0$.  Two components are not a magnetic
field.  This chapter recovers the third one: the divergence constraint determines it
algebraically, $b_3=-\Div_\perp h/(2\pi iK)$, and its own equation supplies the regularity
needed for membership in the operator domain (Lemma~\ref{lem:axial-recovery}).  The
resulting $V^\eps=\e^{2\pi iK_\eps x_3}b^\eps$ is divergence-free \emph{by construction},
which is why the paper's proof never needs a separate solenoidality lemma.  The remaining
two sections are independent of that: the non-zero axial Fourier index lets one
pass to real parts without losing any norm, at any temporal phase, so that the real datum
$B_{\rm in}^\eps=\re V^\eps$ is non-zero and its evolution obeys the exact exponential law;
and a last section records the bounds behind the two-sided $\sqrt\eps$ magnetic scale.

\textbf{Standing conventions of the formalisation.}  Three global conventions apply to every node
of this chapter and are not repeated in each one.  (i) The three-torus is modelled as the product
$\T^2\times\T^1$ (\texttt{ProductTorus3}) with the product measure; a measure-preserving
equivalence with the literal $\T^3=(\R/\Z)^3$ and the induced isometric $\Lp^2$ transport do exist
(\texttt{FastDynamo.productToLiteral}, \texttt{FastDynamo.toLiteralR3L2} in
\texttt{FastDynamo/Analysis/LiteralThreeTorus.lean}), and the real magnetic trajectory and its norm
law are transported along them, but no differential operator is transported, so all operator
statements below live on the product model.  (ii) The velocity is an element of
\texttt{TransverseRealW1Inf}, the Fourier-graph realisation of $\W^{1,\infty}(\T^2;\R^3)$; such a
field is independent of $x_3$ but has a non-trivial third component.  This is exactly the class of
the velocity constructed in the previous chapters, and it is what makes the axial reduction exact.
Axial independence is \emph{not} a narrowing: the only lemma of this section that hypothesises a
velocity at all, Lemma~\ref{lem:axial-recovery}, asks for a divergence-free
$u\in\W^{1,\infty}(\T^3;\R^3)$ \emph{independent of $x_3$} (as in the paper's statement of
Lemma~8.1), which is exactly this class; Lemma~\ref{lem:real-axial-mode} has no velocity, and
Corollary~\ref{cor:magnetic-scale} is about the theorem's own $u$.  What does remain
unformalised is the identification of the Fourier-graph model \texttt{TransverseRealW1Inf} with
the Lipschitz / $\W^{1,\infty}(\T^3)$ class itself (see the scope notes in
\texttt{README.md}).  Divergence-freeness of $u$
is \texttt{IsTransverseDivergenceFree}, the identity $\partial_1u_1+\partial_2u_2=0$ in
$\Lp^\infty$.  (iii) $\Hs^2$ and $\Hs^1$ domains are Fourier graph spaces, and all weak
derivatives, divergences and curls are their Fourier realisations.

\section{Recovery of the axial magnetic component}

\begin{lemma}[Recovery of the axial magnetic component]
  \label{lem:axial-recovery}
  \lean{FastDynamo.axialRecoveredComponent,
    FastDynamo.horizontalEmbedH2, FastDynamo.transverseSnocH2,
    FastDynamo.transverseCoordL2_reducedInductionAction,
    FastDynamo.fixedSectorDivergenceCoeff_reducedInductionAction,
    FastDynamo.axialRecoveredSymbol,
    FastDynamo.transverseH2Weight_le_axialRecoveredSymbol,
    FastDynamo.axialRecoveredCertificateMap,
    FastDynamo.axialRecoveredComponent_coeff_equation,
    FastDynamo.axialRecoveredCertificate_eq,
    FastDynamo.exists_reducedInductionEigenpair_of_horizontalEigenpair}
  \leanok
  \uses{def:operator:horizontal-induction, def:main:induction-operator}
  Let $\eps>0$, $K\in\Z\setminus\{0\}$, $\lambda\in\CC$, and let
  $u\in\W^{1,\infty}(\T^3;\R^3)$ be divergence-free and independent of $x_3$.  Suppose
  $0\neq h\in\Hs^2(\T^2;\CC^2)$ satisfies $A_{\eps,K}h=\lambda h$.  Put
  \[
    b_3=-\frac{\Div_\perp h}{2\pi iK},\qquad b=(h,b_3),
    \qquad V(y,x_3)=\e^{2\pi iKx_3}b(y).
  \]
  Then $b\in\Hs^2(\T^2;\CC^3)$ and
  \[
    0\neq V\in\Hs^2(\T^3;\CC^3)\cap\Lp^2_\sigma(\T^3),\qquad L_{\eps,u}V=\lambda V .
  \]
\end{lemma}

Three points of the Lean form.  (a) The conclusion is packaged as an eigenpair of the
\emph{three}-component reduced operator $A^{(3)}_{\eps,K}$ of
Definition~\ref{def:main:induction-operator} together with
$\pi_\perp b=h$, the formula for the third coordinate, and \emph{both} divergence clauses
$\Div_\perp h+2\pi iKb_3=0$ and $\Div V=0$; the passage to the physical curl form is then
Lemma~\ref{lem:finish:physical-compatibility}.  (b) The identity
\eqref{eq:finish:transport-stretching-divergence} is not re-proved in two dimensions: the
product-torus identity
$\Div[-(w\cdot\nabla)F+(F\cdot\nabla)w]=-w\cdot\nabla\Div F$ is transported into the sector
$K$ along the axial lift, which avoids a second derivative of $w$.  The paper does not
use that identity (see the deviation below).  (c) The elliptic upgrade $b_3\in\Hs^1\Rightarrow b_3\in\Hs^2$ is
\emph{free} at every frequency, with no elliptic-regularity theorem: the multiplier
$(1+4\pi^2\abs n^2)/\bigl(\eps(4\pi^2\abs n^2+(2\pi K)^2)\bigr)$ has modulus at most
$1/\eps$ as soon as $(2\pi K)^2\geq1$, which holds for every $K\neq0$, so the $\Hs^2$
certificate is a bounded Fourier multiplier.

\textbf{Deviation: the scalar route to $b_3$, not the three-component distributional
route.}  The statement above is the paper's (Lemma~8.1, \texttt{lem:axial-recovery}), word for
word.  The proof below is not.  The paper proves the lemma without writing an equation for
$b_3$: for the $\Hs^1$ field $V$, which is divergence-free by construction, it forms
$F=\lambda V-\eps\Delta V-\curl(u\times V)\in\Hs^{-1}(\T^3;\CC^3)$, uses the identity
$\curl(u\times V)=(V\cdot\nabla)u-(u\cdot\nabla)V$ in distributions, observes that $F_\perp=0$ is
the horizontal eigenvalue equation and that $\Div F=0$, so that, every term having axial
Fourier index $K$, $0=\Div F=\partial_{x_3}F_3=2\pi iKF_3$ in distributions forces $F=0$; only
then does it upgrade $V$ to $\Hs^2$ by elliptic regularity from
$\eps\Delta V=\lambda V-\curl(u\times V)\in\Lp^2$.  Lean takes a scalar route instead: the
transport--stretching divergence identity
$\Div[-(w\cdot\nabla)F+(F\cdot\nabla)w]=-w\cdot\nabla\Div F$, transported into the sector
$K$ along the axial lift, turns $A_{\eps,K}h=\lambda h$ into the scalar equation
$(\lambda-S)b_3=(h\cdot\nabla_\perp)u_3$, whose right-hand side is an $\Lp^2$ field;
the $\Hs^2$ certificate for $b_3$ is then a bounded Fourier multiple of it, with no
elliptic-regularity theorem.  The paper's route is not available as stated because Lean's
curl form is realised in $\Hs^{-2}$ against $\Hs^2$ test fields, not in $\Hs^{-1}$ against
$\Hs^1$: forming $F$ at the regularity the paper uses would require a further antidual layer
and the distributional identity $\curl(u\times V)=(V\cdot\nabla)u-(u\cdot\nabla)V$ one
derivative lower.  The two cheap steps of the paper's route -- $\Div F=0\Rightarrow
2\pi iKF_3=0\Rightarrow F=0$ in a single non-zero axial mode, and $\eps\Delta V\in\Lp^2
\Rightarrow V\in\Hs^2$ -- are free in the Fourier model, and are how point (c) above is
proved.  The statement proved is the paper's, including both divergence clauses; only the
route differs.  The paper has no counterpart of the two displays
\eqref{eq:finish:transport-stretching-divergence} and
\eqref{eq:finish:recovered-axial-equation}; their labels are blueprint-local.

\begin{proof}
  \leanok
  \uses{def:operator:horizontal-induction, def:main:induction-operator}
  This is the argument Lean states, not the paper's; see the deviation
  above.  Write $d=\Div_\perp h$ and
  $S=\eps\bigl(\Delta_\perp-(2\pi K)^2\bigr)-u_\perp\cdot\nabla_\perp-2\pi iKu_3$.  For a divergence-free
  $w\in\W^{1,\infty}$ and $F\in\Hs^2$,
  \begin{equation}
    \Div\bigl[-(w\cdot\nabla)F+(F\cdot\nabla)w\bigr]=-w\cdot\nabla\Div F
    \label{eq:finish:transport-stretching-divergence}
  \end{equation}
  in distributions, because the bracket is $\curl(w\times F)-w\Div F$; this needs no second
  derivative of $w$.  Applying it to the horizontal fields and using
  $\Div_\perp(u_3h)=u_3d+h\cdot\nabla_\perp u_3$ turns $A_{\eps,K}h=\lambda h$ into
  $(\lambda-S)d=-2\pi iK(h\cdot\nabla_\perp)u_3$ in $\Hs^{-1}(\T^2)$, that is
  \begin{equation}
    (\lambda-S)b_3=(h\cdot\nabla_\perp)u_3 .
    \label{eq:finish:recovered-axial-equation}
  \end{equation}
  Initially $b_3\in\Hs^1$; rearranging \eqref{eq:finish:recovered-axial-equation} gives
  $\eps\Delta_\perp b_3\in\Lp^2$ and hence $b_3\in\Hs^2$.  The horizontal equation together
  with \eqref{eq:finish:recovered-axial-equation} is exactly
  $\eps\Delta V-(u\cdot\nabla)V+(V\cdot\nabla)u=\lambda V$, and
  $\Div V=\e^{2\pi iKx_3}(d+2\pi iKb_3)=0$, so the left-hand side is $L_{\eps,u}V$.  Finally $h\neq0$
  gives $V\neq0$.
\end{proof}

\section{The eigenfunction is divergence-free}

For the recovered mode, $\Div V=0$ holds by construction (Lemma~\ref{lem:axial-recovery}).
The lemma of this section proves it for every eigenfunction of the auxiliary operator
$\mathcal{A}_\eps$ with $\re\lambda>0$; the Lean proof of
Proposition~\ref{prop:finish:real-strong-mode} takes the solenoidality of the lifted mode
from it, and the dynamics chapter uses its resolvent form
(Lemma~\ref{lem:dynamics:solenoidal-subspace}).  The passage from $\mathcal{A}_\eps$ to the
physical curl form (Lemma~\ref{lem:finish:physical-compatibility}) is used by the axial
recovery and by the dynamics chapter.

\begin{definition}[The physical curl-form operator]
  \label{def:finish:physical-operator}
  \lean{FastDynamo.productL2ToHMinus2, FastDynamo.productWeakCurl,
  FastDynamo.productCrossWithVelocity, FastDynamo.productPhysicalInductionAction,
  FastDynamo.productPhysicalInductionOperator}
  \leanok
  \uses{def:main:induction-operator}
  For $u\in\W^{1,\infty}$ and $\eps\in\R$ the \emph{physical} induction operator is the curl form
  \[
    L_{\eps,u}B=\eps\Delta B+\curl(u\times B),
  \]
  realised as a bounded map from the Fourier graph $\Hs^2(\T^2\times\T^1;\CC^3)$ into a Fourier
  model of $\Hs^{-2}$, in which an $\Lp^2$ coefficient $\hat f(n,k)$ is represented by
  $\hat f(n,k)/(1+\abs\xi^2)$.  In that model the distributional curl of an arbitrary $\Lp^2$
  field is bounded, and $\Div\circ\curl=0$ is an exact identity frequency by frequency.
\end{definition}

In Lean this is \texttt{productPhysicalInductionAction} and the bundled domain operator
\texttt{productPhysicalInductionOperator}, whose ambient space is $\Hs^{-2}$ and whose embedding
is $\Hs^2\hookrightarrow\Lp^2\hookrightarrow\Hs^{-2}$ (module
\texttt{FastDynamo/Operator/Compatibility.lean}).  The paper never names a functional-analytic
realisation of $L_{\eps,u}$; the weak $\Hs^{-2}$ realisation is the formalisation's choice, and
it is weaker than an $\Lp^2$ realisation on solenoidal fields, which does not exist in the
development ($\Lp^2_\sigma$ is a predicate, not a subspace).

\begin{lemma}[Compatibility of the two actions]
  \label{lem:finish:physical-compatibility}
  \lean{FastDynamo.productInductionAction_toHMinus2_eq_physical,
  FastDynamo.productInductionEigenpair_to_physical}
  \leanok
  \uses{def:main:induction-operator, def:finish:physical-operator}
  Let $u$ be divergence-free and let $V\in\Hs^2$ satisfy $\Div V=0$.  Then
  \[
    \curl(u\times V)=-(u\cdot\nabla)V+(V\cdot\nabla)u,
  \]
  so $\mathcal{A}_\eps V$ and $L_{\eps,u}V$ agree in $\Hs^{-2}$.  Consequently an eigenpair of
  $\mathcal{A}_\eps$ whose eigenfunction is divergence-free is an eigenpair of $L_{\eps,u}$.
\end{lemma}

Faithful, modulo the standing conventions.

\begin{proof}
  \leanok
  \uses{def:main:induction-operator}
  The general vector identity for divergence-free $u$ reads
  $-(u\cdot\nabla)V+(V\cdot\nabla)u=\curl(u\times V)-u\,\Div V$, and $\Div V=0$ kills the last
  term.  In Lean the identity is proved coefficientwise from weak Leibniz certificates
  (\texttt{productInductionLowerField\_add\_velocityDiv\_eq\_weakCurl}), and the eigenpair transfer
  is then immediate because the embedding $\Lp^2\to\Hs^{-2}$ is injective.
\end{proof}

\begin{lemma}[The divergence equation and its energy]
  \label{lem:finish:divergence-energy}
  \lean{FastDynamo.productInductionLowerField_add_velocityDiv_eq_weakCurl,
  FastDynamo.productWeakDivergence_productInductionAction,
  FastDynamo.productWeakDivergence_velocityTimesScalar_pair_re_eq_zero,
  FastDynamo.productScalarWeakLaplacian_pair_self_re_nonpos,
  FastDynamo.productScalarWeakEigenfunction_eq_zero}
  \leanok
  \uses{def:main:induction-operator}
  Let $u\in\W^{1,\infty}$ be divergence-free, let $V\in\Hs^2$ and put $d=\Div V\in\Hs^1$.  Then, in
  the sense of distributions,
  \[
    -(u\cdot\nabla)V+(V\cdot\nabla)u+u\,d=\curl(u\times V),
    \qquad
    \Div\left[\mathcal{A}_\eps V\right]=\eps\Delta d-\Div(u\,d).
  \]
  Moreover the transport term is skew and the Laplacian is dissipative,
  \[
    \re\inner{\Div(u\,d)}{d}=0,
    \qquad
    \re\inner{\Delta d}{d}=-\norm{\nabla d}_2^2\leq0,
  \]
  so that a scalar $d\in\Hs^1$ obeying $\lambda d=-\Div(u\,d)+\eps\Delta d$ weakly, with
  $\eps\geq0$ and $\re\lambda>0$, vanishes.
\end{lemma}

In Lean the two pairings are the $\Hs^{-1}$--$\Hs^1$ duality, $\Div$ of an $\Lp^2$ field is
\texttt{productWeakDivergence} valued in $\Hs^{-1}$, and the final implication is the
velocity-free algebraic endpoint \texttt{productScalarWeakEigenfunction\_eq\_zero} (module
\texttt{FastDynamo/Analysis/ProductScalarEnergy.lean}); the vector identity and the divergence of
the induction action are in \texttt{FastDynamo/Analysis/InductionDivergence.lean}, the skewness in
\texttt{FastDynamo/Analysis/ProductTransportEnergy.lean}.

\begin{proof}
  \leanok
  The vector identity is checked on each of the three coordinates from the weak Leibniz rule for an
  $\Lp^\infty$ factor with an $\Lp^2$ factor; taking weak divergences and using
  $\Div\circ\curl=0$ gives the second display, since the diffusion term contributes
  $\eps\Delta d$.  Testing the scalar equation with $\bar d$, using periodicity and $\Div u=0$ for
  the transport term, and taking real parts yields
  $(\re\lambda)\norm d_2^2=-\eps\norm{\nabla d}_2^2\leq0$, whence $d=0$.
\end{proof}

\begin{lemma}[Eigenfunctions with $\re\lambda>0$ are divergence-free]
  \label{lem:positive-eigenfunction-solenoidal}
  \lean{FastDynamo.productInductionEigenpair_divergence_eq_zero,
  FastDynamo.productInductionEigenpair_solenoidal,
  FastDynamo.axiallyIndependent_positive_eigenfunction_solenoidal,
  FastDynamo.axiallyIndependent_eigenpair_solenoidal}
  \leanok
  \uses{def:main:induction-operator, def:finish:physical-operator}
  Let $u\in\W^{1,\infty}(\T^3;\R^3)$ be divergence-free and let $\eps\geq0$.  If
  $0\neq V\in\Hs^2(\T^3;\CC^3)$ satisfies
  \[
    \mathcal{A}_\eps V=\lambda V,
    \qquad\re\lambda>0,
  \]
  then $\Div V=0$.  Consequently $L_{\eps,u}V=\lambda V$.
\end{lemma}

The paper states no such lemma, because its recovered mode is solenoidal by construction
(the definition of $b_3$ in Lemma~8.1 of the paper); this is a result of the
formalisation.  The sign hypothesis on $\eps$ cannot be dropped: the statement is false for
$\eps<0$ (Lemma \ref{lem:finish:negative-diffusivity}).  The Lean statement therefore assumes
$0\leq\eps$, which is weaker than the paper's standing $\eps>0$ and hence covers it; the
positive-diffusivity wrapper is
\texttt{axiallyIndependent\_positive\_eigenfunction\_solenoidal}, and
\texttt{axiallyIndependent\_eigenpair\_solenoidal} also returns the physical eigenpair.  As explained in the
chapter preamble, the Lean velocity is $x_3$-independent; the general $\T^3$-velocity form of the
lemma is not formalised.  The conclusion comes in two equivalent shapes, $\Div V=0$ as an
$\Lp^2$ field and \texttt{IsProductSolenoidal} coefficientwise.

\begin{proof}
  \leanok
  \uses{lem:finish:divergence-energy, lem:finish:physical-compatibility}
  Set $d=\Div V$.  Taking the weak divergence of the eigenvalue equation and using the two displays
  of Lemma \ref{lem:finish:divergence-energy} gives the scalar advection-diffusion equation
  $\lambda d=-\Div(u\,d)+\eps\Delta d$ in $\Hs^{-1}$; since $d\in\Hs^1$ it is an admissible test
  function, and the energy endpoint of the same lemma gives $d=0$.  The last sentence is Lemma
  \ref{lem:finish:physical-compatibility}.  The argument starts from the auxiliary operator, not
  from the curl form whose agreement is only known after solenoidality, so the chain is not
  circular.
\end{proof}

\begin{lemma}[The sign hypothesis is necessary]
  \label{lem:finish:negative-diffusivity}
  \lean{FastDynamo.negativeDiffusivityEigenvalue,
  FastDynamo.negativeDiffusivity_auxiliary_eigenpair,
  FastDynamo.negativeDiffusivityEigenvalue_pos,
  FastDynamo.negativeDiffusivity_mode_not_productSolenoidal,
  FastDynamo.exists_positive_nonsolenoidal_eigenpair_of_negative_diffusivity}
  \leanok
  \uses{def:main:induction-operator}
  For every $\eps<0$ the zero velocity admits an eigenpair of $\mathcal{A}_\eps$ with positive real
  part and non-solenoidal eigenfunction: with $V(y,x_3)=\e^{2\pi ix_3}e_3$ one has
  \[
    \mathcal{A}_\eps V=\lambda V,
    \qquad\lambda=-\eps(2\pi)^2>0,
    \qquad\Div V\neq0 .
  \]
  Hence the hypothesis $\eps\geq0$ cannot be dropped from Lemma
  \ref{lem:positive-eigenfunction-solenoidal}.
\end{lemma}

This node is not in the paper; it is a remark-level Lean result recording that the sign
hypothesis of Lemma~\ref{lem:positive-eigenfunction-solenoidal} is sharp (module \texttt{FastDynamo/Finish/NegativeDiffusivity.lean}).  The main theorem is
unaffected, since it assumes $0<\eps$ throughout.

\begin{proof}
  \leanok
  With $u=0$ the transport and stretching terms vanish, and on the axial sector $K=1$ the constant
  vertical amplitude is annihilated by the transverse Laplacian, so only the axial diffusion symbol
  $-\eps(2\pi)^2$ survives.  Its divergence has the non-zero Fourier coefficient
  $2\pi i\cdot1$ at the zero transverse frequency.
\end{proof}

\section{Passage to a real-valued solution}

\begin{lemma}[Real part of a non-zero $x_3$-Fourier mode]
  \label{lem:real-axial-mode}
  \lean{FastDynamo.realPart_axialLiftPointwise_norm_sq,
  FastDynamo.phaseAxialFieldL2_norm_sq,
  FastDynamo.phaseAxialFieldL2_norm_sq_eq,
  FastDynamo.phaseAxialFieldL2_ne_zero}
  \leanok
  Let $K\in\Z\setminus\{0\}$, $b\in\Lp^2(\T^2;\CC^3)$ and $\varphi\in\R$.  Then, for almost every
  $y$,
  \[
    \int_0^1\abs{\re\bigl(\e^{i\varphi}\e^{2\pi iKx_3}b(y)\bigr)}^2\dd x_3
    =\frac12\abs{b(y)}_{\CC^3}^2 .
  \]
  In particular $\norm{\re(\e^{i\varphi}\e^{2\pi iKx_3}b)}_2^2=\frac12\norm b_2^2$ for every phase,
  this quantity does not depend on $\varphi$, and if $b\neq0$ the real part of the corresponding
  $x_3$-Fourier mode is non-zero for every phase.
\end{lemma}

Faithful.  The Lean statement comes in two layers: the fibrewise identity
\texttt{realPart\_axialLiftPointwise\_norm\_sq} is the displayed equation verbatim, for a
pointwise representative and the Haar measure of the axial circle; the global form
\texttt{phaseAxialFieldL2\_norm\_sq} is stated for the $\Lp^2$-class field
\texttt{phaseAxialFieldL2}~$K\,\varphi\,b=\re(\e^{i\varphi}\e^{2\pi iKx_3}b)$, which lives in the
real space $\Lp^2(\T^2\times\T^1;\R^3)$.  Module
\texttt{FastDynamo/Fourier/AxialL2.lean}.

\begin{proof}
  \leanok
  Write $b=a+ic$ with $a,c\in\R^3$.  Componentwise the integrand is the square of
  $a\cos(2\pi Kx_3+\varphi)-c\sin(2\pi Kx_3+\varphi)$.  Because $K$ is a non-zero integer, the
  squared sine and cosine have average $1/2$ on $[0,1]$ while their product has average zero;
  summing over the components gives the identity.  Integrating in $y$ by Fubini gives the global
  form.
\end{proof}

\begin{definition}[The real datum, the real solution and the strong-mode certificate]
  \label{def:finish:real-solution}
  \lean{FastDynamo.realAxialInitialField, FastDynamo.realAxialTrajectory,
  FastDynamo.productConjH2, FastDynamo.productRealPartH2,
  FastDynamo.complexifyProductR3L2, FastDynamo.realifiedEigenTrajectoryH2,
  FastDynamo.DomainOperator.IsStrongSolution,
  FastDynamo.ConditionalModeCertificate, FastDynamo.RealStrongModeCertificate}
  \leanok
  \uses{def:main:induction-operator, def:finish:physical-operator}
  Let $K\in\Z\setminus\{0\}$, $\lambda\in\CC$ and let $b$ be a transverse amplitude, and write
  $V=\e^{2\pi iKx_3}b$, $\alpha=\re\lambda$, $\omega=\im\lambda$.  Define
  \[
    B_{\rm in}=\re V,
    \qquad
    B(t)=\e^{\alpha t}\,\re\bigl(\e^{i\omega t}V\bigr)
    =\re\bigl(\e^{t\lambda}V\bigr)
    \qquad(t\in\R),
  \]
  both elements of the real space $\Lp^2(\T^2\times\T^1;\R^3)$
  (\texttt{realAxialInitialField}, \texttt{realAxialTrajectory}).  Inside the complex $\Hs^2$ graph
  the same field is represented by $W(t)=\frac12(\id+\mathrm{conj})(\e^{t\lambda}V)$
  (\texttt{realifiedEigenTrajectoryH2}, built from \texttt{productConjH2} and
  \texttt{productRealPartH2}), a fixed point of conjugation whose complexification is $B(t)$.  A
  curve $t\mapsto W(t)$ in the domain of a bundled operator $A$ is a \emph{strong solution} when
  $s\mapsto A.\text{embedding}(W(s))$ is differentiable in the ambient space with derivative
  $A.\text{op}(W(t))$ at every $t$ (\texttt{DomainOperator.IsStrongSolution}).  Finally
  \texttt{RealStrongModeCertificate}~$u\,\eps\,K\,\lambda\,b\,g$ bundles all the conclusions of
  this chapter for one mode: the lifted auxiliary eigenpair and its solenoidality, the physical
  eigenpair, $B_{\rm in}\neq0$, $B(t)\neq0$ for all $t$, the exact norm law, the lower bound with
  rate $g$, the conjugation-fixedness and solenoidality of $W(t)$, the identification of $W(t)$
  with $B(t)$, and the strong-solution equation for $W$.  (\texttt{ConditionalModeCertificate} is
  the sub-bundle not mentioning the real evolution equation.)
\end{definition}

The paper speaks of \emph{the} solution with datum $B_{\rm in}^\eps$, and its proof of
Theorem~1.2 writes that solution down.  This chapter exhibits the explicit strong solution in
$\Hs^2$ of the induction equation in the $\Hs^{-2}$ realisation of
Definition~\ref{def:finish:physical-operator}, with $W(0)=B_{\rm in}$; that it is the only one,
and that it is the orbit of $B_{\rm in}$ under the semigroup $\e^{tL_{\eps,u}}$ on
$\Lp^2_\sigma(\T^3)$, is Lemma~\ref{lem:dynamics:physical-uniqueness} in the appendix chapter.
Modules \texttt{FastDynamo/Finish/RealSolution.lean}, \texttt{FastDynamo/Operator/Evolution.lean},
\texttt{FastDynamo/Analysis/Realification.lean}.

\begin{lemma}[Exact norm law for the real mode]
  \label{lem:finish:norm-law}
  \lean{FastDynamo.realAxialTrajectory_eq_realPart_complexExp,
  FastDynamo.realAxialInitialField_ne_zero,
  FastDynamo.realAxialTrajectory_ne_zero,
  FastDynamo.norm_realAxialTrajectory,
  FastDynamo.realAxialTrajectory_norm_lower_bound}
  \leanok
  \uses{def:finish:real-solution}
  Let $K\neq0$ and $b\neq0$.  Then $B_{\rm in}\neq0$, $B(t)\neq0$ for every $t\in\R$, and
  \[
    \norm{\re\bigl(\e^{i\omega t}V\bigr)}_2^2=\frac12\norm b_2^2=\norm{\re V}_2^2,
    \qquad\text{hence}\qquad
    \norm{B(t)}_2=\e^{\alpha t}\norm{B_{\rm in}}_2 .
  \]
  If moreover $g\leq\alpha$, then $\e^{gt}\norm{B_{\rm in}}_2\leq\norm{B(t)}_2$ for $t\geq0$.
\end{lemma}

Faithful, and slightly broader than the paper in two respects: the norm identity is proved for
\emph{every} real time, not only $t\geq0$, and the same law is available after transport to the
literal three-torus (\texttt{norm\_literalRealTrajectory\_eq\_exp\_mul\_norm\_zero} in
\texttt{FastDynamo/Analysis/LiteralThreeTorus.lean}).  The auxiliary
\texttt{realAxialTrajectory\_eq\_realPart\_complexExp} identifies the definition of $B(t)$ with
$\re(\e^{t\lambda}V)$.

\begin{proof}
  \leanok
  \uses{lem:real-axial-mode}
  The non-zero axial Fourier index removes the oscillatory cross term when one integrates in
  $x_3$, so no temporal phase can make the real part small.  Apply Lemma
  \ref{lem:real-axial-mode} first with $\varphi=\omega t$ and then with $\varphi=0$: both sides
  equal $\frac12\norm b_2^2$, so the phase factor is norm-preserving and only the real
  exponential $\e^{\alpha t}$ survives.  Non-vanishing at every time follows from the same
  identity.
\end{proof}

\begin{lemma}[The realified mode is a strong solution]
  \label{lem:finish:strong-solution}
  \lean{FastDynamo.DomainOperator.isStrongSolution_eigenTrajectory,
  FastDynamo.productRealPartH2_preserves_divergenceFree,
  FastDynamo.productConjH2_realifiedEigenTrajectoryH2,
  FastDynamo.realifiedEigenTrajectoryH2_divergence_eq_zero,
  FastDynamo.complexify_realAxialTrajectory_eq_realified_embedding,
  FastDynamo.realifiedEigenTrajectoryH2_isPhysicalStrongSolution}
  \leanok
  \uses{def:finish:real-solution, def:finish:physical-operator}
  Let $u$ be divergence-free, let $\eps\in\R$, and let $V\in\Hs^2$ be divergence-free with
  $\mathcal{A}_\eps V=\lambda V$.  Then $W(t)=\re(\e^{t\lambda}V)$ is real (fixed by conjugation)
  and divergence-free at every time, it is the $\Hs^2$ representative of $B(t)$ whenever $V$ is the
  axial lift $\e^{2\pi iKx_3}b$, and it is a strong solution of the physical induction equation
  \[
    \partial_tB=L_{\eps,u}B,\qquad B(0)=B_{\rm in} .
  \]
\end{lemma}

This stands in for the paper's one-line justification in the proof of Theorem~1.2 (``since
the coefficients of the induction equation are real-valued, its solution is
$B^\eps(t)=\e^{\alpha_\eps t}\re(\e^{i\omega_\eps t}V^\eps)$''), which in the paper refers to
the solution given by the semigroup.  Here the complex solution $\e^{t\lambda}V$ is produced directly by scalar
exponentiation (\texttt{DomainOperator.isStrongSolution\_eigenTrajectory}), and no semigroup is
used; its identification with the semigroup orbit is
Lemma~\ref{lem:dynamics:physical-uniqueness}.

\begin{proof}
  \leanok
  \uses{lem:finish:physical-compatibility}
  The curve $t\mapsto\e^{t\lambda}V$ is a strong solution of the auxiliary equation, by
  differentiating the scalar exponential and using the eigenvalue equation.  Taking real parts is a
  fixed continuous real-linear map, so it commutes with the time derivative and with the embeddings
  into $\Lp^2$ and $\Hs^{-2}$; it also preserves the divergence-free constraint.  Finally, for
  divergence-free velocity and divergence-free field the auxiliary and physical actions agree
  (Lemma \ref{lem:finish:physical-compatibility}), which converts the auxiliary equation satisfied
  by the real part into the physical one.
\end{proof}

\begin{proposition}[One unstable reduced mode gives the real growing solution]
  \label{prop:finish:real-strong-mode}
  \lean{FastDynamo.conditionalModeCertificate_of_reducedEigenpair,
  FastDynamo.realStrongModeCertificate_of_reducedEigenpair}
  \leanok
  \uses{def:finish:real-solution, def:main:induction-operator, lem:axial-recovery}
  Let $u$ be divergence-free, $\eps\geq0$, $K\in\Z\setminus\{0\}$, and let $(\lambda,b)$ be an
  eigenpair of the reduced operator $A^{(3)}_{\eps,K}$ with $\re\lambda>0$; let
  $g\leq\re\lambda$.  Then
  the strong-mode certificate of Definition \ref{def:finish:real-solution} holds.  Explicitly, with
  $V=\e^{2\pi iKx_3}b$:
  \begin{itemize}
    \item $V\in\Hs^2(\T^3;\CC^3)$, $\norm V_{\Lp^2(\T^3)}=\norm b_{\Lp^2(\T^2)}>0$, and
      $\mathcal{A}_\eps V=\lambda V$;
    \item $\Div V=0$ and $L_{\eps,u}V=\lambda V$;
    \item $B_{\rm in}=\re V\neq0$ is divergence-free, and $B(t)$ is a real divergence-free strong
      solution of the induction equation with $\norm{B(t)}_2=\e^{t\re\lambda}\norm{B_{\rm in}}_2$
      for all $t\in\R$ and $\e^{gt}\norm{B_{\rm in}}_2\leq\norm{B(t)}_2$ for $t\geq0$.
  \end{itemize}
\end{proposition}

This proposition is the Lean form of the paper's proof of Theorem \ref{thm:main}, minus the
growth-rate clause.  It is the single hinge of the finish layer: the main chapter's
\texttt{MainPhysicalClauses} is obtained from the reduced clauses by applying it (see
\texttt{FastDynamo/Main.lean}).  The remaining clause $\gamma(u,\eps)\geq\gamma_0$ of Theorem
\ref{thm:main} is proved elsewhere and by a different route, from the same reduced eigenpair:
it is Theorem \ref{thm:main-growth-rate}, which needs the $\C_0$-semigroup theory built in
\texttt{FastDynamo/Dynamics/HilleYosida/} and \texttt{FastDynamo/Dynamics/Assembly/} rather
than the analytic-semigroup input of Lemma \ref{lem:semigroup-spectral-growth}.  The norm of the
complex eigentrajectory, $\norm{\e^{t\lambda}V}=\e^{t\re\lambda}\norm V$, is
\texttt{norm\_embedding\_eigenTrajectory}
(\texttt{FastDynamo/Finish/GrowingModeEvolution.lean}).

\begin{proof}
  \leanok
  \uses{lem:positive-eigenfunction-solenoidal, lem:finish:physical-compatibility,
  lem:finish:norm-law, lem:finish:strong-solution, lem:real-axial-mode}
  Since $K\neq0$ and every eigenvector of $A_{\eps,K}$ lies in its domain, the lift
  $V=\e^{2\pi iKx_3}b$ is one-periodic and lies in $\Hs^2$, with
  $\norm V_{\Lp^2(\T^3)}=\norm b_{\Lp^2(\T^2)}$; the axial reduction is an equivalence, so
  $\mathcal{A}_\eps V=\lambda V$.  Because $\re\lambda>0$, Lemma
  \ref{lem:positive-eigenfunction-solenoidal} gives $\Div V=0$, and Lemma
  \ref{lem:finish:physical-compatibility} then gives $L_{\eps,u}V=\lambda V$.  Lemma
  \ref{lem:finish:norm-law} gives the non-vanishing and the exact norm law for
  $B_{\rm in}=\re V$ and $B(t)$, and Lemma \ref{lem:finish:strong-solution} gives the evolution
  equation and the solenoidality of the curve.  Unless $\omega=\im\lambda$ vanishes the datum
  $B_{\rm in}$ is not itself a magnetic eigenmode: it evolves in the real invariant plane spanned
  by $\re V$ and $\im V$.
\end{proof}

\begin{lemma}[Zero mean]
  \label{lem:finish:zero-mean}
  \lean{FastDynamo.HasZeroMean, FastDynamo.productCoeffCLM_zero_zero_eq_integral,
    FastDynamo.hasZeroMean_axialLiftH2, FastDynamo.hasZeroMean_realifiedEigenTrajectoryH2,
    FastDynamo.hasZeroMean_complexify_realAxialTrajectory,
    FastDynamo.hasZeroMean_complexify_realAxialInitialField}
  \leanok
  \uses{def:finish:real-solution}
  Let $K\in\Z\setminus\{0\}$.  For every amplitude $b$ and every $\lambda\in\CC$ the mode
  $V=\e^{2\pi iKx_3}b$, the real datum $B_{\rm in}=\re V$ and the real solution
  $B(t)=\re(\e^{t\lambda}V)$ at every time have zero spatial mean,
  $\int_{\T^3}B(t)\dd x=0$.
\end{lemma}

This is the sentence ``Since $K_\eps\neq0$, every magnetic field constructed here has zero
spatial mean'' of Remark~1.5 of the paper (\texttt{rem:scope-intro}).  The mean is the integral
over the product torus with its probability Haar measure, which is Lebesgue measure on $\T^3$,
and a real field is read in $\CC^3$ through the complexification.  Module
\texttt{FastDynamo/Finish/ZeroMean.lean}.

\begin{proof}
  \leanok
  The mean is the zeroth Fourier coefficient, and the coefficients of $\e^{2\pi iKx_3}b$ are
  supported on the axial frequency $K\neq0$.  Scalar multiples, conjugates and real parts of a
  zero-mean field have zero mean.
\end{proof}

\begin{lemma}[The real part is not an eigenmode]
  \label{lem:finish:real-part-not-eigenmode}
  \lean{FastDynamo.productInductionAction_eq_smul_of_physicalEigenpair,
    FastDynamo.not_isEigenpair_productRealPartH2}
  \leanok
  \uses{def:finish:physical-operator, lem:finish:physical-compatibility,
    def:finish:real-solution}
  Let $u$ be real and divergence-free and let $V\in\Hs^2$ be a divergence-free eigenmode of
  $L_{\eps,u}$ with eigenvalue $\lambda$, $\im\lambda\neq0$.  Then $\re V$ is not an
  eigenfunction of $L_{\eps,u}$, for any eigenvalue.
\end{lemma}

This is Remark~1.4 of the paper (unlabelled): if $\im\lambda_\eps\neq0$, the real part of the
complex magnetic eigenmode $V^\eps$ is not itself a magnetic eigenmode, since it is not an
eigenfunction of $L_{\eps,u}$.  The Lean statement is for every such eigenmode, with
$L_{\eps,u}$ the curl form of Definition~\ref{def:finish:physical-operator}; the divergence-free
hypothesis on $V$ is part of the paper's notion of a magnetic eigenmode.  Module
\texttt{FastDynamo/Finish/RealPartNotEigenmode.lean}.

\begin{proof}
  \leanok
  \uses{lem:finish:physical-compatibility}
  On divergence-free fields $L_{\eps,u}$ is the auxiliary action, which commutes with complex
  conjugation because $u$ is real.  If $\re V=\frac12(V+\bar V)$ were an eigenfunction with
  eigenvalue $\mu$, then $\lambda V+\bar\lambda\bar V=\mu(V+\bar V)$; since
  $\lambda\neq\bar\lambda$, comparing the two relations forces $V=0$, contradicting that $V$
  is an eigenmode.
\end{proof}

\section{Axial Fourier index and magnetic length}

\begin{definition}[Magnetic length of an axial mode]
  \label{def:finish:magnetic-length}
  \lean{FastDynamo.productGradientFromH2,
  FastDynamo.norm_sq_productGradientFromH2_axialLiftH2,
  FastDynamo.complexAxialMagneticLength,
  FastDynamo.phaseAxialGradientL2,
  FastDynamo.phaseAxialGradientL2_norm_sq,
  FastDynamo.complexifyProductR3L2_phaseAxialGradientL2,
  FastDynamo.realPhaseAxialMagneticLength}
  \leanok
  \uses{def:finish:real-solution}
  For $B\in\Hs^2$ write $\nabla B=(\partial_1B,\partial_2B,\partial_3B)$ for the full weak
  gradient, an element of $\Lp^2(\T^3;\CC^3)^3$.  On a fixed axial sector $V=\e^{2\pi iKx_3}b$,
  \[
    \norm{\nabla V}_2^2
    =\norm{\partial_1b}_2^2+\norm{\partial_2b}_2^2+(2\pi K)^2\norm b_2^2 .
  \]
  The $\Hs^1$-based magnetic length of the mode and of its phase-rotated real part are
  \[
    L_B=\frac{\norm V_2}{\norm{\nabla V}_2},
    \qquad
    L_B^{\re,\varphi}
    =\frac{\norm{\re(\e^{i\varphi}V)}_2}{\norm{\re(\e^{i\varphi}\nabla V)}_2},
  \]
  the second one taken at $\varphi=0$ for the real datum $B_{\rm in}$.
\end{definition}

In Lean the real gradient is defined as the vector of the phase-rotated real parts of the three
complex derivative amplitudes (\texttt{phaseAxialGradientL2}), and
\texttt{complexifyProductR3L2\_phaseAxialGradientL2} proves that it is the weak gradient of the
real field: complexified, its components are the weak derivatives of the $\Hs^2$ field
$\re(\e^{i\varphi}V)$, because the weak derivatives commute with conjugation and hence with
the real part.  \texttt{phaseAxialGradientL2\_norm\_sq} proves that
its squared norm is half of the complex one, which is the paper's
$\norm{\nabla B_{\rm in}}_2^2=\frac12\norm{\nabla V}_2^2$ (proof of Corollary~8.3).  Module
\texttt{FastDynamo/Finish/MagneticScale.lean}.

\begin{lemma}[Eigenmode energy estimate]
  \label{lem:finish:magnetic-energy}
  \lean{FastDynamo.ReducedInductionMagneticScale.epsilon_mul_norm_sq_productGradient_le_stretching,
  FastDynamo.ReducedInductionMagneticScale.norm_transverseStretchingCoefficient_pos_of_reducedEigenpair,
  FastDynamo.complexAxialMagneticLength_lower_of_energy,
  FastDynamo.magneticLength_lower_of_energy}
  \leanok
  \uses{def:main:induction-operator, def:finish:magnetic-length}
  Let $u$ be divergence-free, $\eps\geq0$, and let $(\lambda,b)$ be an eigenpair of
  $A^{(3)}_{\eps,K}$ with $\re\lambda\geq0$.  Then, with $M_u$ the $\Lp^\infty$ operator norm of the stretching matrix
  field $b\mapsto(b\cdot\nabla_\perp)u$,
  \[
    \eps\norm{\nabla V}_2^2\leq M_u\norm V_2^2,
    \qquad\text{hence}\qquad
    \sqrt{\eps/M_u}\leq L_B\quad\text{when }M_u>0 .
  \]
  If moreover $\eps>0$, $K\neq0$ and $b\neq0$, then $M_u>0$.
\end{lemma}

This is the energy argument in the proof of Corollary~8.3 of the paper
(\texttt{cor:magnetic-scale}), which the paper does not state as a separate lemma.  One
difference: the paper's constant is
$M_u=\mathrm{ess\,sup}_x\norm{\nabla u(x)}_{\mathrm{op}}$; the Lean constant is the norm of
\texttt{transverseStretchingCoefficient}, the $\Lp^\infty$ operator norm of the stretching matrix
field $x\mapsto Du(x)$, which is what the energy identity uses; by its definition it is the
paper's $M_u$, an identification that no Lean statement records.
Like the paper (``In particular, $M_u>0$'', after the same energy estimate), Lean derives
$M_u>0$ from the existence of the eigenpair, with no property of the constructed velocity.  The hypotheses
$\eps\geq0$, $\re\lambda\geq0$ are weaker than the paper's $\eps>0$, $\re\lambda_\eps>0$.

\begin{proof}
  \leanok
  Take the real part of the eigenvalue identity paired with $V$.  Transverse and axial transport
  are skew, so they drop out; diffusion contributes $-\eps\norm{\nabla V}_2^2$ with the exact
  fixed-sector gradient identity of Definition \ref{def:finish:magnetic-length}; stretching is
  bounded by $M_u\norm V_2^2$.  Since $\re\lambda\geq0$ this gives the energy inequality, and the
  elementary algebraic step $\eps G^2\leq MF^2\Rightarrow\sqrt{\eps/M}\leq F/G$ gives the lower
  bound.  If $\eps>0$ and the gradient is non-zero, the left-hand side is positive, so $M_u>0$.
\end{proof}

\begin{corollary}[Axial Fourier index and $\Hs^1$-based magnetic length of the growing fields]
  \label{cor:magnetic-scale}
  \lean{FastDynamo.complexAxialMagneticLength_le,
  FastDynamo.realPhaseAxialMagneticLength_eq_complex,
  FastDynamo.realPhaseAxialMagneticLength_le,
  FastDynamo.ReducedInductionMagneticScale.complexAxialMagneticLength_lower,
  FastDynamo.ScaleFamily.matching_fourierIndex_sq_bounds,
  FastDynamo.ScaleFamily.matching_fourierIndex_lower,
  FastDynamo.ScaleFamily.matching_fourierIndex_upper,
  FastDynamo.MainPhysicalClausesAt, FastDynamo.mainPhysicalClausesAt_of_reducedEigenpair,
  FastDynamo.main_theorem_magnetic_scale}
  \leanok
  \uses{def:finish:magnetic-length, def:finish:real-solution, def:velocity:scale-family,
    thm:main-spectral, def:main:witnesses}
  Define the $\Hs^1$-based magnetic length of the magnetic eigenmode $V^\eps$ of Theorem
  \ref{thm:main} by $L_B^\eps=\norm{V^\eps}_2/\norm{\nabla V^\eps}_2$.  There are constants
  $c,C>0$, independent of $\eps$, such that
  \[
    c\eps^{-1/2}\leq\abs{K_\eps}\leq C\eps^{-1/2},
    \qquad
    c\sqrt\eps\leq L_B^\eps\leq C\sqrt\eps
    \qquad(0<\eps\leq\eps_0),
  \]
  the range in which Theorem~\ref{thm:main} produces $V^\eps$.  The real datum
  $B_{\rm in}^\eps=\re V^\eps$ has the same magnetic length.
\end{corollary}

The Lean statement is \texttt{main\_theorem\_magnetic\_scale}
(\texttt{FastDynamo/Finish/MagneticScaleMain.lean}): there are $c,C>0$ such that for every
$0<\eps\leq\eps_0$ there is a triple $(K,\lambda,b)$ satisfying the clauses of
Theorem~\ref{thm:main} at the explicit witnesses (\texttt{MainPhysicalClausesAt}, the body of
the sixteen clauses of Definition~\ref{def:main:physical-clauses}) together with the two
two-sided estimates for $\abs K$ and for the magnetic length of the complex mode
$V=\e^{2\pi iKx_3}b$, and such that the real part of $V$ at every phase, in particular the real
datum $B_{\rm in}^\eps$, has the same magnetic length.  The triple is the eigenpair in the axial
sector $K=K_{n(\eps)}$ of the matching cell of the parameter chain.  The general bounds it is
assembled from are also stated on their own:
\begin{itemize}
  \item the index estimate, for every cell $n$ of an arbitrary scale family and every
    $\eps>0$ with $\eps/\ell_n^2\in I$,
    $k_*\sqrt{\delta_-/\eps}\leq K_n\leq k_*\sqrt{\delta_+/\eps}$
    (\texttt{ScaleFamily.matching\_fourierIndex\_lower} and \texttt{\_upper}, with the squared
    form \texttt{matching\_fourierIndex\_sq\_bounds});
  \item the upper magnetic-length bound $L_B\leq1/(2\pi\abs K)$ for every $K\neq0$ and every
    non-zero amplitude, for the complex mode (\texttt{complexAxialMagneticLength\_le}) and at
    every real phase (\texttt{realPhaseAxialMagneticLength\_le});
  \item the lower bound $\sqrt{\eps/M_u}\leq L_B$ of the paper's proof, for every non-zero
    eigenpair of $A^{(3)}_{\eps,K}$ with $\eps>0$, $K\neq0$ and $\re\lambda\geq0$
    (\texttt{ReducedInductionMagneticScale.complexAxialMagneticLength\_lower},
    Lemma~\ref{lem:finish:magnetic-energy});
  \item the identity of magnetic lengths of the real and the complex mode,
    \texttt{realPhaseAxialMagneticLength\_eq\_complex}, for every $K\neq0$, every non-zero
    amplitude and every phase.
\end{itemize}
Deviations: the norm $\norm{\nabla V}_2$ is that of the weak gradient on the product torus and
$\norm{V}_2$ the product-measure norm; the lower bound uses
$M_u=\norm{\texttt{transverseStretchingCoefficient}\ u}$ in place of
$\mathrm{ess\,sup}\norm{\nabla u}_{\rm op}$ (Lemma~\ref{lem:finish:magnetic-energy}); the
constants are $c=\min\{\sqrt{\delta_-},M_u^{-1/2}\}$ and
$C=\max\{\sqrt{\delta_+},(2\pi\sqrt{\delta_-})^{-1}\}$ at $k_*=1$.  The triple $(K,\lambda,b)$
has its own existential in this statement, as in \texttt{MainPhysicalClauses}: that it is the
triple of \texttt{MainSemigroupOrbitClauses} is not stated, and the clause ``its solution'' (the
semigroup orbit and the uniqueness of item (8) of Theorem~\ref{thm:main-strengthened}) is not
part of it.  The remark after the
corollary in the paper (Remark~8.4, unlabelled), which only observes that this magnetic length is
comparable to the classical resistive scale $\sqrt\eps$ of Moffatt and Proctor, has no Lean
counterpart and is not a node.

\begin{proof}
  \leanok
  \uses{lem:real-axial-mode, lem:finish:magnetic-energy, prop:finish:real-strong-mode}
  For the matching index $n=n(\eps)$ one has $\ell_n=\sqrt{\eps/\delta_n}$, $2\pi K_n\ell_n=N$
  and $\delta_n\in[\delta_-,\delta_+]$, so $K_\eps=K_n=k_*\sqrt{\delta_n/\eps}$, which is the
  index estimate.  (The paper writes this arithmetic at $N=2\pi$, where the middle identity
  is $K_n\ell_n=1$ and the index is $K_\eps=\sqrt{\delta_n/\eps}$; Lean keeps
  $k_*$ free, so its constants read $k_*\sqrt{\delta_\mp/\eps}$.)  Since
  $\partial_{x_3}V^\eps=2\pi iK_nV^\eps$, the full gradient dominates
  $2\pi\abs{K_n}\norm{V^\eps}_2$, whence
  $L_B^\eps\leq1/(2\pi K_n)=\ell_n/N=\ell_n/(2\pi k_*)$ -- the paper's
  $L_B^\eps\leq\ell_n/(2\pi)\leq\sqrt\eps/(2\pi\sqrt{\delta_-})$ at $k_*=1$.
  The converse inequality is the energy estimate of Lemma
  \ref{lem:finish:magnetic-energy}.  Finally, applying
  Lemma \ref{lem:real-axial-mode} to $V^\eps$ and to each of its three derivatives gives
  $\norm{B_{\rm in}^\eps}_2^2=\frac12\norm{V^\eps}_2^2$ and
  $\norm{\nabla B_{\rm in}^\eps}_2^2=\frac12\norm{\nabla V^\eps}_2^2$, so the real datum has the
  same magnetic length.  The constants are the minimum and the maximum of the constants of
  the two steps.
\end{proof}


\chapter{Appendix A: the semigroup, and the Lean generation chain}
This chapter is the blueprint's appendix.  It records Lemma~A.1 of the paper
(\texttt{lem:semigroup-spectral-growth}, in its Appendix~A, \texttt{sec:semigroup-growth}) ---
Lemma~\ref{lem:semigroup-spectral-growth} here, the analytic input that turns the spectrum of
$L_{\eps,u}$ into a statement about the growth rate $\gamma(u,\eps)$ --- together with the
facts about $\gamma(u,\eps)$ that §1.1 of the paper states with it (the limit
\texttt{eq:semigroup-growth-rate} exists, and the spectral supremum is attained), and the Lean
chain that proves them.  The other nodes of this chapter are the formalisation's own: the paper
states none of them, and they are compared with the paper only through
Lemma~\ref{lem:semigroup-spectral-growth}, the growth-rate clause of Theorem~\ref{thm:main}, the
energy identity \texttt{eq:energy-competition} and the phrase ``its solution'' of
Theorem~\ref{thm:main}.  The particle flow, its entropy, the ideal induction group and the
vanishing-diffusivity limit are Section~4 and Remark~1.7 of the paper and have their own
chapter.

The paper proves its Lemma~A.1 on $\Lp^2_\sigma$ directly: $\eps\Delta$ with domain
$\Hs^2\cap\Lp^2_\sigma$ generates an analytic semigroup there, $C_uB=\curl(u\times B)$ is
infinitesimally bounded relative to it, so $L_{\eps,u}$ generates an analytic semigroup by
the perturbation theorem [Engel--Nagel, Theorem III.2.10]; the graph norm is equivalent to
the $\Hs^2$ norm, so the resolvent is compact; analyticity makes the semigroup immediately
compact, which gives the spectral mapping theorem and the equality of the growth bound with
the spectral bound; and $0\in\sigma(L_{\eps,u})$, together with the finiteness of the
spectrum in any right half-plane, makes the spectral supremum a maximum.  The formalisation
proves every one of these statements (Lemma~\ref{lem:semigroup-spectral-growth},
\texttt{SemigroupSpectralGrowth}), along a different route.  Mathlib has no
$\C_0$-semigroup theory, so the development builds what it needs from scratch.
\begin{itemize}
  \item \emph{Generation} (Sections~\ref{sec:dynamics:hille-yosida}
    to~\ref{sec:dynamics:solenoidal}).  A resolvent estimate $\norm{R(r)}\leq(r-M)^{-1}$ on a
    real ray (Definition~\ref{def:dynamics:hille-yosida-ray}) drives a Yosida approximation
    scheme whose limit is a strongly continuous semigroup
    (Theorem~\ref{thm:dynamics:hille-yosida}).  Feeding it the dissipativity of the reduced
    operator $A_{\eps,K}$ and of its formal adjoint
    (Lemma~\ref{lem:dynamics:reduced-adjoint},
    Proposition~\ref{prop:dynamics:reduced-half-plane-resolvent}) gives each axial sector a
    Hille--Yosida ray (Theorem~\ref{thm:dynamics:sector-semigroup}); summing the sector
    resolvents over the axial Fourier index (Lemmas~\ref{lem:dynamics:axial-hilbert-sum}
    and~\ref{lem:dynamics:uniform-sector-estimate},
    Proposition~\ref{prop:dynamics:product-resolvent}) gives a ray for, and hence generates
    the semigroup of, the all-frequency auxiliary operator $\mathcal A_\eps$
    (\texttt{productInductionOperator u $\eps$}, domain $\Hs^2(\T^3;\CC^3)$) on the whole of
    $\Lp^2(\T^3;\CC^3)$ (Theorem~\ref{thm:dynamics:product-semigroup}); restricting it to the
    closed solenoidal subspace (Lemma~\ref{lem:dynamics:solenoidal-subspace},
    Theorem~\ref{thm:dynamics:solenoidal-semigroup}) puts it on the paper's state space
    $\Lp^2_\sigma(\T^3)$.
  \item \emph{The generator} (Section~\ref{sec:dynamics:generator}).  The generator of a
    $\C_0$-semigroup and the spectrum of an unbounded operator are defined
    (Definition~\ref{def:dynamics:generator}); the generator of the Hille--Yosida semigroup is
    the operator it was built from (Lemma~\ref{lem:dynamics:hille-yosida-generator}); and the
    generator of the semigroup on $\Lp^2_\sigma$ is the paper's $L_{\eps,u}$ with domain
    $\Hs^2\cap\Lp^2_\sigma$, which determines the semigroup
    (Theorem~\ref{thm:dynamics:solenoidal-generator}).  So the semigroup whose growth bound is
    the Lean value of $\gamma(u,\eps)$ (Definition~\ref{def:main:growth-rate},
    \texttt{FastDynamo.gamma}) is the paper's $\e^{tL_{\eps,u}}$.
  \item \emph{Analyticity, compactness, spectral mapping} (Section~\ref{sec:dynamics:analytic}).
    A sector condition on the numerical range of $\mathcal A_\eps$
    (Lemma~\ref{lem:dynamics:sectorial}) passes to the Yosida approximants in complex time and
    makes the semigroup analytic of angle $\pi/4$
    (Theorem~\ref{thm:dynamics:analytic-generation}); the resolvent factors through the compact
    embedding $\Hs^2\hookrightarrow\Lp^2$ (Lemma~\ref{lem:dynamics:compact-embedding}); an
    analytic semigroup with compact resolvent is immediately compact
    (Lemma~\ref{lem:dynamics:immediate-compactness}); and for such semigroups a direct
    point-spectrum argument gives the spectral mapping theorem and the equality of the growth
    bound with the spectral bound, attained at an eigenvalue
    (Theorem~\ref{thm:dynamics:spectral-mapping}).  Mean conservation gives
    $0\in\sigma(L_{\eps,u})$ and $\norm{\e^{tL_{\eps,u}}}\geq1$
    (Lemma~\ref{lem:dynamics:mean-conservation}), and the growth bound of every
    $\C_0$-semigroup is a limit (Lemma~\ref{lem:dynamics:growth-bound-limit}).
    Theorem~\ref{thm:dynamics:induction-spectrum} collects the instances at the induction
    semigroup.
  \item \emph{Solutions} (Section~\ref{sec:dynamics:solutions}).  The energy identity
    \texttt{eq:energy-competition} (Lemma~\ref{lem:dynamics:energy-identity}), and uniqueness
    of strong solutions of the curl-form induction equation, which identifies every solution
    with divergence-free datum with the semigroup orbit
    (Lemma~\ref{lem:dynamics:physical-uniqueness}).
\end{itemize}

The growth-rate clause of Theorem~\ref{thm:main} needs only part of this: the paper bounds
$\norm{\e^{tL_{\eps,u}}}$ from below by $\norm{B^\eps(t)}_2/\norm{B_{\rm in}^\eps}_2$, and Lean
bounds the growth bound from below by the real part of an eigenvalue whose eigenfunction lies in
$\Lp^2_\sigma$ (Theorem~\ref{thm:main-growth-rate}); analyticity, compactness and spectral
mapping are not used there.  Both the paper and the Lean clause (\texttt{gamma\_ge\_gamma0})
state the growth-rate clause for $0<\eps\leq\eps_0$.

Besides the semigroup chain, the first section records the Lean-side nodes: the
$\C_0$-semigroup interface (Definition~\ref{def:dynamics:semigroup-interface}); the a priori
energy bound $\norm{B(t)}_2\leq\e^{M_ut}\norm{B(0)}_2$ for forward strong solutions of the
reduced operator and the uniqueness of those solutions
(Lemmas~\ref{lem:dynamics:energy-bound} and~\ref{lem:dynamics:strong-solution-uniqueness});
the exponential growth bound of an \emph{arbitrary} strongly continuous semigroup, its two
elementary norm bounds and its description as a limit
(Definition~\ref{def:dynamics:growth-bound},
Lemmas~\ref{lem:dynamics:growth-bound-norm-bounds} and~\ref{lem:dynamics:growth-bound-limit});
the exact exponential law of an explicit eigen-trajectory
(Lemma~\ref{lem:dynamics:eigen-trajectory-law}); and the growth floor of the initial-value
problem on the operator domain, sandwiched at the explicit witnesses by
$\gamma_0\leq\gamma_{\rm str}(A_{\eps,K_\eps})\leq M_u$
(Lemma~\ref{lem:dynamics:strong-solution-growth-floor}).  Throughout,
$M_u:=\norm{\nabla_\perp u}_\infty$ is the $\Lp^\infty$ norm of the stretching matrix field,
independent of $\eps$ and of the axial index.  The map \texttt{docs/MODULE\_MAP.md} contains
the same accounting on the Lean side.

\section{The growth rate and the analytic semigroup}

\begin{lemma}[Semigroup generation and spectral growth]
  \label{lem:semigroup-spectral-growth}
  \lean{FastDynamo.SemigroupSpectralGrowth, FastDynamo.semigroupSpectralGrowth_holds}
  \leanok
  \uses{def:main:induction-operator, def:main:growth-rate, def:dynamics:generator,
    def:dynamics:analytic-semigroup}
  Let $u\in\W^{1,\infty}(\T^3;\R^3) $ be divergence-free and let $\eps>0 $.  Then the
  following hold.
  \begin{enumerate}
    \item The operator $L_{\eps,u} $, with domain $\Hs^2\cap\Lp^2_\sigma $, generates an
      analytic $\C_0 $-semigroup on $\Lp^2_\sigma $.
    \item The resolvent of $L_{\eps,u} $ is compact, $\e^{tL_{\eps,u}} $ is compact for
      every $t>0 $, and
      \begin{equation}
        \sigma(\e^{tL_{\eps,u}})\setminus\{0\}
        =\left\{\e^{t\lambda}:\lambda\in\sigma(L_{\eps,u})\right\},
        \qquad t>0.
        \label{eq:dynamics:compact-semigroup-spectral-mapping}
      \end{equation}
    \item Furthermore,
      \begin{equation}
        \gamma(u,\eps)
        =\max\{\re\lambda:\lambda\in\sigma(L_{\eps,u})\}.
        \label{eq:dynamics:growth-rate-spectral-bound}
      \end{equation}
  \end{enumerate}
\end{lemma}

The paper states the lemma in two paragraphs, the second with the two displays; the
enumeration into items (1)--(3) is the blueprint's, for reference.  The Lean statement is the
structure \texttt{SemigroupSpectralGrowth u hu h$\eps$}, proved for every divergence-free
transverse velocity and every $\eps>0$ by \texttt{semigroupSpectralGrowth\_holds}
(\texttt{FastDynamo/Dynamics/Assembly/SemigroupSpectralGrowth.lean}).  Its fields, in order:
\texttt{generator\_eq}, \texttt{generator\_curlForm}, \texttt{generator\_unique} and
\texttt{analytic} (item~(1)); \texttt{resolventSet\_nonempty},
\texttt{isCompactOperator\_resolvent}, \texttt{isCompactOperator\_semigroup} and
\texttt{spectralMapping} (item~(2)); \texttt{gamma\_tendsto} (the limit
\texttt{eq:semigroup-growth-rate} exists and is a real number, as §1.1 of the paper states);
\texttt{gamma\_eq\_iSup} ($\gamma(u,\eps)$ is the supremum);
\texttt{zero\_mem\_spectrum}; and \texttt{exists\_eigenvalue\_re\_eq\_gamma} (the supremum is
attained, at an eigenvalue with an eigenvector, as §1.1 of the paper states);
\texttt{gamma\_eq\_iSup} and \texttt{exists\_eigenvalue\_re\_eq\_gamma} together are
\eqref{eq:dynamics:growth-rate-spectral-bound}.  It differs from the paper's
statement as follows.
\begin{itemize}
  \item \emph{Model.}  The torus is the product $\T^2\times\T$, the state space is the
    Fourier-defined $\Lp^2_\sigma$ of that model (\texttt{productSolenoidalSubspace}), and the
    velocity is axially independent (\texttt{TransverseRealW1Inf}), as everywhere in this
    blueprint.  The clauses are stated for the semigroup \texttt{solenoidalInductionSemigroup}
    on the product model; \texttt{gamma} is the growth bound of its isometric transport to the
    literal three-torus, which has the same operator norms and the same growth bound
    (\texttt{growthBound\_literalInductionSemigroup}).
  \item \emph{The operator.}  The generator is \texttt{solenoidalInductionOperator}: domain
    the divergence-free $\Hs^2$ fields, values the auxiliary action $\mathcal A_\eps B$ of
    Definition~\ref{def:main:induction-operator}, which lies in $\Lp^2_\sigma$.  On that
    domain the action is the curl form: its image in the Fourier model of $\Hs^{-2}$ is
    $\curl(u\times B)+\eps\Delta B$ (\texttt{generator\_curlForm}), and the embedding of
    $\Lp^2$ into that model is injective.  This is an equivalent formulation of the paper's
    operator.
  \item \emph{Analyticity} is \texttt{IsAnalyticSemigroup} with the explicit angle $\pi/4$
    (Definition~\ref{def:dynamics:analytic-semigroup}); the paper names no angle.  The
    spectrum is \texttt{LinearPMap.spectrum}, the complement of the set of $z$ at which
    $z-L$ has a bounded two-sided inverse (Definition~\ref{def:dynamics:generator}).
  \item \emph{Additional clauses.}  \texttt{generator\_unique} (every $\C_0$-semigroup on
    $\Lp^2_\sigma$ with this generator is this one), \texttt{resolventSet\_nonempty} and
    \texttt{zero\_mem\_spectrum} are not in the paper's statement; the last is a step of the
    paper's proof.  Attainment is stated with an eigenvector.
\end{itemize}
The route differs from the paper's at every step; the proof below gives both.

\begin{proof}
  \leanok
  \uses{def:main:induction-operator, def:main:growth-rate,
    thm:dynamics:product-semigroup, thm:dynamics:solenoidal-semigroup,
    thm:dynamics:solenoidal-generator, thm:dynamics:induction-spectrum,
    lem:dynamics:growth-bound-limit, lem:dynamics:mean-conservation}
  \emph{The paper's proof.}  Work on $\Lp^2_\sigma $ throughout.  The operator $\eps\Delta $
  with domain $\Hs^2\cap\Lp^2_\sigma $ generates an analytic semigroup there by its Fourier
  representation, and $C_uB=\curl(u\times B) $ belongs to $\Lp^2_\sigma $ for every
  $B\in\Hs^2\cap\Lp^2_\sigma $.  For every $\eta>0 $, interpolation gives
  $\norm{C_uB}_2\leq C\norm u_{\W^{1,\infty}}\norm B_{\Hs^1}\leq\eta\norm{\Delta B}_2+C_\eta\norm B_2 $,
  so $C_u $ is infinitesimally bounded relative to $\eps\Delta $ on $\Lp^2_\sigma $, and the
  perturbation theorem for analytic semigroups [Engel--Nagel, Theorem III.2.10] gives item
  (1) with the stated domain.  By the same relative bound and the elliptic estimate
  $\norm B_{\Hs^2}\leq C(\norm{\Delta B}_2+\norm B_2) $ the graph norm of $L_{\eps,u} $ is
  equivalent to the $\Hs^2 $-norm, so the compact embedding
  $\Hs^2(\T^3)\hookrightarrow\Lp^2(\T^3) $ gives a compact resolvent; analyticity plus
  [Engel--Nagel, Theorem II.4.29] makes the semigroup immediately compact, and
  [Engel--Nagel, Corollaries IV.3.11 and IV.3.12] give
  \eqref{eq:dynamics:compact-semigroup-spectral-mapping} and the equality of the growth
  bound with the spectral bound, which is
  \eqref{eq:dynamics:growth-rate-spectral-bound}.  For attainment: every
  $B\in D(L_{\eps,u}) $ has $\int_{\T^3}L_{\eps,u}B\dd x=0 $, and $\Lp^2_\sigma $ contains
  non-zero constant fields, so $L_{\eps,u} $ is not surjective and $0\in\sigma(L_{\eps,u}) $.
  By [Engel--Nagel, Corollary V.3.2] only finitely many spectral points lie in any fixed
  right half-plane; applied to $\{\re z\geq0\} $, this makes the spectral supremum a maximum.
  The paper's proof notes in addition that the limit in \texttt{eq:semigroup-growth-rate} is
  this growth bound.

  \emph{The Lean proof.}  Item~(1): the semigroup is generated on all of
  $\Lp^2(\T^3;\CC^3)$ by the auxiliary operator through the Hille--Yosida theorem
  (Theorem~\ref{thm:dynamics:product-semigroup}) and restricted to the invariant closed
  subspace $\Lp^2_\sigma$ (Theorem~\ref{thm:dynamics:solenoidal-semigroup}); its generator is
  identified with $L_{\eps,u}$ on $\Hs^2\cap\Lp^2_\sigma$ and determines it
  (Theorem~\ref{thm:dynamics:solenoidal-generator}); analyticity comes from a sector condition
  on the numerical range through the Yosida approximants in complex time, not from
  [Engel--Nagel, Theorem III.2.10] (Theorem~\ref{thm:dynamics:analytic-generation}).
  Item~(2): the resolvent factors through the compact embedding, as in the paper;
  compactness of $\e^{tL_{\eps,u}}$ for $t>0$ is proved from analyticity directly
  (Lemma~\ref{lem:dynamics:immediate-compactness}), where the paper cites
  [Engel--Nagel, Theorem II.4.29]; and the spectral mapping theorem is proved by a direct
  point-spectrum argument (Theorem~\ref{thm:dynamics:spectral-mapping}), where the paper cites
  [Engel--Nagel, Corollaries IV.3.11 and IV.3.12].  Item~(3): the same point-spectrum argument,
  with Gelfand's formula for $\e^{L_{\eps,u}}$ in place of [Engel--Nagel, Corollary V.3.2],
  gives the growth bound as the spectral bound, attained at an eigenvalue; the spectrum is
  non-empty because $0\in\sigma(L_{\eps,u})$, which is the paper's mean argument
  (Lemma~\ref{lem:dynamics:mean-conservation}).  The limit is
  Lemma~\ref{lem:dynamics:growth-bound-limit}, and it is a real number because
  $\norm{\e^{tL_{\eps,u}}}\geq1$ (Lemma~\ref{lem:dynamics:mean-conservation}).  The instances
  at the induction semigroup are Theorem~\ref{thm:dynamics:induction-spectrum}.
\end{proof}

\begin{definition}[The Lean $\C_0$-semigroup interface]
  \label{def:dynamics:semigroup-interface}
  \lean{FastDynamo.StronglyContinuousSemigroup, FastDynamo.DomainOperator.IsForwardStrongSolution, FastDynamo.DomainOperator.ForwardStrongSolutionsUnique, FastDynamo.DomainOperatorSemigroupRealization}
  \leanok
  \uses{def:main:induction-operator}
  For a complex normed space $X $, a \emph{strongly continuous semigroup} is a family
  $S(t)\in\mathcal L(X) $, $t\geq0 $, with $S(0)=\id $,
  $S(s+t)=S(s)S(t) $, and $t\mapsto S(t)x $ continuous for every $x\in X $.  For an
  operator $A $ presented as a bounded map out of a Banach domain $D $ with continuous
  injection $\iota:D\to X $, a curve $B:\R\to D $ is a \emph{forward strong solution}
  when $t\mapsto\iota(B(t)) $ has one-sided derivative $A(B(t)) $ at every $t\geq0 $, and
  $A $ has \emph{unique forward strong solutions} when two forward strong solutions with
  the same initial ambient state agree for all $t\geq0 $.  A \emph{realization} of $S $ by
  $A $ is the data of a domain-valued orbit $B\mapsto\Theta(B,\cdot) $ with
  $\Theta(B,0)=B $, each $\Theta(B,\cdot) $ a forward strong solution, and
  $S(t)\iota(B)=\iota(\Theta(B,t)) $ for every $t\geq0 $.
\end{definition}

In Lean the four items are \texttt{StronglyContinuousSemigroup} (a structure bundling the
three semigroup axioms, with time in \texttt{NNReal}),
\texttt{DomainOperator.IsForwardStrongSolution},
\texttt{DomainOperator.ForwardStrongSolutionsUnique} and
\texttt{DomainOperatorSemigroupRealization}, all in
\texttt{FastDynamo/Dynamics/StronglyContinuousSemigroup.lean}.  The realization structure
packages the \emph{conclusion} of item~(1) of Lemma~\ref{lem:semigroup-spectral-growth} as a
hypothesis, and \emph{both} of the hypotheses this interface isolates are discharged for
the operators of this development, so the theorems quantified over a realization are not
vacuous:
Lemma~\ref{lem:dynamics:strong-solution-uniqueness} and
Lemma~\ref{lem:dynamics:product-dissipative} inhabit
\texttt{ForwardStrongSolutionsUnique} for $A_{\eps,K} $ and for
$\mathcal A_\eps $, at every divergence-free $u $ and $\eps\geq0 $, and \texttt{HilleYosidaRay.realization}
(Theorem~\ref{thm:dynamics:hille-yosida}) and
\texttt{productInductionRealization} (Theorem~\ref{thm:dynamics:product-semigroup})
inhabit \texttt{DomainOperatorSemigroupRealization}.  One reading of the interface should
be resisted: a realization records that the semigroup \emph{solves} the Cauchy
problem, not that it is the only family doing so, and it is forward uniqueness, not the
structure, that supplies the definite article
(Lemma~\ref{lem:dynamics:semigroup-unique}).

\begin{lemma}[A priori energy bound for forward strong solutions]
  \label{lem:dynamics:energy-bound}
  \lean{FastDynamo.DomainOperator.norm_embedding_le_exp_of_isForwardStrongSolution, FastDynamo.reducedInductionOperator_norm_embedding_le_exp, FastDynamo.re_inner_reducedInductionAction_le}
  \leanok
  \uses{def:main:induction-operator}
  Let $A $ be an operator with bundled domain $D $ and continuous injection $\iota:D\to X $
  into a complex inner-product space, and suppose that the numerical range of $A $ lies in
  a half-plane,
  \[
    \re\inner{\iota(d)}{A d}\leq M\norm{\iota(d)}^2
    \qquad(d\in D).
  \]
  Then every forward strong solution $B $ of $A $ obeys
  \[
    \norm{\iota(B(t))}\leq\e^{Mt}\norm{\iota(B(0))}
    \qquad(t\geq0).
  \]
  The hypothesis holds for the reduced fixed-sector operator $A_{\eps,K} $ of
  \eqref{eq:main:torus-sector} with $M=M_u:=\norm{\nabla_\perp u}_\infty $, for every
  divergence-free transverse $u\in\W^{1,\infty} $, every $\eps\geq0 $ and every $K\in\Z $;
  hence every forward strong solution of $A_{\eps,K} $ satisfies
  $\norm{B(t)}_2\leq\e^{M_ut}\norm{B(0)}_2 $ for $t\geq0 $, with a ceiling independent of
  the diffusivity and of the axial index.
\end{lemma}

This is the operator-level form of the paper's energy competition
(\texttt{eq:energy-competition}): stretching contributes at most
$M_u\norm B_2^2 $,
diffusion contributes a nonpositive term and transport contributes nothing.  In Lean the
abstract bound is
\texttt{DomainOperator.norm\_embedding\_le\_exp\_of\_isForwardStrongSolution} and its
instance for $A_{\eps,K} $ is
\texttt{reducedInductionOperator\_norm\_embedding\_le\_exp}, both in
\texttt{FastDynamo/Dynamics/StrongSolutionUniqueness.lean}; the half-plane hypothesis is
discharged there by \texttt{re\_inner\_reducedInductionAction\_le}, which combines the
exact energy identity and the stretching estimate of
\texttt{FastDynamo/Global/ReducedInductionEnergyIdentity.lean}.  The constant $M_u $ is the
$\Lp^\infty $ operator norm of the stretching matrix field
$b\mapsto(b_\perp\cdot\nabla_\perp)u $, and it does not depend on $\eps $ or $K $.

\begin{proof}
  \leanok
  \uses{def:main:induction-operator}
  Along a forward strong solution the squared ambient energy
  $\varphi(t)=\norm{\iota(B(t))}^2 $ is differentiable at every $t>0 $ with
  $\varphi'(t)=2\re\inner{\iota(B(t))}{A(B(t))}\leq2M\varphi(t) $, so
  $t\mapsto\e^{-2Mt}\varphi(t) $ has nonpositive derivative on $(0,T) $ and is antitone on
  $[0,T] $; evaluating at the endpoints and taking square roots gives the bound.  (This is
  Gronwall's lemma in exponential-multiplier form; Lean carries it out by hand, with
  mean-value theory only, rather than through \texttt{Mathlib.Analysis.ODE.Gronwall}.)
  For $A_{\eps,K} $ the
  half-plane hypothesis is the exact identity
  \[
    \re\inner{b}{A_{\eps,K}b}
    =-\eps\bigl(\norm{\nabla_\perp b}_2^2+(2\pi K)^2\norm b_2^2\bigr)
    +\re\inner{b}{(b_\perp\cdot\nabla_\perp)u},
  \]
  in which the transport terms are skew and drop out, together with the stretching bound
  $\re\inner{b}{(b_\perp\cdot\nabla_\perp)u}\leq M_u\norm b_2^2 $ and $\eps\geq0 $.
\end{proof}

\begin{lemma}[Uniqueness of forward strong solutions]
  \label{lem:dynamics:strong-solution-uniqueness}
  \lean{FastDynamo.DomainOperator.forwardStrongSolutionsUnique_of_re_inner_le, FastDynamo.reducedInductionOperator_forwardStrongSolutionsUnique}
  \leanok
  \uses{lem:dynamics:energy-bound}
  With $A $, $\iota $ and $M $ as in Lemma~\ref{lem:dynamics:energy-bound}, two forward
  strong solutions of $A $ with the same initial ambient state agree in $X $ at every
  $t\geq0 $.  In particular, for every divergence-free transverse $u\in\W^{1,\infty} $,
  every $\eps\geq0 $ and every $K\in\Z $, the reduced operator $A_{\eps,K} $ has unique
  forward strong solutions, with no further hypothesis.
\end{lemma}

This is the uniqueness half of an energy well-posedness statement, for the
fixed-sector reduced operator; existence is not formalised.  The paper states no such
proposition, and Theorem~\ref{thm:main} does not need one.  The Lean statement is weaker
than a full energy well-posedness proposition in three ways -- the velocity is
autonomous, the data lie in the bundled $\Hs^2 $ domain rather than in $\Lp^2_\sigma $, and
the solution concept is the strong one of Definition~\ref{def:dynamics:semigroup-interface}
rather than the Lions--Magenes energy one -- and it holds for every divergence-free $u $, every
$\eps\geq0 $ and every $K $.  In Lean:
\texttt{DomainOperator.forwardStrongSolutionsUnique\_of\_re\_inner\_le} and
\texttt{reducedInductionOperator\_forwardStrongSolutionsUnique}
(\texttt{FastDynamo/Dynamics/StrongSolutionUniqueness.lean}).  It inhabits one of the two hypotheses isolated by
Definition~\ref{def:dynamics:semigroup-interface}.

\begin{proof}
  \leanok
  \uses{lem:dynamics:energy-bound}
  The difference of two forward strong solutions is again a forward strong solution, by
  linearity of $\iota $ and of $A $, and its initial ambient state vanishes.  The bound of
  Lemma~\ref{lem:dynamics:energy-bound} applied to the difference gives
  $\norm{\iota(B_1(t)-B_2(t))}\leq\e^{Mt}\cdot0=0 $ for every $t\geq0 $.
\end{proof}

\begin{lemma}[Exact exponential law of an eigen-trajectory]
  \label{lem:dynamics:eigen-trajectory-law}
  \lean{FastDynamo.DomainOperator.isStrongSolution_eigenTrajectory, FastDynamo.DomainOperator.norm_embedding_eigenTrajectory}
  \leanok
  \uses{def:main:induction-operator}
  Let $A $ be an operator with bundled Banach domain $D $ and continuous injection
  $\iota:D\to X $, and let $\lambda\in\CC $, $0\neq B\in D $ satisfy $AB=\lambda\iota(B) $.
  Then the curve $t\mapsto\e^{t\lambda}B $ is a strong solution of
  $\frac{\dd}{\dd t}\iota(B(t))=A(B(t)) $ on all of $\R $, and
  \[
    \norm{\iota(\e^{t\lambda}B)}
    =\e^{t\re\lambda}\norm{\iota(B)}
    \qquad(t\in\R).
  \]
\end{lemma}

In Lean the curve is \texttt{FastDynamo.DomainOperator.eigenTrajectory}, and the two
displayed facts are \texttt{isStrongSolution\_eigenTrajectory}
(\texttt{FastDynamo/Operator/Evolution.lean}) and
\texttt{norm\_embedding\_eigenTrajectory}
(\texttt{FastDynamo/Finish/GrowingModeEvolution.lean}); the same module derives that the
norm tends to infinity when $\re\lambda>0 $ and crosses every threshold at a nonnegative
time.

This is \emph{not} the growth bound of Lemma~\ref{lem:semigroup-spectral-growth}, and the
distinction is the whole content of this chapter.  The lemma describes a single
trajectory, written down by hand by scalar exponentiation; it asserts nothing about the
existence of an evolution operator, and the quantity it controls is
$\norm{\iota(\e^{t\lambda}B)} $ for one datum, not
$\norm{\e^{tL_{\eps,u}}}_{\Lp^2_\sigma\to\Lp^2_\sigma} $.  To convert it into a lower
bound for $\gamma(u,\eps) $ one needs a semigroup to exist and this curve to be its orbit,
which is exactly the realization hypothesis of
Definition~\ref{def:dynamics:semigroup-interface}; that hypothesis is a theorem
(Theorem~\ref{thm:dynamics:product-semigroup}), and the conversion is carried out in
Theorem~\ref{thm:main-growth-rate}.  It is this lemma, at the level of the physical
real-valued field, that supplies the exact norm law in Theorem~\ref{thm:main}; the
growth-rate clause is a separate statement, proved separately and by a different route.

\begin{proof}
  \leanok
  \uses{def:main:induction-operator}
  Differentiating $s\mapsto\e^{s\lambda} $ in the real variable $s $ and applying the
  bounded map $\iota $ gives
  $\frac{\dd}{\dd t}\iota(\e^{t\lambda}B)=\lambda\,\iota(\e^{t\lambda}B)$, which equals
  $A(\e^{t\lambda}B) $ by the eigenvalue equation and linearity.  For the norm,
  $\norm{\iota(\e^{t\lambda}B)}=\abs{\e^{t\lambda}}\norm{\iota(B)}=\e^{t\re\lambda}\norm{\iota(B)} $.
  Neither step needs a hypothesis or an analytic input.
\end{proof}

\begin{definition}[The growth bound of a strongly continuous semigroup]
  \label{def:dynamics:growth-bound}
  \lean{FastDynamo.StronglyContinuousSemigroup.growthBound}
  \leanok
  \uses{def:dynamics:semigroup-interface}
  For a strongly continuous semigroup $S $ on a complex normed space $X $, the
  \emph{exponential growth bound} is
  \[
    \omega_0(S)
    :=\limsup_{t\to\infty}\frac1t\log\norm{S(t)}
    \in[-\infty,+\infty],
  \]
  with the convention that a time at which $S(t)=0 $ contributes $-\infty $.
\end{definition}

In Lean this is \texttt{StronglyContinuousSemigroup.growthBound}
(\texttt{FastDynamo/Dynamics/GrowthBound.lean}), valued in \texttt{EReal}.  Two points of
the encoding matter.  First, it is defined as a $\limsup $, so that it is defined for every
semigroup with no preliminary work; Lemma~\ref{lem:dynamics:growth-bound-limit} proves that
for every $\C_0$-semigroup on a Banach space it is a limit, so it is the limit of
Definition~\ref{def:main:growth-rate}.  Second, the integrand is $-\infty $ rather than
Mathlib's junk value $\log0=0 $ at the times where $\norm{S(t)}=0 $, so the definition is
faithful on the nilpotent semigroups the interface permits.  This is the growth bound of an
arbitrary \emph{given} semigroup; applied to the semigroup of
Theorem~\ref{thm:dynamics:solenoidal-semigroup} it is the Lean value of $\gamma(u,\eps) $
(Definition~\ref{def:main:growth-rate}).

\begin{lemma}[Exponential majorants and minorants control the growth bound]
  \label{lem:dynamics:growth-bound-norm-bounds}
  \lean{FastDynamo.StronglyContinuousSemigroup.growthBound_le_of_norm_le, FastDynamo.StronglyContinuousSemigroup.le_growthBound_of_norm_ge}
  \leanok
  \uses{def:dynamics:growth-bound}
  Let $S $ be a strongly continuous semigroup on $X $ and let $\omega\in\R $.  If
  $\norm{S(t)}\leq C\e^{\omega t} $ for every $t\geq0 $ and some $C>0 $, then
  $\omega_0(S)\leq\omega $.  If $\e^{\omega t}\leq\norm{S(t)} $ for all large $t $, then
  $\omega\leq\omega_0(S) $.
\end{lemma}

\begin{proof}
  \leanok
  \uses{def:dynamics:growth-bound}
  For the first bound, at the times where $\norm{S(t)}=0 $ the integrand is $-\infty $ and
  there is nothing to prove; elsewhere take logarithms, divide by $t>0 $ and use
  monotonicity of the $\limsup $ against the majorant $t^{-1}\log C+\omega $, which tends
  to $\omega $.  For the second bound the hypothesis makes $\norm{S(t)} $ positive along
  $t\to\infty $, so the same computation shows that the integrand is eventually at least
  $\omega $, which is below the $\limsup $.
\end{proof}

\begin{lemma}[The growth bound is a limit]
  \label{lem:dynamics:growth-bound-limit}
  \lean{FastDynamo.StronglyContinuousSemigroup.growthExponent,
    FastDynamo.StronglyContinuousSemigroup.exists_norm_operator_le,
    FastDynamo.StronglyContinuousSemigroup.growthBound_eq_iInf,
    FastDynamo.StronglyContinuousSemigroup.tendsto_growthExponent,
    FastDynamo.StronglyContinuousSemigroup.growthBound_ne_top,
    FastDynamo.StronglyContinuousSemigroup.growthBound_ne_bot_iff,
    FastDynamo.StronglyContinuousSemigroup.exists_growthBound_eq_and_tendsto}
  \leanok
  \uses{def:dynamics:growth-bound}
  Let $S $ be a strongly continuous semigroup on a complex Banach space $X $.  Then
  \[
    \frac1t\log\norm{S(t)}\xrightarrow[t\to\infty]{}\omega_0(S)
    =\inf_{t>0}\frac1t\log\norm{S(t)}\in[-\infty,\infty),
  \]
  with $\log0=-\infty$.  The limit is finite exactly when there is $\omega\in\R$ with
  $\e^{\omega t}\leq\norm{S(t)}$ for every $t>0$, and then it is a real number $L$ with
  $t^{-1}\log\norm{S(t)}\to L$ in $\R$.  In particular the $\limsup$ of
  Definition~\ref{def:dynamics:growth-bound} is the limit of
  \eqref{eq:main:semigroup-growth-rate}.
\end{lemma}

This is the paper's justification of \texttt{eq:semigroup-growth-rate} (the sentence after that
display: subadditivity and local boundedness, [Engel--Nagel, Proposition IV.2.2]), for every
$\C_0$-semigroup, with no hypothesis beyond completeness of $X$.  In Lean:
\texttt{FastDynamo/Dynamics/GrowthBoundLimit.lean}.  \texttt{growthExponent} is the integrand of
\texttt{growthBound}, valued in \texttt{EReal} and equal to $\bot=-\infty$ where $S(t)=0$;
\texttt{growthBound\_eq\_iInf} and \texttt{tendsto\_growthExponent} are the display;
\texttt{growthBound\_ne\_bot\_iff} is the criterion for finiteness, and
\texttt{exists\_growthBound\_eq\_and\_tendsto} the real-valued form.  The limit lies in
$[-\infty,\infty)$, as in the paper; it is $-\infty$ for nilpotent semigroups and also for
semigroups whose norm decays faster than every exponential without vanishing.  Completeness of
$X$ enters through the Banach--Steinhaus theorem (\texttt{exists\_norm\_operator\_le}).

\begin{proof}
  \leanok
  \uses{def:dynamics:growth-bound}
  Fix $s>0$ and put $n=\lfloor t/s\rfloor$.  The semigroup law gives
  $\norm{S(t)}\leq\norm{S(t-ns)}\norm{S(s)}^n\leq K\norm{S(s)}^n$, where $K$ bounds
  $\norm{S(r)}$ for $0\leq r\leq s$ (strong continuity on a compact interval and the
  Banach--Steinhaus theorem).  Taking logarithms,
  $t^{-1}\log\norm{S(t)}\leq t^{-1}\log K+(n/t)\log\norm{S(s)}$, whose right side tends to
  $s^{-1}\log\norm{S(s)}$; if $S(s)=0$ then $S(t)=0$ for $t\geq s$.  Hence the $\limsup$ is
  at most every value $s^{-1}\log\norm{S(s)}$, so at most their infimum, which is at most the
  $\liminf$.  So $\limsup=\inf=\liminf$, and the limit exists.  It is finite exactly when the
  infimum is, which is the exponential minorant.
\end{proof}

\begin{lemma}[Spectral lower bound for the growth bound of a realizing semigroup]
  \label{lem:dynamics:conditional-growth-bound}
  \lean{FastDynamo.DomainOperatorSemigroupRealization.exp_re_le_operatorNorm, FastDynamo.DomainOperatorSemigroupRealization.eigenpair_orbit_and_growth, FastDynamo.DomainOperatorSemigroupRealization.re_le_growthBound_of_eigenpair}
  \leanok
  \uses{def:dynamics:growth-bound, def:dynamics:semigroup-interface, lem:dynamics:strong-solution-uniqueness, lem:dynamics:eigen-trajectory-law}
  Let $S $ be a strongly continuous semigroup on $X $, let $A $ be an operator with
  bundled domain $D $, assume a realization of $S $ by $A $ in the sense of
  Definition~\ref{def:dynamics:semigroup-interface}, and assume that $A $ has unique
  forward strong solutions.  Then for every eigenpair $AB=\lambda\iota(B) $ with
  $\iota(B)\neq0 $ and every $t\geq0 $,
  \[
    S(t)\iota(B)=\iota(\e^{t\lambda}B),
    \qquad
    \norm{S(t)\iota(B)}=\e^{t\re\lambda}\norm{\iota(B)},
    \qquad
    \e^{t\re\lambda}\leq\norm{S(t)},
  \]
  and consequently $\re\lambda\leq\omega_0(S) $.
\end{lemma}

The statement carries its own hypotheses, and in that form it is proved in Lean
(\texttt{FastDynamo/Dynamics/StronglyContinuousSemigroup.lean} for the three displayed
identities, \texttt{FastDynamo/Dynamics/GrowthBound.lean} for
\texttt{re\_le\_growthBound\_of\_eigenpair}).  It isolates the easy half of the growth
clause: the implication ``an eigenvalue with $\re\lambda\geq\gamma_0 $
$\Rightarrow$ $\norm{S(t)}\geq\e^{t\gamma_0} $, hence $\gamma(u,\eps)\geq\gamma_0 $''.  Both
of the hypotheses the lemma carries are theorems: forward
uniqueness for the reduced operator, for the all-frequency operator and for the curl form
(Lemmas~\ref{lem:dynamics:strong-solution-uniqueness},
\ref{lem:dynamics:product-dissipative} and~\ref{lem:dynamics:physical-uniqueness}), and the realization by the generation theorem
(Theorem~\ref{thm:dynamics:product-semigroup}).
The lemma is therefore not conditional in substance: it is the engine of the lower
bound of Lemma~\ref{lem:dynamics:hille-yosida-growth-sandwich} and, through it, of
Theorem~\ref{thm:main-growth-rate}.

\begin{proof}
  \leanok
  \uses{def:dynamics:growth-bound, lem:dynamics:growth-bound-norm-bounds, lem:dynamics:eigen-trajectory-law}
  The realization supplies a forward strong solution starting at $B $ whose embedding is
  the semigroup orbit; the eigen-trajectory of
  Lemma~\ref{lem:dynamics:eigen-trajectory-law} is a second forward strong solution with
  the same initial state.  Forward uniqueness identifies the two, which gives the first
  identity; the norm law of Lemma~\ref{lem:dynamics:eigen-trajectory-law} gives the
  second; and dividing $\norm{S(t)\iota(B)}\leq\norm{S(t)}\norm{\iota(B)} $ by
  $\norm{\iota(B)}>0 $ gives the third.  The third inequality holds for every $t\geq0 $, so
  the minorant half of Lemma~\ref{lem:dynamics:growth-bound-norm-bounds} with
  $\omega=\re\lambda $ gives $\re\lambda\leq\omega_0(S) $.  Forward uniqueness enters as a
  \emph{hypothesis} of this lemma, not as a lemma its proof invokes; it is discharged
  downstream at Lemma~\ref{lem:dynamics:strong-solution-uniqueness}, whose Lean module
  \texttt{Dynamics/StrongSolutionUniqueness.lean} imports
  \texttt{Dynamics/StronglyContinuousSemigroup.lean} and not the other way round, so this
  proof does not \texttt{\textbackslash uses} it.
\end{proof}

\begin{lemma}[The growth floor of the initial-value problem]
  \label{lem:dynamics:strong-solution-growth-floor}
  \lean{FastDynamo.DomainOperator.strongSolutionGrowthFloor, FastDynamo.DomainOperator.re_le_strongSolutionGrowthFloor_of_eigenpair, FastDynamo.DomainOperator.strongSolutionGrowthFloor_le_of_norm_le, FastDynamo.reducedInductionOperator_strongSolutionGrowthFloor_le, FastDynamo.mainGrowthFloor_le_strongSolutionGrowthFloor, FastDynamo.mainGrowthFloor_le_strongSolutionGrowthFloor_le}
  \leanok
  \uses{lem:dynamics:eigen-trajectory-law, lem:dynamics:energy-bound, thm:main-spectral, def:main:witnesses}
  For an operator $A $ with bundled domain $D $ and continuous injection $\iota:D\to X $
  put
  \[
    \gamma_{\rm str}(A)
    :=\sup_B\ \limsup_{t\to\infty}\frac1t
    \log\frac{\norm{\iota(B(t))}}{\norm{\iota(B(0))}}
    \in[-\infty,+\infty],
  \]
  the supremum being over all forward strong solutions $B $ of $A $ with
  $\iota(B(0))\neq0 $ (and $-\infty $ if there are none; a time at which $\iota(B(t)) $
  vanishes contributes $-\infty $ to the inner $\limsup $).  Then
  \[
    \re\lambda\leq\gamma_{\rm str}(A)
  \]
  for every eigenpair $AB=\lambda\iota(B) $ with $\iota(B)\neq0 $, and
  \[
    \gamma_{\rm str}(A)\leq M
  \]
  whenever every forward strong solution of $A $ obeys
  $\norm{\iota(B(t))}\leq\e^{Mt}\norm{\iota(B(0))} $ for $t\geq0 $.  Both halves apply to
  the reduced fixed-sector operator: with $M_u=\norm{\nabla_\perp u}_\infty $ as in
  Lemma~\ref{lem:dynamics:energy-bound} one has $\gamma_{\rm str}(A_{\eps,K})\leq M_u $
  for every divergence-free transverse $u\in\W^{1,\infty} $, every $\eps\geq0 $ and every
  $K\in\Z $; and at the explicit witnesses $u $, $\eps_0 $, $\gamma_0 $ of
  Definition~\ref{def:main:witnesses} the eigenpairs of Theorem~\ref{thm:main-spectral}
  give, for every $0<\eps\leq\eps_0 $, an axial index $K_\eps\neq0 $ with
  \[
    \gamma_0\leq\gamma_{\rm str}(A_{\eps,K_\eps})\leq M_u.
  \]
\end{lemma}

This has no hypothesis: no semigroup, no realization and no uniqueness hypothesis.  It is
\emph{not} the paper's $\gamma(u,\eps) $, and the Lean name
\texttt{strongSolutionGrowthFloor} is chosen so that it cannot be mistaken for it.  The
difference is the one this whole chapter is about: $\gamma_{\rm str} $ is the growth rate
of the initial-value problem on the operator's domain, a supremum over curves that happen
to solve the equation, whereas $\gamma(u,\eps) $ of Definition~\ref{def:main:growth-rate}
is a supremum over \emph{all} $\Lp^2_\sigma $ data taken inside an operator norm.  Given a
strongly continuous semigroup $S $ realizing $A $, the inequality
$\gamma_{\rm str}(A)\leq\omega_0(S) $ follows from
$\norm{S(t)x}\leq\norm{S(t)}\norm x $, and the converse from density of the domain
together with the locally uniform bound of Lemma~\ref{lem:dynamics:energy-bound}; neither
direction is proved.  In Lean:
\texttt{DomainOperator.strongSolutionGrowthFloor} and
\texttt{DomainOperator.re\_le\_strongSolutionGrowthFloor\_of\_eigenpair}
and the ceiling \texttt{DomainOperator.strongSolutionGrowthFloor\_le\_of\_norm\_le}
(\texttt{FastDynamo/Dynamics/GrowthBound.lean}).  The sandwich at the explicit witnesses is
\texttt{mainGrowthFloor\_le\_strongSolutionGrowthFloor\_le}, assembled from
\texttt{mainGrowthFloor\_le\_strongSolutionGrowthFloor} (the eigenpairs of
Theorem~\ref{thm:main-spectral} at the witnesses of Definition~\ref{def:main:witnesses})
and \texttt{reducedInductionOperator\_strongSolutionGrowthFloor\_le} (the ceiling, fed by
the a priori bound of Lemma~\ref{lem:dynamics:energy-bound}), both in
\texttt{FastDynamo/Dynamics/MainGrowthFloor.lean}.  So
$\gamma_0\leq\gamma_{\rm str}(A_{\eps,K_\eps})\leq M_u $ holds for every $0<\eps\leq\eps_0 $:
it is the counterpart of \eqref{eq:main-growth-conclusion} for $\gamma_{\rm str} $, the lower bound carrying the growth conclusion and the
upper bound showing that $\gamma_{\rm str} $ is a real quantity and not a vacuous
$+\infty $.  It is a statement about the initial-value problem on the domain; the
corresponding statement about $\gamma(u,\eps) $ is Theorem~\ref{thm:main-growth-rate}, and
it needs the whole generation chain of the next three sections, not this lemma.

\begin{proof}
  \leanok
  \uses{lem:dynamics:eigen-trajectory-law, lem:dynamics:energy-bound, thm:main-spectral, def:main:witnesses}
  The eigen-trajectory $t\mapsto\e^{t\lambda}B $ of
  Lemma~\ref{lem:dynamics:eigen-trajectory-law} is a forward strong solution with
  $\iota(B(0))=\iota(B)\neq0 $ and
  $\norm{\iota(\e^{t\lambda}B)}/\norm{\iota(B)}=\e^{t\re\lambda} $ for every $t $.  Taking
  logarithms and dividing by $t>0 $ makes the inner integrand identically $\re\lambda $
  along $t\to\infty $, so the $\limsup $ for this single solution is $\re\lambda $, and the
  supremum over all solutions is at least that.  For the ceiling, the hypothesis bounds
  the inner integrand of every admissible solution by $t^{-1}\log\e^{Mt}=M $ along
  $t\to\infty $ (the times where $\iota(B(t)) $ vanishes contribute $-\infty $), so every
  inner $\limsup $, and hence the supremum, is at most $M $.  The instance for
  $A_{\eps,K} $ is Lemma~\ref{lem:dynamics:energy-bound} with $M=M_u $, and the lower
  bound at the witnesses of Definition~\ref{def:main:witnesses} is the eigenpair clause of
  Theorem~\ref{thm:main-spectral} fed to the displayed eigenpair bound.
\end{proof}

\section{The Hille--Yosida generation theorem}
\label{sec:dynamics:hille-yosida}

Mathlib has no $\C_0 $-semigroup theory, so the generation theorem that item~(1) of
Lemma~\ref{lem:semigroup-spectral-growth} needs is built here from scratch, in
\texttt{FastDynamo/Dynamics/HilleYosida/}.  The paper generates an \emph{analytic}
semigroup by an infinitesimal perturbation of $\eps\Delta $; the formalisation instead uses
the dissipative route, whose only analytic input is a resolvent estimate on a real ray and
whose conclusion is a $\C_0 $-semigroup with an exponential norm bound.  Its generator is
identified in Lemma~\ref{lem:dynamics:hille-yosida-generator}, and analyticity is proved from a
sector condition in Theorem~\ref{thm:dynamics:analytic-generation}, on the same Yosida
approximants.

\begin{definition}[Hille--Yosida ray]
  \label{def:dynamics:hille-yosida-ray}
  \lean{FastDynamo.HilleYosidaRay, FastDynamo.HilleYosidaRay.ofHalfPlane,
    FastDynamo.HilleYosidaRay.norm_smul_stateResolvent_le,
    FastDynamo.HilleYosidaRay.tendsto_smul_stateResolvent}
  \leanok
  \uses{def:main:induction-operator}
  Let $A $ be an operator presented as a bounded map out of a Banach domain $D $ with
  continuous injection $\iota:D\to X $, and let $M\in\R $.  A \emph{Hille--Yosida ray of
  constant $M $} for $A $ is a set $S\subset\CC $ together with a two-sided inverse
  $(z-A)^{-1}:X\to D $ of the pencil $z\iota-A $ at every $z\in S $, such that $S $ contains
  every real $r>M $ and the state resolvent $R(r)=\iota\circ(r-A)^{-1} $ obeys
  \[
    \norm{R(r)}_{X\to X}\leq(r-M)^{-1}
    \qquad(r\in\R,\ r>M).
  \]
\end{definition}

In Lean this is the structure \texttt{FastDynamo.HilleYosidaRay A M}
(\texttt{FastDynamo/Dynamics/HilleYosida/ResolventRay.lean}), with fields
\texttt{spectralSet}, a \texttt{HasResolventOn} certificate, membership of the real ray and
the displayed bound.  The certificate records only the two inverse equations: no closed
graph, no completeness of $D $ and no continuity of $z\mapsto R(z) $ is assumed, and every
algebraic identity the construction needs (the resolvent identity, $AR(z)=zR(z)-\id $,
commutation of the resolvents) is derived from them.  \texttt{HilleYosidaRay.ofHalfPlane}
builds a ray from a resolvent on the half-plane $\{\re z>M\} $, which is how the two
instances below arise.  The two quantitative consequences the approximation scheme consumes
are \texttt{norm\_smul\_stateResolvent\_le} ($\norm{rR(r)}\leq r/(r-M) $, a bound uniform for
large $r $) and, when $\iota $ has dense range, the strong convergence
$rR(r)x\to x $ of \texttt{tendsto\_smul\_stateResolvent}.

\begin{theorem}[Hille--Yosida generation]
  \label{thm:dynamics:hille-yosida}
  \lean{FastDynamo.exists_stronglyContinuousSemigroup_of_hilleYosidaRay,
    FastDynamo.HilleYosidaRay.semigroup, FastDynamo.HilleYosidaRay.realization,
    FastDynamo.HilleYosidaRay.semigroup_spec,
    FastDynamo.HilleYosidaRay.norm_semigroup_le,
    FastDynamo.HilleYosidaRay.growthBound_le}
  \leanok
  \uses{def:dynamics:hille-yosida-ray, def:dynamics:semigroup-interface,
    def:dynamics:growth-bound}
  Let $X $ be a complex Banach space and let $A $ carry a Hille--Yosida ray of constant $M $
  whose domain injection $\iota:D\to X $ has dense range.  Then $A $ generates a strongly
  continuous semigroup $S $ on $X $: there is a realization of $S $ by $A $ in the sense of
  Definition~\ref{def:dynamics:semigroup-interface} --- every $d\in D $ evolves along a
  domain-valued forward strong solution whose embedding is the orbit $t\mapsto S(t)\iota(d) $
  --- and
  \[
    \norm{S(t)}_{X\to X}\leq\e^{Mt}
    \qquad(t\geq0),
    \qquad\text{hence}\qquad
    \omega_0(S)\leq M .
  \]
\end{theorem}

The existential form is
\texttt{exists\_stronglyContinuousSemigroup\_of\_hilleYosidaRay}; because the consumers need
the orbit and not only its existence, the same content is also stated about the named
objects \texttt{HilleYosidaRay.semigroup} and \texttt{HilleYosidaRay.realization}, as
\texttt{HilleYosidaRay.semigroup\_spec} (\texttt{FastDynamo/Dynamics/HilleYosida/Generation.lean}).
This is the analogue of the classical generation theorem for contraction-type semigroups
[Engel--Nagel, Theorem II.3.5] in this development's \texttt{DomainOperator} formulation.
Analyticity is not part of this theorem: it needs the additional sector condition of
Theorem~\ref{thm:dynamics:analytic-generation}.

\begin{proof}
  \leanok
  \uses{def:dynamics:hille-yosida-ray, def:dynamics:semigroup-interface,
    lem:dynamics:growth-bound-norm-bounds}
  The Yosida route.  The approximants $A_r:=r^2R(r)-r $
  (\texttt{HilleYosidaRay.yosida}) are bounded operators with
  $\norm{\e^{tA_r}}\leq\e^{Mrt/(r-M)} $, so their exponentials
  $\e^{tA_r} $ (\texttt{HilleYosidaRay.approx}) are uniformly bounded on compact time
  intervals; $A_r\iota(d)\to A d $ for every $d\in D $, and the two approximants commute,
  which turns the telescoping identity
  $\e^{tA_r}x-\e^{tA_s}x=\int_0^t\frac{\dd}{\dd\tau}\bigl(\e^{(t-\tau)A_s}\e^{\tau A_r}x\bigr)\dd\tau $
  into the Cauchy estimate of \texttt{cauchySeq\_approx}, first on the dense range of $\iota $
  and then, by the uniform bound, on all of $X $.  Completeness of $X $ gives the strong
  limit $S(t) $ (\texttt{HilleYosidaRay.semigroupOperator}); the semigroup law and
  $S(0)=\id $ pass to the limit, and strong continuity at $t=0 $ comes from the convergence
  being uniform on compact time intervals (\texttt{tendstoUniformlyOn\_approx}).  The
  exponential bound survives the limit, and the majorant half of
  Lemma~\ref{lem:dynamics:growth-bound-norm-bounds} turns it into $\omega_0(S)\leq M $.  For
  the realization, fix a real $\lambda_0>M $ and transport the ambient orbit of
  $(\lambda_0-A)d $ back to $D $ along $(\lambda_0-A)^{-1} $; because $S(t) $ commutes with
  $R(\lambda_0) $, the embedding of that domain-valued curve is $S(t)\iota(d) $ and the
  operator applied to it is $S(t)(Ad) $.  Its one-sided derivative at each $t\geq0 $ then
  comes from the integral identity
  $S(t)\iota(d)-\iota(d)=\int_0^tS(\tau)(Ad)\dd\tau $ by the second fundamental theorem of
  calculus, the endpoint $t=0 $ included.
\end{proof}

\begin{lemma}[The generated semigroup is the only one]
  \label{lem:dynamics:semigroup-unique}
  \lean{FastDynamo.StronglyContinuousSemigroup.operator_eq_of_realization,
    FastDynamo.HilleYosidaRay.semigroup_unique}
  \leanok
  \uses{def:dynamics:semigroup-interface, thm:dynamics:hille-yosida}
  Let $A $ have a densely embedded domain and unique forward strong solutions.  Then two
  strongly continuous semigroups admitting a realization by $A $ have the same evolution
  operator at every time.  In particular, if $A $ carries a Hille--Yosida ray, every
  strongly continuous semigroup realizing $A $ equals the one built in
  Theorem~\ref{thm:dynamics:hille-yosida}.
\end{lemma}

This is what earns the definite article in ``\emph{the} semigroup of $A $'' below, and it
removes any content from the choices made inside the construction: the semigroup does not
depend on which resolvent certificate the ray carries.  In Lean:
\texttt{StronglyContinuousSemigroup.operator\_eq\_of\_realization}
(\texttt{FastDynamo/Dynamics/StronglyContinuousSemigroup.lean}) and
\texttt{HilleYosidaRay.semigroup\_unique}
(\texttt{FastDynamo/Dynamics/HilleYosida/Generation.lean}).

\begin{proof}
  \leanok
  \uses{def:dynamics:semigroup-interface}
  On the image of $\iota $ the two realizations give two forward strong solutions with the
  same initial state, so forward uniqueness makes the two orbits agree there; both evolution
  operators are bounded and the image of $\iota $ is dense, so they agree on $X $.  As in
  Lemma~\ref{lem:dynamics:conditional-growth-bound}, forward uniqueness is a hypothesis of
  the statement rather than a lemma the proof calls: the concrete instance
  Lemma~\ref{lem:dynamics:strong-solution-uniqueness} lives in a Lean module that
  \emph{imports} this one, so it is not a \texttt{\textbackslash uses} of this proof.
\end{proof}

\begin{lemma}[Growth-bound sandwich for a Hille--Yosida generator]
  \label{lem:dynamics:hille-yosida-growth-sandwich}
  \lean{FastDynamo.re_le_growthBound_of_hilleYosidaRay_of_eigenpair}
  \leanok
  \uses{thm:dynamics:hille-yosida, def:dynamics:growth-bound}
  Let $A $ carry a Hille--Yosida ray of constant $M $, let its domain embed densely, and let
  its forward strong solutions be unique.  Then for every eigenpair $Ab=\lambda\iota(b) $
  with $\iota(b)\neq0 $,
  \[
    \re\lambda\ \leq\ \omega_0(S)\ \leq\ M ,
  \]
  $S $ being the semigroup of Theorem~\ref{thm:dynamics:hille-yosida}.
\end{lemma}

This is the shape in which the rest of the chapter uses the generation theorem: it converts
an eigenvalue of the operator into a two-sided pin on the growth bound of a
constructed semigroup.  It is
\texttt{re\_le\_growthBound\_of\_hilleYosidaRay\_of\_eigenpair}.

\begin{proof}
  \leanok
  \uses{lem:dynamics:conditional-growth-bound, thm:dynamics:hille-yosida}
  The upper bound is the exponential norm bound of Theorem~\ref{thm:dynamics:hille-yosida}.
  For the lower bound, the realization and forward uniqueness put the semigroup orbit of the
  eigenvector equal to the eigen-trajectory $t\mapsto\e^{t\lambda}b $, which is exactly the
  hypothesis discharged in Lemma~\ref{lem:dynamics:conditional-growth-bound}; its conclusion
  $\re\lambda\leq\omega_0(S) $ applies verbatim.
\end{proof}

\begin{lemma}[Invariant closed subspaces and the restricted semigroup]
  \label{lem:dynamics:invariant-restriction}
  \lean{FastDynamo.HilleYosidaRay.semigroup_mem_of_closed_resolvent_invariant,
    FastDynamo.HilleYosidaRay.restrictedSemigroup,
    FastDynamo.StronglyContinuousSemigroup.restrict,
    FastDynamo.StronglyContinuousSemigroup.growthBound_restrict_le,
    FastDynamo.StronglyContinuousSemigroup.congr,
    FastDynamo.StronglyContinuousSemigroup.growthBound_congr,
    FastDynamo.StronglyContinuousSemigroup.growthBound_mono_of_norm_le}
  \leanok
  \uses{thm:dynamics:hille-yosida, def:dynamics:growth-bound}
  Let $A $ carry a Hille--Yosida ray of constant $M $ and a densely embedded domain, and let
  $W\subset X $ be a closed subspace such that $R(r)W\subset W $ for every real $r>M $.  Then
  the generated semigroup leaves $W $ invariant, and its restriction $S|_W $ is a strongly
  continuous semigroup on $W $ with $\norm{S|_W(t)}\leq\e^{Mt} $ and
  $\omega_0(S|_W)\leq\omega_0(S) $.  Moreover, conjugating a strongly continuous semigroup by
  a surjective linear isometry changes no operator norm and hence no growth bound, and
  $\norm{S(t)}\leq\norm{T(t)} $ for all $t\geq0 $ implies $\omega_0(S)\leq\omega_0(T) $.
\end{lemma}

The invariance hypothesis is checked on the resolvent -- an explicit operator -- and
propagated to the semigroup along the approximation scheme,
$R(r)\rightsquigarrow A_r\rightsquigarrow\e^{tA_r}\rightsquigarrow S(t) $; closedness of $W $
is used exactly once, to pass the exponential series to the limit.  This is the only
available route, since the generated semigroup is a limit and not a formula.  Note the
direction: passing to an invariant subspace can strictly \emph{lower} a growth bound, so the
reverse inequality is false in general and is never claimed; a lower bound on
$\omega_0(S|_W) $ has to be proved on $W $ itself.  In Lean:
\texttt{FastDynamo/Dynamics/HilleYosida/Invariance.lean} for the invariance and
\texttt{HilleYosidaRay.restrictedSemigroup}, and
\texttt{FastDynamo/Dynamics/SemigroupRestriction.lean} for
\texttt{StronglyContinuousSemigroup.restrict}, \texttt{congr} and the two growth-bound
comparisons.

\begin{proof}
  \leanok
  \uses{thm:dynamics:hille-yosida, def:dynamics:growth-bound}
  A closed subspace invariant under a bounded operator $B $ is invariant under $\e^B $,
  because the exponential series converges in operator norm and every partial sum lies in
  $W $; the Yosida approximant $A_r=r^2R(r)-r $ preserves $W $ as soon as $R(r) $ does on the
  whole ray, hence so does $\e^{tA_r} $, and $S(t)x $ is the limit of $\e^{tA_r}x $.  The
  restricted family is the ambient one read through the subtype inclusion, so the three
  semigroup laws and strong continuity are inherited and $\norm{S|_W(t)}\leq\norm{S(t)} $.
  The comparison of growth bounds is monotonicity of the $\limsup $, the times with
  $\norm{S(t)}=0 $ contributing $-\infty $; for a surjective isometry both conjugations
  preserve norms, so the comparison holds in both directions and the growth bounds are equal.
\end{proof}

\section{The induction operator on one axial sector}

The Hille--Yosida ray of each axial sector comes from the dissipativity proved for the
reduced fixed-sector operator $A_{\eps,K} $ of \eqref{eq:main:torus-sector}.  What dissipativity of
$A_{\eps,K} $ alone gives is the lower bound $(\re z-M_u)\norm{\iota b}\leq\norm{(z-A)b} $,
that is, injectivity with closed range; surjectivity of the pencil is the Lumer--Phillips
half, and it comes from the same bound for the formal adjoint.

\begin{lemma}[The displayed formal adjoint of the reduced operator]
  \label{lem:dynamics:reduced-adjoint}
  \lean{FastDynamo.reducedInductionAdjointOperator,
    FastDynamo.reducedInductionStretchingTranspose,
    FastDynamo.reducedInductionOperator_formalAdjointPairing,
    FastDynamo.re_inner_reducedInductionAdjointAction_le,
    FastDynamo.reducedInductionAdjointOperator_forwardStrongSolutionsUnique,
    FastDynamo.reducedInductionOperator_isAdjointRealization}
  \leanok
  \uses{def:main:induction-operator, lem:dynamics:energy-bound}
  For a transverse divergence-free $u\in\W^{1,\infty} $, $\eps\geq0 $ and $K\in\Z $ put
  \[
    A^*_{\eps,K}b
    =\eps\bigl(\Delta_\perp-(2\pi K)^2\bigr)b
    +u_\perp\cdot\nabla_\perp b+2\pi iKu_3b
    +(\nabla_\perp u)^{\mathsf T}b,
    \qquad D(A^*_{\eps,K})=\Hs^2(\T^2;\CC^3),
  \]
  on the same domain and with the same $\Lp^2 $ embedding as $A_{\eps,K} $.  Then
  $\inner{A_{\eps,K}b}{c}=\inner{b}{A^*_{\eps,K}c} $ for all $b,c $ in the domain, the
  numerical range of $A^*_{\eps,K} $ lies in the same half-plane,
  \[
    \re\inner{b}{A^*_{\eps,K}b}\leq M_u\norm{b}_2^2,
  \]
  and $A^*_{\eps,K} $ has unique forward strong solutions.  If moreover $\eps>0 $, then
  $A^*_{\eps,K} $ is a maximal adjoint realization of $A_{\eps,K} $, and $A_{\eps,K} $ has a
  closed graph.
\end{lemma}

Two devices let this rest on facts proved elsewhere.  The transposed stretching matrix
is the pointwise adjoint $\star $ of the literal stretching matrix field of
\texttt{FastDynamo/Global/TorusLocalCoreMeasurability.lean}, so no new $\W^{1,\infty} $ datum is
introduced; and the three bilinear integration-by-parts identities (diffusion, transverse
transport, axial transport) are obtained from the \emph{quadratic} identities proved
in \texttt{FastDynamo/Global/ReducedInductionEnergyIdentity.lean} by complex polarization,
which over $\CC $ loses nothing.  In particular the divergence-freeness of $u $ enters
exactly where it enters the energy identity of Lemma~\ref{lem:dynamics:energy-bound}.  The
passage from the formal pairing to a \emph{maximal} adjoint realization, and with it the
closed graph, needs one far real reference point at which both $A_{\eps,K} $ and
$A^*_{\eps,K} $ are invertible; that is where $\eps>0 $ is used.  Module
\texttt{FastDynamo/Global/ReducedInductionAdjoint.lean}, with the reference point and the
maximality in \texttt{FastDynamo/Global/ReducedInductionHalfPlaneResolvent.lean}.

\begin{proof}
  \leanok
  \uses{lem:dynamics:energy-bound}
  A sesquilinear form over $\CC $ vanishing on the diagonal vanishes identically (test with
  $b+c $ and $b+ic $), so each quadratic identity
  $\im\inner{b}{\eps\Delta_\perp b}=0 $, $\re\inner{b}{u_\perp\cdot\nabla_\perp b}=0 $,
  $\re\inner{b}{2\pi iKu_3b}=0 $ upgrades to the corresponding bilinear pairing identity, and
  the stretching term contributes the pointwise adjoint relation
  $\inner{(b_\perp\cdot\nabla_\perp)u}{c}=\inner{b}{(\nabla_\perp u)^{\mathsf T}c} $
  integrated over $\T^2 $.  Summing the four gives the formal pairing.  The half-plane bound
  for $A^*_{\eps,K} $ is then the same computation as for $A_{\eps,K} $, the diffusion term
  being nonpositive and the transposed stretching matrix having the same $\Lp^\infty $ norm;
  uniqueness of forward strong solutions follows from
  Lemma~\ref{lem:dynamics:strong-solution-uniqueness} applied to $A^*_{\eps,K} $.  For the
  last claim, the exponential-ray Neumann construction of
  \texttt{FastDynamo/Operator/FixedSectorLowerOrderDecay.lean} inverts the diffusion pencil
  far out on the positive real axis and absorbs the lower-order part; it applies verbatim to
  $A^*_{\eps,K} $, whose diffusion pencil is the same and whose lower-order part decays
  relative to the diffusion inverse in the same way, the transposed stretching matrix being
  an $\Lp^\infty $ multiplier composed with the $\Hs^2\to\Lp^2 $ embedding.
\end{proof}

\begin{proposition}[The reduced operator has a resolvent on the right half-plane]
  \label{prop:dynamics:reduced-half-plane-resolvent}
  \lean{FastDynamo.reducedInductionRightHalfPlane,
    FastDynamo.reducedInductionOperator_pencil_bijective_of_lt_re,
    FastDynamo.reducedInductionRightHalfPlaneResolvent,
    FastDynamo.norm_stateResolvent_reducedInductionOperator_rightHalfPlane_le,
    FastDynamo.norm_stateResolvent_reducedInductionOperator_realRay_le}
  \leanok
  \uses{lem:dynamics:reduced-adjoint, lem:dynamics:energy-bound}
  Let $u $ be transverse divergence-free, $\eps>0 $ and $K\in\Z $.  Then for every
  $z\in\CC $ with $\re z>M_u $ the pencil $z-A_{\eps,K}:\Hs^2(\T^2;\CC^3)\to\Lp^2 $ is
  bijective, and the resolvents assemble into a single continuous family on the open
  half-plane $\{\re z>M_u\} $ with
  \[
    \norm{(z-A_{\eps,K})^{-1}}_{\Lp^2\to\Lp^2}\leq(\re z-M_u)^{-1} .
  \]
  In particular the bound holds at every real $r>M_u $.
\end{proposition}

This is the Lumer--Phillips input of Theorem~\ref{thm:dynamics:sector-semigroup}: the real-ray
specialization \texttt{norm\_stateResolvent\_reducedInductionOperator\_realRay\_le} is
exactly the hypothesis of Definition~\ref{def:dynamics:hille-yosida-ray}.  Module
\texttt{FastDynamo/Global/ReducedInductionHalfPlaneResolvent.lean}; the set
$\{\re z>M_u\} $ is \texttt{reducedInductionRightHalfPlane u}.  The coherent family is
assembled from the pointwise inverses by the repository's rational-propagation construction
\texttt{DomainOperator.hasResolventOnOfPointwiseTwoSidedInverses}, following the template
already used for the periodic background, so no continuity in $z $ is proved by hand.

\begin{proof}
  \leanok
  \uses{lem:dynamics:reduced-adjoint, lem:dynamics:energy-bound}
  The numerical-range bound of Lemma~\ref{lem:dynamics:energy-bound} places the numerical
  range of $A_{\eps,K} $ in $\{\re\nu\leq M_u\} $, which gives the ambient pencil lower
  bound $(\re z-M_u)\norm{\iota b}\leq\norm{(z-A_{\eps,K})b} $; hence the pencil is injective
  with closed range, and the same bound at $\bar z $ for the adjoint
  (Lemma~\ref{lem:dynamics:reduced-adjoint}) makes the range dense.  Closed plus dense is
  everything, so the pencil is bijective, and the open mapping theorem gives a bounded
  inverse; dividing the lower bound by $\norm{(z-A_{\eps,K})b} $ turns it into the displayed
  estimate.  The closed graph of $A_{\eps,K} $ and the maximality of the adjoint, both from
  the far real reference point of Lemma~\ref{lem:dynamics:reduced-adjoint}, are what make
  the range closed and identify the displayed formal adjoint with the Hilbert adjoint.
\end{proof}

\begin{theorem}[The Hille--Yosida ray of one axial sector]
  \label{thm:dynamics:sector-semigroup}
  \lean{FastDynamo.reducedInductionHilleYosidaRay,
    FastDynamo.reducedInductionHilleYosidaRay_spectralSet,
    FastDynamo.norm_stateResolvent_reducedInductionHilleYosidaRay_le}
  \leanok
  \uses{prop:dynamics:reduced-half-plane-resolvent, def:dynamics:hille-yosida-ray}
  Let $u $ be transverse divergence-free, $\eps>0 $ and $K\in\Z $.  Then $A_{\eps,K} $
  carries a Hille--Yosida ray of constant $M_u $ in the sense of
  Definition~\ref{def:dynamics:hille-yosida-ray}, whose spectral set is the whole open
  half-plane $\{\re z>M_u\} $ and whose state resolvent obeys
  $\norm{(r-A_{\eps,K})^{-1}}_{\Lp^2\to\Lp^2}\leq(r-M_u)^{-1} $ at every real $r>M_u $.
\end{theorem}

The ray is the input that the next section sums over the axial sectors: its resolvent
certificate is the summand of the all-frequency resolvent of
Proposition~\ref{prop:dynamics:product-resolvent}.  No semigroup is generated at the level
of a single sector; the semigroup of the development is that of the all-frequency operator
(Theorem~\ref{thm:dynamics:product-semigroup}).  The spectral set being the whole
half-plane is definitional: \texttt{HilleYosidaRay.ofHalfPlane} keeps the half-plane
certificate and only adds the ray estimate.  Module
\texttt{FastDynamo/Dynamics/ReducedInductionHilleYosidaRay.lean}.

\begin{proof}
  \leanok
  \uses{prop:dynamics:reduced-half-plane-resolvent, def:dynamics:hille-yosida-ray}
  Restrict the half-plane resolvent of
  Proposition~\ref{prop:dynamics:reduced-half-plane-resolvent} to the real ray by
  \texttt{HilleYosidaRay.ofHalfPlane}; the ray estimate is the half-plane estimate at real
  $z=r $.
\end{proof}

\section{From the axial sectors to the solenoidal semigroup}
\label{sec:dynamics:solenoidal}

The paper's semigroup acts on $\Lp^2_\sigma(\T^3) $.  Three steps separate it from the
fixed-sector operator $A_{\eps,K} $ and its Hille--Yosida ray
(Theorem~\ref{thm:dynamics:sector-semigroup}): summing the axial
sectors into the all-frequency operator $\mathcal A_\eps $ on $\Lp^2(\T^3;\CC^3) $,
restricting to the closed solenoidal subspace, and transporting to the literal three-torus.
The first step is the only one with analytic content, and it is done at the level of
resolvents rather than of semigroups: the sector resolvents sum to a resolvent as soon as
one elliptic estimate is uniform in the axial index, and the generation theorem is then
applied once, to the sum.

\begin{lemma}[The axial sectors decompose product $\Lp^2$ and diagonalise the operator]
  \label{lem:dynamics:axial-hilbert-sum}
  \lean{FastDynamo.isHilbertSum_axialLiftL2, FastDynamo.axialDecomposition,
    FastDynamo.axialSectorProjL2, FastDynamo.hasSum_axialLiftL2_axialSectorProjL2,
    FastDynamo.hasSum_norm_sq_axialSector, FastDynamo.hasSum_re_inner_axialSector,
    FastDynamo.axialSectorProjH2, FastDynamo.norm_axialSectorProjH2_le,
    FastDynamo.hasSum_axialLiftH2_axialSectorProjH2,
    FastDynamo.axialSectorProjL2_productInductionAction}
  \leanok
  \uses{def:main:induction-operator}
  The isometries $b\mapsto\e^{2\pi iKx_3}b $, $K\in\Z $, make $\Lp^2(\T^3;\CC^3) $ the
  Hilbert sum of copies of $\Lp^2(\T^2;\CC^3) $: they have pairwise orthogonal ranges, and
  the closed span of those ranges is everything.  Writing $P_K $ for the resulting
  orthogonal projection onto the $K $-th sector, read as a transverse amplitude,
  \[
    \norm{V}_2^2=\sum_{K\in\Z}\norm{P_KV}_2^2
    \qquad(V\in\Lp^2(\T^3;\CC^3)).
  \]
  The same decomposition lifts to the $\Hs^2 $ domain: there are projections
  $\Pi_K:\Hs^2(\T^3;\CC^3)\to\Hs^2(\T^2;\CC^3) $ of norm at most one, uniformly in $K $,
  such that every $B\in\Hs^2(\T^3;\CC^3) $ is the $\Hs^2 $-convergent sum
  $B=\sum_K\e^{2\pi iKx_3}\Pi_KB $.  The all-frequency operator is block diagonal for this
  decomposition,
  \[
    P_K\bigl(\mathcal A_\eps B\bigr)=A_{\eps,K}\bigl(\Pi_KB\bigr)
    \qquad(K\in\Z,\ B\in\Hs^2(\T^3;\CC^3)).
  \]
\end{lemma}

The ambient statement is Fourier analysis on the product torus together with Mathlib's
\texttt{IsHilbertSum} API: orthogonality is Fubini plus orthonormality of the circle
characters, and totality is completeness of the product Fourier system.  The domain-level
projection is not the adjoint of the $\Hs^2 $ lift -- that lift is only
$1+(2\pi K)^2 $-bounded -- but the ambient sector projection of the stored field paired with
the $K $-th row of the stored Fourier certificate, rescaled by a ratio bounded by one; this
is what makes $\norm{\Pi_K}\leq1 $ uniform in $K $.  Modules
\texttt{FastDynamo/Dynamics/Assembly/AxialHilbertSum.lean} and
\texttt{FastDynamo/Dynamics/Assembly/AxialSectorDomain.lean}.

\begin{proof}
  \leanok
  \uses{def:main:induction-operator}
  For orthogonality, $\inner{\e^{2\pi ikx_3}b}{\e^{2\pi iKx_3}c} $ splits by Fubini into the
  transverse inner product times $\int_{\T}\e^{2\pi i(K-k)x_3}\dd x_3 $, which is $1 $ when
  $k=K $ and $0 $ otherwise.  For totality, a field orthogonal to every sector has every
  product Fourier coefficient zero, by testing against a transverse scalar field placed in
  one vector coordinate.  Parseval is then Mathlib's, and the domain-level sum follows from
  the $\ell^2 $ expansion of the stored certificate together with closedness of the Fourier
  graph submodule.  Block diagonality holds on a single lifted sector by the exact sector
  identity of Definition~\ref{def:main:induction-operator}, and both sides are continuous in
  $B $, so the $\Hs^2 $-convergent sector expansion extends it to the whole domain.
\end{proof}

\begin{lemma}[The all-frequency operator is dissipative up to $M_u$]
  \label{lem:dynamics:product-dissipative}
  \lean{FastDynamo.re_inner_productInductionAction_le,
    FastDynamo.productInductionOperator_forwardStrongSolutionsUnique}
  \leanok
  \uses{lem:dynamics:axial-hilbert-sum, lem:dynamics:energy-bound}
  For every transverse divergence-free $u\in\W^{1,\infty} $, every $\eps\geq0 $ and every
  $B\in\Hs^2(\T^3;\CC^3) $,
  \[
    \re\inner{B}{\mathcal A_\eps B}\leq M_u\norm{B}_2^2 ,
  \]
  and consequently $\mathcal A_\eps $ has unique forward strong solutions, with no further
  hypothesis.
\end{lemma}

The energy competition is never redone at the product level: transport skewness and the sign
of the diffusion are consumed only through the reduced operator, sector by sector, which is
why the constant is the sector constant $M_u $ and is independent of $\eps $ and of the axial
index.  The statement is about the auxiliary operator $\mathcal A_\eps $; uniqueness for the
$\Hs^{-2}$ curl-form realisation in which clause~(16) of
Definition~\ref{def:main:physical-clauses} is stated is
Lemma~\ref{lem:dynamics:physical-uniqueness}, which reduces to this one.  Module
\texttt{FastDynamo/Dynamics/Assembly/AxialSectorDomain.lean}.

\begin{proof}
  \leanok
  \uses{lem:dynamics:axial-hilbert-sum, lem:dynamics:energy-bound,
    lem:dynamics:strong-solution-uniqueness}
  Both sides are sums over the axial sectors: the left-hand side by block diagonality and
  Parseval for the real part of the inner product, the right-hand side by Parseval for the
  norm.  Termwise the inequality is the fixed-sector bound of
  Lemma~\ref{lem:dynamics:energy-bound} with the same constant $M_u $, so it survives the
  summation.  Uniqueness is then
  Lemma~\ref{lem:dynamics:strong-solution-uniqueness} applied to $\mathcal A_\eps $.
\end{proof}

\begin{lemma}[An elliptic estimate uniform in the axial index]
  \label{lem:dynamics:uniform-sector-estimate}
  \lean{FastDynamo.productH2Weight_le_diffusion_bound_times_symbol,
    FastDynamo.norm_axialMultiplier_mul_inverseSymbol_le,
    FastDynamo.exists_uniform_sigma_norm_lowerOrder_comp_diffusionInverse_le_half,
    FastDynamo.norm_axialLiftH2_le_graph}
  \leanok
  \uses{def:main:induction-operator}
  Let $\eps>0 $.  There is a threshold $\sigma $, as large as one likes, at which the
  lower-order part of $A_{\eps,K} $ composed with the inverse of the diffusion pencil at
  $\sigma $ has norm at most $\tfrac12 $ \emph{simultaneously for every} $K\in\Z $; and there
  is $C>0 $ with
  \[
    \norm{\e^{2\pi iKx_3}b}_{\Hs^2(\T^3)}
    \leq C\bigl(\norm{b}_{\Lp^2(\T^2)}+\norm{A_{\eps,K}b}_{\Lp^2(\T^2)}\bigr)
    \qquad(K\in\Z,\ b\in\Hs^2(\T^2;\CC^3)).
  \]
\end{lemma}

This is the analytic crux of the assembly, and the only place where anything beyond algebra
is needed.  The fixed-sector chain of \texttt{FastDynamo/Operator/FixedSectorLowerOrderDecay.lean}
produces both a threshold and an $\Hs^2 $ bound that degrade as $\abs K\to\infty $, because
the axial transport term carries $\norm{2\pi K} $ and the sector lift is only
$1+(2\pi K)^2 $-bounded.  Two $K $-free scalar estimates remove both dependencies: the
product-weight elliptic gain
$\texttt{productH2Weight}\,(n,K)\leq(\sigma^{-1}+\eps^{-1})\,\texttt{symbol}_{\eps,\sigma,K}(n) $,
the three-dimensional analogue of the transverse gain already used to land the diffusion
inverse in $\Hs^2 $; and the arithmetic--geometric bound
$2\pi\abs K\cdot\texttt{inverseSymbol}_{\eps,\sigma,K}(n)\leq(2\sqrt{\eps\sigma})^{-1} $,
which absorbs the axial factor into the diffusion symbol.  The uniform threshold is a separate construction alongside the fixed-sector one, which
the fixed-sector statements do not depend on.
Module \texttt{FastDynamo/Dynamics/Assembly/UniformSectorResolvent.lean}.

\begin{proof}
  \leanok
  \uses{def:main:induction-operator}
  The lower-order part of $A_{\eps,K} $ is transverse transport plus axial transport plus
  stretching.  The transverse-transport and stretching bounds relative to the diffusion
  inverse are already $K $-free; the axial one is made $K $-free by the
  arithmetic--geometric estimate above, and all three tend to zero as $\sigma\to\infty $, so
  a single $\sigma $ works for every sector.  At such a $\sigma $ the Neumann series
  converges, which bounds $\norm{b}_{\Hs^2} $ by the diffusion-pencil norm and hence, through
  the product-weight gain applied to the lifted field, the $\Hs^2(\T^3) $ norm of
  $\e^{2\pi iKx_3}b $ by the graph norm of $A_{\eps,K} $ with a constant independent of $K $.
\end{proof}

\begin{proposition}[The all-frequency resolvent and its Hille--Yosida ray]
  \label{prop:dynamics:product-resolvent}
  \lean{FastDynamo.productResolventInv,
    FastDynamo.productInductionHalfPlaneResolvent,
    FastDynamo.norm_stateResolvent_productInduction_le,
    FastDynamo.productHilleYosidaRay}
  \leanok
  \uses{prop:dynamics:reduced-half-plane-resolvent, lem:dynamics:axial-hilbert-sum,
    lem:dynamics:uniform-sector-estimate, def:dynamics:hille-yosida-ray}
  Let $u $ be transverse divergence-free and $\eps>0 $.  On the half-plane $\{\re z>M_u\} $
  the operator
  \[
    R(z)x:=\sum_{K\in\Z}\e^{2\pi iKx_3}\,(z-A_{\eps,K})^{-1}(P_Kx)
  \]
  converges in $\Hs^2(\T^3;\CC^3) $, is a two-sided inverse of the pencil
  $z-\mathcal A_\eps $, depends continuously on $z $, and obeys
  $\norm{R(z)}_{\Lp^2\to\Lp^2}\leq(\re z-M_u)^{-1} $.  Hence $\mathcal A_\eps $ carries a
  Hille--Yosida ray of constant $M_u $.
\end{proposition}

The constant is not recomputed at the product level: it is the sector constant $M_u $, which
is independent of $\eps $ and of the axial index, and Parseval turns the sector estimates
into the same estimate for the sum.  Module
\texttt{FastDynamo/Dynamics/Assembly/ProductResolvent.lean}.

\begin{proof}
  \leanok
  \uses{prop:dynamics:reduced-half-plane-resolvent, lem:dynamics:axial-hilbert-sum,
    lem:dynamics:uniform-sector-estimate}
  Square summability of the summands is where the uniform estimate is used: by
  Lemma~\ref{lem:dynamics:uniform-sector-estimate} and the identity
  $A_{\eps,K}(z-A_{\eps,K})^{-1}=z(z-A_{\eps,K})^{-1}-\id $, the $\Hs^2 $ norm of the
  $K $-th summand is at most a constant depending only on $z $, $\eps $ and $M_u $ times
  $\norm{P_Kx} $, and $\sum_K\norm{P_Kx}^2=\norm{x}^2 $.  A family of sector lifts with square-summable graph
  norms sums in $\Hs^2(\T^3;\CC^3) $, the physical halves through the Hilbert-sum
  decomposition of Lemma~\ref{lem:dynamics:axial-hilbert-sum} and the stored certificates
  through their $\ell^2 $ expansion, the two limits being glued by closedness of the Fourier
  graph submodule.  Both inverse equations hold sector by sector, by block diagonality, and
  transport to the sum by uniqueness of limits.  The Parseval bound on the sum gives the
  estimate.  Continuity in $z $ is not proved by hand: the pointwise inverses are fed to the
  same rational-propagation construction as in
  Proposition~\ref{prop:dynamics:reduced-half-plane-resolvent}, whose stored inverse is
  continuous and, two-sided pencil inverses being unique, agrees with the sector sum.
  Restricting to the real ray gives the Hille--Yosida ray.
\end{proof}

\begin{theorem}[The all-frequency induction semigroup]
  \label{thm:dynamics:product-semigroup}
  \lean{FastDynamo.productInductionSemigroup, FastDynamo.productInductionRealization,
    FastDynamo.norm_productInductionSemigroup_le,
    FastDynamo.growthBound_productInductionSemigroup_le,
    FastDynamo.re_le_growthBound_productInductionSemigroup_of_eigenpair,
    FastDynamo.productInductionSemigroup_unique}
  \leanok
  \uses{prop:dynamics:product-resolvent, thm:dynamics:hille-yosida,
    lem:dynamics:product-dissipative, lem:dynamics:hille-yosida-growth-sandwich,
    lem:dynamics:semigroup-unique}
  Let $u $ be transverse divergence-free and $\eps>0 $.  Then the auxiliary operator
  $\mathcal A_\eps $ with domain $\Hs^2(\T^3;\CC^3) $ generates a strongly continuous
  semigroup on $\Lp^2(\T^3;\CC^3) $, realizing it in the sense of
  Definition~\ref{def:dynamics:semigroup-interface}, with
  $\norm{\e^{t\mathcal A_\eps}}\leq\e^{M_ut} $, hence
  $\omega_0\leq M_u $, and with $\re\lambda\leq\omega_0 $ for every eigenpair of
  $\mathcal A_\eps $.  It is the only strongly continuous semigroup realizing
  $\mathcal A_\eps $ in this sense.
\end{theorem}

This is the generation step of item~(1) of Lemma~\ref{lem:semigroup-spectral-growth}, for the
\emph{auxiliary} operator on all of $\Lp^2(\T^3;\CC^3)$.  Uniqueness is
Lemma~\ref{lem:dynamics:semigroup-unique} applied to the realization
(\texttt{productInductionSemigroup\_unique}), so the choice of resolvent family inside the
construction is immaterial.  The paper's operator is the curl form on $\Lp^2_\sigma$; the two
ambient operators are different on $\Lp^2(\T^3;\CC^3)$ (they differ by $-u\,\Div B $, which
vanishes only on solenoidal fields), and their actions agree on divergence-free $\Hs^2 $
fields (\texttt{productInductionAction\_toHMinus2\_eq\_physical},
\texttt{FastDynamo/Operator/Compatibility.lean}).  The generator of this semigroup is
$\mathcal A_\eps$ with domain $\Hs^2$ (\texttt{generator\_productInductionSemigroup},
Lemma~\ref{lem:dynamics:hille-yosida-generator}), and the generator of its restriction to
$\Lp^2_\sigma$ is the part of $\mathcal A_\eps$ there, which is the paper's $L_{\eps,u}$
(Theorem~\ref{thm:dynamics:solenoidal-generator}).  Module
\texttt{FastDynamo/Dynamics/Assembly/ProductSemigroup.lean}.

\begin{proof}
  \leanok
  \uses{prop:dynamics:product-resolvent, thm:dynamics:hille-yosida,
    lem:dynamics:product-dissipative, lem:dynamics:hille-yosida-growth-sandwich}
  The Hille--Yosida ray is Proposition~\ref{prop:dynamics:product-resolvent}, the
  $\Hs^2\to\Lp^2 $ inclusion has dense range
  (\texttt{FastDynamo/Analysis/ProductH2DenseEmbedding.lean}), and forward strong solutions
  are unique by Lemma~\ref{lem:dynamics:product-dissipative}.  Apply
  Theorem~\ref{thm:dynamics:hille-yosida} and
  Lemma~\ref{lem:dynamics:hille-yosida-growth-sandwich}.
\end{proof}

\begin{lemma}[The solenoidal subspace is closed and resolvent-invariant]
  \label{lem:dynamics:solenoidal-subspace}
  \lean{FastDynamo.productSolenoidalSubspace, FastDynamo.mem_productSolenoidalSubspace_iff,
    FastDynamo.isClosed_productSolenoidalSubspace,
    FastDynamo.stateResolvent_mem_productSolenoidalSubspace_of_re_pos,
    FastDynamo.productH2Embedding_mem_productSolenoidalSubspace_of_isEigenpair}
  \leanok
  \uses{def:main:induction-operator, lem:positive-eigenfunction-solenoidal}
  The set $\Lp^2_\sigma(\T^3)=\{V\in\Lp^2(\T^3;\CC^3):\Div V=0\} $ is a closed subspace of
  $\Lp^2(\T^3;\CC^3) $, hence a Hilbert space in its own right.  It is invariant under every
  state resolvent of $\mathcal A_\eps $ at a spectral parameter with positive real part: if
  $\eps\geq0 $, $u $ is transverse divergence-free, $(z-\mathcal A_\eps)^{-1} $ exists,
  $\re z>0 $ and $\Div x=0 $, then $\Div\bigl((z-\mathcal A_\eps)^{-1}x\bigr)=0 $.  Moreover
  every eigenfunction of $\mathcal A_\eps $ with eigenvalue of positive real part lies in
  $\Lp^2_\sigma(\T^3) $.
\end{lemma}

Both halves are cheap here, and neither is a new analytic input.  Solenoidality is a family
of scalar linear equations in the Fourier coefficients -- for each transverse frequency $n $
and axial frequency $k $ the combination
$\sum_{j<2}2\pi in_j\widehat V_j(n,k)+2\pi ik\widehat V_3(n,k) $ is a continuous linear
functional -- so $\Lp^2_\sigma $ is an intersection of kernels, exactly as the Fourier graph
domain is.  Invariance is the resolvent form of
Lemma~\ref{lem:positive-eigenfunction-solenoidal}: taking the weak divergence of the
resolvent equation turns it into a scalar advection--diffusion equation for $d=\Div B $ whose
forcing is the weak divergence of the solenoidal datum and therefore vanishes, and the same
energy endpoint forces $d=0 $.  The statement is quantified over an \emph{arbitrary}
resolvent certificate, so it does not depend on the construction of
Proposition~\ref{prop:dynamics:product-resolvent}.  Module
\texttt{FastDynamo/Dynamics/Assembly/SolenoidalSubspace.lean}.

\begin{proof}
  \leanok
  \uses{lem:positive-eigenfunction-solenoidal}
  A submodule cut out by a family of continuous linear functionals is closed, and a closed
  subspace of a Hilbert space is complete.  For invariance, write $B=(z-\mathcal A_\eps)^{-1}x $
  and apply the weak divergence to $zB-\mathcal A_\eps B=x $: incompressibility of $u $ makes
  the transport and stretching terms combine into $-(u\cdot\nabla)\Div B $, so
  $d=\Div B $ satisfies $(z-\eps\Delta+u\cdot\nabla)d=\Div x=0 $ weakly, and the energy
  endpoint of Lemma~\ref{lem:positive-eigenfunction-solenoidal} with $\re z>0 $ gives $d=0 $.
  The eigenfunction statement is the same argument with $x=0 $.
\end{proof}

\begin{theorem}[The induction semigroup on the solenoidal space]
  \label{thm:dynamics:solenoidal-semigroup}
  \lean{FastDynamo.productInductionSemigroup_mem_productSolenoidalSubspace,
    FastDynamo.solenoidalInductionSemigroup,
    FastDynamo.growthBound_solenoidalInductionSemigroup_le,
    FastDynamo.re_le_growthBound_solenoidalInductionSemigroup_of_eigenpair,
    FastDynamo.literalSolenoidalSubspace, FastDynamo.isClosed_literalSolenoidalSubspace,
    FastDynamo.exists_ne_zero_literalSolenoidalSubspace,
    FastDynamo.solenoidalToLiteral, FastDynamo.literalInductionSemigroup,
    FastDynamo.norm_literalInductionSemigroup_le,
    FastDynamo.growthBound_literalInductionSemigroup,
    FastDynamo.productInductionRealization_domainOrbit_mem_productSolenoidalSubspace,
    FastDynamo.productL2ToHMinus2_domainOrbit_eq_physical}
  \leanok
  \uses{thm:dynamics:product-semigroup, lem:dynamics:solenoidal-subspace,
    lem:dynamics:invariant-restriction, lem:main:literal-torus}
  Let $u $ be transverse divergence-free and $\eps>0 $.  Then $\e^{t\mathcal A_\eps} $ leaves
  $\Lp^2_\sigma(\T^3) $ invariant, and its restriction
  $T(t):=\e^{t\mathcal A_\eps}|_{\Lp^2_\sigma} $ is a strongly continuous semigroup on
  $\Lp^2_\sigma(\T^3) $ with $\norm{T(t)}\leq\e^{M_ut} $ and $\omega_0(T)\leq M_u $.  For
  every eigenpair $\mathcal A_\eps V=\lambda V $ whose eigenfunction lies in
  $\Lp^2_\sigma(\T^3) $ one has $\re\lambda\leq\omega_0(T) $.  All of this transports, with
  no change of any operator norm and hence no change of the growth bound, to the literal
  three-torus $\T^3=(\R/\Z)^3 $, where the solenoidal subspace is again closed.  Moreover
  the orbit of every solenoidal $\Hs^2 $ datum is a solenoidal $\Hs^2 $ field at every
  $t\geq0 $ and solves the paper's curl-form induction equation
  $\partial_tB=\operatorname{curl}(u\times B)+\eps\Delta B $ in $\Hs^{-2} $ there.
\end{theorem}

Three points of the passage to the subspace matter.

First, invariance is proved on the resolvent, at every real $r>M_u\geq0 $, and propagated
through the approximation scheme by Lemma~\ref{lem:dynamics:invariant-restriction}; there is
no other route, the generated semigroup being a limit and not a formula.  Second, the
spectral lower bound has to be \emph{re-proved} on the subspace: restriction can strictly
lower a growth bound, so the ambient bound of
Theorem~\ref{thm:dynamics:product-semigroup} is no evidence about the restricted one.  It is
re-proved by evaluating the restricted operator on an eigenvector that lies in
$\Lp^2_\sigma $ -- the subspace carries the ambient norm, so nothing is lost -- which is why
Lemma~\ref{lem:dynamics:solenoidal-subspace} is needed twice.  That same eigenvector shows
$\Lp^2_\sigma(\T^3)\neq0$ (\texttt{exists\_ne\_zero\_literalSolenoidalSubspace}): it is
non-zero by the definition of an eigenpair and solenoidal by the auxiliary solenoidality
theorem.  The statement is recorded because every clause of the ideal half of
Theorem~\ref{thm:main} is quantified over $B\in\Lp^2_\sigma(\T^3) $ and would hold verbatim
on the zero space; see Corollary~\ref{cor:ideal-growth}.

Third, \emph{$T(t)$ is the paper's $\e^{tL_{\eps,u}}$}: its generator is $L_{\eps,u}$ with
domain $\Hs^2\cap\Lp^2_\sigma$, and a $\C_0$-semigroup is determined by its generator
(Theorem~\ref{thm:dynamics:solenoidal-generator}).  What remains a deviation is the route: the
paper generates on $\Lp^2_\sigma$ directly, by perturbation of $\eps\Delta$, whereas this node
restricts the semigroup of the ambient auxiliary operator, and the generator is identified
afterwards.  The orbit of every solenoidal $\Hs^2 $ datum stays solenoidal and satisfies the
curl-form equation in $\Hs^{-2} $ at every $t\geq0 $
(\texttt{productL2ToHMinus2\_domainOrbit\_eq\_physical}), and it is the only solution of that
equation with that datum (Lemma~\ref{lem:dynamics:physical-uniqueness}).  So
$\gamma(u,\eps)$ (Definition~\ref{def:main:growth-rate}) is the growth bound of $T $.

Modules \texttt{FastDynamo/Dynamics/Assembly/ProductSemigroup.lean} and
\texttt{FastDynamo/Dynamics/Assembly/MainGrowthRate.lean}; the isometric identification of
the product model with the literal three-torus is the one of
Lemma~\ref{lem:main:literal-torus}.

\begin{proof}
  \leanok
  \uses{thm:dynamics:product-semigroup, lem:dynamics:solenoidal-subspace,
    lem:dynamics:invariant-restriction}
  The state resolvent of the product Hille--Yosida ray at a real $r>M_u $ preserves
  $\Lp^2_\sigma $ by Lemma~\ref{lem:dynamics:solenoidal-subspace}, the positivity of $r $
  being free since $M_u\geq0 $; the subspace is closed, so
  Lemma~\ref{lem:dynamics:invariant-restriction} propagates the invariance to the semigroup
  and produces the restriction with the same exponential bound.  For the lower bound,
  evaluate the restricted operator at time $t $ on an eigenvector $V\in\Lp^2_\sigma $: its
  orbit is $\e^{t\lambda}V $ by Theorem~\ref{thm:dynamics:product-semigroup}, so
  $\e^{t\re\lambda}\norm{V}\leq\norm{T(t)}\norm{V} $ and the minorant half of
  Lemma~\ref{lem:dynamics:growth-bound-norm-bounds} applies.  The transport is conjugation by
  a surjective linear isometry, which preserves operator norms exactly, so the growth bound
  is unchanged (Lemma~\ref{lem:dynamics:invariant-restriction}).
\end{proof}

\section{The generator}
\label{sec:dynamics:generator}

The semigroup of the previous section is built from an operator; this section proves that the
operator is its generator, which identifies the semigroup on $\Lp^2_\sigma$ with the paper's
$\e^{tL_{\eps,u}}$.

\begin{definition}[Generator, resolvent set and spectrum]
  \label{def:dynamics:generator}
  \lean{FastDynamo.StronglyContinuousSemigroup.diffQuotient,
    FastDynamo.StronglyContinuousSemigroup.generator,
    FastDynamo.StronglyContinuousSemigroup.mem_generator_domain_iff,
    FastDynamo.LinearPMap.IsResolvent, FastDynamo.LinearPMap.resolventSet,
    FastDynamo.LinearPMap.spectrum, FastDynamo.DomainOperator.toLinearPMap}
  \leanok
  \uses{def:dynamics:semigroup-interface}
  The \emph{generator} of a strongly continuous semigroup $S$ on $X$ is the operator
  \[
    Gx:=\lim_{h\downarrow0}\frac{S(h)x-x}{h},
    \qquad
    D(G)=\{x\in X:\text{the limit exists}\}.
  \]
  For a linear operator $T$ with domain $D(T)\subset X$, the \emph{resolvent set}
  $\rho(T)$ is the set of $z\in\CC$ at which $z-T:D(T)\to X$ has a bounded two-sided inverse
  $X\to X$ with range in $D(T)$, and the \emph{spectrum} is $\sigma(T)=\CC\setminus\rho(T)$.
  An operator presented by a bounded map out of a graph domain with an injective embedding
  (\texttt{DomainOperator}) is read as an operator on the range of the embedding.
\end{definition}

The definitions are the standard ones, which the paper uses without stating them.  In Lean the
generator is a \texttt{LinearPMap} (\texttt{FastDynamo/Dynamics/SemigroupGenerator.lean}), and
\texttt{LinearPMap.IsResolvent}, \texttt{resolventSet}, \texttt{spectrum} and
\texttt{DomainOperator.toLinearPMap} are in \texttt{FastDynamo/Operator/UnboundedResolvent.lean}.
The resolvent set asks for a bounded inverse explicitly; no closedness of $T$ is assumed.

\begin{lemma}[Properties of the generator]
  \label{lem:dynamics:generator-basic}
  \lean{FastDynamo.StronglyContinuousSemigroup.generator_eq_iff,
    FastDynamo.StronglyContinuousSemigroup.apply_mem_generator_domain,
    FastDynamo.StronglyContinuousSemigroup.hasDerivWithinAt_apply,
    FastDynamo.StronglyContinuousSemigroup.dense_generator_domain,
    FastDynamo.StronglyContinuousSemigroup.apply_eq_cexp_smul_of_generator_eq_smul,
    FastDynamo.StronglyContinuousSemigroup.ext_of_generator_eq,
    FastDynamo.StronglyContinuousSemigroup.mem_generator_restrict_domain_iff,
    FastDynamo.StronglyContinuousSemigroup.generator_restrict_apply,
    FastDynamo.LinearPMap.mem_spectrum_of_eigenvector,
    FastDynamo.LinearPMap.exists_eigenvector_of_mem_spectrum,
    FastDynamo.LinearPMap.IsResolvent.isCompactOperator_of_isResolvent,
    FastDynamo.DomainOperator.isResolvent_stateResolvent}
  \leanok
  \uses{def:dynamics:generator, lem:dynamics:invariant-restriction}
  Let $S$ be a strongly continuous semigroup on a complex Banach space $X$ with generator $G$.
  Then $D(G)$ is dense and invariant under $S(t)$, $GS(t)x=S(t)Gx$ on $D(G)$, and
  $t\mapsto S(t)x$ is differentiable on $[0,\infty)$ with derivative $GS(t)x$ for $x\in D(G)$;
  an eigenvector $Gx=\lambda x$ has $S(t)x=\e^{t\lambda}x$.  Two strongly continuous semigroups
  with the same generator are equal.  If $W\subset X$ is a closed subspace invariant under
  $S$, the generator of $S|_W$ is the part of $G$ in $W$: $D(G|_W)=D(G)\cap W$ with the same
  values.  Every eigenvalue of an operator lies in its spectrum; if the resolvent at one point
  is compact, it is compact at every point of the resolvent set, and then every point of the
  spectrum is an eigenvalue.
\end{lemma}

The semigroup facts are [Engel--Nagel, Lemma~II.1.3, Theorem~II.1.4, Proposition~II.2.3]; the
last two sentences are the resolvent identity and the Fredholm alternative for compact
operators.  The paper uses all of them without stating them.

\begin{proof}
  \leanok
  \uses{def:dynamics:generator}
  The derivative of $t\mapsto S(t)x$ is computed from the difference quotient on both sides;
  the integrals $\int_0^tS(s)x\dd s$ lie in $D(G)$ and approximate $x$, which gives density.
  For uniqueness, if $S$ and $T$ have the same generator then
  $s\mapsto T(t-s)S(s)x$ has derivative zero for $x\in D(G)$, so $S(t)x=T(t)x$ on a dense set
  and hence everywhere.  The restriction statement compares the difference quotients in $W$
  and in $X$.  For the spectrum: the resolvent identity
  $R(w)=R(z)+(z-w)R(w)R(z)$ transfers compactness, and for $z\in\sigma(T)$ and a compact
  resolvent $R(r)$ the operator $z-T$ is $(\id-(r-z)R(r))(r-T)$, so non-invertibility of
  $z-T$ forces a non-trivial kernel of $\id-(r-z)R(r)$ by the Fredholm alternative.
\end{proof}

\begin{lemma}[The generator of the Hille--Yosida semigroup]
  \label{lem:dynamics:hille-yosida-generator}
  \lean{FastDynamo.HilleYosidaRay.generator_semigroup,
    FastDynamo.HilleYosidaRay.tendsto_diffQuotient_semigroup_embedding,
    FastDynamo.generator_productInductionSemigroup}
  \leanok
  \uses{thm:dynamics:hille-yosida, def:dynamics:generator, lem:dynamics:generator-basic,
    thm:dynamics:product-semigroup}
  Let $A$ carry a Hille--Yosida ray whose domain embeds densely in the Banach space $X$, and let
  $S$ be the semigroup of Theorem~\ref{thm:dynamics:hille-yosida}.  Then the generator of $S$
  is $A$, with domain the range of the embedding.  In particular the generator of the
  all-frequency semigroup of Theorem~\ref{thm:dynamics:product-semigroup} is $\mathcal A_\eps$
  with domain $\Hs^2(\T^3;\CC^3)$.
\end{lemma}

This is [Engel--Nagel, Theorem~II.3.5] in the generator direction.  Modules
\texttt{FastDynamo/Dynamics/HilleYosida/GeneratorIdentification.lean} and
\texttt{FastDynamo/Dynamics/Assembly/SolenoidalGenerator.lean}.

\begin{proof}
  \leanok
  \uses{thm:dynamics:hille-yosida, lem:dynamics:generator-basic}
  The range of the embedding lies in $D(G)$ with $G=A$ there: the realization of
  Theorem~\ref{thm:dynamics:hille-yosida} differentiates the orbit at $t=0$.  Conversely, let
  $x\in D(G)$ and fix a real $r$ beyond the ray constant and with $\norm{S(1)}<\e^r$.  Put
  $w=R(r)(rx-Gx)-x$; then $w\in D(G)$ and $Gw=rw$, so $S(1)w=\e^rw$ by
  Lemma~\ref{lem:dynamics:generator-basic}, and $\norm{S(1)}<\e^r$ forces $w=0$.  Hence $x$
  is in the range of the embedding.
\end{proof}

\begin{theorem}[The induction semigroup is generated by $L_{\eps,u}$]
  \label{thm:dynamics:solenoidal-generator}
  \lean{FastDynamo.solenoidalH2, FastDynamo.productInductionAction_mem_productSolenoidalSubspace,
    FastDynamo.solenoidalInductionOperator,
    FastDynamo.generator_solenoidalInductionSemigroup,
    FastDynamo.mem_generator_solenoidalInductionSemigroup_domain_iff,
    FastDynamo.coe_generator_solenoidalInductionSemigroup,
    FastDynamo.productL2ToHMinus2_generator_solenoidalInductionSemigroup,
    FastDynamo.eq_solenoidalInductionSemigroup_of_generator_eq,
    FastDynamo.gamma_eq_growthBound_of_generator_eq,
    FastDynamo.exists_isResolvent_generator_solenoidalInductionSemigroup_eq}
  \leanok
  \uses{thm:dynamics:solenoidal-semigroup, lem:dynamics:hille-yosida-generator,
    lem:dynamics:generator-basic, lem:finish:physical-compatibility}
  Let $u$ be transverse divergence-free and $\eps>0$.  For every divergence-free
  $B\in\Hs^2$ the field $\mathcal A_\eps B$ lies in $\Lp^2_\sigma$, and it is the curl form
  $L_{\eps,u}B=\curl(u\times B)+\eps\Delta B$.  The generator of the semigroup $T$ of
  Theorem~\ref{thm:dynamics:solenoidal-semigroup} is $L_{\eps,u}$ with domain
  $\Hs^2\cap\Lp^2_\sigma$; every strongly continuous semigroup on $\Lp^2_\sigma$ with that
  generator is $T$, and its growth bound is $\gamma(u,\eps)$.
\end{theorem}

This is the generation statement of item~(1) of Lemma~\ref{lem:semigroup-spectral-growth},
without analyticity.  The Lean operator \texttt{solenoidalInductionOperator} has domain the
divergence-free $\Hs^2$ fields (\texttt{solenoidalH2}) and values the auxiliary action
$\mathcal A_\eps B$; the identification with the curl form is
\texttt{productL2ToHMinus2\_generator\_solenoidalInductionSemigroup}, an identity in the Fourier
model of $\Hs^{-2}$, into which $\Lp^2$ embeds injectively.  So the curl form of the paper's
\texttt{eq:physical-generator}, computed distributionally, is the $\Lp^2$ field
$\mathcal A_\eps B$: an equivalent formulation.  The route differs from the paper's: the
semigroup is generated on all of $\Lp^2$ by Hille--Yosida, its generator identified
(Lemma~\ref{lem:dynamics:hille-yosida-generator}), and restricted to $\Lp^2_\sigma$
(Lemma~\ref{lem:dynamics:generator-basic}, [Engel--Nagel, Proposition~II.2.3]), not generated
on $\Lp^2_\sigma$ by perturbation of $\eps\Delta$.  Module
\texttt{FastDynamo/Dynamics/Assembly/SolenoidalGenerator.lean}.

\begin{proof}
  \leanok
  \uses{lem:dynamics:hille-yosida-generator, lem:dynamics:generator-basic,
    lem:finish:physical-compatibility}
  On a divergence-free $\Hs^2$ field the lower-order part of $\mathcal A_\eps B$ is the weak
  curl $\curl(u\times B)$, which is solenoidal, and $\Delta B$ is solenoidal coefficientwise;
  the two actions agree in $\Hs^{-2}$ by Lemma~\ref{lem:finish:physical-compatibility}.  The
  generator of the restriction is the part in $\Lp^2_\sigma$ of the ambient generator
  $\mathcal A_\eps$ (Lemmas~\ref{lem:dynamics:hille-yosida-generator}
  and~\ref{lem:dynamics:generator-basic}), whose domain is
  $\Hs^2\cap\Lp^2_\sigma$.  Uniqueness is uniqueness of the semigroup with a given generator.
\end{proof}

\section{Analyticity, compactness and spectral mapping}
\label{sec:dynamics:analytic}

\begin{definition}[Analytic semigroup]
  \label{def:dynamics:analytic-semigroup}
  \lean{FastDynamo.openSector, FastDynamo.AnalyticExtension, FastDynamo.IsAnalyticSemigroup,
    FastDynamo.isOpen_openSector, FastDynamo.isPreconnected_openSector}
  \leanok
  \uses{def:dynamics:semigroup-interface}
  For $0<\delta\leq\pi/2$ let $\Sigma_\delta=\{z\in\CC\setminus\{0\}:\abs{\arg z}<\delta\}$.  A
  strongly continuous semigroup $S$ is \emph{analytic of angle $\delta$} if there is a map
  $T:\Sigma_\delta\cup[0,\infty)\to\mathcal L(X)$ that agrees with $S$ on $[0,\infty)$, obeys
  $T(z+w)=T(z)T(w)$ on $\Sigma_\delta$, is holomorphic on $\Sigma_\delta$ in the operator norm,
  and converges strongly to the identity as $z\to0$ within every smaller sector
  $\Sigma_{\delta'}$.
\end{definition}

This is [Engel--Nagel, Definition~II.4.5], the notion of the paper's
Lemma~\ref{lem:semigroup-spectral-growth}, with the angle explicit.  Module
\texttt{FastDynamo/Dynamics/AnalyticSemigroup.lean}.

\begin{lemma}[The sector condition]
  \label{lem:dynamics:sectorial}
  \lean{FastDynamo.productInductionTransport, FastDynamo.productInductionAction_eq_transport,
    FastDynamo.re_inner_productInductionAction_le_sub,
    FastDynamo.exists_norm_productInductionTransport_le,
    FastDynamo.re_inner_productInductionTransport_skew,
    FastDynamo.exists_norm_productInductionStretching_le,
    FastDynamo.exists_productInduction_sectorial}
  \leanok
  \uses{def:main:induction-operator, lem:dynamics:energy-identity}
  Let $u$ be transverse divergence-free and $\eps>0$.  There is $\omega\in\R$ with
  \[
    \abs{\im\inner{B}{\mathcal A_\eps B}}\leq\omega\norm B_2^2-\re\inner{B}{\mathcal A_\eps B}
    \qquad(B\in\Hs^2(\T^3;\CC^3)).
  \]
\end{lemma}

The paper states no such estimate; it is the numerical-range condition from which the Lean
proof derives analyticity, in place of the paper's relative bound for $\curl(u\times B)$.  The
constant is $\omega=M_u+d+c^2/(4\eps)$, with $c$ a bound of the transport term by the $\Hs^1$
seminorm and $d$ a bound of the stretching term in $\Lp^2$.  Module
\texttt{FastDynamo/Dynamics/Assembly/ProductInductionEnergy.lean}.

\begin{proof}
  \leanok
  \uses{lem:dynamics:energy-identity}
  The diffusion contributes the real number $-\eps\norm{\nabla B}_2^2$, transport is skew and
  contributes only to the imaginary part, bounded by $c\norm{\nabla B}_2\norm B_2$, and
  stretching contributes at most $d\norm B_2^2$ to either part.  Young's inequality absorbs
  $c\norm{\nabla B}_2\norm B_2$ into $\eps\norm{\nabla B}_2^2+\tfrac{c^2}{4\eps}\norm B_2^2$,
  and the energy identity bounds $\re\inner{B}{\mathcal A_\eps B}$ by
  $M_u\norm B_2^2-\eps\norm{\nabla B}_2^2$.
\end{proof}

\begin{theorem}[Analytic generation from the sector condition]
  \label{thm:dynamics:analytic-generation}
  \lean{FastDynamo.HilleYosidaRay.sectorial_yosida,
    FastDynamo.HilleYosidaRay.sectorExtension,
    FastDynamo.HilleYosidaRay.isAnalyticSemigroup_of_sectorial,
    FastDynamo.IsAnalyticSemigroup.restrict,
    FastDynamo.differentiableOn_of_forall_apply}
  \leanok
  \uses{thm:dynamics:hille-yosida, def:dynamics:analytic-semigroup,
    lem:dynamics:invariant-restriction}
  Let $X$ be a complex Hilbert space and let $A$ carry a Hille--Yosida ray with densely
  embedded domain.  If there is $\omega\in\R$ with
  $\abs{\im\inner{\iota d}{Ad}}\leq\omega\norm{\iota d}^2-\re\inner{\iota d}{Ad}$ for every
  $d$ in the domain, then the semigroup of Theorem~\ref{thm:dynamics:hille-yosida} is analytic
  of angle $\pi/4$.  The restriction of an analytic semigroup of angle $\delta$ to a closed
  invariant subspace is analytic of angle $\delta$.  A family of bounded operators on an open
  set of $\CC$ that is locally bounded and holomorphic at every vector is holomorphic in the
  operator norm.
\end{theorem}

The paper obtains analyticity from [Engel--Nagel, Theorem~III.2.10] (relative boundedness with
respect to $\eps\Delta$); this theorem replaces that step, and exhibits the angle.  Modules
\texttt{FastDynamo/Dynamics/HilleYosida/AnalyticGeneration.lean},
\texttt{FastDynamo/Dynamics/AnalyticSemigroup.lean} and
\texttt{FastDynamo/Operator/StrongHolomorphy.lean}.

\begin{proof}
  \leanok
  \uses{thm:dynamics:hille-yosida, lem:dynamics:invariant-restriction}
  The sector condition passes to the Yosida approximants $A_s=s^2R(s)-s$ of the proof of
  Theorem~\ref{thm:dynamics:hille-yosida}, with $\abs\omega+1$ in place of $\omega$, for all
  large $s$; for a bounded operator with numerical range in that sector,
  $\norm{\e^{wA_s}}\leq\e^{(\abs\omega+1)\abs w}$ on $\{\abs{\im w}\leq\re w\}$.  The
  exponentials $\e^{wA_s}$ are then Cauchy uniformly on bounded parts of that closed sector, by
  the complex form of the telescoping estimate of Theorem~\ref{thm:dynamics:hille-yosida}; the
  strong limit is the extension, holomorphic in $w$ at every vector as a locally uniform limit
  of holomorphic maps, and locally bounded, hence holomorphic in the operator norm by the
  last statement (Cauchy estimates for the first two derivatives and the mean value
  inequality).  The open sector $\Sigma_{\pi/4}$ lies in $\{\abs{\im w}\leq\re w\}$.  For the
  restriction, the map $z\mapsto(\id-P_W)T(z)x$, $x\in W$, vanishes on $[0,\infty)$ and is
  holomorphic on the connected open sector, so it vanishes there by the identity theorem.
\end{proof}

\begin{lemma}[The compact embedding]
  \label{lem:dynamics:compact-embedding}
  \lean{FastDynamo.isCompactOperator_lpDiagonal,
    FastDynamo.isCompactOperator_productH2Embedding}
  \leanok
  \uses{def:main:induction-operator}
  The embedding $\Hs^2(\T^3;\CC^3)\hookrightarrow\Lp^2(\T^3;\CC^3)$ is a compact operator.
\end{lemma}

This is the compact embedding the paper's proof of Lemma~A.1 uses.  Module
\texttt{FastDynamo/Operator/ProductH2Compact.lean}.

\begin{proof}
  \leanok
  In Fourier variables the embedding is division of the $\Hs^2$-weighted coefficients by the
  weight $1+4\pi^2\abs\xi^2$, which tends to infinity; a diagonal operator on $\ell^2$ with
  multipliers tending to zero is the operator-norm limit of its finite truncations, which have
  finite rank.
\end{proof}

\begin{lemma}[Immediate compactness]
  \label{lem:dynamics:immediate-compactness}
  \lean{FastDynamo.IsAnalyticSemigroup.exists_generator_comp,
    FastDynamo.IsAnalyticSemigroup.isCompactOperator}
  \leanok
  \uses{def:dynamics:analytic-semigroup, def:dynamics:generator,
    lem:dynamics:generator-basic}
  Let $S$ be an analytic semigroup on a Banach space with generator $G$.  For every $t>0$,
  $S(t)$ maps into $D(G)$ and $GS(t)$ is bounded.  If moreover $G$ has a compact resolvent,
  then $S(t)$ is compact for every $t>0$.
\end{lemma}

The paper cites [Engel--Nagel, Theorem~II.4.29] for immediate compactness; the Lean proof is
the direct factorisation below.  Module
\texttt{FastDynamo/Dynamics/AnalyticSemigroupGenerator.lean}.

\begin{proof}
  \leanok
  \uses{lem:dynamics:generator-basic}
  The holomorphic extension $T$ is differentiable at the real point $t$, and its complex
  derivative there is a bounded operator $D$; the real difference quotient of $S$ at $S(t)x$
  converges to $Dx$, so $S(t)x\in D(G)$ and $GS(t)=D$.  Then
  $S(t)=R(r)\bigl(rS(t)-GS(t)\bigr)$ is a compact operator composed with a bounded one.
\end{proof}

\begin{theorem}[Spectral mapping and the growth bound]
  \label{thm:dynamics:spectral-mapping}
  \lean{FastDynamo.StronglyContinuousSemigroup.exp_re_mul_le_norm,
    FastDynamo.StronglyContinuousSemigroup.operator_natCast,
    FastDynamo.StronglyContinuousSemigroup.cexp_mem_spectrum,
    FastDynamo.IsAnalyticSemigroup.exists_cexp_eq_of_mem_spectrum,
    FastDynamo.IsAnalyticSemigroup.spectrum_diff_zero,
    FastDynamo.IsAnalyticSemigroup.growthBound_eq_iSup}
  \leanok
  \uses{lem:dynamics:immediate-compactness, lem:dynamics:generator-basic,
    lem:dynamics:growth-bound-limit, def:dynamics:growth-bound}
  Let $S$ be an analytic semigroup on a complex Banach space whose generator $G$ has a compact
  resolvent.  Then for every $t>0$
  \[
    \sigma(S(t))\setminus\{0\}=\{\e^{t\lambda}:\lambda\in\sigma(G)\},
  \]
  and if $\sigma(G)\neq\emptyset$, then
  $\omega_0(S)=\sup\{\re\lambda:\lambda\in\sigma(G)\}$ and the supremum is attained at an
  eigenvalue of $G$, with an eigenvector.
\end{theorem}

The paper cites [Engel--Nagel, Corollaries~IV.3.11 and~IV.3.12] for the spectral mapping and
the spectral bound, and [Engel--Nagel, Corollary~V.3.2] for the attainment.  The Lean proof is a
direct point-spectrum argument.  Module \texttt{FastDynamo/Dynamics/SemigroupSpectralMapping.lean}.

\begin{proof}
  \leanok
  \uses{lem:dynamics:immediate-compactness, lem:dynamics:generator-basic,
    lem:dynamics:growth-bound-limit}
  Every point of $\sigma(G)$ is an eigenvalue (compact resolvent), and an eigenvector
  $Gx=\lambda x$ gives $S(t)x=\e^{t\lambda}x$, so $\e^{t\lambda}\in\sigma(S(t))$ and
  $\e^{t\re\lambda}\leq\norm{S(t)}$.  Conversely $S(t)$, $t>0$, is compact
  (Lemma~\ref{lem:dynamics:immediate-compactness}), so a non-zero $\mu\in\sigma(S(t))$ is an
  eigenvalue with finite-dimensional eigenspace; that eigenspace is invariant under the bounded
  operator $GS(t)$, which therefore has an eigenvector there, and it is an eigenvector of $G$
  whose eigenvalue $\lambda$ has $\e^{t\lambda}=\mu$.  For the growth bound: the eigenvector
  bound gives $\re\lambda\leq\omega_0(S)$; by Lemma~\ref{lem:dynamics:growth-bound-limit} and
  Gelfand's formula $\e^{\omega_0(S)}$ is the spectral radius of the compact operator $S(1)$,
  which is attained at a non-zero spectral point $\e^\lambda$, $\lambda\in\sigma(G)$, with
  $\re\lambda=\omega_0(S)$.
\end{proof}

\begin{lemma}[Mean conservation]
  \label{lem:dynamics:mean-conservation}
  \lean{FastDynamo.HasZeroMean, FastDynamo.productCoeffCLM_zero_zero_eq_integral,
    FastDynamo.hasZeroMean_productInductionAction,
    FastDynamo.exists_productSolenoidal_mean_eq_one,
    FastDynamo.one_le_norm_solenoidalInductionSemigroup,
    FastDynamo.gamma_eq_lim, FastDynamo.zero_le_gamma,
    FastDynamo.zero_mem_spectrum_generator_solenoidalInductionSemigroup}
  \leanok
  \uses{thm:dynamics:solenoidal-generator, lem:dynamics:growth-bound-limit,
    lem:dynamics:generator-basic}
  Let $u$ be transverse divergence-free and $\eps>0$.  For every divergence-free
  $B\in\Hs^2$, $\int_{\T^3}L_{\eps,u}B\dd x=0$; $\Lp^2_\sigma$ contains a constant field of norm
  one; hence $0\in\sigma(L_{\eps,u})$ and $\norm{\e^{tL_{\eps,u}}}\geq1$ for every $t\geq0$.
  Consequently $\gamma(u,\eps)$ is a real number $L$ with
  $0\leq L\leq M_u$ and $t^{-1}\log\norm{\e^{tL_{\eps,u}}}\to L$.
\end{lemma}

The first three facts are the paper's argument for $0\in\sigma(L_{\eps,u})$ in the proof of
Lemma~A.1; the paper does not state $\norm{\e^{tL_{\eps,u}}}\geq1$ or $\gamma\geq0$.  The mean
is the integral over the product torus with its probability Haar measure, which is Lebesgue
measure on $\T^3$ (\texttt{HasZeroMean}, \texttt{FastDynamo/Finish/ZeroMean.lean}).  Modules
\texttt{FastDynamo/Dynamics/Assembly/InductionMeanConservation.lean} and
\texttt{FastDynamo/Dynamics/Assembly/InductionSemigroupSpectrum.lean}.

\begin{proof}
  \leanok
  \uses{lem:dynamics:growth-bound-limit, lem:dynamics:generator-basic}
  The zeroth Fourier coefficient is the mean; that of $\eps\Delta B$ and of the weak curl
  vanishes.  So $L_{\eps,u}$ maps into the zero-mean fields and is not surjective onto
  $\Lp^2_\sigma$, which contains the constant field $c=(1,0,0)$: $0\in\sigma(L_{\eps,u})$.
  The mean of $\e^{tL_{\eps,u}}c$ has right derivative zero, so it stays $1$, and
  $\norm{\e^{tL_{\eps,u}}c}\geq1=\norm c$.  The minorant $1\leq\norm{\e^{tL_{\eps,u}}}$ makes
  the limit of Lemma~\ref{lem:dynamics:growth-bound-limit} finite and non-negative, and the
  Hille--Yosida bound gives $L\leq M_u$.
\end{proof}

\begin{theorem}[The spectral theory of the induction semigroup]
  \label{thm:dynamics:induction-spectrum}
  \lean{FastDynamo.isAnalyticSemigroup_solenoidalInductionSemigroup,
    FastDynamo.exists_isResolvent_generator_solenoidalInductionSemigroup,
    FastDynamo.isCompactOperator_of_isResolvent_generator_solenoidalInductionSemigroup,
    FastDynamo.isCompactOperator_solenoidalInductionSemigroup,
    FastDynamo.spectrum_solenoidalInductionSemigroup_diff_zero,
    FastDynamo.gamma_eq_iSup_re_spectrum,
    FastDynamo.exists_eigenvalue_re_eq_gamma}
  \leanok
  \uses{thm:dynamics:solenoidal-generator, lem:dynamics:sectorial,
    thm:dynamics:analytic-generation, lem:dynamics:compact-embedding,
    lem:dynamics:immediate-compactness, thm:dynamics:spectral-mapping,
    lem:dynamics:mean-conservation, prop:dynamics:product-resolvent}
  Let $u$ be transverse divergence-free and $\eps>0$.  The semigroup $\e^{tL_{\eps,u}}$ on
  $\Lp^2_\sigma$ is analytic of angle $\pi/4$; its generator $L_{\eps,u}$ has a non-empty
  resolvent set and a compact resolvent at every point of it; $\e^{tL_{\eps,u}}$ is compact
  for $t>0$; \eqref{eq:dynamics:compact-semigroup-spectral-mapping} holds; and
  $\gamma(u,\eps)=\sup\{\re\lambda:\lambda\in\sigma(L_{\eps,u})\}$, attained at an
  eigenvalue.
\end{theorem}

These are the instances, at the induction semigroup, of the abstract statements of this
section; Lemma~\ref{lem:semigroup-spectral-growth} collects them with the generator and the
limit.  Module \texttt{FastDynamo/Dynamics/Assembly/InductionSemigroupSpectrum.lean}.

\begin{proof}
  \leanok
  \uses{lem:dynamics:sectorial, thm:dynamics:analytic-generation,
    lem:dynamics:compact-embedding, lem:dynamics:immediate-compactness,
    thm:dynamics:spectral-mapping, lem:dynamics:mean-conservation}
  Analyticity: Theorem~\ref{thm:dynamics:analytic-generation} at the product Hille--Yosida ray
  with the sector condition of Lemma~\ref{lem:dynamics:sectorial}, then restriction to
  $\Lp^2_\sigma$.  Compact resolvent: at a real $r$ beyond the ray constant the resolvent of
  the generator is the restricted product resolvent
  (\texttt{exists\_isResolvent\_generator\_solenoidalInductionSemigroup\_eq}), which maps into
  $\Hs^2$ and is compact by Lemma~\ref{lem:dynamics:compact-embedding}; the resolvent identity
  transfers compactness to every point of the resolvent set.  The rest is
  Lemma~\ref{lem:dynamics:immediate-compactness} and
  Theorem~\ref{thm:dynamics:spectral-mapping}, the spectrum being non-empty by
  Lemma~\ref{lem:dynamics:mean-conservation}, and the growth bound of the product semigroup
  being that of its transport to the literal torus.
\end{proof}

\section{Solutions of the induction equation}
\label{sec:dynamics:solutions}

\begin{lemma}[The energy identity]
  \label{lem:dynamics:energy-identity}
  \lean{FastDynamo.productGradientEnergy,
    FastDynamo.inner_productH2Embedding_laplacian_eq_neg_gradientEnergy,
    FastDynamo.re_inner_productInductionTransverseTransport_eq_zero,
    FastDynamo.re_inner_productInductionAxialTransport_eq_zero,
    FastDynamo.re_inner_productInductionAction_eq,
    FastDynamo.hasDerivAt_half_norm_sq_of_isForwardStrongSolution,
    FastDynamo.hasDerivAt_half_norm_sq_of_isStrongSolution,
    FastDynamo.realStrongSolution_hasDerivAt_half_norm_sq}
  \leanok
  \uses{def:main:induction-operator, def:dynamics:semigroup-interface}
  Let $u$ be transverse divergence-free and $\eps\in\R$.  For every $B\in\Hs^2(\T^3;\CC^3)$,
  \[
    \re\inner{B}{\mathcal A_\eps B}=\re\inner{B}{(B\cdot\nabla)u}-\eps\norm{\nabla B}_2^2,
  \]
  and along every strong solution $B$ of $\partial_tB=\mathcal A_\eps B$ (at every positive
  time, for a forward strong solution)
  \[
    \frac12\frac{\dd}{\dd t}\norm{B(t)}_2^2
    =\re\inner{B(t)}{(B(t)\cdot\nabla)u}-\eps\norm{\nabla B(t)}_2^2 .
  \]
  In particular this holds along the real solution $B^\eps$ of Theorem~\ref{thm:main}.
\end{lemma}

This is the paper's \texttt{eq:energy-competition}, which the paper displays formally (``this
competition is expressed by the energy identity'').  The Lean identity is for complex fields with
the real part of the pairing; for a real field the pairing is the paper's
$\int_{\T^3}B\cdot(B\cdot\nabla)u\dd x$.  The proof uses no divergence condition on $B$:
both transport terms of $\mathcal A_\eps$ drop out because $\Div u=0$ and $u$ is independent of
$x_3$.  Off divergence-free fields, however, $\mathcal A_\eps$ is not $L_{\eps,u}$, so the
paper's identity is the case $\Div B=0$, where the auxiliary and the curl forms agree; along
solutions of the curl-form equation itself with divergence-free datum the identity is part of
Lemma~\ref{lem:dynamics:physical-uniqueness}.  Module
\texttt{FastDynamo/Dynamics/EnergyIdentity.lean}.

\begin{proof}
  \leanok
  \uses{def:main:induction-operator}
  In Fourier variables $\re\inner{B}{\eps\Delta B}=-\eps\norm{\nabla B}_2^2$; the transverse
  transport term is skew because $\Div_\perp u_\perp=0$, and the axial transport term
  $u_3\partial_3B$ is skew because $u_3$ does not depend on $x_3$.  Along a strong solution,
  $\frac{\dd}{\dd t}\norm{B}_2^2=2\re\inner{B}{\mathcal A_\eps B}$.
\end{proof}

\begin{lemma}[Uniqueness of the solution of the induction equation]
  \label{lem:dynamics:physical-uniqueness}
  \lean{FastDynamo.productHMinus2Synthesis,
    FastDynamo.exists_norm_stateResolvent_le_productL2ToHMinus2,
    FastDynamo.productHMinus2DivergenceFunctional,
    FastDynamo.productDivergenceFromH2_eq_zero_of_isForwardStrongSolution,
    FastDynamo.productPhysicalInductionOperator_forwardStrongSolutionsUnique,
    FastDynamo.isForwardStrongSolution_productPhysicalInductionOperator_domainOrbit,
    FastDynamo.productH2Embedding_eq_productInductionSemigroup_of_isForwardStrongSolution,
    FastDynamo.productH2Embedding_eq_solenoidalInductionSemigroup_of_isForwardStrongSolution,
    FastDynamo.hasDerivAt_half_norm_sq_of_isForwardStrongSolution_physical}
  \leanok
  \uses{def:finish:physical-operator, lem:dynamics:product-dissipative,
    thm:dynamics:product-semigroup, thm:dynamics:solenoidal-semigroup,
    lem:finish:physical-compatibility, lem:dynamics:energy-identity}
  Let $u$ be transverse divergence-free and $\eps>0$.  Two forward strong solutions
  $B_1,B_2:[0,\infty)\to\Hs^2(\T^3;\CC^3)$ of
  $\partial_tB=\curl(u\times B)+\eps\Delta B$, the time derivative taken in $\Hs^{-2}$, with
  $B_1(0)=B_2(0)$ agree for all $t\geq0$.  A forward strong solution with divergence-free
  initial value stays divergence-free, and it is the orbit
  $B(t)=\e^{tL_{\eps,u}}B(0)$ of the semigroup on $\Lp^2_\sigma$; conversely the orbit of
  every divergence-free $\Hs^2$ datum is a forward strong solution.  Along a forward strong
  solution with divergence-free initial value the energy identity
  \texttt{eq:energy-competition} holds at every $t>0$:
  $\frac12\frac{\dd}{\dd t}\norm{B(t)}_2^2
  =\re\inner{B(t)}{(B(t)\cdot\nabla)u}-\eps\norm{\nabla B(t)}_2^2$.
\end{lemma}

This is what the definite article in ``\emph{its} solution'' of Theorem~\ref{thm:main} asks
for.  The paper has no separate statement: Theorem~\ref{thm:main} asserts the norm law for
``its solution'', and its proof writes that solution down.  The solution class is the Lean
one: $\Hs^2$-valued curves on the product torus, with one-sided derivative at $t=0$ and
two-sided derivative at $t>0$ in the Fourier model of $\Hs^{-2}$, and no divergence condition
on the data; conclusions are for $t\geq0$.  It contains the class of clause~(16) of
Definition~\ref{def:main:physical-clauses}, whose curve is a strong solution at every real
time.  Modules \texttt{FastDynamo/Dynamics/Assembly/NegativeSobolevResolvent.lean},
\texttt{FastDynamo/Dynamics/Assembly/PhysicalInductionUniqueness.lean} and
\texttt{FastDynamo/Dynamics/Assembly/SemigroupOrbit.lean}.

\begin{proof}
  \leanok
  \uses{lem:dynamics:product-dissipative, lem:finish:physical-compatibility,
    lem:dynamics:energy-identity}
  The divergence of a forward strong solution satisfies, Fourier mode by Fourier mode, the
  heat equation $\dot c=-4\pi^2\eps\abs\xi^2c$, because the curl term has no divergence; so a
  divergence-free datum stays divergence-free.  On divergence-free fields the curl form is the
  $\Hs^{-2}$ image of $\mathcal A_\eps$ (Lemma~\ref{lem:finish:physical-compatibility}).  For the
  difference $w$ of two solutions, the resolvent $R(r)$ of $\mathcal A_\eps$ at a real $r>M_u$ is
  bounded in the $\Hs^{-2}$ norm (by the energy inequality and the skewness of transport), so
  $R(r)w(t)$ is an $\Lp^2$ forward strong solution of $\partial_t=\mathcal A_\eps$ with initial
  value $0$, hence $0$ by Lemma~\ref{lem:dynamics:product-dissipative}; $R(r)$ is injective, so
  $w=0$.  The same argument with the orbit of the realization of
  Theorem~\ref{thm:dynamics:product-semigroup} as second solution identifies a solution with
  divergence-free datum with the semigroup orbit.  That orbit is an $\Lp^2$ forward strong
  solution of $\partial_t=\mathcal A_\eps$, so Lemma~\ref{lem:dynamics:energy-identity} gives
  the energy identity along it, and the two curves agree on $[0,\infty)$.
\end{proof}

